\documentclass[UKenglish,acmsmall,authorversion]{acmart}

%\documentclass[UKenglish,acmsmall,authorversion]{acmart}
%\usepackage[utf8]{inputenc} %input font encoding
%\usepackage[latin1]{inputenc} 
%\usepackage[T1]{fontenc} % output font encoding


\usepackage{virginialake}
\usepackage{xypic}
\vlsmallopstrue


\usepackage{booktabs} % For formal tables
\usepackage{proof}
\usepackage{enumitem}

\newcommand{\tohide}[2]{#1} %{#1} % Change to #1 to show full proofs
%comments
\newcommand{\fullproof}[1]{} % Cases only for tech report, not journal.

\usepackage{lipsum}                     % Dummytext
\usepackage{xargs}                      % Use more than one optional parameter in a new commands
%\usepackage[dvipsnames]{xcolor}  % Coloured text etc.
%
\usepackage[colorinlistoftodos,prependcaption,textsize=small]{todonotes}
\newcommandx{\ross}[2][1=]
{\todo[linecolor=blue,backgroundcolor=blue!25,bordercolor=blue,#1]{Ross: #2}} 
% Uncomment to show my comments.
\newcommandx{\bogdan}[2][1=]{\todo[linecolor=red,backgroundcolor=red!25,bordercolor=red,#1]{Bogdan: #2}}
\newcommandx{\gabriel}[2][1=]{\todo[linecolor=green,backgroundcolor=green!25,bordercolor=green,#1]{Gabriel: #2}}
\newcommandx{\alwen}[2][1=]{\todo[linecolor=purple,backgroundcolor=purple!25,bordercolor=purple,#1]{Alwen: #2}}




\usepackage[T2A,T1]{fontenc}
%\usepackage[russian,english]{babel}

\usepackage{wrapfig}
\usepackage{etex}
\usepackage{amsmath,amsfonts,amsthm}
\theoremstyle{plain}
\usepackage{verbatim}
%\usepackage{txfonts}
\usepackage{fancyhdr}
%\usepackage{epsfig}
%\usepackage{latexsym}
\usepackage{microtype}
%\usepackage{lastpage} 
%\usepackage{hyperref}
%\usepackage{breakurl}
\usepackage{cmll}
\usepackage{prooftree}
\usepackage{rotating}
%\usepackage{bussproofs}
%\usepackage[table]{xcolor}
%\usepackage[noxy]{virginialake}
%\usepackage{xypic}
\usepackage{syndmeta}

\usepackage{flushend}

%\usepackage{graphicx}
%\usepackage{amssymb}
\usepackage{fancybox,tikz}
%\usepackage{makeidx}
\usetikzlibrary{positioning}

%\usepackage[dvipsnames*,svgnames,table]{xcolor}
%\hypersetup{colorlinks=true,allcolors=DarkBlue,linkcolor=black}

\newcommand{\AV}{\textsf{AV}}
\newcommand{\SAV}{\textsf{SAV}}
\newcommand{\terminates}[1]{#1 \checkmark}


\def\proof{{\noindent \bf Proof.~}}
\def\endproof{\hspace{\fill}{\boldmath $\Box$}\par}
%\newtheorem{counterexample}{Counterexample}

\title{De Morgan Dual 
 Nominal Quantifiers Modelling Private Names in Non-Commutative Logic}
%\subtitle{A Proof System with Polarised Nominal Quantifiers}

\acmJournal{TOCL}
% Metadata Information
\acmVolume{20}
\acmNumber{4}
\acmArticle{22}
\acmYear{2019}
\acmMonth{7}
\copyrightyear{2019}
%\acmArticleSeq{9}

% Copyright
%\setcopyright{acmcopyright}
\setcopyright{acmlicensed}
%\setcopyright{rightsretained}
%\setcopyright{usgov}
%\setcopyright{usgovmixed}
%\setcopyright{cagov}
%\setcopyright{cagovmixed}

% DOI
\acmDOI{10.1145/3325821}

% Paper history
\received{November 2017}
\received[revised]{September 2018}
\received[accepted]{April 2019}


\author[R.~Horne]{Ross Horne}
\affiliation{
Computer Science and Communications,
University of Luxembourg
}
\email{ross.horne@uni.lu}


\author[A.~Tiu]{Alwen Tiu}
\affiliation{
Research School of Computer Science,
The Australian National University, Australia
}
\email{alwen.tiu@anu.edu.au}

\author[B.~Aman]{Bogdan Aman}
\affiliation{
Alexandru Ioan Cuza University of Ia\c si, Romania
%Romanian Academy, Institute of Computer Science,
%Blvd.~Carol I no.8, 700505,
%Ia\c si, Romania
}
\email{bogdan.aman@iit.academiaromana-is.ro}

\author[G.~Ciobanu]{Gabriel Ciobanu}
\affiliation{
Alexandru Ioan Cuza University of Ia\c si, Romania
%Romanian Academy, Institute of Computer Science,
%Blvd.~Carol I no.8, 700505,
%Ia\c si, Romania
}
\email{gabriel@info.uaic.ro}



% --- begin abstract.tex ---
\begin{abstract}
This paper explores the proof theory necessary for recommending an expressive but decidable first-order system, named \textsf{MAV1}, featuring a de Morgan dual pair of nominal quantifiers. These nominal quantifiers called `new' and `wen' are distinct from the self-dual Gabbay-Pitts and Miller-Tiu nominal quantifiers. The novelty of these nominal quantifiers is they are polarised in the sense that `new' distributes over positive operators while `wen' distributes over negative operators. This greater control of bookkeeping enables private names to be modelled in processes embedded as formulae in \textsf{MAV1}. The technical challenge is to establish a cut elimination result, from which essential properties including the transitivity of implication follow.
Since the system is defined using the calculus of structures, a generalisation of the sequent calculus, novel techniques are employed. The proof relies on an intricately designed multiset-based measure of the size of a proof, which is used to guide a normalisation technique called \textit{splitting}. The presence of equivariance, which swaps successive quantifiers, induces complex inter-dependencies between nominal quantifiers, additive conjunction and multiplicative operators in the proof of splitting. Every rule is justified by an example demonstrating why the rule is necessary for soundly embedding processes and ensuring that cut elimination holds.
\end{abstract}



% --- end abstract.tex ---

\begin{document}


%
% The code below should be generated by the tool at
% http://dl.acm.org/ccs.cfm
% Please copy and paste the code instead of the example below. 
%
\begin{CCSXML}
<ccs2012>
<concept>
<concept_id>10003752.10003790.10003792</concept_id>
<concept_desc>Theory of computation~Proof theory</concept_desc>
<concept_significance>500</concept_significance>
</concept>
<concept>
<concept_id>10003752.10003753.10003761.10003764</concept_id>
<concept_desc>Theory of computation~Process calculi</concept_desc>
<concept_significance>500</concept_significance>
</concept>
<concept>
<concept_id>10003752.10003790.10003801</concept_id>
<concept_desc>Theory of computation~Linear logic</concept_desc>
<concept_significance>500</concept_significance>
</concept>
</ccs2012>
\end{CCSXML}

\ccsdesc[500]{Theory of computation~Proof theory}
\ccsdesc[500]{Theory of computation~Process calculi}
\ccsdesc[500]{Theory of computation~Linear logic}

%
% End generated code
%

% We no longer use \terms command
%\terms{Design, Algorithms, Performance}

\keywords{calculus of structures, nominal logic, non-commutative logic}


\thanks{
R.\ Horne\ and A.\ Tiu\ acknowledge support from Ministry of Education, Singapore, Tier 2 grant MOE2014-T2-2-076.
%A.\ Tiu\ acknowledges support from Nanayang Technological University Start Up grant M4081190.020.
%B.\ Aman and G.\ Ciobanu\ are supported by National Authority for Scientific Research, Romania, grant number PN-II-ID-PCE-2011-3-0919.

Authors' addresses: 
R.~Horne,
Faculty of Science, Technology and Communication,
6 avenue de la Fonte,
L-4364 Esch-sur-Alzette,
Luxembourg,
%50 Nanyang Avenue, Singapore 639798;
email:
\texttt{ross.horne@uni.lu};
A.~Tiu,
Research School of Computer Science,
The Australian National University,
Building 108, North Road,
Canberra, ACT 2061, Australia;
email: \texttt{alwen.tiu@anu.edu.au};
B.~Aman and G.~Ciobanu,
Alexandru Ioan Cuza University of Ia\c si, Romania;
email: \texttt{bogdan.aman@iit.academiaromana-is.ro}
and
\texttt{gabriel@info.uaic.ro}.
}


\maketitle



% --- begin introduction.tex ---
\section{Introduction}
\label{sec:intro}
%\ross{We should cut to the chase straight away, due to lack of space to motivate fully.}
%This paper introduces a calculus that straddles the fields of proof theory and process calculi. From process calculi we draw inspiration from applied process calculi for modelling protocols notably the applied $\pi$-calculus~\cite{Abadi2001}. We use calculi that are restricted to finite processes, i.e.\ without replication or fixed points, that have the advantage of being decidable while permitting complex protocols to be modelled.

This paper investigates the proof theory of a novel pair of de Morgan dual nominal quantifiers.
These quantifiers are motivated by the desire to model private name binders in processes by embedding the processes directly as formulae in a suitable logical system.
The logical system in which this investigation is conducted is sufficiently expressive to soundly embed the finite fragment of several process calculi.

A requirement of directly embedding processes as formulae is that the logic should be able to capture causal dependencies.
To do so, we employ a non-commutative multiplicative operator, which can be used to model the fact that `$a$ happens before $b$' is not equivalent to `$b$ happens before $a$'.
Such non-commutative operators are problematic for traditional proof frameworks such as the sequent calculus;
hence we adopt a formalism called the \textit{calculus of structures}~\cite{Guglielmi2007,Tiu2006,Guglielmi2011,Strassburger2011,Strassburger2003thesis}.
%The calculus of structures is effectively a term rewriting system modulo a congruence satisfying some desirable proof theoretic properties (in particular cut elimination).
The calculus of structures permits more proofs than the sequent calculus, by allowing inference rules to be applied in any context; while still satisfying proof theoretic properties, notably cut elimination.
An advantage of the calculus of structures is that it can express proof systems combining connectives for sequentiality and parallelism.
% not expressible in natural deduction or the sequent calculus.
%We assume the reader is comfortable with basic term rewriting, and the $\pi$-calculus.
% in order to obtain the key results.
%
The calculus of structures was motivated by a need for understanding
why pomset logic~\cite{Retore1997} could not be expressed in the
sequent calculus. Pomset logic is inspired by
pomsets~\cite{Pratt1986} and linear logic~\cite{Girard1987}, the
former being a model of concurrency respecting
causality, while the latter can be interpreted in
various ways as a logic of resources and
concurrency~\cite{Caires2010,Kobayashi1993,Wadler2014}.
%Acknowledging connections between pomsets, linear logic and concurrency, it is natural to consider the calculus of structures in the context of concurrency theory.
%This paper is part of a line of work clarifying this question.
%% Initial investigations have been conducted in
%% a basic process calculus based on \textsf{BV}~\cite{Bruscoli2002}; where it is
%% demonstrated that completed traces have a logical characterisation
%% using the calculus of structures; consequently logical equivalence
%% implies completed trace equivalence.


These observations lead to the propositional system \textsf{MAV}~\cite{Horne2015} and its first-order extension presented in this work, named \textsf{MAV1}.
Related work establishes that linear implication in such logical systems is sound with respect to both pomset ideals~\cite{Horne2017}
and weak simulation~\cite{MSCS}.
These results tighten results in initial investigations concerning a minimal calculus \textsf{BV} and trace inclusion~\cite{Bruscoli2002}.
Hence reasoning using linear implication is sound with respect to most useful (weak) preorders over processes,
for a range of languages not limited to \textsf{CCS}~\cite{Milner1982} and $\pi$-calculus~\cite{Milner1992}.


This paper resolves the fundamental logical problem of whether cut elimination holds for \textsf{MAV1}.
Cut elimination, the corner stone of a proof system, is essential for confidently recommending a proof system.
In the setting of the calculus of structures, cut elimination is formalised quite differently compared to traditional proof frameworks;
hence the proof techniques employed in this paper are of considerable novelty.
Furthermore, this paper is the first paper to establish cut elimination for a de Morgan dual pair of nominal quantifiers in any proof framework.
These nominal quantifiers introduce intricate interdependencies between other operators in the calculus, reflected in the technique of \textit{splitting} (Lemma~\ref{lemma:split-times}) which is the key lemma required to establish cut elimination (Theorem~\ref{theorem:cut}).



\begin{comment}
 relate a basic process calculus with parallel composition ($P \cpar Q$) and action prefix ($\alpha.P$)
to a proof system in the calculus of structures called \textsf{BV}~\cite{Guglielmi2007}.
In that work, a logical characterisation of completed traces is established; thereby, linear implication is strictly finer than completed trace inclusion.
Strictness follows since some processes related by linear implication are not related by trace inclusion. 
For example, a desirable property of linear implication is that \textit{autoconcurrency}~\cite{Bednarczyk1991,Glabbeek2001,Glabbeek2015} is avoided, since the embedding of $\alpha \cpar \alpha$ does \textbf{not} logically imply the embedding of $\alpha.\alpha$.
Avoiding autoconcurrency indicates that linear implication fully respects the causal order of events.
%
Preorders that respect causality are significant for various applications, including soundly reasoning about processes deployed on large distributed systems where guaranteeing a consensus on interleavings 
is infeasible but causality is respected~\cite{Lamport1978}.

Recently, the authors extended \textsf{BV} with additives to obtain the system
\textsf{MAV}~\cite{Horne2015}, enabling choice in processes to be modelled.
%% \alwen{What do you mean by `branching'?
%%   Do you mean the choice operator, or branching-time equivalence?} 
Results concerning multi-party compatibility and subtyping in a
session type system inspired by \textsf{Scribble}~\cite{Honda2011}
have been established using \textsf{MAV}~\cite{Ciobanu2015}.
The current paper continues this line of enquiry by modelling name passing processes in a conservative extension of \textsf{MAV} with first-order quantifiers,
named \textsf{MAV1}.  The system \textsf{MAV1} is also a conservative
extension of first-order multiplicative additive linear logic
\textsf{MALL1}~\cite{Lincoln1994b} with mix. 
A novelty is that \textsf{MAV1} includes a pair of de Morgan dual nominal quantifiers
pronounced \textit{new} and \textit{wen} and written $\new$ and $\wen$ respectively.
%These quantifiers were concerning name passing processes modelled as predicates.
In a processes-as-predicates embedding~\cite{Guglielminote},  $\new$ models new name restriction in the $\pi$-calculus.
The dual nominal quantifier $\wen$ is essential for defining linear implication and models the input of private names in a processes-as-predicates embedding for the $\pi I$-calculus~\cite{Sangiorgi1996}. %and is significant can define tests that only a fresh name can specify tests pass.
%Such embeddings allow us to study what sorts of process equivalence that logical equivalence gives rise to.
\end{comment}

Logically speaking, nominal quantifiers $\new$ and $\wen$, pronounced `new' and `wen' respectively, sit between $\forall$ and $\exists$ such that
%, i.e., we have the following implications:
$
\forall x. P \multimap \new x. P
$
and
$
\new x. P \multimap \wen x. P
$ and
$
\wen x. P \multimap \exists x. P
$, where $\multimap$ is linear implication.
The quantifier $\new$ is similar in some respects to $\forall$, whereas $\wen$
is similar to $\exists.$ A crucial difference between $\exists x. P$
and $\wen x. P$ is that variable $x$ in the latter cannot be instantiated with arbitrary terms,
but only `fresh' names introduced by $\new.$
%
Our \textit{new} quantifier $\new$, distinct from the Gabbay-Pitts quantifier, addresses limitations of established self-dual nominal quantifiers for modelling private names in embeddings of processes as formulae. 
In particular, our $\new$ quantifier does not distribute over parallel composition in either direction.
In \textsf{MAV1}, the formulae
$
\mathopen{\new x.}({\mathrm{event}(x)} \cpar {\mathrm{event}(x)})
$ and
$
\new x.{\mathrm{event}(x)} \cpar \new x.{\mathrm{event}(x)}
$
are unrelated by linear implication. This property is essential for soundly modelling private name binders in processes.
%This leads us to the discovery of the ``new'' operator and its de Morgan dual ``wen''.


\begin{comment}
\item $\exists$: Using the two processes $R_1$ and $R_2$ above as an example,
an encoding using $\exists$ are: 
$$
G_1 = \exists x.( \embed{P(x)} || \embed{Q(x)})
\qquad
G_2 = \exists x. \embed{P(x)} || \exists x.\embed{Q(x)}.
$$
In this case, we have that $G_1 \multimap G_2.$
This is also unsound, since $R_1$ has a trace that outputs two identical names,
whereas $R_2$ admits so such traces.

\item Self-dual nominal quantifiers: 
Rather surprisingly, encoding restriction with either Gabbay-Pitts quantifier $\new$~\cite{Pitts2003} or
Miller-Tiu's $\nabla$ quantifier~\cite{Tiu2009,Gacek2011} would not work either. Both quantifiers are self-dual
and for our purpose here, they are interchangeable, so
we illustrate our example using the $\new$ quantifier.
Processes $R_1$ and $R_2$ in this case are encoded, respectively, as predicates $H_1$ and $H_2$ below:
$$
H_1 = \new x.(\embed{P(x)} \cpar \embed{Q(x)})
\qquad
H_2 = \new x.\embed{P(x)} \cpar \new x.\embed{Q(x)}.
$$
One property of $\nabla$ (and $\new$)
is that it commutes with all propositional connectives.
\footnote{
  Self-dual nominal quantifiers have so far been formulated in classical or intuitionistic
  systems, but it is straightforward to introduce them in linear logic. 
}
So in this case we actually have a logical equivalence: $H_1 \multimap H_2$ and $H_2 \multimap H_1$, which is clearly unsound with respect
to the trace semantics.
\end{comment}

\begin{comment}
We illustrate why neither universal quantification nor an established self-dual nominal quantifier~\cite{Pitts2003,Tiu2009,Gacek2011} are capable of soundly modelling name restriction in a processes-as-predicates embedding.
%If we are to interpret the restriction operator as a logical operator, naturally the only possibilities are the quantifiers. We discuss possible choices of quantifiers below.
In the following, observe that $R_1 = \mathopen{\nu x.}(\co{a}x \cpar \co{b}x)$ is a $\pi$-calculus process that can output a fresh name twice, once on channel
$a$ and once on channel $b$; but cannot output two distinct names in any execution.
In contrast, observe that $R_2 = \nu x.\co{a}x \cpar \nu x.\co{b}x$ is a $\pi$-calculus process that outputs
two distinct fresh names before terminating, but cannot output the same name twice in any execution. The processes $R_1$ and $R_2$ must not be related by any sound preorder.
%We also assume parallel composition of processes is encoded using the \textit{par} operator $\cpar$ of linear logic.

For an encoding using universal quantifiers,
%$\forall$,
% a process $\nu x.P$ would be encoded as a predicate $\forall x \embed{P}$, where $\embed{P}$ is an encoding of $P$.
%Under this encoding scheme,
processes $R_1$ and $R_2$ are respectively encoded as predicates 
$P_1 = \mathopen{\forall x.}\left(\co{ax} \cpar \co{bx}\right)$
and
$P_2 = \forall x.\co{ax} \cpar \forall x.\co{bx}$,
%
%where $\cpar$ denotes the linear logic par connective.
where operator $\cpar$ is overloaded to connect parallel composition and \textit{par} from linear logic.
Unfortunately, the implication $P_2 \multimap P_1$ is provable. However, $R_1$ can output the same fresh name on two channels but $R_2$ cannot, so implication would not be sound with respect to trace inclusion.
%
Additionally, we must also avoid the following \textit{diagonalisation} property~\cite{Tiu2009}: $\forall x .\forall y .P(x,y) \multimap \forall z .P(z,z)$.
%\quad\mbox{and}\quad
%\exists z P(z,z) \multimap \exists x \exists y P(x,y)

The self-dual nominal quantifiers of either Gabbay-Pitts~\cite{Pitts2003} or Miller-Tiu~\cite{Tiu2009,Gacek2011}, as recently investigated in the calculus of structures~\cite{Roversi14}, do successfully avoid the above diagonalisation property.
Unfortunately, rather surprisingly, encoding private names using any of these self-dual nominal quantifiers, say $\nabla$, leads to the following problem.
%Both quantifiers are self-dual and for our purpose here, they are interchangeable, so we illustrate our example using the $\nabla$ quantifier.
Suppose processes $R_1$ and $R_2$ are encoded by the respective predicates 
$
Q_1 = \mathopen{\nabla x.} (\co{ax} \cpar \co{bx})
$ and
$
Q_2 = \nabla x.\co{ax} \cpar \nabla x.\co{bx}.
$
In this case, the linear implication $Q_1 \multimap Q_2$ is provable.
This implication is also unsound, since $R_1$ has a trace that outputs two identical names, whereas $R_2$ admits no such trace.
%The scenario of an encoding using $\exists$ also suffers from the above problem, as well as the dual `diagonalization' property.

Our \textit{new} quantifier $\new$, distinct from the Gabbay-Pitts operator, addresses the above limitations of universal quantification and established self-dual nominal quantifiers. In addition to avoiding diagonalisation, our $\new$ quantifier does not distribute over parallel composition in either direction.
In \textsf{MAV1}, the predicates
$
\mathopen{\new x.}(\co{ax} \cpar \co{bx})
$ and
$
\new x.\co{ax} \cpar \new x.\co{bx}
$
are, correctly, unrelated by linear implication.
%This leads us to the discovery of the ``new'' operator and its de Morgan dual ``wen''.
\end{comment}

%% Note that there are existing work on a deep-embedding of pi-calculus in logic,
%% e.g., as done in \cite{tiu,bengtson}. In those works, processes are encoded as terms,
%% and judgments about processes (e.g., transitions, bisimulation) are encoded as
%% atomic formulas in logic. In such a setting, self-dual nominal quantifiers are sufficient,
%% but the logical correspondence between logic and process calculus is not as strong
%% as in the shallow embedding approach.

%% To model fresh names a nominal quantifier is required. The difficulty is that, in this linear setting, established self-dual nominal quantifiers~\cite{Pitts03,Gacek2011,Tiu2009}, as recently investigated in the calculus of structures~\cite{Roversi14}, have the weakness that any two fresh names can be unified during proof search. A self-dual nominal quantifier would thereby wrongly permit communications that are forbidden in process calculi with a name restriction operator, such as $\pi$-calculus~\cite{Milner1992}. Examples demonstrate how this weakness is avoided by employing the dual pair of nominal quantifiers `new' and `wen'.


%significant in the development of methods for reasoning about processes that executing on systems where  on systems where respecting causality is essential, which is becomming the norm for scalable distributed systems.
%The long term aim is to demonstrate the applications of \textsf{MAV1} to the verification of protocols.
%and its relation to formalisms such as the applied $\pi$-calculus~\cite{Abadi2001} and dependent session types~\cite{Yoshida2010,Toninho2011}.
\textbf{Outline.}
For a new logical system it is necessary to justify correctness, which we approach in proof theoretic style by cut elimination.
Section~\ref{section:self-dual} illustrates why an established self-dual nominal quantifier~\cite{Gabbay2002,Pitts2003,Tiu2009,Gacek2011} is incapable of soundly modelling name restriction in a processes-as-formulae embedding.
Section~\ref{section:syntax} defines \textsf{MAV1}, explains cut elimination and discusses rules.
%Section~\ref{section:process} illustrates a processes-as-predicates embedding in \textsf{MAV1} and explains that linear implication defines a branching-time preorder that respects causality.
Section~\ref{section:nominal} presents an explanation of the rules for the nominal quantifiers.
Section~\ref{section:splitting} presents technical lemmas and the \textit{splitting} technique which is key to cut elimination.
Section~\ref{section:context} presents a context lemma which is used to eliminate \textit{co-rules} that form a cut; thereby establishing cut elimination.
Section~\ref{section:complexity} explains the complexity classes for various fragments of \textsf{MAV1}.
%Note that existing cut elimination proofs involving first-order quantifiers in the calculus of structures~\cite{Brunnler2006} rely on a correspondence with the sequent calculus that is impossible in the presence of the self-dual non-commutative operator \textit{seq} that models sequential composition.

The cut elimination result in this article was announced at CONCUR 2016~\cite{Horne2016}, without full proofs.
This journal version of the paper explains the cut elimination proof, elaborates on the motivating discussion, and highlights further corollaries of cut elimination.
%The symbol $\wen$ follows the tradition of Gabbay and Pitts of using Cyrillic vowels for nominal quantifiers. This Cyrillic vowel is pronounced as the hard e in `wen' and reminds the reader of its existential nature.
Since $\new$ is a Cyrillic vowel, we use another Cyrillic vowel $\wen$ for nominal quantifier 'wen'. This Cyrillic vowel is pronounced as the hard e in `wen' and reminds the reader of its existential nature.

Due to the space limitation, some proofs are omitted in the printed version of this article, but are available in the accompanying Electronic Appendix.



%The result is a calculus that can model many scenarios that the finite applied $\pi$-calculus~\cite{Abadi2004,Kremer2005,Smyth2011} can model, and furthermore introduces connectives that are useful for the static analysis of protocols, including external choice and synchronisation.


%In this paper we consider an expressive calculus that extends the finite (exponential-free) first-order linear logic, known as \textsf{MALL1}~\cite{Lincoln1994}, with a self-dual non-commutative operator corresponding to sequential composition in the theory of pomsets. 
%For discussion, we also include a self-dual renaming quantifier~\cite{Retore1997}.

%To illustrate the versatility of our calculus we present several examples of security protocols and fault-tolerant protocols. We adapt and extend several examples on security protocols modelled in the applied $\pi$-calculus, while the fault-tolerant protocol examples draw inspiration from scalable databases~\cite{Corbett2012,Ghemawat2003} that are critical components for Cloud providers.

%For a more objective justification of the presented deductive system, we provide a cut elimination result for our calculus. This result supports the claim that the finite applied calculus of structure is a consistent logical system in its own right. Among the consequences of cut elimination, are that we can employ the semantics of our calculus to compare the behaviour of processes modelled by predicates in the calculus.

\begin{comment}
\subparagraph*{Structure of paper.}
%\tohide{}{\footnote{For review purposes proofs are available at \url{http://iit.iit.tuiasi.ro/TR/reports/fml1502.pdf}.}}}
Section~\ref{section:syntax} introduces the syntax and semantics of \textsf{MAV1}.
%In Section~\ref{section:applied}, we exemplify the ``applied'' part of the name of our calculus by considering a range of protocols to be modelled using the predicates of our calculus.
Section~\ref{section:implication} %we justify the ``calculus of structures'' part of the name of our calculus, by
settles cut elimination for \textsf{MAV1}. %that provides an objective justification for the correctness of our term rewriting system.
%Section~\ref{section:complexity} summarises complexity classes for sub-systems of \textsf{MAV1}.
%In Section~\ref{section:consequences}, we explore some useful consequences of the cut elimination result. Conclusions and references end the~paper.
%Section~\ref{section:process}, discusses embeddings of processes in \textsf{MAV1}.
\end{comment}

%\tohide{
%
%}
%{}





\begin{comment}
\subsection{Notes for motivation}

Since trace equivalence is among the weakest equivalence relations for processes~\cite{Glabbeek1990},
it makes sense to impose a minimum requirement that logical implication should be
sound with respect to trace inclusion, and that logical equivalence implies trace equivalence.
%Bruscoli~\cite{Bruscoli2002} shows that CCS processes can be encoded as propositions in BV such that implication implies trace inclusion.
To extend results for completed trace equivalence~\cite{Bruscoli2002} to the $\pi$-calculus, we require an encoding of the
restriction operator $\nu$ of the $\pi$-calculus. 




\begin{comment}
\item $\exists$: Using the two processes $R_1$ and $R_2$ above as an example,
an encoding using $\exists$ are: 
$$
G_1 = \exists x.( \embed{P(x)} || \embed{Q(x)})
\qquad
G_2 = \exists x. \embed{P(x)} || \exists x.\embed{Q(x)}.
$$
In this case, we have that $G_1 \multimap G_2.$
This is also unsound, since $R_1$ has a trace that outputs two identical names,
whereas $R_2$ admits so such traces.

\item Self-dual nominal quantifiers: 
Rather surprisingly, encoding restriction with either Gabbay-Pitts quantifier $\new$~\cite{Pitts2003} or
Miller-Tiu's $\nabla$ quantifier~\cite{Tiu2009,Gacek2011} would not work either. Both quantifiers are self-dual
and for our purpose here, they are interchangeable, so
we illustrate our example using the $\new$ quantifier.
Processes $R_1$ and $R_2$ in this case are encoded, respectively, as predicates $H_1$ and $H_2$ below:
$$
H_1 = \new x.(\embed{P(x)} \cpar \embed{Q(x)})
\qquad
H_2 = \new x.\embed{P(x)} \cpar \new x.\embed{Q(x)}.
$$
One property of $\nabla$ (and $\new$)
is that it commutes with all propositional connectives.
\footnote{
  Self-dual nominal quantifiers have so far been formulated in classical or intuitionistic
  systems, but it is straightforward to introduce them in linear logic. 
}
So in this case we actually have a logical equivalence: $H_1 \multimap H_2$ and $H_2 \multimap H_1$, which is clearly unsound with respect
to the trace semantics.

If we are to interpret the restriction operator as a logical operator, naturally the only
possibilities are the quantifiers.
We discuss possible choices of quantifiers below.
In the following illustration, suppose $\co{a}x$ is a process that outputs $x$ on a channel $a$ and then terminates,
and $\co{b}x$ is another process that does the same but on channel $b$.
Observe that $R_1 = \nu x.(\co{a}x ~|~ \co{b}x)$ is a process that can output a fresh name twice --- once on channel
$a$ and once on channel $b$; and, more importantly, it cannot output two distinct names in any execution.
On the other hand, observe that the process $R_2 = \nu x.\co{a}x ~|~ \nu x.\co{b}x$ is a process that outputs
two distinct fresh names before terminating, but cannot output the same name twice in any execution. We also assume parallel composition of processes is encoded using the
\textit{par} operator $\cpar$ of linear logic.


For an encoding using universal quantifiers,
%$\forall$,
% a process $\nu x.P$ would be encoded as a predicate $\forall x \embed{P}$, where $\embed{P}$ is an encoding of $P$.
%Under this encoding scheme,
the processes $R_1$ and $R_2$ are respectively encoded as predicates 
$F_1 = \forall x.\left(\co{ax} \cpar \co{bx}\right)$
and
$F_2 = \forall x.\co{ax} \cpar \forall x.\co{bx}$.
%
%where $\cpar$ denotes the linear logic par connective.
For these predicates, we have the implication $F_2 \multimap F_1$. But clearly $R_1$  has a trace that cannot
be produced by $R_2$, so implication would not be sound with respect to trace inclusion.
%
Additionally, we must also avoid the following `diagonalization' property~\cite{Tiu2009} that is again unsound with respect to trace semantics:
$
\forall x .\forall y .P(x,y) \multimap \forall z .P(z,z)$.
%\quad\mbox{and}\quad
%\exists z P(z,z) \multimap \exists x \exists y P(x,y)

The self-dual nominal quantifiers of either Gabbay-Pitts~\cite{Pitts2003} or Miller-Tiu~\cite{Tiu2009,Gacek2011}, as recently investigated in the calculus of structures~\cite{Roversi14}, do successfully avoid the above `diagonalization' property.
Unfortunately, rather surprisingly, encoding private names using any of these self-dual nominal quantifiers, say $\nabla$, leads to the following problem.
%Both quantifiers are self-dual and for our purpose here, they are interchangeable, so we illustrate our example using the $\nabla$ quantifier.
Suppose processes $R_1$ and $R_2$ are encoded by the respective predicates 
$
G_1 = \nabla x.(\co{ax} \cpar \co{bx})
$ and
$
G_2 = \nabla x.\co{ax} \cpar \nabla x.\co{bx}.
$
In this case, we have that $G_1 \multimap G_2.$
This implication is also unsound, since $R_1$ has a trace that outputs two identical names,
whereas $R_2$ admits no such trace. The scenario of an encoding using $\exists$ also suffers from the above problem, as well as the dual `diagonalization' property.
Thus, for a sound embedding of $\pi$-processes as predicates, the ``new'' quantifier is necessary since it does not distribute over the \textit{par} operator.
%This leads us to the discovery of the ``new'' operator and its de Morgan dual ``wen''.


\end{comment}
% --- end introduction.tex ---
% --- begin self-dual.tex ---
\section{Why not a self-dual nominal quantifier?}
\label{section:self-dual}

Nominal quantifiers in the literature are typically self-dual in the sense of de Morgan dualities.
That is, for a nominal quantifier, say $\nabla$, ``not $\nabla x$ $P$'' is equivalent to ``$\nabla x$ not $P$.''
Such self-dual nominal quantifiers have been successfully introduced in classical and intuitionistic frameworks, typically used to reason about higher-order abstract syntax with name binders.
Such nominal frameworks are therefore suited to program analysis, where the semantics of a programming language are encoded as a theory over terms in the logical framework.

Rather surprisingly, when processes themselves are directly embedded as formulae in a logic, where constructs are mapped directly to primitive logical connectives 
(as opposed to terms inside a logical encoding of the semantics of processes),
self-dual quantifiers do not exhibit typical properties expected of name binders.
To understand this problem, in this section we recall an established calculus \textsf{BVQ}~\cite{Roversi14} that can directly embed processes but features a self-dual nominal quantifier.
We explain that such a self-dual quantifier provides an unsound semantics for name binders.
This motivates the need for a finer polarised nominal quantifier, which leads to the calculus introduced in subsequent sections.

We assume the reader has a basic understanding of 
%term-rewriting systems~\cite{Meseguer1992} and
the semantics of the $\pi$-calculus~\cite{Milner1992} and $\textsf{CCS}$~\cite{Milner1982}.
%but may not necessarily be experts in deep inference or nominal logic.
This section provides necessary preliminaries for the calculus of structures.

\subsection{An established extension of \textsf{BV} with a self-dual quantifier}

An abstract syntax for formulae and the rules of \textsf{BVQ} are defined in Fig~\ref{figure:BVQ}.
In an inference rule, the formula appearing above the horizontal line is the premise and the formula below the horizontal line is the conclusion.
%  can be applied to rewrite a formula of the form on the left of the long right arrow to the formula on the right (note rewriting --- the direction of proof search --- is in the opposite direction to linear implication).
The key feature of the calculus of structures is \textit{deep inference}, which is the ability to apply
all rules in any context, i.e.\ formulae with a hole of the following form: $\context{} \Coloneqq \ehole{\,\cdot\,} \mid \context{} \odot P \mid P \odot \context{} \mid \nabla x . \context{}$, 
where $\odot \in \left\{ \cseq, \cpar, \tensor \right\}$.

% and $\quantifier \in \left\{ \exists, \forall, \new, \wen \right\}$. %, \rename{} \right\}$.


\begin{figure}
 \begin{center}
  %\vspace{-30pt}
\[
\begin{array}{l}
\mbox{\footnotesize Structural rules}
\\[4pt]
\begin{tabular}{c}
%$\left(P, \cpar, \cunit\right)\mbox{ is a commutative monoid}$\\$\left(P, \tensor, \cunit\right)\mbox{ is a commutative monoid}$
$\left(P, \cpar, \cunit\right)\mbox{ and }\left(P, \tensor, \cunit\right)\mbox{ are commutative monoids}$
%\quad \quantifier x. P \equiv \mathopen{\quantifier y .}\left(P\sub{x}{y}\right) \mbox{ if $\nfv{y}{\quantifier x. P}$} ~~\mbox{ ($\alpha$-conversion)}
\\[9pt]
$\left(P, \cseq, \cunit\right)\mbox{ is a monoid}$
\qquad
\mbox{$\alpha$-conversion for $\nabla$ quantifier}
\\[9pt]
$\nabla x .\nabla y. P \equiv \nabla y .\nabla x. P
~~\mbox{(equivariance)}$
\\[9pt]
$\nabla x . P \equiv P
~~\mbox{only if $\nfv{x}{P}$}
~~\mbox{(vacuous)}$
\end{tabular}
\end{array}
\qquad\quad
\begin{array}{l}
\mbox{\footnotesize Syntax}
\\[4pt]
\begin{tabular}{l@{\hspace{1ex}}l@{\hspace{1ex}}l@{\hspace{1ex}}l}
$P$&$::=$ & $\cunit$ & (unit)\\
&& $\alpha$ & (atom)\\
&& $\overline{\alpha}$ & (co-atom)\\
&& $\nabla{x}. P$ & (nabla)\\
&& $P \cpar P$ & (par)\\
&& $P \tensor P$ & (times)\\
&& $P \cseq P$ & (seq)
\end{tabular}
\end{array}
\]
\[
\begin{array}{l}
\mbox{\footnotesize Inference rules}
\\[4pt]
\begin{array}{c}
\infer[\mbox{(atomic interaction)}]{
\context{ \co{\alpha} \cpar \alpha }
}{
\context{ \cunit }
}
\qquad\qquad
\infer[\mbox{(switch)}]{
\context{P \cpar \left(Q \tensor S\right)}
}{
\context{(P \cpar Q) \tensor S}
}
\\[9pt]
\infer[\mbox{(sequence)}]{
\context{(P \cseq Q) \cpar (R \cseq S)}
}{
\context{(P \cpar R) \cseq (Q \cpar S)}
}
\qquad\qquad
\infer[\mbox{(unify)}]{
\context{ \nabla x. P \cpar \nabla x. Q }
}{
\context{ \mathopen{\nabla x .}\left( P \cpar Q \right) }
}
\end{array}\end{array}
\]
\end{center}
  %\vspace{-20pt}
\caption{Syntax and rules of system \textsf{BVQ}~\cite{Roversi14}: which is \textsf{BV} extended with a self-dual nominal quantifier.
}\label{figure:BVQ}
\end{figure}


Inference rules are defined \textit{modulo a structural congruence}, where a congruence is an equivalence relation that holds in any context.
A \textit{derivation} is a sequence of rules from Fig.~\ref{figure:BVQ}, where the structural congruence can be applied at any point in a derivation.
The length of a derivation involving only the structural congruence is zero. 
The length of a derivation involving one inference rule instance is one.
Given a derivation $\vcenter{\infer{Q}{P}}$ of length $m$ and another $\vcenter{\infer{R}{Q}}$ of length $n$, the derivation $\vcenter{\infer{R}{P}}$ is of length $m+n$.
Unless we make it clear in the context that we refer to a specific rule, this horizontal line notation is generally used to represent derivations of any length.
For example, since $\nabla x.\cunit \equiv \cunit$, derivation $\infer{\nabla x.\cunit}{\cunit}$  of length 0, and derivation 
$\vcenter{\infer{\left(P \tensor Q\right) \cpar R \cpar S}{\left(P \cpar R\right) \tensor \left(Q \cpar S\right)}}$ is of length 2, since two instances of \textit{switch} are applied.


%The structural congruence may be applied at any point in a   be applied In this work, we define use $\infer{Q}{P}$ to mean any 

The congruence, $\equiv$ in Fig.~\ref{figure:BVQ}, makes \textit{par} and \textit{times} commutative and \textit{seq} non-commutative in general. %Specifically, $\left(A,\cpar,\cunit\right)$ and $\left(A,\tensor,\cunit\right)$ are commutative monoids and $(A,\cseq,\cunit)$ is a monoid.
For the nominal quantifier $\nabla$, the congruence enables: $\alpha$-conversion for renaming bound names;
\textit{equivariance} which allows names bound by successive nominal quantifiers to be swapped; and \textit{vacuous} that allows the nominal quantifier to be introduced or removed whenever the bound variable does not appear in the formula. 
As standard, we define a freshness predicate such that a variable $x$ is fresh for a formulae $P$, written $\nfv{x}{P}$, if and only if $x$ is not a member of the set of free variables of $P$, where $\nabla x.P$ binds occurrences of $x$ in $P$.

Consider the syntax and rules of $\textsf{BVQ}$ in Figure~\ref{figure:BVQ}.
The three rules \textit{atomic interaction} and $\textit{switch}$ and \textit{sequence} define the basic system \textsf{BV}~\cite{Guglielmi2007} that also forms the core of the system \textsf{MAV1} investigated in later sections.
The only additional inference rule for $\nabla$ is called \textit{unify}.

\textbf{Atomic interaction.}
The atomic interaction rule should remind the reader of the classical tautology $\neg \alpha \vee \alpha$ or intuitionistic axiom $\alpha \Rightarrow \alpha$, applied only to the predicates forming the atoms of the calculus.
%\ross{atomic predicates or predicates}
Since there is no contraction rule for $\cpar$, once atoms are consumed by \textit{atomic interaction} they cannot be reused. Thus \textit{atomic interaction} is useful for modelling communication in process, where $\alpha$ models a receive action or event and $\co{\alpha}$ is the complementary send, which cancel each other out.
% by committing to an interaction.

\textbf{Switch and sequence.}
The \textit{atomic interaction} and $\textit{switch}$ rules together provide a model for multiplicative linear logic (with \textit{mix})~\cite{Girard1987}.
The difference between $\cpar$ and $\tensor$ is that $\cpar$ allows interaction, but $\tensor$ does not. In this sense the switch rule restricts where which atoms may interact.
The \textit{seq} rule also restricts where interactions can take place, but, since \textit{seq} is non-commutative, it can be used to capture causal dependencies between atoms.
The \textit{sequence} rule preserves these causal dependencies, while permitting new causal dependencies. In terms of process models, the \textit{sequence} rule appears in the theory of pomsets~\cite{Gischer1988} and can refine parallel composition to its interleavings.

\textbf{Unify.}
The novel rule for \textsf{BVQ} is \textit{unify} for nominal quantifier $\nabla$.
The \textit{unify} rule should be admissible in a well-designed extension of linear logic with a self-dual quantifier. To see why, consider the following auxiliary definitions.
Observe that the following definition of linear implication ensures that 
$\nabla$ is self-dual in the sense that the de Morgan dual of $\nabla$ is $\nabla$ itself. Similarly, \textit{seq} and the unit are self-dual, while $\tensor$ and $\cpar$ are a de Morgan dual pair of operators. 
\begin{definition}
\textit{Linear negation} is defined by the following function over formulae.
\begin{gather*}
\co{\cunit} = \cunit
\qquad
\co{\co{\alpha}} = \alpha
\qquad
\co{ P \tensor Q } = \co{P} \cpar \co{Q}
\qquad
\co{ P \cpar Q } = \co{P} \tensor \co{Q}
\qquad
\co{P \cseq Q} = \co{P} \cseq \co{Q}
\qquad
\co{\nabla x. P} = \nabla x. \co{P}
\end{gather*}
\textit{Linear implication}, written $P \multimap Q$, is defined as $\co{P} \cpar Q$.
\end{definition}

%A derivation is a sequence of zero or more rewrite rules from Fig.~\ref{figure:rewrite}, where the congruence in Fig.~\ref{figure:congruence} can be applied at any point.
We are particularly interested in special derivations, called proofs.
\begin{definition}
A \textit{proof} is a derivation of any length with conclusion $P$ and premise $\cunit$. When such a derivation exists, we say that $P$ is provable, and write $\vdash P$ holds.
\end{definition}

As a basic property of linear implication $\vdash P \multimap P$ must hold for any $P$. Now assume that $\vdash Q \multimap Q$ is provable in \textsf{BVQ} (hence, by the above definitions, there exists a derivation with conclusion $\co{Q} \cpar Q$ and premise $\cunit$), and consider formula $\nabla x. Q$. Using the \textit{unify} rule and the definition of linear implication, we can construct the following proof of $\vdash \nabla x.Q \multimap \nabla x.Q$.
\[
%\infer[\mbox{by definition}]{
%\nabla x.Q \multimap \nabla x.Q
%}{
\infer[\mbox{by the \textit{unify} rule}]{
\nabla x.\co{Q} \cpar \nabla x.Q 
}{
\infer[\mbox{by the assumption $\vdash \co{Q} \cpar Q$}]{
\mathopen{\nabla x.}\left(\co{Q} \cpar Q\right)
}{
\infer[\mbox{by the \textit{vacuous} rule}]{
\nabla x.\cunit
}{
\cunit
}}}
%}
\]
The above illustrates why \textit{unify} should be admissible in order to guarantee \textit{reflexivity} --- the most basic property of implication --- for an extension of \textsf{BV} with a self-dual nominal quantifier.
In the next section, we explain why the \textit{unify} rule is problematic for modelling processes as formulae.


\subsection{Fundamental problems with a self-dual nominal for embeddings of processes}

Initially, it seems that desirable properties of name binding, typical of process calculi, are achieved in \textsf{BVQ}.
For example, we expect that if $\nfv{x}{Q}$ then $\vdash \nabla x.\left( P \cpar Q \right) \multimap \nabla x.P \cpar Q$,
indicating that the scope of a name can be \textit{extruded} as long as another name is not captured,
which is provable using the \textit{vacuous} and \textit{unify} rules.
The \textit{equivariance} rule that swaps name binders is also a property preserved by most equivalences over processes.

Another strong property of \textsf{BVQ}, expected of all nominal quantifiers, is that we avoid the \textit{diagonalisation} property.
Diagonalisation $\vdash \forall x .\forall y .P(x,y) \multimap \forall z .P(z,z)$ holds in any system with universal quantifiers, as does the converse for existential quantifiers.
However, for nominals such at $\nabla$, \textbf{neither} $\nabla x .\nabla y .P(x,y) \multimap \nabla z .P(z,z)$ \textbf{nor} its converse $\nabla z .P(z,z) \multimap \nabla x .\nabla y .P(x,y)$ hold.
This is a critical feature of all nominal quantifiers that ensures that distinct fresh names in the same scope never collapse to the same name, and explains why universal and existential quantifiers are not suited modelling fresh name binders.
It is precisely the absence of diagonalisation for nominals that is used in classical~\cite{Pitts2003,Gabbay2002} and intuitionistic frameworks~\cite{Tiu2009,Gacek2011}
to logically manage the bookkeeping of fresh name in, so called, \textit{deep embeddings} of processes as terms in a theory.
Avoiding diagonalisation is sufficient in such deep embeddings since nominal quantifiers cannot appear inside a term representation of a process, so are always pushed to the outermost level where formulae are used to define the operational semantics of processes as a theory over process terms.

\textbf{Soundness criterion.}
The problem with \textsf{BVQ} is that when processes are directly embedded as formulae $\nabla$ quantifiers may appear inside embeddings of processes, which can result in unsound behaviours.
To see why the \textit{unify} rule induces unsound behaviours consider the following $\pi$-calculus terms. 
% $\mathopen{\nu x.}(\co{a}x \cpar \co{b}x)$ and $R_2$.
$\mathopen{\nu x.}(\co{z}x \mathrel{|} \co{y}x)$ is a $\pi$-calculus process that can output a fresh name twice, once on channel
$z$ and once on channel $y$; but cannot output two distinct names in any execution.
In contrast, observe that $\nu x.\co{z}x \mathrel{|} \nu x.\co{y}x$ is a $\pi$-calculus process that outputs
two distinct fresh names before terminating, but cannot output the same name twice in any execution.
As a soundness criterion, since the processes $\mathopen{\nu x.}(\co{z}x \mathrel{|} \co{y}x)$ and $\nu x.\co{z}x \mathrel{|} \nu x.\co{y}x$ do not have any complete traces in common, these processes must not be related by any sound preorder over processes.

Now consider an embedding of these processes in \textsf{BVQ}, where the parallel composition operator of the $\pi$-calculus is encoded as \emph{par} and $\nu$ is encoded as $\nabla$.
This gives us the formulae
$\mathopen{\nabla x.} \left(\co{\mathrm{act}(z, x)} \cpar \co{\mathrm{act}(y, x)}\right)$
and
$\nabla x.\co{\mathrm{act}(z, x)} \cpar \nabla x.\co{\mathrm{act}(y, x)}$.
Note that output action prefixes are encoded as negated predicates, e.g., $\co{z}x$ is
encoded $\co{\mathrm{act}(z,x)}$. 
%However, to simplify presentation, in the following discussion and for the rest of the paper, 
%we shall omit the $act$ predicate symbol in the process embedding
%and would simply write $zx$ as an atomic predicate in logic.

Observe that 
$\vdash
\mathopen{\nabla x.} \left(\co{\mathrm{act}(z,x)} \cpar \co{\mathrm{act}(y,x)}\right)
\multimap
\nabla x.\co{\mathrm{act}(z,x)} \cpar \nabla x.\co{\mathrm{act}(y,x)}
$
is provable, as follows. %is unsound for trace inclusion.
\[
\infer[\mbox{by \textit{unify}}]{
\mathopen{\nabla x.} \left({\mathrm{act}(z,x)} \tensor {\mathrm{act}(y,x)}\right)
\cpar
\nabla x.\co{\mathrm{act}(z,x)} \cpar \nabla x.\co{\mathrm{act}(y,x)}
}{
\infer[\mbox{by \textit{unify}}]{
\mathopen{\nabla x.} \left({\mathrm{act}(z,x)} \tensor {\mathrm{act}(y,x)}\right)
\cpar
\mathopen{\nabla x.}\left(\co{\mathrm{act}(z,x)} \cpar \co{\mathrm{act}(y,x)}\right)
}{
\infer[\mbox{by \textit{switch}}]{
\mathopen{\nabla x.} \left(\left({\mathrm{act}(z,x)} \tensor {\mathrm{act}(y,x)}\right)
\cpar
\co{\mathrm{act}(z,x)} \cpar \co{\mathrm{act}(y,x)}\right)
}{
\infer[\mbox{by \textit{switch}}]{
\mathopen{\nabla x.} \left(\left(\left({\mathrm{act}(z,x)} \cpar \co{\mathrm{act}(z,x)}\right) \tensor {\mathrm{act}(y,x)}\right)
 \cpar \co{\mathrm{act}(y,x)}\right)
}{
\infer[\mbox{by \textit{atomic interaction}}]{
\mathopen{\nabla x.} \left(\left({\mathrm{act}(z,x)} \cpar \co{\mathrm{act}(z,x)}\right)
\tensor
\left({\mathrm{act}(y,x)} \cpar \co{\mathrm{act}(y,x)}\right)\right)
}{
\infer[\mbox{by \textit{atomic interaction}}]{
\mathopen{\nabla x.} \left({\mathrm{act}(y,x)} \cpar \co{\mathrm{act}(y,x)}\right)
}{
\infer[\mbox{by \textit{vacuous}}]{
\mathopen{\nabla x.} \cunit
}{
\cunit
}}}}}}}
\]
The above implication is \textbf{unsound} with respect to trace inclusion for the $\pi$-calculus. 
The implication wrongly suggests that the process $\nu x.\co{z}x \mathrel{|} \nu x.\co{y}x$, that cannot output the same names twice, can be refined to a process $\mathopen{\nu x.}(\co{z}x \mathrel{|} \co{y}x)$, that outputs the same name twice.
This is exactly the contradiction that we avoid by using polarised nominal quantifiers investigated in subsequent sections.

As a further example of unsoundness issues for a self-dual nominal, consider the following criterion: an embedding of a process is provable if and only if there is a series of internal transitions leading to a successful termination state.
A successful termination state is a state without any unconsumed actions.
Now consider the process $\mathopen{\nu x.} \left(x.y\right) \mathrel{|} \nu z.\co{z} \mathrel{|} \co{y}$ in process calculus \textbf{CCS}~\cite{Milner1982}.
We can attempt to embed this process in \textsf{BVQ} as $\mathopen{\nabla x.} \left(\mathrm{event}(x) \cseq \mathrm{event}(y)\right) \cpar \nabla z.\co{\mathrm{event}(z)} \cpar \co{\mathrm{event}(y)}$, where $\mathrm{event}(x)$ is a unary predicate representing an event identified by variable $x$.
This embedding \textbf{violates} our soundness criterion. Under the semantics of \textsf{CCS} the process is immediately deadlocked; hence none of the four actions are consumed.
However, the embedding is a provable formula, by the following derivation.
% using the rules of Figure~\ref{figure:BVQ}.
\[
\infer[\mbox{by \textit{unify} and $\alpha$-conversion}]{
\mathopen{\nabla x.} \left(\mathrm{event}(x) \cseq \mathrm{event}(y)\right) \cpar \nabla z.\co{\mathrm{event}(z)} \cpar \co{\mathrm{event}(y)}
}{
\infer[\mbox{by  \textit{vacuous} and \textit{unify}}]{
\mathopen{\nabla x.} \left(\left(\mathrm{event}(x) \cseq \mathrm{event}(y)\right) \cpar \co{\mathrm{event}(x)}\right) \cpar \co{\mathrm{event}(y)}
}{
\infer[\mbox{by  \textit{sequence}}]{
\mathopen{\nabla x.} \left(\left(\mathrm{event}(x) \cseq \mathrm{event}(y)\right) \cpar \co{\mathrm{event}(x)} \cpar \co{\mathrm{event}(y)}\right)
}{
\infer[\mbox{by  \textit{sequence}}]{
\mathopen{\nabla x.} \left(\left(\mathrm{event}(x) \cseq \mathrm{event}(y)\right) \cpar \left(\co{\mathrm{event}(x)} \cseq \co{\mathrm{event}(y)}\right)\right)
}{
\infer[\mbox{by  \textit{atomic interaction}}]{
\mathopen{\nabla x.} \left(\left(\mathrm{event}(x) \cpar \co{\mathrm{event}(x)}\right) \cseq \left(\mathrm{event}(y) \cpar \co{\mathrm{event}(y)}\right)\right)
}{
\infer[\mbox{by  \textit{atomic interaction} and \textit{vacuous}}]{
\mathopen{\nabla x.} \left(\mathrm{event}(y) \cpar \co{\mathrm{event}(y)}\right)
}{
\cunit
}}}}}}
\]
%For a sound process embedding, we desire the property: 
% This is because the bound names $\nabla a$ and $\nabla c$ should always be kept distinct and hence $a$ and $\co{c}$ cannot interact; in contrast the proof above which should indicate that the process may reach a successfully terminated state.

The above observations lead to a specification of the properties desired for a nominal quantifier suitable for direct embeddings of processes as formulae.
We desire a nominal quantifier, say $\new$, such that properties such as \textit{no diagonalisation}, \textit{equivariance} and \textit{extrusion} hold except that
also 
\textbf{neither} 
$
\mathopen{\new x.} \left(P \cpar Q \right)
\multimap
\new x.P \cpar \new x.Q
$
\textbf{nor}
$
\new x.P \cpar \new x.Q
\multimap
\mathopen{\new x.} \left(P \cpar Q \right)
$
hold in general.
Also, by the arguments above the quantifier cannot be self-dual;
and hence, as a side effect, we expose another nominal quantifier, called ``wen'', denoted $\wen$, that is de Morgan dual to $\new$.
The rest of this paper is devoted to establishing that indeed there does exist a logical system with such a pair of nominal quantifiers.




% --- end self-dual.tex ---
% --- begin syntax.tex ---
\section{Introducing a Proof System with a Pair of Nominal Quantifiers}
\label{section:syntax}

Soundness issues associated with a self-dual nominal quantifier in embeddings of processes as formulae, can be resolved by instead using a pair of de Morgan dual nominal quantifiers.
This section introduces a proof system for such a pair of nominal quantifiers, 
building on the core system \textsf{BV}, further extended with: additives useful for expressing non-deterministic choice; and first-order quantifiers which range over terms not only fresh names.
%Further to this pair of nominal quantifiers, the system extends the core system \textsf{BV}, with: additives useful for expressing non-deterministic choice; and first-order quantifiers which range over terms not only fresh names.
Investigating the pair of nominal quantifiers in the presence of these operators is essential for understanding  the interplay between nominal quantifiers and other operators, showing that this pair of nominal quantifiers can exist in a system sufficiently expressive to embed rich process models.
This section also summarises the main proof theoretic result, although lemmas are postponed until later sections.

%The first half of the paper introduces sufficient terminology to use the calculus, without knowledge of results from the literature on the calculus of structures. The intention is to engage with scientists enjoying process calculi.

\subsection{The inference rules and structural rules}

\begin{figure}[t]
\begin{tabular}{c}
\begin{tabular}{c}
$x$ a variable
\\[9pt]
$c$ a constant
\\[9pt]
$f$ a function symbol
\\[9pt]
$\mathrm{p}$ a predicate symbol
\\[9pt]
\end{tabular}
\\
\begin{tabular}{l@{\hspace{1ex}}l@{\hspace{1ex}}l@{\hspace{1ex}}l}
$t$&$::=$ & $x$ & (variable)\\
&& $c$ & (constant)\\
&& $f(t, \hdots t)$ & ($n$-ary function)
%&& $\top$ & (top)\\
%&& $0$ & (bottom)\\
%&& $\rename{x}. P$ & (rename)\\
\\[9pt]
$\alpha$&$::=$ & $\mathrm{p}(t, \hdots t)$ & ($n$-ary predicate)\\
\end{tabular}
\qquad\qquad\qquad
\end{tabular}
\begin{tabular}{l@{\hspace{1ex}}l@{\hspace{1ex}}l@{\hspace{1ex}}l}
$P$&$::=$ & $\cunit$ & (unit)\\
&& $\alpha$ & (atom)\\
&& $\overline{\alpha}$ & (co-atom)\\
%&& $\top$ & (top)\\
%&& $0$ & (bottom)\\
%&& $\rename{x}. P$ & (rename)\\
&& $\forall x. P$ & (all)\\
&& $\exists x. P$ & (some)\\
&& $\new x. P$ & (new)\\
&& $\wen x. P$ & (wen)\\
&& $P \wwith P$ & (with)\\
&& $P \ooplus P$ & (plus) \\
&& $P\cpar P$ & (par)\\
&& $P \tensor P$ & (times)\\
&& $P \cseq P$ & (seq)
\end{tabular}
  %\vspace{-20pt}
\caption{Syntax for \textsf{MAV1} formulae.}\label{figure:syntax}
\end{figure}
%\end{wrapfigure}


\begin{figure}[t]
\begin{tabular}{c}
$\left(P, \cpar, \cunit\right)\mbox{ and }\left(P, \tensor, \cunit\right)\mbox{ are commutative monoids}
\mbox{ and }
\left(P, \cseq, \cunit\right)\mbox{ is a monoid.}$
%\quad \quantifier x. P \equiv \mathopen{\quantifier y .}\left(P\sub{x}{y}\right) \mbox{ if $\nfv{y}{\quantifier x. P}$} ~~\mbox{ ($\alpha$-conversion)}
\\[9pt]
$\new x .\new y. P \equiv \new y .\new x. P
%~~\mbox{(equivariance)}$
\qquad
\wen x .\wen y .P \equiv \wen y .\wen x .P
\qquad\mbox{(equivariance)}$
\end{tabular}
\caption{Structural congruence ($\equiv$)
 for \textsf{MAV1} formulae,
plus $\alpha$-conversion for all quantifiers.
}\label{figure:congruence}
\end{figure}

\begin{figure}[t]
\[
\begin{array}{c}
\infer[\mbox{(atomic interaction)}]{
\context{ \co{\alpha} \cpar \alpha }
}{
\context{ \cunit }
}
\qquad\quad\quad
\infer[\mbox{(switch)}]{
\context{P \cpar \left(Q \tensor S\right)}
}{
\context{(P \cpar Q) \tensor S}
}
\\ [10pt]
\infer[\mbox{(sequence)}]{
\context{(P \cseq Q) \cpar (U \cseq V)}
}{
\context{(P \cpar U) \cseq (Q \cpar V)}
}
\\[10pt]
\hline
\\
%\context{ \rename{x} P \cpar \rename{x} Q }
%\longrightarrow
%\context{ \rename{x} \left( P \cpar Q \right) }
%\quad
%\mbox{only if }
%P \not \equiv \cunit \mbox{ and } Q \not \equiv \cunit
%~~\mbox{(rename)}
%\\[17pt]
\infer[\mbox{(external)}]{
\context{ (P \wwith Q) \cpar S }
}{
\context{ (P \cpar S) \wwith (Q \cpar S) }
}
\qquad\qquad
\infer[\mbox{(medial)}]{
\context{ \left(P \cseq Q\right) \wwith \left(U \cseq V\right) }
}{
\context{ \left(P \wwith U\right) \cseq \left(Q \wwith V\right) }
}
\\[10pt]
\infer[\mbox{(tidy)}]{
\context{ \cunit \wwith \cunit }
}{
\context{ \cunit }
}
\qquad\qquad
\infer[\mbox{(left)}]{
\context{ P \ooplus Q }
}{
\context{ P}
}
\qquad\qquad
\infer[\mbox{(right)}]{
\context{ P \ooplus Q }
}{
\context{ Q}
}
\\[10pt]
\hline
\\
\infer[\mbox{(extrude1)}]{
\context{ \forall x. P \cpar R }
}{
\context{ \mathopen{\forall x.}\left( P \cpar R \right) }
}
%\mbox{ only if $\nfv{x}{R}$}
\qquad\qquad\quad
\infer[\mbox{(medial1)}]{
\context{ \mathopen{\forall x .}\left(P \cseq S \right) }
}{
\context{ \forall x .P \cseq \forall x .S }
}
\\[10pt]
\infer[\mbox{(tidy1)}]{
\context{ \forall x. \cunit }
}{
\context{ \cunit }
}
\qquad\qquad\quad
%\context{ \top \cpar P } \longrightarrow \context{ \top } \quad \mbox{only if } P \not \equiv \cunit ~~\mbox{(top)}
%\\[19pt]
\infer[\mbox{(select1)}]{
\context{ \exists x. P }
}{
\context{ P \sub{x}{t} }
}
\\[10pt]
\hline
\\
\infer[\mbox{(extrude new)}]{
\context{ \new x .P \cpar R }
}{
\context{ \mathopen{\new x .}\left( P \cpar R \right) }
}
\qquad
\infer[\mbox{(medial new)}]{
\context{ \mathopen{\new x .}\left( P \cseq S \right) }
}{
\context{ \new x. P \cseq \new x. S }
}
\\[10pt]
\infer[\mbox{(tidy name)}]{
\context{ \new x. \cunit }
}{
\context{ \cunit }
}
\qquad
\infer[\mbox{(close)}]{
\context{ \new x. P \cpar \wen x. Q }
}{
\context{ \mathopen{\new x .}\left( P \cpar Q \right) }
}
\\[10pt]
\infer[\mbox{(fresh)}]{
\context{ \wen x. P }
}{
\context{ \new x. P }
}
\qquad
\infer[\mbox{(new wen)}]{
\context{ \new x .\wen y .P }
}{
\context{ \wen y. \new x .P }
}
\qquad
\infer[\mbox{(all name)}]{
\context{ \forall x .\quantifier y. P }
}{
\context{ \quantifier y. \forall x .P }
}
\\[10pt]
\infer[\mbox{(suspend)}%\mbox{(suspend)}
]{
\context{ \wen x. P \odot \wen x. S }
}{
\context{ \mathopen{\wen x .}\left( P \odot S \right) }
}
\qquad
\infer[\mbox{(left wen)}%\mbox{(left wen)}
]{
\context{ \wen x .P \odot R }
}{
\context{ \mathopen{\wen x .}\left( P \odot R \right) }
}
%\mbox{ where $\odot \in \left\{ \cpar, \cseq \right\}$}
%\mbox{ only if $\nfv{x}{R}$}
\qquad
\infer[\mbox{(right wen)}%\mbox{(right wen)}
]{
\context{ R \odot \wen x. Q }
}{
\context{ \mathopen{\wen x .}\left( R \odot Q \right) }
}
%\mbox{ where $\odot \in \left\{ \cpar, \cseq \right\}$}
%\mbox{ only if $\nfv{x}{R}$}
\\[10pt]
\infer[\mbox{(with name)}
]{
\context{ \quantifier x .P \wwith \quantifier x. S }
}{
\context{ \mathopen{\quantifier x .}\left( P \wwith S \right) }
}
%\mbox{ for $\quantifier \in \left\{ \new, \wen \right\}$}
\quad
\infer[\mbox{(left name)}
]{
\context{ \quantifier x .P \wwith R }
}{
\context{ \mathopen{\quantifier x .}\left( P \wwith R \right) }
}
%\mbox{ only if $\nfv{x}{R}$}
%\mbox{ for $\quantifier \in \left\{ \new, \wen \right\}$}
%~~\mbox{(left name)}
\quad
\infer[\mbox{(right name)}]{
\context{ R \wwith \quantifier x. Q }
}{
\context{ \mathopen{\quantifier x .}\left( R \wwith Q \right) }
}
%\mbox{ only if $\nfv{x}{R}$}
%\mbox{ for $\quantifier \in \left\{ \new, \wen \right\}$}
\\\\
\end{array}
\]
\begin{flushright}
where $\protect{\quantifier \in \left\{ \new, \wen \right\}}$, $\odot \in \left\{ \cpar, \cseq \right\}$ and $\nfv{x}{R}$, in all rules containing $R$
\end{flushright}
\caption[Rules for formulae in system \textsf{MAV1}.]{Rules for formulae in system \textsf{MAV1}. %In the \textit{atomic interact} rule, $\vdash \alpha \leq \beta$ defines any preorder closed under substitution.
%The first, second and fourth parts together are called subsystem \textsf{nominal-MAV}.
%The first, third and fourth parts together are called sub-system \textsf{BV1}.
Notice the figure is divided into four parts. The first part defines sub-system \textsf{BV}~\cite{Guglielmi2007}. The first and second parts define sub-system \textsf{MAV}~\cite{Horne2015}.
}
\label{figure:rewrite}
\end{figure}




We present the syntax and rules of a first-order system expressed in the calculus of structures, with the technical name \textsf{MAV1}.
%expressed as a term rewriting system modulo a structural congruence. 
The derivations of the system are defined by the \textit{abstract syntax} in Fig.~\ref{figure:syntax}, \textit{structural congruence} in Fig.~\ref{figure:congruence},
and the \textit{inference rules}, in Fig~\ref{figure:rewrite}.
We emphasise that, in contrast to the sequent calculus, 
rules can be applied in any context, i.e.\ \textsf{MAV1} formulae from Fig.~\ref{figure:syntax} with a hole of the form 
\[
\context{} \Coloneqq \ehole{\,\cdot\,} \mid \context{} \odot P \mid P \odot \context{} \mid \quantifier x . \context{}, 
\mbox{ where $\odot \in \left\{ \cseq, \cpar, \tensor, \wwith, \ooplus \right\}$ and $\quantifier \in \left\{ \exists, \forall, \new, \wen \right\}$. }
\] 
%In addition to the structural rules of \textsf{BV}, the congruence in Fig.~\ref{figure:syntax} enables $\alpha$-conversion for all quantifiers and \textit{equivariance} which allows successive nominal quantifiers to be swapped.
% where $\alpha$-conversion is such that, if $\nfv{y}{P}$, then $\quantifier{x}P \equiv \quantifier{y}(P\sub{x}{y})$.
We also assume the standard notion of capture avoiding substitution of a variable for a term.
Terms may be constructed from variables, constants and function symbols.

To explore the theory of proofs, two auxiliary definitions are introduced: linear negation and linear implication. 
Notice in the syntax in Fig.~\ref{figure:syntax} linear negation applies only to atoms.
% thus all rules are defined  always consider formulae in \textit{negation-normal-form}.
%Instead complementation is a function from predicates to predicates, that maps a predicate to its dual predicate in , where complementation is applied only to atoms, as follows.
\begin{definition}
\textit{Linear negation} is defined by the following function from formulae to formulae.
\begin{gather*}
\co{\co{\alpha}} = \alpha
\qquad
\co{ P \tensor Q } = \co{P} \cpar \co{Q}
\qquad
\co{ P \cpar Q } = \co{P} \tensor \co{Q}
\qquad
%\co{\top} = 0
%\qquad
%\co{0} = \top
%\qquad
\co{ P \ooplus Q } = \co{P} \wwith \co{Q}
\qquad
\co{ P \wwith Q } = \co{P} \ooplus \co{Q}
\\%[10pt]
\co{\cunit} = \cunit
\quad
\co{P \cseq Q} = \co{P} \cseq \co{Q}
\qquad
\co{\forall x. P} = \exists x. \co{P}
\qquad
\co{\exists x. P} = \forall x. \co{P}
\qquad
\co{\new x. P} = \wen x. \co{P}
\qquad
\co{\wen x. P} = \new x. \co{P}
%\co{\rename x P} = \rename x \co{P}
\end{gather*}
\textit{Linear implication}, written $P \multimap Q$, is defined as $\co{P} \cpar Q$.
\end{definition}
Linear negation defines de Morgan dualities.
As in linear logic, the multiplicatives $\tensor$ and $\cpar$ are de Morgan dual; as are the additives $\wwith$ and $\ooplus$, the first-order quantifiers $\exists$ and $\forall$, and the nominal quantifiers $\new$ and $\wen$. As in \textsf{BV}, \textit{seq} and the unit are self-dual.

A basic, but essential, property of implication can be established immediately.
The following proposition is simply a reflexivity property of linear implication in \textsf{MAV1}.
\begin{proposition}[Reflexivity]\label{proposition:reflexivity}
For any formula $P$, $\vdash \co{P} \cpar P$ holds, i.e., $\vdash P \multimap P$.
\end{proposition}
The proof of the above follows by a straightforward induction over the structure of $P$.
%\tohide{
%%\ross{This can be cut down to the case for nominal only.}
%\begin{proof}
%\input{proof_reflexivity}
%\end{proof}}{}



\subsection{Intuitive explanations for the rules of \textsf{MAV1}.}

We briefly recall the established system \textsf{MAV},
before explaining the rules for quantifiers.
This paper focuses on necessary proof theoretical prerequisites, and hints at result for process embeddings in \textsf{MAV1}.
Details on the soundness of process embeddings appear in a companion paper~\cite{MSCS}.

\textbf{The additives.}
The rules of the basic system \textsf{BV} in the top part of Fig.~\ref{figure:rewrite} are as described previously in Section~\ref{section:self-dual}.
The first and second parts of Fig.~\ref{figure:rewrite} define multiplicative-additive system \textsf{MAV}~\cite{Horne2015}.
The additives are useful for modelling non-deterministic choice in processes~\cite{Abramsky1993}: the \textit{left} rule $\vcenter{\infer{P \ooplus Q}{P}}$ suggests we chose the left branch $P$ \textbf{or} alternatively the right branch $Q$ by using the \textit{right} rule; the \textit{external} rule 
$\vcenter{\infer{\left(P \wwith Q\right) \cpar R}{\left(P  \cpar R\right) \wwith \left(Q \cpar R\right)}}$
suggests that we try both branches $P \cpar R$ \textbf{and} $Q \cpar R$ separately; and the \textit{tidy} rule indicates a derivation is successfully only if both branches explored are successful.
The \textit{medial} rule is a partial distributivity property between the additives and \textit{seq} (in concurrency theory, this is a property expected of most preorders over processes).
The role of the additives as a form of \textit{internal} and \textit{external} choice has been investigated in related work~\cite{Ciobanu2015}.

\textbf{The first-order quantifiers.}
The rules for the first-order quantifiers in the third part of Fig.~\ref{figure:rewrite} follow a similar pattern to the additives.
The \textit{select1} rule allows a variable to be replaced by any term. Notice we stick to the first-order case, since variables only appear in atomic formulae and may only be replaced by terms.
% and not an arbitrary predicate.
The \textit{extrude1}, \textit{tidy1} and \textit{medial1} rules follow a similar pattern to the rules for the additives \textit{external}, \textit{tidy} and \textit{medial} respectively.
In process embeddings, first-order quantifiers are useful as input binders.
For example we can soundly embed the $\pi$-calculus process $\co{y}z \pipar y(x).\co{x}w \pipar z(x)$ as the following provable formula:
% $\co{ab} \cpar \exists x.\left(ax\cseq \co{xc}\right) \cpar \exists y.b(y)$.
\[
\infer[\mbox{by \textit{select1}}]{
\co{\mathrm{act}(y, z)} \cpar \mathopen{\exists x.}\left(\mathrm{act}(y, x)\cseq \co{\mathrm{act}(x, w)}\right) \cpar \exists v.\mathrm{act}(z, v)
}{
\infer[\mbox{by \textit{sequence}}]{
\co{\mathrm{act}(y, z)} \cpar \left(\mathrm{act}(y, z)\cseq \co{\mathrm{act}(z, w)}\right) \cpar \exists v.\mathrm{act}(z, v)
}{
\infer[\mbox{by \textit{atomic interaction}} ]{
\left(\left(\co{\mathrm{act}(y, z)} \cpar \mathrm{act}(y, z)\right)\cseq \co{\mathrm{act}(z, w)}\right) \cpar \exists v.\mathrm{act}(z, v) 
}{
\infer[\mbox{by \textit{select1}}]{
\co{\mathrm{act}(z, w)} \cpar \exists v.\mathrm{act}(z, v)
}{
\infer[\mbox{by \textit{atomic interaction}}]{
\co{\mathrm{act}(z, w)} \cpar \mathrm{act}(z, w)
}{
\cunit
}}}}}
\]
Notice, that the above process can also reach a successfully terminated state using $\tau$ transitions in the $\pi$-calculus semantics.
Indeed the cut elimination result established in this paper is a prerequisite in order to prove this soundness criterion holds for finite $\pi$-calculus processes. %(paper submitted).
% (to appear in a companion paper).
%This paper focuses on the proof theory of \textsf{MAV1} rather than the soundness of process embeddings.

\textbf{The polarised nominal quantifiers.}
The rules for the de Morgan dual pair of nominal quantifiers are more intricate.
For first-order quantifiers many properties are derivable, e.g., the following implications hold (appealing to Prop.~\ref{proposition:reflexivity}):
$\vdash \forall x. \forall y. P \multimap \forall y. \forall x. P$,
$\vdash \exists x. \forall y. P \multimap \forall y. \exists x. P$
and
$\vdash \mathopen{\forall x.}\left( P \cpar Q \right) \multimap \forall x. P \cpar \exists x. Q$.
The three proofs proceed as follows.
%\quad
%\vdash \mathopen{\forall x.} \left( P \wwith Q \right) \multimap \forall x. P \wwith \forall x. Q
%\begin{comment}
\[
\infer[]{
\exists x. \exists y. \co{P} \cpar \forall y. \forall x. P
}{
\infer[]{
\mathopen{\forall y. \forall x.}\left(\exists x. \exists y. \co{P} \cpar  P\right)
}{
\infer[]{
\mathopen{\forall y. \forall x.}\left(\co{P} \cpar  P\right)
}{
\infer[]{
\mathopen{\forall y. \forall x.}\cunit
}{
\cunit
}}}}
\qquad\qquad%%%%%%%%%%%%%%%%%%
\infer[]{
\forall x. \exists y. \co{P} \cpar \forall y. \exists x. P
}{
\infer[]{
\mathopen{\forall x. \forall y.} \left(\exists y. \co{P} \cpar \exists x. P\right)
}{
\infer[]{
\mathopen{\forall x. \forall y.} \left(\co{P} \cpar P\right)
}{
\infer[]{
\mathopen{\forall x. \forall y.} \cunit
}{
\cunit
}}}}
\qquad\qquad%%%%%%%%%%%%%%%%%
\infer[]{
\mathopen{\exists x.}\left( \co{P} \tensor \co{Q} \right) \cpar \forall x. P \cpar \exists x. Q
}{
\infer[]{
\mathopen{\forall x.}\left(\mathopen{\exists x.}\left( \co{P} \tensor \co{Q} \right) \cpar P \cpar \exists x. Q\right)
}{
\infer[]{
\mathopen{\forall x.}\left( \co{P \cpar Q} \cpar P \cpar Q\right)
}{
\infer[]{
\mathopen{\forall x.}\cunit
}{
\cunit
}}}}
\]
%\end{comment}
We desire analogous properties for the nominals $\new$ and $\wen$.
However, in contrast to first-order quantifiers, these properties must be induced for our pair of nominals. 
The first property is induced for $\new$ and $\wen$ by \textit{equivariance} in the structural congruence. 
The other rules analogous to the above derived implications are induced by the rules: \textit{new wen}, which allow a weaker quantifier $\wen$ to commute over a stronger quantifier $\new$;
and \textit{close} which models that $\wen$ can select a name as long as it is fresh as indicated by $\new$.
%and \textit{with name} which combines the bookkeeping of quantifiers in separate branches.

We avoid \textit{new} distributing over $\cpar$, i.e., in general \textbf{neither} 
$
\mathopen{\new x.} \left(P \cpar Q \right)
\multimap
\new x.P \cpar \new x.Q
$
\textbf{nor}
$
\new x.P \cpar \new x.Q
\multimap
\mathopen{\new x.} \left(P \cpar Q \right)
$
hold.
Hence $\new$ is suitable for embedding the name binder $\nu$ of the $\pi$-calculus.
Interestingly, the dual quantifier $\wen$ is also useful for embedding a variant of the $\pi$-calculus called the $\pi I$-calculus, where every communication creates a new fresh name.
For example, $\pi I$-calculus process $\mathopen{\oprivate{v}{x}}.\iprivate{x}{y} \pipar \mathopen{\iprivate{v}{z}}.\oprivate{z}{w}$ can be embedded as the following provable formula.\footnote{
To disambiguate from the $\pi$-calculus we use square brackets as binders for the $\pi I$-calculus.
So $\mathopen{\oprivate{v}{x}}.P$ denotes a process that outputs a fresh name $x$
and $\mathopen{\iprivate{v}{x}}.P$ denotes a process that receives a name $x$ only if it is fresh.
}
\[
\infer[\mbox{by \textit{close} and $\alpha$-conversion}]{
\mathopen{\new x.}\left( \co{\mathrm{act}(v, x)} \cseq \wen y.\mathrm{act}(x, y)\right) \cpar \mathopen{\wen z.}\left(\mathrm{act}(v, z) \cseq \new w.\co{\mathrm{act}(z, w)}\right)
}{
\infer[\mbox{by \textit{sequence}}]{
\mathopen{\new x.}\left(\left( \co{\mathrm{act}(v, x)} \cseq \wen y.\mathrm{act}(x, y)\right) \cpar \left(\mathrm{act}(v, x) \cseq \new w.\co{\mathrm{act}(x, w)}\right)\right)
}{
\infer[\mbox{by \textit{atomic interaction}}]{
\mathopen{\new x.}\left(\left( \co{\mathrm{act}(v, x)} \cpar \mathrm{act}(v, x)\right) \cseq \left(\wen y.\mathrm{act}(x, y) \cpar \new w.\co{\mathrm{act}(x, w)}\right)\right)
}{
\infer[\mbox{by \textit{close}}]{
\mathopen{\new x.}\left(\wen y.\mathrm{act}(x, y) \cpar \new w.\co{\mathrm{act}(x, w)}\right)
}{
\infer[\mbox{by \textit{atomic interaction}}]{
\mathopen{\new x. \new w.}\left(\mathrm{act}(x, w) \cpar \co{\mathrm{act}(x, w)}\right)
}{
\infer[\mbox{by \textit{tidy name}}]{
\mathopen{\new x. \new w.}\cunit
}{
\cunit
}}}}}}
\]
Note that $\alpha$-renaming is implicitly applied in the derivation above. 

%\ross{This can go where explaining tidy new, extrude name, etc.,.}
There is no \textit{vacuous} rule in Fig.~\ref{figure:syntax}, in contrast to the presentation of \textsf{BVQ} in Fig.~\ref{figure:BVQ}.
This is because the \textit{vacuous} rule creates problems for proof search, since arbitrarily many nominal quantifiers can be introduced at any point in the proof leading to unnecessary infinite search paths. Instead we build the introduction and elimination of fresh names into rules only where required.
For example, \textit{extrude new} is like \textit{close} with a vacuous $\wen$ implicitly introduced;
similarly, for \textit{left wen}, \textit{right wen}, \textit{left name} and \textit{right name} a vacuous $\wen$ is implicitly introduced.
Also the \textit{tidy name} allows vacuous $\new$ operators to be removed from a successful proof in order to terminate with $\cunit$ only.
The reason why the rules \textit{medial new}, \textit{suspend}, \textit{all name} and \textit{with name} are required are in order to make cut elimination work; hence we postpone their explanation until after the statement of the cut elimination result.
%Notice that another problem with \textsf{BVQ} is that in a sense the proof system  is not analytic. 



\begin{comment}
\begin{figure}
\[
\begin{array}{c}
\quantifier x P \equiv \quantifier y \left(P\sub{x}{y}\right)
~~\mbox{only if $\nfv{y}{\quantifier x P}$}
\quad
\mbox{($\alpha$-conversion)}
\\ \qquad\qquad\qquad\qquad\qquad\mbox{for $\quantifier \in \left\{ \exists, \forall, \new, \wen \right\}$}
\\
\wen x \wen y P \equiv \wen y \wen x P
~~\mbox{(swap)}
\\[9pt]
\left(P, \cpar, \cunit\right) %{\rm ~is~a~commutative~monoid}
\mbox{ and }
\left(P, \tensor, \cunit\right) {\rm ~are~commutative~monoids}
~
\left(P, \cseq, \cunit\right) {\rm ~is~a~monoid}
%\quad
%(P,=) {\rm ~is~an~equivalence~relation}
\end{array}
\]
\caption{The structural congruence ($\equiv$) for \textsf{MAV1}.}
\label{figure:congruence}
\end{figure}
\end{comment}

%\paragraph{Abstract syntax technicalities.}
%We introduce some standard technical devices: a fresh name predicate and capture avoiding substitution.
%The freshness predicate is such that a variable $x$ is fresh for a predicate $P$, written $\nfv{x}{P}$, if and only if $x$ is not a member of the set of free variables of $P$.
%Definitions are standard and clarify that quantifiers are variable binders.






In addition to forbidding the \textit{vacuous} rule, 
the following restrictions are placed on the rules to avoid meaningless infinite paths in proof search.
\begin{itemize}
\item For the \textit{switch}, \textit{sequence}, \textit{medial1}, \textit{medial new} and \textit{extrude new} rules, $P \not \equiv \cunit$ and $S \not \equiv \cunit$.
\item The \textit{medial} rule is such that either $P \not\equiv \cunit$ or $R \not\equiv \cunit$
and also either $Q \not\equiv \cunit$ or $S \not\equiv \cunit$.
\item The rules \textit{external}, \textit{extrude1}, \textit{extrude new}, \textit{left wen} and \textit{right wen} are such that $R \not \equiv \cunit$.
\end{itemize}
Avoiding infinite search paths is important for the termination of our cut elimination procedure.
Essentially, we desire that our system for \textsf{MAV1} is in a sense \textit{analytic}~\cite{Bruscoli2009}.


\begin{comment}
%%%%%%%%%%%%% V2 %%%%%%%%%%%%%%%%%%%%%%%
\hline
\\
\context{ \new x P \cpar \wen x Q }
\longrightarrow
\context{ \new x \left( P \cpar Q \right) }
~~\mbox{(close)}
\qquad
\context{ \new x \cunit } \longrightarrow \context{ \cunit }
~~\mbox{(tidy new)}
\\[9pt]
\context{ \new x P \cpar R }
\longrightarrow
\context{ \new x \left( P \cpar R \right) }
\mbox{ only if $\nfv{x}{R}$}
~~\mbox{(extrude new)}
\\[9pt]
\context{ \wen x P }
\longrightarrow
\context{ \new x P }
~~\mbox{(fresh)}
\qquad
\context{ \new x \left( P \cseq S \right) }
\longrightarrow
\context{ \new x P \cseq \new x S }
~~\mbox{(medial new)}
\\[9pt]
\context{ \forall x \wen y P }
\longrightarrow
\context{ \wen y \forall x P }
~~\mbox{(all wen)}
\qquad
\context{ \new x \wen y P }
\longrightarrow
\context{ \wen y \new x P }
~~\mbox{(new wen)}
\\[9pt]
\context{ \wen x P \odot \wen x S }
\longrightarrow
\context{ \wen x \left( P \odot S \right) }
\mbox{ where $\odot \in \left\{ \cpar, \cseq, \wwith \right\}$}
~~\mbox{(suspend)}
\\[9pt]
\context{ \wen x P \odot R }
\longrightarrow
\context{ \wen x \left( P \odot R \right) }
\mbox{ only if $\nfv{x}{R}$}
\mbox{ where $\odot \in \left\{ \cpar, \cseq, \wwith \right\}$}
~~\mbox{(left wen)}
\\[9pt]
\context{ R \odot \wen x Q }
\longrightarrow
\context{ \wen x \left( R \odot Q \right) }
\mbox{ only if $\nfv{x}{R}$}
\mbox{ where $\odot \in \left\{ \cpar, \cseq, \wwith \right\}$}
~~\mbox{(right wen)}
\\[9pt]
\new x P \wwith \new x Q \longrightarrow \new x \left( P \wwith Q \right)
~~\mbox{(propogate new)}
\\[9pt]
\new x P \wwith R \longrightarrow \new x \left( P \wwith R \right)
\mbox{ only if $\nfv{x}{R}$}
~~\mbox{(propogate new)}
\\[9pt]
R \wwith \new x Q \longrightarrow \new x \left( R \wwith Q \right)
\mbox{ only if $\nfv{x}{R}$}
~~\mbox{(propogate new)}
\\[9pt]
\forall x \new y P \longrightarrow \forall x \new y P
~~\mbox{(propogate new)}
\end{comment}


\begin{comment}
\context{ \wen x P \cseq \wen x Q }
\longrightarrow
\context{ \wen x \left( P \cseq Q \right) }
~~\mbox{(seq/wen)}
\\[9pt]
\context{ \wen x P \cpar \wen x Q }
\longrightarrow
\context{ \wen x \left( P \cpar Q \right) }
~~\mbox{(par/wen)}
\\[9pt]
\context{ \wen x P \wwith \wen x Q }
\longrightarrow
\context{ \wen x \left( P \wwith Q \right) }
~~\mbox{(with/wen)}
\\[9pt]
\context{ \wen x P \cpar R }
\longrightarrow
\context{ \wen x \left( P \cpar R \right) }
\mbox{ only if $\nfv{x}{R}$}
~~\mbox{(extrude wen)}
\\[9pt]
\context{ \wen x P \cseq R }
\longrightarrow
\context{ \wen x \left( P \cseq R \right) }
\mbox{ only if $\nfv{x}{R}$}
~~\mbox{(seq/wen left)}
\\[9pt]
\context{ R \cseq \wen x Q }
\longrightarrow
\context{ \wen x \left( R \cseq Q \right) }
\mbox{ only if $\nfv{x}{R}$}
~~\mbox{(seq/wen right)}
\\[9pt]\\[9pt]
\context{ \wen x P \wwith R }
\longrightarrow
\context{ \wen x \left( P \wwith R \right) }
\mbox{ only if $\nfv{x}{R}$}
~~\mbox{(with/wen left)}
\\[9pt]
\context{ R \wwith \wen x Q }
\longrightarrow
\context{ \wen x \left( R \wwith Q \right) }
\mbox{ only if $\nfv{x}{R}$}
~~\mbox{(with/wen right)}
\end{comment}





%Note that, by adopting HOAS~\cite{Pfenning1988}, these technicalities would be handled automatically.




%We define a freshness predicate, where a variable $x$ is fresh for a predicate $P$, written $\nfv{x}{P}$, if and only if $x$ is not a member of the set of free variables of $P$. Note that this notation for the freshness predicate is borrowed from nominal logic~\cite{Pitts03}, and is just and abbreviation for $x \not\in \fv{P}$. %Future work includes building freshness and substitution into the formal proof system by developing ideas from nominal logic and explicit substitutions~\cite{Abadi1989}. Freshness is defined as follows for all binary connectives over predicates, terms and tests $\odot \in \left\{ \cseq, \tensor, \cpar, \ooplus, \wwith, =, 
%\wedge, 
%\Rightarrow \right\}$ and variable binding quantifiers $\quantifier\ \in \left\{ \forall, \exists, \rename{} \right\}$. For quantifiers, if $x = y$ or $\nfv{x}{P}$ then $\nfv{x}{\left(\quantifier{y}P\right)}$. For binary connectives, if $\nfv{x}{P}$ and $\nfv{x}{Q}$ then $\nfv{x}{\left(P \odot Q\right)}$.
%For channels, if $\nfv{x}{t}$ then $\nfv{x}{(\mathopen{a}t)}$.
%For unary operators $\Box \in \left\{ \co{(\cdot)}, \neg \right\}$, if $\nfv{x}{\alpha}$, then $\nfv{x}{(\Box \alpha)}$.
%For constant terms and predicates $c \in \left\{ a, \top, 0, \cunit \right\}$, $\nfv{x}{c}$ always hold. For variables, if $x \not= y$, then $\nfv{x}{y}$.

\begin{comment}
\begin{gather*}
\fv{x} = \left\{ x \right\}
\qquad
\fv{f(v_1,\dots,v_n)} = \bigcup_{1 \leq i \leq n} \fv{v_i}
\qquad
\fv{\cunit} = \fv{\top} = \fv{0} = \emptyset
\qquad
\fv{\co{\alpha}} = \fv{\alpha}
\\[9pt]
\fv{av} = \fv{v}
\qquad
\fv{u=v} = \fv{u} \cup \fv{v}
\qquad
\fv{\neg \phi} = \fv{\phi}
\qquad
\fv{\phi \Rightarrow \psi} = \fv{\phi} \cup \fv{\psi}
\\[9pt]
\begin{array}{rl}
\fv{\quantifier{x} B} &= \fv{B} \setminus \left\{ x \right\}
\\
&\mbox{where $\quantifier \in \left\{ \forall, \exists, \rename{} \right\}$}
\end{array}
\qquad
\fv{A \odot B} = \fv{A} \cup \fv{B}
\ 
\mbox{where $\odot \in \{;,\cpar,\tensor,\wwith,\ooplus\}$}
\end{gather*}
\end{comment}

%Capture avoiding substitution is defined for all binary connectives over predicates, terms and tests $\odot \in \left\{ \cseq, \tensor, \cpar, \ooplus, \wwith, =, %\wedge, 
%\Rightarrow \right\}$ and variable binding quantifiers $\quantifier\ \in \left\{ \forall, \exists, \rename{} \right\}$. For binary connectives $(P \odot Q)\sub{x}{t} = P\sub{x}{t} \odot Q\sub{x}{t}$. 
%For quantifiers, $(\quantifier{x}P)\sub{x}{t} = \quantifier{x}P$ and, if $x \not= y$ and $\nfv{y}{t}$, then $(\quantifier{y}P)\sub{x}{t} = \quantifier{y}\left(P\sub{x}{t}\right)$.
%For quantifiers if $\nfv{z}{\quantifier{y}P}$ and $\nfv{z}{t}$ then $(\quantifier{y}P)\sub{x}{t} \equiv \quantifier{z}\left(P\sub{y}{z}\sub{x}{t}\right)$.
%For unary connectives $\Box \in \left\{ \co{(\cdot)}, \neg\right\}$, $\left(\Box{\alpha}\right)\sub{x}{t} = \Box\left({\alpha\sub{x}{t}}\right)$.
%For channels, $\left(\mathopen{a}v\right)\sub{x}{t} = \mathopen{a}\left(v\sub{x}{t}\right)$.
%For equality tests, $(u = v)\sub{x}{t} = u\sub{x}{t} = v\sub{x}{t}$.
%For function terms, $\mathopen{f}(t_0, \hdots, t_n)\sub{x}{t} = \mathopen{f}(t_0\sub{x}{t}, \hdots, t_n\sub{x}{t})$.
%For constant terms and predicates $c \in \left\{ a, \top, 0, \cunit \right\}$, $c\sub{x}{t} = c$.
%For variables, if $x \not= y$, then $y\sub{x}{t} = y$, otherwise $x\sub{x}{t} = t$. Note that we use equality ($=$) for substitutions to emphasise that substitutions are functions over terms and not part of the syntax. We assume that $\alpha$-conversion can be applied to any bound variables to avoid name capture, where $\alpha$-conversion is such that, if $\nfv{y}{P}$, then $\quantifier{x}P \equiv \quantifier{y}(P\sub{x}{y})$.
%Whereas, for $\alpha$-conversion, we use equivalence ($\equiv$). $However, of course, abstract syntax terms that differ only with respect to $\alpha$-conversion are identified by de Brujin indices.
\begin{comment}
\begin{gather*}
y \sub{x}{v} =
\left\{
\begin{array}{lr}
v & \mbox{if $x = y$} \\
y & \mbox{if $x \not = y$}
\end{array}\right.
\qquad
f\left(v_1,\dots,v_{n}\right)\sub{x}{v} = f\left(v_1\sub{x}{v},\dots,v_{n}\sub{x}{v}\right)
\\[9pt]
\cunit \sub{x}{v} = \cunit
\qquad
\top \sub{x}{v} = \top
\qquad
0 \sub{x}{v} = 0
\qquad
\left(au\right) \sub{x}{v} = a(u\sub{x}{v})
\qquad
\co{\alpha}\sub{x}{v} = \co{\alpha\sub{x}{v}}
\\[9pt]
(t=u) \sub{x}{v} = (t\sub{x}{v} = u\sub{x}{v})
\qquad
(\phi \Rightarrow \psi) \sub{x}{v} = \phi\sub{x}{v} \Rightarrow \psi\sub{x}{v}
\qquad
(\neg \phi) \sub{x}{v} = \neg\left(\phi\sub{x}{v}\right)
%\mbox{y \not \in \fv{v}}
\\[9pt]
\begin{array}{rl}
(\quantifier{y} T) \sub{x}{v} &=
\left\{
\begin{array}{lr}
\quantifier{y}T                      & \mbox{if $x = y$} \\
\quantifier{y}(T\sub{x}{v})          & \mbox{if $x \not = y$ and $y \not\in \fv{v}$}
\end{array}\right.
\\
&\mbox{where $\quantifier \in \left\{ \forall, \exists, \rename{} \right\}$}
\end{array}
\qquad
\begin{array}{rl}
(T \odot U)\sub{x}{v} &= (T\sub{x}{v}) \odot (U\sub{x}{v}) \\
                       & \mbox{where $\odot \in \left\{;,\cpar,\tensor, \wwith, \ooplus \right\}$}
\end{array}
\end{gather*}
\end{comment}

\begin{comment}
\begin{array}{c}
v~~\mbox{value}
\qquad
x~~\mbox{variable}
\\[9pt]
\lambda \Coloneqq x \mid v ~~\mbox{variable or value}
\\[9pt]
\begin{array}{rl}
\atom \Coloneqq & p(\lambda) \\
           \mid & p(\lambda,\lambda) \\
           \mid & p(\lambda,\lambda,\lambda) \\
           \mid & \hdots 
\end{array}~~\mbox{atomic predicates}
\end{array}
\end{comment}


%The atoms are dependent on the application. In this case atoms are tuples of values, which are the basic units of a loosely structured data format, such as RDF triples~\cite{Carroll2004} 


\paragraph{Note on term ``medial''} Medials were introduced, historically, to make contraction local (reducing contraction to a rule acting only over atoms)~\cite{Brunnler2001}.
Although the rules in Fig.~\ref{figure:rewrite} do not define such a local system, we discovered these rules by first
defining a local system, and then designing a more controlled system retaining only the medials of the local system that are not admissible.
Related work~\cite{Tubella2018} shows that medials are a ubiquitous recipe underlying the rules of proof systems.





% --- end syntax.tex ---
% --- begin implication.tex ---
\subsection{Cut elimination and its consequences}
\label{section:implication}

This section confirms that the rules of \textsf{MAV1} indeed define a logical system, as established by a cut elimination theorem.
Surprisingly, prior to this work, the only direct proof of cut elimination involving quantifiers in the calculus of structures was for \textsf{BVQ}~\cite{Roversi14}.
Related cut elimination results involving first-order quantifiers in the calculus of structures relied on a correspondence with the sequent calculus~\cite{Brunnler2006,Strassburger09}.
However, due to the presence of the non-commutative operator \textit{seq} there is no sequent calculus presentation~\cite{Tiu2006} for \textsf{MAV1}; hence we pursue here a direct proof.

\begin{comment}
A \textit{derivation} is a sequence of rewrites, where the congruence in Fig.~\ref{figure:syntax} can be applied at any point in a derivation.
The length of a derivation involving only the congruence is zero. 
The length of a derivation involving one rule from Fig.~\ref{figure:rewrite} is one.
Given a derivation $P \longrightarrow Q$ of length $m$ and another $Q \longrightarrow R$ of length $n$, the derivation $P \longrightarrow R$ is of length $m+n$.
Unless we make it clear in the context that we refer to a specific rule, $\longrightarrow$ is generally used to represent derivations of any length.
%A derivation is a sequence of zero or more rewrite rules from Fig.~\ref{figure:rewrite}, where the congruence in Fig.~\ref{figure:congruence} can be applied at any point.
We are particularly interested in special derivations, called proofs.
\begin{definition}
A \textit{proof} in \textsf{MAV1} is a derivation $P \longrightarrow \cunit$ from a predicate $P$ to the unit $\cunit$. When such a derivation exists, we say that $P$ is provable, and write $\vdash P$.
\end{definition}
%% When a predicate models a protocol, there is the following useful intuition behind provability: 
%% %If a predicate models a protocol, then
%% a provable predicate represents a protocol that may be executed such that all parties successfully complete their contract (as opposed to halting due to deadlock for example).
%% %Must properties can also be tested by use of dual operators.
%% %There may exist bad execution paths, but there exists at least one good execution path.
%% %Several examples of successfully terminating, hence provable, protocols were provided in the first half of the paper.
%% In related work~\cite{Carbone2015,Ciobanu2015}, provability is employed to characterise multi-party compatibility in session types. Given a multiset of local session types specifying the local behaviour of a party in a protocol, multi-party compatibility determines whether it is possible for the protocols to execute together to complete a session.

%Proofs have properties that do not hold for all derivations, as explored using proof theoretic techniques.
%In particular, proofs have a proof theory explored in this section.
To explore the theory of proofs, two auxiliary definitions are introduced: linear negation and linear implication. 
Notice in the syntax in Fig.~\ref{figure:syntax} linear negation applies only to atoms.
%Instead complementation is a function from predicates to predicates, that maps a predicate to its dual predicate in \textit{negation-normal-form}, where complementation is applied only to atoms, as follows.
\begin{definition}
\textit{Linear negation} is defined by the following function from predicates to predicates.
\begin{gather*}
\co{\co{\alpha}} = \alpha
\qquad
\co{ P \tensor Q } = \co{P} \cpar \co{Q}
\qquad
\co{ P \cpar Q } = \co{P} \tensor \co{Q}
\qquad
%\co{\top} = 0
%\qquad
%\co{0} = \top
%\qquad
\co{ P \ooplus Q } = \co{P} \wwith \co{Q}
\qquad
\co{ P \wwith Q } = \co{P} \ooplus \co{Q}
\\%[10pt]
\co{\cunit} = \cunit
\quad
\co{P \cseq Q} = \co{P} \cseq \co{Q}
\qquad
\co{\forall x. P} = \exists x. \co{P}
\qquad
\co{\exists x. P} = \forall x. \co{P}
\qquad
\co{\new x. P} = \wen x. \co{P}
\qquad
\co{\wen x. P} = \new x. \co{P}
%\co{\rename x P} = \rename x \co{P}
\end{gather*}
\textit{Linear implication}, written $P \multimap Q$, is defined as $\co{P} \cpar Q$.
\end{definition}
Linear negation defines de Morgan dualities.
As in linear logic, the multiplicatives $\tensor$ and $\cpar$ are de Morgan dual; as are the additives $\wwith$ and $\ooplus$, the first-order quantifiers $\exists$ and $\forall$, and the nominal quantifiers $\new$ and $\wen$. As in \textsf{BV}, sequential composition and the unit are self-dual.
% As with complementation, linear implication is not part of the syntax of predicates, but is a function from pairs of predicates to predicates.
%A corollary of the main result in this paper establishes that the relation defined by all provable linear implications is a precongruence, i.e.\ a reflexively and transitively closed relation.
%Reflexivity of linear implication can be established by a straightforward induction.
%Since $P \multimap P$ is defined as $\co{P} \cpar P$, 
\end{comment}




%\paragraph{Cut elimination.}
The main result of this paper is the following, which is a generalisation of \textit{cut elimination} to the setting of the calculus of structures.
\begin{theorem}[Cut elimination]\label{theorem:cut}
For any formula $P$, if $\vdash \context{ P \tensor \co{P} }$ holds,
then $\vdash \context{ \cunit }$ holds.
\end{theorem}
The above theorem can be stated alternatively by supposing that there is a proof in \textsf{MAV1} extended with the extra inference rule: 
$\vcenter{\infer[\mbox{(cut)}]{\context{\cunit}}{\context{P \tensor \co{P}}}}.$
Given such a proof, a new proof can be constructed that uses only the rules of \textsf{MAV1}.
In this formulation, we say that \textit{cut} is \textit{admissible}.

%The proof of Theorem~\ref{theorem:cut} involves a technique known as \textit{splitting}~\cite{Guglielmi2007,Strassburger2003}.
Cut elimination for the propositional sub-system \textsf{MAV} has been previously established~\cite{Horne2015}.
The current paper advances cut-elimination techniques to tackle first-order system \textsf{MAV1}, as achieved by the lemmas in later sections.
Before proceeding with the necessary lemmas, we provide a corollary that demonstrates that one of many consequences of cut elimination is indeed that linear implication defines a precongruence --- a reflexive transitive relation that holds in any context.
%A stronger statement is proven: since implication is preserved in any context, according to the syntax in Fig.~\ref{figure:syntax}, linear implication is a precongruence.
\begin{corollary}~\label{corollary:precongruence}
Linear implication defines a precongruence.
\end{corollary}
\begin{proof}
For transitivity, if $\vdash P \multimap Q$ and $\vdash Q \multimap R$ hold, we have the following.
%the following proof can be constructed.
\[
\infer[\mbox{by the \textit{switch} rule}]{
 \left(\co{P} \cpar \left( Q \tensor \co{Q} \right) \cpar R \right)
}{
\infer[\mbox{by the assumptions $\vdash \co{P} \cpar Q$ and $\vdash \co{Q} \cpar R$}]{
 \left(\co{P} \cpar Q \right) \tensor \left( \co{Q} \cpar R \right)
}{
\cunit
}}
\]
Hence, by Theorem~\ref{theorem:cut}, $\vdash P \multimap R$ as required.

For contextual closure, if $\vdash P \multimap Q$ holds, we have the following.
\[
\infer[\mbox{by the \textit{switch} rule}]{
\co{\context{P}} \cpar \context{ \left( P \tensor \co{P} \right) \cpar Q}
}{
\infer[\mbox{by the assumption $\vdash P \multimap Q$}]{
\co{\context{P}} \cpar \context{ P \tensor \left( \co{P}  \cpar Q \right) } 
}{
\infer[\mbox{by Proposition~\ref{proposition:reflexivity}}]{
\co{\context{P}} \cpar \context{ P } 
}{
 \cunit
}}}
\]
Hence by Theorem~\ref{theorem:cut}, $\vdash \context{P} \multimap \context{Q}$ as required.
Reflexivity holds by Proposition~\ref{proposition:reflexivity}.
\end{proof}


\begin{comment}
Another important consequence of cut elimination is that \textfs{MAV1} is a consistent logical system.
This can be stated by saying that for any provable predicate containing at least one atom, the linear negation of that predicate is not provable.
Notice that the predicate must contain at least on atom since any predicate formed using $\cunit$, binary operators and quantifiers is equivalent to $\cunit$, which is self-dual.
\begin{corollary}
If $\vdash \context{\alpha}$ then $\not\vdash \context{\alpha}$.
\end{corollary}
\begin{proof}
Assume for contradiction that both $\vdash \context{\alpha}$ and $\vdash \co{\context{\alpha}}$ hold.
By Lemma~\ref{lemma:context},
there exists killing context $\tcontext{}$ and $Q$ such that
$\context{\alpha} \longrightarrow \tcontext{\alpha \cpar Q}$
and $\vdash \alpha \cpar Q$.
Hence, by Lemma~\ref{lemma:split-atom},
there exists killing context $\tcontextp{}$ such that $Q \longrightarrow \tcontextp{\co{\alpha}}$.
Hence the following derivation can be constructed.
\[
\context{\alpha} \longrightarrow \tcontext{\alpha \cpar Q}
\longrightarrow
\tcontext{\alpha \cpar \tcontextp{\co{\alpha}}}
\longrightarrow
\tcontext{\tcontextp{\alpha \cpar \co{\alpha}}}
\]
Consequently, 
$
\vdash
\tcontext{\tcontextp{\alpha \cpar \co{\alpha}}}
\multimap
\context{\alpha}
$
holds
and hence
$
\vdash
\co{\context{\alpha}}
\multimap
\co{\tcontext{\tcontextp{\alpha \cpar \co{\alpha}}}}
$ holds.
So, since we assumed that $\vdash \co{\context{\alpha}}$ holds, by Theorem~\ref{theorem:cut-elimination},
$\vdash \co{\tcontext{\tcontextp{\alpha \cpar \co{\alpha}}}}$ must hold.
But in $\co{\tcontext{\tcontextp{\alpha \cpar \co{\alpha}}}}$ there is no 


\end{proof}
Notice that, in contrast to classical logic, the converse of the above does not hold in general it is possible that neither the predicate nor its complement is provable, e.g. $\alpha$.
\end{comment}
% --- end implication.tex ---
% --- begin nominal.tex ---

\subsection{Discussion on logical properties of the rules for nominal quantifiers}
\label{section:nominal}

The rules for the nominal quantifiers \textit{new} and \textit{wen} require justification.
The \textit{close} and \textit{tidy name} rules ensure the reflexivity of implication for nominal quantifiers.
% and are admissible for all dual pairs of quantifiers.
Using the \textit{extrude new} rule (and Proposition~\ref{proposition:reflexivity}) we can establish the following proof of $\vdash \wen x. P \multimap \exists x. P$.
\[
\infer[\mbox{by the \textit{extrude new} rule}]{
 \exists x. P \cpar \new x. \co{P}
}{
\infer[\mbox{by the \textit{select1} rule}]{
 \mathopen{\new x.} \left(\exists x .P \cpar \co P\right)
}{
\infer[\mbox{by Proposition~\ref{proposition:reflexivity}}]{
 \mathopen{\new x.} \left( P \cpar \co P \right)
}{
\infer[\mbox{by the \textit{tidy name} rule}]{
 \new x. \cunit
}{
 \cunit
}}}}
\]
The above also serves as a proof of the dual statement $\vdash \forall x. P \multimap \new x. P$.

Using the \textit{fresh} rule we can establish the following implication $\vdash \new x. P \multimap \wen x. P$, as follows.
%between \textit{new} and \textit{wen}.
\[
\infer[\mbox{by the \textit{fresh} rule}]{
 \wen x .\co{P} \cpar \wen x .P
}{
\infer[\mbox{by Proposition~\ref{proposition:reflexivity}}]{
 \new x .\co{P} \cpar \wen x .P
}{
 \cunit
}}
\]
\begin{comment}
\[
\infer[\mbox{by the \textit{fresh} rule}]{
 \wen x .\co{P} \cpar \wen x .P
}{
\infer[\mbox{by the \textit{close} rule}]{
 \new x .\co{P} \cpar \wen x .P
}{
\infer[\mbox{by Proposition~\ref{proposition:reflexivity}}]{
 \mathopen{\new x .}\left( \co{P} \cpar P \right)
}{
\infer[\mbox{by the \textit{tidy} rule}]{
 \new x .\cunit
}{
 \cunit
}}}}
\]
\end{comment}
This completes the chain $\vdash \forall x. P \multimap \new x. P$, $\vdash \new x. P \multimap \wen x .P$ and $\vdash \wen x. P \multimap \exists x. P$. These linear implications are strict unless $\nfv{x}{P}$, in which case, for $\quantifier \in \left\{ \forall, \exists, \new, \wen \right\}$, $\quantifier x .P$ is logically equivalent to $P$. For example, using the \textit{fresh} rule followed by the \textit{extrude new} and \textit{tidy name} rules, $\vdash
 \new x .P \multimap P$
 holds, whenever $\nfv{x}{P}$.
%\[
%\infer[]{
% \wen x. \co{P} \cpar P
%}{
%\infer[]{
% \new x .\co{P} \cpar P
%}{
%\infer[]{
%\mathopen{\new x.} \left( \co{P} \cpar P \right)
%}{
%\infer[]{
% \new x .\cunit
%}{
% \cunit
%}}}}
%\]
Thus the implication corresponding to the \textit{vacuous} rule as in Fig.~\ref{figure:BVQ} is provable for any quantifier.
% in \textsf{MAV1}.
%Alternatively, the \textit{extrude new} and \textit{fresh} rules could be replaced by extending the congruence with $\new x .P \equiv P \equiv \wen x. P$, where $\nfv{x}{P}$. However, this has the disadvantage that an arbitrary number of nominal quantifiers can be introduced during proof search thereby jeopardising analyticity~\cite{Bruscoli2009}.

\textbf{The medial rules for nominals.}
The \textit{medial new} rule is particular to handling nominals in the presence of the self-dual non-commutative operator \textit{seq}. To see why this medial rule cannot be excluded, consider the following
% provable predicates.
formulae, where $x$ is free for atoms $\bbb$, $\ccc$, $\eee$ and $\ffff.$
%\ross{Choose symbols for atoms.}
\[
\begin{array}{rl}
 \left( \aaa \cseq \mathopen{\wen x.}\left( \bbb \cseq \ccc \right)\right) \tensor \left( \ddd \cseq \mathopen{\wen x.}\left( \eee \cseq \ffff \right) \right)
 \!\!\!&\multimap
 \left( \aaa \cseq \exists x. \bbb \cseq \exists x. \ccc \right) \tensor \left( \ddd \cseq \exists x. \eee \cseq \exists x .\ffff \right) 
\\
 \left( \aaa \cseq \exists x. \bbb \cseq \exists x .\ccc \right) \tensor \left( \ddd \cseq \exists x. \eee \cseq \exists x .\ffff \right) 
 \!\!\!&\multimap
 \left(\left( \aaa \cseq \exists x .\bbb \right) \tensor \left( \ddd \cseq \exists x. \eee \right)\right) \cseq \left( \exists x. \ccc \tensor \exists x. \ffff \right)
\end{array}
\]
Without using the \textit{medial new} rule, the above formulae are provable.
The first is as follows.
\[
\infer[\mbox{\textit{switch}}]{
 \left( \co{\aaa} \cseq \mathopen{\new x.}\left( \co{\bbb} \cseq \co{\ccc} \right)\right) \cpar \left( \co{\ddd} \cseq \mathopen{\new x.}\left( \co{\eee} \cseq \co{\ffff} \right) \right)
 \cpar
 \left( \aaa \cseq \exists x. \bbb \cseq \exists x. \ccc \right) \tensor \left( \ddd \cseq \exists x. \eee \cseq \exists x .\ffff \right) 
}{
\infer[\mbox{\textit{sequence}}]{
 \left(
  \left( \co{\aaa} \cseq \mathopen{\new x.}\left( \co{\bbb} \cseq \co{\ccc} \right)\right) \cpar \left( \aaa \cseq \exists x. \bbb \cseq \exists x. \ccc \right)
 \right)
 \tensor
 \left(
  \left( \co{\ddd} \cseq \mathopen{\new x.}\left( \co{\eee} \cseq \co{\ffff} \right) \right) \cpar \left( \ddd \cseq \exists x. \eee \cseq \exists x .\ffff \right) 
 \right)
}{
\infer[\mbox{\textit{extrude}}]{
 \left(
  \left( \co{\aaa} \cpar \aaa \right) \cseq \left( \mathopen{\new x.}\left( \co{\bbb} \cseq \co{\ccc} \right) \cpar \left( \exists x. \bbb \cseq \exists x. \ccc \right)\right)
 \right)
 \tensor
 \left(
  \left( \co{\ddd} \cpar \ddd \right) \cseq \left(\mathopen{\new x.}\left( \co{\eee} \cseq \co{\ffff} \right) \cpar \left( \exists x. \eee \cseq \exists x .\ffff \right) \right)
 \right)
}{
\infer[\mbox{\textit{select1}}]{
 \left(
  \left( \co{\aaa} \cpar \aaa \right) \cseq  \mathopen{\new x.}\left(\left( \co{\bbb} \cseq \co{\ccc} \right) \cpar \left( \exists x. \bbb \cseq \exists x. \ccc \right)\right)
 \right)
 \tensor
 \left(
  \left( \co{\ddd} \cpar \ddd \right) \cseq \mathopen{\new x.}\left(\left( \co{\eee} \cseq \co{\ffff} \right) \cpar \left( \exists x. \eee \cseq \exists x .\ffff \right) \right)
 \right)
}{
\infer[\mbox{by Proposition~\ref{proposition:reflexivity}}]{
 \left(
  \left( \co{\aaa} \cpar \aaa \right) \cseq  \mathopen{\new x.}\left(\left( \co{\bbb \cseq \ccc} \right) \cpar \left( \bbb \cseq \ccc \right)\right)
 \right)
 \tensor
 \left(
  \left( \co{\ddd} \cpar \ddd \right) \cseq \mathopen{\new x.}\left(\co{\eee \cseq \ffff} \cpar \left( \eee \cseq \ffff \right) \right)
 \right)
}{
\infer[\mbox{by \textit{tidy name}}]{
 \left(
  \mathopen{\new x.}\cunit
 \right)
 \tensor
 \left(
  \mathopen{\new x.}\cunit
 \right)
}{
\cunit
}}}}}}
\]
The proof of the second formula above is as follows.
\[
\infer[\mbox{by \textit{sequence}}]{
 \left( \co{\aaa} \cseq \forall x. \co{\bbb} \cseq \forall x .\co{\ccc} \right) \cpar \left( \co{\ddd} \cseq \forall x. \co{\eee} \cseq \forall x .\co{\ffff} \right) 
 \cpar
 \left(\left( \aaa \cseq \exists x .\bbb \right) \tensor \left( \ddd \cseq \exists x. \eee \right)\right) \cseq \left( \exists x. \ccc \tensor \exists x. \ffff \right)
\\\quad%\!\!\!\!\!
}{
\infer[\mbox{by \textit{sequence}}]{
 \left(\left( \co{\aaa} \cseq \forall x. \co{\bbb}\right) \cpar \left(\co{\ddd} \cseq \forall x. \co{\eee}\right) \right) \cseq \left( \forall x .\co{\ccc} \cpar \forall x .\co{\ffff} \right) 
 \cpar
 \left(\left( \aaa \cseq \exists x .\bbb \right) \tensor \left( \ddd \cseq \exists x. \eee \right)\right) \cseq \left( \exists x. \ccc \tensor \exists x. \ffff \right)
}{
\infer[\mbox{by Prop.~\ref{proposition:reflexivity}}]{
 \left( \co{\left( \aaa \cseq \exists x .\bbb \right) \tensor \left( \ddd \cseq \exists x. \eee \right)} \cpar \left(\left( \aaa \cseq \exists x .\bbb \right) \tensor \left( \ddd \cseq \exists x. \eee \right)\right)\right) \cseq \left( \co{\exists x. \ccc \tensor \exists x. \ffff} \cpar \left(\exists x. \ccc \tensor \exists x. \ffff \right)\right)
}{
 \cunit
}}}
\]


However, the issue is that
the following formula would not be provable without using the \textit{medial new} rule; hence cut elimination cannot hold without the \textit{medial new} rule.
\[
\begin{array}{l}
 \left( \aaa \cseq \mathopen{\wen x.}\left( \bbb \cseq \ccc \right)\right) \tensor \left( \ddd \cseq \mathopen{\wen x.}\left( \eee \cseq \ffff \right) \right)
 %\\\qquad\qquad\qquad\qquad\qquad
 \multimap
 \left(\left( \aaa \cseq \exists x. \bbb \right) \tensor \left( \ddd \cseq \exists x. \eee \right)\right) \cseq \left( \exists x .\ccc \tensor \exists x. \ffff \right)
\end{array}
\]
In contrast, with the \textit{medial new} rule the above formula is provable, as verified by the proof in Figure~\ref{fig:medial-ex}.
Notice the above proofs use only the \textit{medial new}, \textit{extrude new} and \textit{tidy name} rules for nominals. These rules are of the same form as rules \textit{medial1}, \textit{extrude1} and \textit{tidy1} for universal quantifiers, hence the same argument holds for the necessity of the \textit{medial1} rule by replacing $\new$ with $\forall$.

\begin{figure}[h!]
\[
\infer[]{
 \left( \co{\aaa} \cseq \new x. \left( \co{\bbb} \cseq \co{\ccc} \right)\right) \cpar \left( \co{\ddd} \cseq \new x. \left( \co{\eee} \cseq \co{\ffff} \right) \right)
 \cpar 
 \left(\left( \aaa \cseq \exists x .\bbb \right) \tensor \left( \ddd \cseq \exists x .\eee \right)\right) \cseq
 \left( \exists x .\ccc \tensor \exists x. \ffff \right)
}{
\infer[]{
  \left( \co{\aaa} \cseq \new x. \co{\bbb} \cseq \new x. \co{\ccc} \right) \cpar \left( \co{\ddd} \cseq \new x. \co{\eee} \cseq \new x .\co{\ffff} \right)
  %\\
 \cpar 
 \left( \left( \aaa \cseq \exists x. \bbb \right) \tensor \left( \ddd \cseq \exists x. \eee \right) \right) \cseq
 \left( \exists x.\ccc \tensor \exists x. \ffff \right)
}{
\infer[]{
  \left( \left( \co{\aaa} \cseq \new x. \co{\bbb} \right) \cpar \left( \co{\ddd} \cseq \new x. \co{\eee} \right)\right)
  \cseq
  \left( \new x. \co{\ccc} \cpar \new x .\co{\ffff} \right)
%  \\
  \cpar
  \left( \left( \aaa \cseq \exists x. \bbb \right) \tensor \left( \ddd \cseq \exists x. \eee \right) \right) \cseq
  \left( \exists x .\ccc \tensor \exists x. \ffff \right)
}{
\infer[]{
  \left( \left( \co{\aaa} \cseq \new x .\co{\bbb} \right) \cpar \left( \co{\ddd} \cseq \new x. \co{\eee} \right)
  \cpar
  \left( \left( \aaa \cseq \exists x. \bbb \right) \tensor \left( \ddd \cseq \exists x .\eee \right) \right) \right) 
%  \\
  \cseq 
  \left( \new x. \co{\ccc} \cpar \new x. \co{\ffff} \cpar \left( \exists x.\ccc \tensor \exists x. \ffff \right) \right)
}{
\infer[]{
 \left( \left( \left( \co{\aaa} \cseq \new x. \co{\bbb} \right) \cpar \left( \aaa \cseq \exists x .\bbb \right) \right)
  \tensor
  \left(  \left( \co{\ddd} \cseq \new x .\co{\eee} \right) \cpar \left( \ddd \cseq \exists x. \eee \right) \right) \right) 
% \\
 \cseq 
  \left(  \left( \new x. \co{\ccc} \cpar \exists x .\ccc \right) \tensor \left( \new x. \co{\ffff} \cpar \exists x. \ffff \right) \right)
}{
\infer[]{
  \left( \left( \left( \co{\aaa} \cpar \aaa \right) \cseq \left( \new x .\co{\bbb} \cpar \exists x .\bbb \right) \right)
  \tensor
  \left(  \left( \co{\ddd} \cpar \ddd \right) \cseq \left( \new x .\co{\eee} \cpar \exists x .\eee \right) \right) \right) 
 % \\
  \cseq 
  \left(  \left( \new x. \co{\ccc} \cpar \exists x .\ccc \right) \tensor \left( \new x. \co{\ffff} \cpar \exists x. \ffff \right) \right)
}{
\infer[]{
  \left( \left( \left( \co{\aaa} \cpar \aaa \right) \cseq \mathopen{\new x.}\left(  \co{\bbb} \cpar \exists x .\bbb \right) \right)
  \tensor
  \left(  \left( \co{\ddd} \cpar \ddd \right) \cseq  \mathopen{\new x.} \left(\co{\eee} \cpar \exists x. \eee \right) \right) \right) 
  %\\
  \cseq 
  \left(  \mathopen{\new x.} \left( \co{\ccc} \cpar \exists x. \ccc \right) \tensor \mathopen{\new x.}\left( \co{\ffff} \cpar \exists x .\ffff \right) \right)
}{
\infer[]{
  \left( \left( \left( \co{\aaa} \cpar \aaa \right) \cseq \mathopen{\new x.} \left(  \co{\bbb} \cpar \bbb \right) \right)
  \tensor
  \left(  \left( \co{\ddd} \cpar \ddd \right) \cseq  \mathopen{\new x.}\left(\co{\eee} \cpar \eee \right) \right) \right) 
  %\\
  \cseq
  \left(  \mathopen{\new x.}\left( \co{\ccc} \cpar \ccc \right) \tensor \mathopen{\new x.}\left( \co{\ffff} \cpar \ffff \right) \right)
}{
\infer[]{
  \left(
  \new x .\cunit
  \tensor
  \new x .\cunit \right) 
  \cseq
  \left(  \new x. \cunit \tensor \new x. \cunit \right)    
}{
  \cunit
}}}}}}}}}
%\mbox{by \textit{atomic interaction} and \textit{tidy name}}
 \]
 \caption{A proof of  $ \left( \aaa \cseq \mathopen{\wen x.}\left( \bbb \cseq \ccc \right)\right) \tensor \left( \ddd \cseq \mathopen{\wen x.}\left( \eee \cseq \ffff \right) \right)
 %\\\qquad\qquad\qquad\qquad\qquad
 \multimap
 \left(\left( \aaa \cseq \exists x. \bbb \right) \tensor \left( \ddd \cseq \exists x. \eee \right)\right) \cseq \left( \exists x .\ccc \tensor \exists x. \ffff \right)$
}
 \label{fig:medial-ex}
 \end{figure}


Including the \textit{medial new} rule forces the \textit{suspend} rule to be included. To see why, observe that the following linear implications are provable.
\[
\begin{array}{rl}
\left( \new x. \aaa \cseq \new x. \bbb \right) \tensor \left( \new x .\ccc \cseq \new x .\ddd \right)
\multimap&
\!\!\!\!
\mathopen{\new x.}\left( \aaa \cseq \bbb \right) \tensor \mathopen{\new x.}\left( \ccc \cseq \ddd \right)
\\
\mathopen{\new x.}\left( \aaa \cseq \bbb \right) \tensor \mathopen{\new x.}\left( \ccc \cseq \ddd \right)
\multimap&
\!\!\!\!
\mathopen{\new x.}\left( \left( \aaa \cseq \bbb \right) \tensor \left( \ccc \cseq \ddd \right) \right)
\end{array}
\]
However, without the \textit{suspend} rule the following implication is not provable, which would contradict the cut elimination result of this paper.
\[
\left( \new x .\aaa \cseq \new x .\bbb \right) \tensor \left( \new x. \ccc \cseq \new x. \ddd \right)
\multimap
\mathopen{\new x.}\left( \left( \aaa \cseq \bbb \right) \tensor \left( \ccc \cseq \ddd \right) \right)
\]
Fortunately, including the \textit{suspend} rule ensures that the above implication is provable as follows.
\[
\infer[\mbox{by \textit{suspend}}]{
\left( \wen x .\co{\aaa} \cseq \wen x .\co{\bbb} \right) \cpar \left( \wen x. \co{\ccc} \cseq \wen x. \co{\ddd} \right)
\cpar
\mathopen{\new x.}\left( \left( \aaa \cseq \bbb \right) \tensor \left( \ccc \cseq \ddd \right) \right)
}{
\infer[\mbox{by \textit{suspend}}]{
\mathopen{\wen x.}\left( \co{\aaa} \cseq \co{\bbb} \right) \cpar \mathopen{\wen x.}\left( \co{\ccc} \cseq \co{\ddd} \right)
\cpar
\mathopen{\new x.}\left( \left( \aaa \cseq \bbb \right) \tensor \left( \ccc \cseq \ddd \right) \right)
}{
\infer[\mbox{by Proposition~\ref{proposition:reflexivity}}]{
\mathopen{\wen x.}\left(\left( \co{\aaa} \cseq \co{\bbb} \right) \cpar \left( \co{\ccc} \cseq \co{\ddd} \right)\right)
\cpar
\mathopen{\new x.}\left( \left( \aaa \cseq \bbb \right) \tensor \left( \ccc \cseq \ddd \right) \right)
}{
\cunit
}}}
\]
A similar argument justifies the inclusion of the \textit{left wen} and \textit{right wen} rules.

\textbf{Rules induced by equivariance.}
Interestingly, \textit{equivariance} is a design decision in the sense that cut elimination still holds if we drop the \textit{equivariance} rule from the structural congruence.
For such a system without \textit{equivariance}, also the rules \textit{all name}, \textit{with name}, \textit{left name} and \textit{right name} could also be dropped.
Perhaps there may be interesting applications for a non-equivariant nominal quantifiers; however, for embedding of process such as $\nu$ in the $\pi$-calculus, \textit{equivariance} is an essential property for scope extrusion. For example, \textit{equivariance} is used when proving the embedding of labelled transition $\nu x. \nu y.\co{z}y.p \lts{\co{z}{(y)}} \nu x. p$, assuming $z \not= x$ and $z \not= y$.

In our embedding of the $\pi$-calculus in $\mathsf{MAV1}$, addressed thoroughly in a companion paper~\cite{MSCS}, we assume process $p$ is embedded as formula $P$.
In this case, process $\nu x. \nu y.\co{z}y.p$ maps to $Q = \mathopen{\new x. \new y.}\left(\co{\textrm{act}(z, y)} \cseq P\right)$, process $\nu x. p$ maps to $R = \new x. P$.
In this embedding of processes as formulae, we can prove that whenever the above labelled transition is enabled,
we can prove the following implication
$
\mathopen{\new y.}\left( \co{\textrm{act}(z, y)} \cseq {R} \right)
\multimap
Q
$, where the binder $\new y$ and atom $\textrm{act}(z, y)$ indicate that the process can commit to a bound output.
Indeed this formula is provable, as follows, by using \textit{equivariance}.
\[
\infer[\mbox{by \textit{equivariance}}]{
\mathopen{\wen y.}\left( \textrm{act}(z, y) \cseq \wen x.\co{P} \right)
\cpar
\mathopen{\new x. \new y.}\left(\co{\textrm{act}(z, y)} \cseq P\right)
}{
\infer[\mbox{by \textit{close}}]{
\mathopen{\wen y.}\left( \textrm{act}(z, y) \cseq \wen x.\co{P} \right)
\cpar
\mathopen{\new y. \new x.}\left(\co{\textrm{act}(z, y)} \cseq P\right)
}{
\infer[\mbox{by \textit{medial new}}]{
\mathopen{\new y.} \left(\left( \textrm{act}(z, y) \cseq \wen x.\co{P}\right)
\cpar
\mathopen{\new x.} \left(\co{\textrm{act}(z, y)} \cseq P\right) \right)
}{
\infer[\mbox{by \textit{sequence}}]{
\mathopen{\new y.} \left(\left( \textrm{act}(z, y) \cseq \wen x.\co{P}\right)
\cpar
 \left(\mathopen{\new x.}\co{\textrm{act}(z, y)} \cseq \mathopen{\new x.}P\right) \right)
}{
\infer[\mbox{by \textit{extrude new}}]{
\mathopen{\new y.} \left(
\left( \textrm{act}(z, y) \cpar \mathopen{\new x.}\co{\textrm{act}(z, y)} \right)
\cseq
 \left(\wen x.\co{P} \cpar \mathopen{\new x.}P\right) \right)
}{
\infer[\mbox{by Proposition~\ref{proposition:reflexivity}}]{
\mathopen{\new y.} \left(
\mathopen{\new x.}\left( \textrm{act}(z, y) \cpar \co{\textrm{act}(z, y)} \right)
\cseq
 \left(\co{\new x.P} \cpar \mathopen{\new x.}P\right) \right)
}{
\infer[\mbox{by \textit{tidy name}}]{
\mathopen{\new y.} 
\mathopen{\new x.}\cunit
}{
\cunit
}}}}}}}
\]
In response to the above problem, modelling the $\pi$-calculus, \textsf{MAV1} includes equivariance.



The \textit{equivariance} rule forces additional distributivity properties for $\new$ and $\wen$ over $\wwith$ and $\forall$, given by the \textit{all name}, \textit{with name}, \textit{left name}, \textit{right name} rules.
These rules allow $\new$ and $\wen$ quantifiers to propagate to the front of certain contexts.
To see why these rules are necessary consider the following implications,
with matching formulae, respectively, after and before the implication.
\[
\vdash
\mathopen{\new x.}\left(\new y. \forall z. \aaa \cpar \mathopen{\wen y.}\left( \bbb \wwith \ccc \right)\right)
\multimap
\new x.\new y. \forall z. \aaa \cpar \mathopen{\wen x. \wen y.}\left( \bbb \wwith \ccc \right)
\]
\[
\vdash
\new x.\new y. \forall z. \aaa \cpar \mathopen{\wen x. \wen y.}\left( \bbb \wwith \ccc \right)
\multimap
\new y. \forall z. \new x. \aaa \cpar \mathopen{\wen y.}\left( \wen x. \bbb \wwith \wen x. \ccc \right)
\]
Any proof of the second implication does involve \textit{equivariance}; but neither proof requires
\textit{all name} or \textit{with name}.
A proof of the first implication above is as follows.
\[
\infer[\mbox{by \textit{close}}]{
\mathopen{\wen x.}\left(\wen y. \exists z. \co{\aaa} \tensor \mathopen{\new y.}\left( \co{\bbb} \ooplus \co{\ccc} \right)\right)
\cpar
\new x.\new y. \forall z. \aaa \cpar \mathopen{\wen x. \wen y.}\left( \bbb \wwith \ccc \right)
}{
\infer[\mbox{by Proposition~\ref{proposition:reflexivity}}]{
\mathopen{\wen x.}\left(
\wen y. \exists z. \co{\aaa} \tensor \mathopen{\new y.}\left( \co{\bbb} \ooplus \co{\ccc} \right)\right)
\cpar
\mathopen{\new x.}\left(\new y. \forall z. \aaa \cpar \mathopen{\wen y.}\left( \bbb \wwith \ccc \right)\right)
}{
\cunit
}}
\]
\begin{comment}
\infer[\mbox{by \textit{switch}}]{
\mathopen{\new x.}\left(
\wen y. \exists z. \co{\aaa} \tensor \mathopen{\new y.}\left( \co{\bbb} \ooplus \co{\ccc} \right)
\cpar
\new y. \forall z. \aaa \cpar \mathopen{\wen y.}\left( \bbb \wwith \ccc \right)
\right)
}{
\infer[\mbox{by Proposition~\ref{proposition:reflexivity}}]{
\mathopen{\new x.}\left(
\left(\co{\new y. \forall z. \aaa} \cpar \new y. \forall z. \aaa\right)
\tensor
\left(\co{\mathopen{\wen y.}\left( \bbb \wwith \ccc \right)} \cpar \mathopen{\wen y.}\left( \bbb \wwith \ccc \right)\right)
\right)
}{
\infer[\mbox{by \textit{tidy name}}]{
\mathopen{\new x.}\cunit
}{
\cunit
}}}}}
\]
\end{comment}
A proof of the second implication above is given in Figure~\ref{fig:equiv-ex}.

\begin{figure}[h!]
\[
\infer[\mbox{by \textit{switch}}]{
\left(
\wen x.\wen y. \exists z. \co{\aaa} \tensor \mathopen{\new x. \new y.}\left( \co{\bbb} \ooplus \co{\ccc} \right)
\right)
\cpar
\new y. \forall z. \new x. \aaa \cpar \mathopen{\wen y.}\left( \wen x. \bbb \wwith \wen x.\ccc \right)
}{
\infer[\mbox{by \textit{equivariance} and \textit{close}}]{
\left(\wen x.\wen y. \exists z. \co{\aaa} \cpar \new y. \forall z. \new x. \aaa \right)
\tensor
\left(
\mathopen{\new x. \new y.}\left( \co{\bbb} \ooplus \co{\ccc} \right)
\cpar \mathopen{\wen y.}\left( \wen x. \bbb \wwith \wen x.\ccc \right)
\right)
}{
\infer[\mbox{by \textit{extrude1}}]{
\mathopen{\new y.}
\left(\wen x. \exists z. \co{\aaa} \cpar \forall z. \new x. \aaa \right)
\tensor
\left(
\mathopen{\new x. \new y.}\left( \co{\bbb} \ooplus \co{\ccc} \right)
\cpar \mathopen{\wen y.}\left( \wen x. \bbb \wwith \wen x.\ccc \right)
\right)
}{
\infer[\mbox{by \textit{close}}]{
\mathopen{\new y.\forall z. }
\left(\wen x. \exists z. \co{\aaa} \cpar \new x. \aaa \right)
\tensor
\left(
\mathopen{\new x. \new y.}\left( \co{\bbb} \ooplus \co{\ccc} \right)
\cpar \mathopen{\wen y.}\left( \wen x. \bbb \wwith \wen x.\ccc \right)
\right)
}{
\infer[\mbox{by \textit{select1}}]{
\mathopen{\new y.\forall z. \new x.}
\left(\exists z. \co{\aaa} \cpar  \aaa \right)
\tensor
\left(
\mathopen{\new x. \new y.}\left( \co{\bbb} \ooplus \co{\ccc} \right)
\cpar \mathopen{\wen y.}\left( \wen x. \bbb \wwith \wen x.\ccc \right)
\right)
}{
\infer[\mbox{by \textit{equivariance} and \textit{close}}]{
\mathopen{\new y.\forall z. \new x.}
\left(\co{\aaa} \cpar  \aaa \right)
\tensor
\left(
\mathopen{\new x. \new y.}\left( \co{\bbb} \ooplus \co{\ccc} \right)
\cpar \mathopen{\wen y.}\left( \wen x. \bbb \wwith \wen x.\ccc \right)
\right)
}{
\infer[\mbox{by \textit{external}}]{
\mathopen{\new y.\forall z. \new x.}
\left(\co{\aaa} \cpar \aaa \right)
\tensor
\mathopen{\new y.}
\left(
\mathopen{\new x.}\left( \co{\bbb} \ooplus \co{\ccc} \right)
\cpar \left( \wen x. \bbb \wwith \wen x.\ccc \right)
\right)
}{
\infer[\mbox{by \textit{close}}]{
\mathopen{\new y.\forall z. \new x.}
\left(\co{\aaa} \cpar \aaa \right)
\tensor
\mathopen{\new y.}
\left(
\left(
\mathopen{\new x.}\left( \co{\bbb} \ooplus \co{\ccc} \right) \cpar \wen x. \bbb
\right) \wwith
\left(
\mathopen{\new x.}\left( \co{\bbb} \ooplus \co{\ccc} \right) \cpar \wen x.\ccc
\right)
\right)
}{
\infer[\mbox{by \textit{left} and \textit{right}}]{
\mathopen{\new y.\forall z. \new x.}
\left(\co{\aaa} \cpar  \aaa \right)
\tensor
\mathopen{\new y.}
\left(
\mathopen{\new x.}\left(\left( \co{\bbb} \ooplus \co{\ccc} \right) \cpar \bbb\right)
 \wwith
\mathopen{\new x.}\left(\left( \co{\bbb} \ooplus \co{\ccc} \right) \cpar \ccc\right)
\right)
}{
\infer[\mbox{by \textit{atomic interaction}}]{
\mathopen{\new y.\forall z. \new x.}
\left(\co{\aaa} \cpar \aaa \right)
\tensor
\mathopen{\new y.}
\left(
\mathopen{\new x.}\left( \co{\bbb} \cpar \bbb\right)
 \wwith
\mathopen{\new x.}\left(\co{\ccc}  \cpar \ccc\right)
\right)
}{
\infer[\mbox{by \textit{tidy name} and \textit{tidy1}}]{
\mathopen{\new y.\forall z. \new x.}
\cunit
\tensor
\mathopen{\new y.}
\left(
\mathopen{\new x.}\cunit
 \wwith
\mathopen{\new x.}\cunit
\right)
}{
\cunit
}}}}}}}}}}}
\]
\caption{A proof of $\new x.\new y. \forall z. \aaa \cpar \mathopen{\wen x. \wen y.}\left( \bbb \wwith \ccc \right) \multimap
\new y. \forall z. \new x. \aaa \cpar \mathopen{\wen y.}\left( \wen x. \bbb \wwith \wen x. \ccc \right)$}
\label{fig:equiv-ex}
\end{figure}


%proceeds as follows.
By the implications above, if cut elimination holds, it must be the case that the following is provable.
\[
\mathopen{\new x.}\left(\new y. \forall z. \aaa \cpar \mathopen{\wen y.}\left( \bbb \wwith \ccc \right)\right)
\multimap
\new y. \forall z. \new x. \aaa \cpar \mathopen{\wen y.}\left( \wen x. \bbb \wwith \wen x. \ccc \right)
\]
However, without the \textit{all name} and \textit{with name} rules,
the above implication is not provable and hence cut elimination would not hold in the presence of \textit{equivariance}.
Fortunately, using both the \textit{all name} and \textit{with name} rules the above implication is provable, as follows.
\[
\infer[\mbox{\textit{all name} and \textit{equivariance}}]{
\mathopen{\wen x.}\left(\wen y. \exists z. \co{\aaa} \tensor \mathopen{\new y.}\left( \co{\bbb} \ooplus \co{\ccc} \right)\right)
\cpar
\new y. \forall z. \new x. \aaa \cpar \mathopen{\wen y.}\left( \wen x. \bbb \wwith \wen x. \ccc \right)
}{
\infer[\mbox{\textit{with name} and \textit{equivariance}}]{
\mathopen{\wen x.}\left(\wen y. \exists z. \co{\aaa} \tensor \mathopen{\new y.}\left( \co{\bbb} \ooplus \co{\ccc} \right)\right)
\cpar
\new x. \new y. \forall z. \aaa \cpar \mathopen{\wen y.}\left( \wen x. \bbb \wwith \wen x. \ccc \right)
}{
\infer[\mbox{by \textit{close}}]{
\mathopen{\wen x.}\left(\wen y. \exists z. \co{\aaa} \tensor \mathopen{\new y.}\left( \co{\bbb} \ooplus \co{\ccc} \right)\right)
\cpar
\new x. \new y. \forall z. \aaa \cpar \mathopen{\wen x. \wen y.}\left( \bbb \wwith \ccc \right)
}{
\infer[\mbox{by Proposition~\ref{proposition:reflexivity}}]{
\mathopen{\wen x.}\left(\wen y. \exists z. \co{\aaa} \tensor \mathopen{\new y.}\left( \co{\bbb} \ooplus \co{\ccc} \right)\right)
\cpar
\mathopen{\new x.} \left(\new y. \forall z. \aaa \cpar \mathopen{\wen y.}\left( \bbb \wwith \ccc \right)\right)
}{
\cunit
}}}}
\]
\begin{comment}
\infer[\mbox{by \textit{switch}}]{
\mathopen{\new x.}\left(
\left(\wen y. \exists z. \co{\aaa} \tensor \mathopen{\new y.}\left( \co{\bbb} \ooplus \co{\ccc} \right)\right)
\cpar
\new y. \forall z. \aaa \cpar \mathopen{\wen y.}\left( \bbb \wwith \ccc \right)
\right)
}{
\infer[\mbox{by Proposition~\ref{proposition:reflexivity}}]{
\mathopen{\new x.}\left(
\left(\co{\new y. \forall z. \aaa} \cpar \new y. \forall z. \aaa\right)
\tensor
\left( \co{\mathopen{\wen y.}\left( \bbb \wwith \ccc \right)} \cpar \mathopen{\wen y.}\left( \bbb \wwith \ccc \right) \right)
\right)
}{
\infer[\mbox{by \textit{tidy name}}]{
\mathopen{\new x.}\cunit
}{
\cunit
}}}}}}}
\] 
\end{comment}
A similar argument justifies the necessity of the \textit{left name} and \textit{right name} rules. 
%in the presence of \textit{equivariance}.
%Also the above argument justifies the \textit{with name} and \textit{all name} for the \textit{new} quantifier by switching the roles of $\wen$ and $\new$ in the above proofs.


\textbf{Polarities of the nominals.}
As with focussed proof search~\cite{Andreoli1992,Chaudhuri2011}, assigning a positive or negative polarity to operators explains certain distributivity properties.
Consider $\cpar$, $\wwith$, $\forall$ and $\new$ to be negative operators, %in the sense of focussing~\cite{Andreoli1992};
and $\tensor$, $\ooplus$, $\exists$ and $\wen$ to be positive operators, where \textit{seq} is both positive and negative. The negative quantifier $\new$ distributes over all positive operators. 
Considering positive operator \textit{tensor} %\textit{times}
for example, $\vdash \new x .\alpha \tensor \new x. \beta  \multimap \new x. \left( \alpha \tensor \beta \right)$ holds but the converse implication does not hold. Furthermore, $\wen x. \alpha \tensor \wen x. \beta$ and $\wen x .\left(\alpha \tensor \beta\right)$ are unrelated by linear implication in general.
%(assuming $x$ is a free variable in $\alpha$ and $\beta$).
Dually, for the negative operator \textit{par} the only distributivity property that holds for nominal quantifiers is $\vdash \wen x .\left( \alpha \cpar \beta \right) \multimap \wen x .\alpha \cpar \wen x .\beta$.
%Considering positive quantifier $\exists$, where $x$ and $y$ appear free in $\gamma$, $\vdash \exists x \new y \gamma \multimap \new y \exists x \gamma$ holds but the converse does not hold.
The \textit{new wen} rule completes this picture of \textit{new} distributing over positive operators and \textit{wen} distributing over negative operators.
From the perspective of embedding name-passing process calculi in logic, the above distributivity properties of \textit{new} and \textit{wen} suggest that processes should be
encoded using negative operators $\new$ and $\cpar$ for private names and parallel composition
(or perhaps dually, using positive operators $\wen$ and $\tensor$),
so as to avoid private names distributing over parallel composition,
%the distributivity property $\new$ over $\cpar$ (respectively, distributivity of $\wen$ over $\tensor$),
which we have shown to be problematic in Section~\ref{section:self-dual}.
%(see Section~\ref{sec:intro}).



The control of distributivity exercised by \textit{new} and \textit{wen} contrasts with the situation for universal and existential quantifiers, where $\exists$ commutes in one direction over all operators and $\forall$ commutes with all operators in the opposite direction, similarly to the additive $\ooplus$ and $\wwith$ which are also  insensitive to the polarity of operators with which they commute.
In the sense of control of distributivity~\cite{Blute2012}, \textit{new} and \textit{wen} behave more like multiplicatives than additives, but are unrelated to multiplicative quantifiers in the logic of bunched implications~\cite{BI}.
%considering $\exists$ for example, $\vdash \exists x \quantifier y \gamma \multimap \quantifier y \exists x \gamma$ and $\vdash \exists x \left( P \odot Q \right) \multimap \exists x P \odot \exists x Q$ hold for any quantifier $\quantifier \in \left\{ \new, \wen, \exists, \forall\right\}$ and any binary operator $\odot \in \left\{ \cpar, \tensor, \cseq, \wwith, \ooplus \right\}$ whether positive or negative, as exploited to push freshness checking to the atoms~\cite{Strassburger09}.


% --- end nominal.tex ---
% --- begin lemmas.tex ---
\section{The Splitting Technique for Renormalising Proofs}
\label{section:splitting}

This section presents the \textit{splitting} technique that is central to the cut elimination proof for \textsf{MAV1}.
%Cut elimination (Theorem~\ref{theorem:cut}) is achieved by advancing methods developed in the propositional sub-system \textsf{MAV}~\cite{Horne2015}.
Splitting is used to recover a syntax directed approach for sequent-like contexts.
%Thereby, despite rules being applied deep withing any context,
Recall that in the sequent calculus rules are always applied to the root connective of a formula in a sequent, whereas deep inference rules can be applied deep within any context.
The technique is used to guide proof normalisation
%(Lemma~\ref{lemma:context} and Lemmas~\ref{lemma:co-close} to~\ref{lemma:co-atoms})
leading to the cut elimination result at the end of Section~\ref{section:context}.

There are complex inter-dependencies between the nominals \textit{new} and \textit{wen} and other operators, particularly the multiplicatives \textit{times} and \textit{seq} and additive \textit{with}.
As such, the splitting proof is tackled as follows, as illustrated in Fig.~\ref{figure:strategy}:
\begin{itemize}
\item Splitting for the first-order universal quantifier $\forall$ can be treated independently of the other operators; hence a direct proof of splitting for this operator is provided first as a simple induction over the length of a derivation in Lemma~\ref{lemma:universal}. Splitting for all other operators are dependent on this lemma.

\item Due to inter-dependencies between $\new$, $\wen$, $\tensor$, $\cseq$ and $\wwith$, splitting for these operators are proven simultaneously by a (huge) mutual induction in Lemma~\ref{lemma:split-times}.
The induction is guided by an intricately designed multiset-based measure of the size of a proof in Definition~\ref{definition:size}. 
The balance of dependencies between operators in this lemma is, by far, the most challenging aspect of this paper.

\item Having established 
Lemma~\ref{lemma:universal} and Lemma~\ref{lemma:split-times},
splitting for the remaining operators $\exists$ and $\ooplus$ and the atoms can each be established independently of each other in Lemmas~\ref{lemma:split-exists},~\ref{lemma:split-plus} and~\ref{lemma:split-atoms} respectively.
\end{itemize}
%A direct proof of \textit{co-rule elimination} for universal quantifiers (Lemma~\ref{lemma:universal}) simplifies \textit{splitting} (Lemma~\ref{lemma:split-times}), and 
%Lemmas check the interplay between nominal quantifiers and other connectives, as illustrated by the discussion in the following subsection. %Examples explain the impact of the polarities of $\new$ and $\wen$.
\begin{figure}[h!]
\[
\xymatrix@C=1.5pc@R=2pc{
&& 
\mbox{\txt{
Splitting $\exists$
\\ (Lemma~\ref{lemma:split-exists})}}
\ar[dr]
%&&\mbox{(Lemmas~\ref{lemma:co-close}}
\\
\mbox{
\txt{
Splitting $\forall$
\\ (Lemma~\ref{lemma:universal})}}
\ar[r]
&
\mbox{\txt{
Splitting $\new, \wen, \tensor, \cseq, \wwith$
\\
 (Lemma~\ref{lemma:split-times})}}
\ar[r]
\ar[ru]
\ar[rd]
&
\mbox{\txt{
Splitting $\ooplus$\\
 (Lemma~\ref{lemma:split-plus})}}
\ar[r]
&
\mbox{Section~\ref{section:context}}
%\context{}
%\mbox{(Lemma~\ref{lemma:context})}
%\ar[r]\ar[ur]\ar[dr]
%&\vdots
\\
&&
\mbox{Splitting $\alpha, \co{\alpha}$ (Lemma~\ref{lemma:split-atoms})}
\ar[ur]
%&&\mbox{Lemma~\ref{lemma:co-atoms}}
}
\]
\caption{
The proof strategy:
dependencies between splitting lemmas leading to cut elimination.}
\label{figure:strategy}
\end{figure}

\subsection{Elimination of universal quantifiers from a proof}

%\begin{comment}
%\begin{lemma}
%The following properties hold for substitutions.
%\begin{itemize}
%\item
%For variables $x$, $y$ and value $v$, such that $\nfv{y}{P}$,
%$P \sub{x}{y} \sub{y}{v} \equiv P \sub{x}{v}$.
%
%\item
%For $x \not = y$ and $\nfv{y}{v}$,
%$P \sub{y}{t}\sub{x}{v}
%\equiv
%P \sub{x}{v} \sub{y}{t\sub{x}{v}}$.
%\end{itemize}
%\end{lemma}
%%\begin{proof}
%%Simple induction on the size of the predicate.
%%\end{proof}
%\end{comment}

%At several points in proofs we work up to substitutions.
We employ a trick 
%employed for \textsf{MAV}~\cite{Horne2015} to \textsf{MAV1} where %\textit{with} $\wwith$ and
where universal quantification $\forall$ receives a more direct treatment than other operators.
%The proof for \textit{with} has a knock on effect for the nominal quantifiers requiring some vacuous \textit{new} and \textit{wen} quantifiers to be removed; while 
The proof requires closure of rules under substitution of terms for variables, established as follows directly by induction over the length of a derivation using a function over formulae.
%(suggesting explicit substitutions are possible).
\begin{lemma}[Substitution]
\label{lemma:substitution}
If we have derivation $\vcenter{\infer{Q}{P}}$,
then we have derivation $\vcenter{\infer{Q\sub{x}{v}}{P\sub{x}{v}}}$.
\end{lemma}
\fullproof{
\begin{proof}
% --- begin sub-proof.tex ---
Consider the case of the \textit{switch} rule. The following rule can be derived.
\[
\left(\left(P \tensor Q\right) \cpar V\right) \sub{x}{v}
\equiv
\left(P\sub{x}{v} \tensor Q\sub{x}{v}\right) \cpar V\sub{x}{v}
\longrightarrow
P\sub{x}{v} \tensor \left(Q\sub{x}{v} \cpar V\sub{x}{v}\right)
\equiv
\left(P \tensor \left(Q \cpar V\right)\right)\sub{x}{v}
\]
The cases are similar for the rules \textit{sequence}, \textit{left}, \textit{right}, \textit{external}, \textit{atomic interaction}, \textit{tidy}, \textit{tidy1} and \textit{tidy name}.


Consider the case of the \textit{extrude new} rule. The following rule can be derived, where $\nfv{y}{Q}$, $z \not = x$, $\nfv{z}{v}$ and $\nfv{x}{(\new y. P \cpar Q)}$.
\[
\begin{array}{rl}
\left(\new y. P \cpar Q\right)\sub{x}{v}
\equiv&
\mathopen{\new z.} \left(P\sub{y}{z}\sub{x}{v}\right) \cpar Q
\\
\longrightarrow&
\mathopen{\new z.} \left(P\sub{y}{z}\sub{x}{v} \cpar Q \right)
\equiv
\mathopen{\new z.} \left(\left(P \cpar Q \right)\sub{y}{z}\sub{x}{v}\right)
\equiv
\mathopen{\new y.} \left(\left(P \cpar Q \right)\right)\sub{x}{v}
\end{array}
\]
The cases for \textit{extrude1}, \textit{left wen}, and \textit{right wen} rules are similar.


In the case of the \textit{select} rule. The following rule can be derived, where $z \not = x$, $\nfv{z}{v}$ and $\nfv{x}{(\exists y. P)}$.
\[
\begin{array}{rl}
\left( \exists y. P \right) \sub{x}{v}
\equiv
\left( \exists z. P\sub{y}{z}\sub{x}{v} \right) 
\longrightarrow
P\sub{y}{z}\sub{x}{v}\sub{z}{t\sub{x}{v}}
\equiv
P\sub{y}{z}\sub{z}{t}\sub{x}{v}
\equiv
P\sub{y}{t}\sub{x}{v}
\end{array}
\]

Consider the case of the \textit{close} rule. The following rule can be derived, where $z \not = x$, $\nfv{z}{v}$ and $\nfv{x}{(\wen{y}. P \cpar \new{y}. Q)}$.
\[
\begin{array}{rl}
\left( \wen{y}. P \cpar \new{y}. Q \right) \sub{x}{v}
\equiv&
\mathopen{\wen{z}}.\left(P\sub{y}{z}\sub{x}{v}\right) \cpar \mathopen{\new{z}.} \left( Q\sub{y}{z}\sub{x}{v} \right)
\\
\longrightarrow&
\mathopen{\new{z}.}\left(P\sub{y}{z}\sub{x}{v} \cpar Q\sub{y}{z}\sub{x}{v} \right)
\equiv
\mathopen{\new{z}.}\left(\left(P \cpar Q \right) \sub{y}{z}\sub{x}{v}\right)
\equiv
\left(\mathopen{\new{y}.}\left(P \cpar Q \right) \right)\sub{x}{v}
\end{array}
\]
The cases for \textit{medial1}, \textit{medial new}, \textit{suspend} are similar.

Consider the case of the \textit{new wen} rule. Pick fresh names $y'$ and $z'$ such that $x \not= y'$, $x \not=z'$, $y' \not= z'$, $\nfv{y'}{v}$, $\nfv{z'}{v}$, $\nfv{y'}{P}$, and $\nfv{z'}{P}$, hence $\sub{y}{y'}\sub{z}{z'} = \sub{z}{z'}\sub{y}{y'}$; and thereby the following rule can be construced.
\[
\left( \new y. \wen z. P \right)\sub{x}{v} 
\equiv
\mathopen{\new y'. \wen z'.} \left( P \sub{z}{z'}\sub{y}{y'}\sub{x}{v}  \right)
\longrightarrow
\mathopen{\wen z'. \new y'.} \left( P \sub{y}{y'}\sub{z}{z'}\sub{x}{v}  \right)
\equiv
\left(\wen z. \new y. P \right)\sub{x}{v}
\] 
The lemma follows by induction on the length of a derivation.
% --- end sub-proof.tex ---
\end{proof}
}
%

%Full proofs for the above lemma for the \textit{with} operator appears in a journal paper~\cite{Horne2015}. 
%The same techniques extend to first-order quantifiers, as demonstrated in the following lemma.

We can now establish, the following lemma directly, which is a \textit{co-rule} elimination lemma. 
By a co-rule, we mean that, for \textit{select} rule $\vcenter{\infer[]{\context{\exists x.P}}{\context{P\sub{x}{v}}}}$, 
there is complementary rule $\vcenter{\infer[]{\context{P\sub{x}{v}}}{\context{\forall x.P}}}$ where the direction of inference is reversed and the formulae are complemented. 
Such a co-rule can always be eliminated from a proof, in which case we say \textit{co-select1} is \textit{admissible}, as established by the following lemma.
\begin{lemma}[Universal]\label{lemma:universal}
If $\vdash \context{ \forall x. P }$ holds
 then, for all terms~$v$, $\vdash \context{ P\sub{x}{v} }$ holds.
\end{lemma}

%\begin{comment}
%\begin{lemma}[Universal]
%\label{lemma:universal}
%If $\vdash \context{ \forall x P }$ holds
% then for all terms~$v$, $\vdash \context{ P\sub{x}{v} }$ holds.
%\end{lemma}
%\begin{proof}
%We require a function over predicates $s_v(P)$ that replaces a certain universal quantifier in $P$ with a substitution for a value $v$. 
%The universal quantifiers to be replaced are highlighted in bold $\boldsymbol\forall$. Note that during a proof the bold external choice may be duplicated by the distribute rule for external choice, hence there may be multiple bold occurrences in a formula.
%The function is defined as follows, where $\odot \in \left\{ \cseq, \cpar, \tensor, \ooplus, \wwith \right\}$ is any binary connective and $\quantifier \in \left\{ \forall, \exists, \new, \wen \right\}$ is any quantifier except bold universal quantification.
%\begin{gather*}
%s_v(\quantifier x P) =
%\quantifier x s_v(P)
%\qquad
%s_v(\boldsymbol\forall x P) = s_v(P) \sub{x}{v}
%\qquad
%s_v(\cunit) = \cunit 
%\\[6pt]
%s_v\left( P \odot Q \right) = 
%s_v(P) \odot s_v(Q)
%\qquad
%%s_v(0) = 0 \qquad s_v(\top) = \top \qquad
%s_v(a) = a
%\qquad
%s_v(\co{a}) = \co{a}
%\end{gather*}
%Firstly, observe that if $P \equiv Q$, then $s_v(P) \equiv s_v(Q)$. The only case to consider is when $\alpha$-conversion is applied to a bold universal quantifier. In this case, $s_v\left( \boldsymbol\forall x P \right) = P\sub{x}{v} = P\sub{x}{y}\sub{y}{v} = s_v\left( \boldsymbol\forall y P\sub{x}{y} \right)$, where $\nfv{y}{\forall x P}$.
%
%
%
%For the base case, $s_v(P) = P$, in which case trivially if $\vdash P$ then $\vdash s_v(P)$, for example where $P \equiv \cunit$.
%\smallskip
%
%Consider the case when the bottommost rule in a proof involves the bold universal quantifier, as follows, where $\nfv{x}{Q}$ and $\vdash \context{ \boldsymbol\forall x \left( P \cpar Q \right) }$ holds.
%\[
%\context{ \boldsymbol\forall x P \cpar Q }
%\longrightarrow
%\context{ \boldsymbol\forall x \left( P \cpar Q \right) }
%\]
%
%By the induction hypothesis, $\vdash s_v\left(\context{ \boldsymbol\forall x \left( P \cpar Q \right) }\right)$ holds.
%Now the following equalities hold.
%\[
%\begin{array}{rl}
%s_v\left(\context{ \boldsymbol\forall x \left( P \cpar Q \right) }\right)
%&= s_v\left(\context{ \left( s_v(P) \cpar s_v(Q) \right) \sub{x}{v}}\right) \\
%&= s_v\left(\context{ s_v(P)\sub{x}{v} \cpar s_v(Q) }\right) \\
%&= s_v\left(\context{ \boldsymbol\forall x P \cpar Q }\right)
%\end{array}
%\]
%Hence $\vdash s_v\left(\context{ \boldsymbol\forall x P \cpar Q }\right)$ holds as required.
%\smallskip
%
%Consider the case where the bottommost rule of a proof is of the following form, where $\vdash \context{\cunit}$ holds.
%\[
%\context{ \boldsymbol \forall x \cunit } \longrightarrow \context{ \cunit }
%\]
%By the induction hypothesis, $\vdash s_v\left( \context{ \cunit } \right)$ holds.
%Furthermore, $s_v\left( \context{ \boldsymbol \forall x \cunit } \right) = s_v\left( \context{ \cunit } \right)$. Hence $\vdash s_v \left( \context{ \forall x \cunit } \right)$ holds, as required.
%\smallskip
%
%Consider when a rule appears within the scope of a bold universal quantifier, of the following form, where $\vdash \context{ \boldsymbol \forall x \contextb{ Q } }$ holds.
%\[
%\context{ \boldsymbol \forall x \contextb{ P } }
%\longrightarrow
%\context{ \boldsymbol \forall x \contextb{ Q } }
%\]
%By the induction hypothesis, $\vdash s_v\left( \context{ \boldsymbol \forall x \contextb{ Q } } \right)$ holds.
%Now both of the following equivalences hold.
%\[
%\begin{array}{rl}
%s_v\left( \context{ \boldsymbol \forall x \contextb{ Q } } \right)
%&\equiv
%s_v\left( \context{ \contextb{ Q }\sub{x}{v} } \right)
%\\
%s_v\left( \context{ \boldsymbol \forall x \contextb{ P } } \right)
%&\equiv
%s_v\left( \context{ \contextb{ P }\sub{x}{v} } \right)
%\end{array}
%\]
%Hence, by Lemma~\ref{lemma:substitution}, 
%$
%\vdash s_v\left( \context{ \boldsymbol \forall x \contextb{ P } } \right)
%$
%holds as required.
%\smallskip
%
%Otherwise, the bottommost rule does not directly involve a bold universal quantifier nor is it applied inside the scope of a bold universal quantifier. In this case, the bottommost rule is of the form.
%\[
%\context{ P } \longrightarrow \context{ Q }
%\]
%where $\vdash \context{ Q }$ holds.
%By the induction hypothesis, $\vdash s_v\left(\context{ Q }\right)$ holds. 
%Now $s_v\left(\context{ Q }\right) = \contextp{ s_v\left(Q\right) }$ for some context and %\bogdan{Who is $V$? Shouldn't be $T$ here and in what follows (3 times in total)?}
%$s_v\left(\context{ P }\right) = \contextp{ s_v\left(P\right) }$ for the same context. By Lemma~\ref{lemma:substitution}, since $P \longrightarrow Q$, $s_v\left(P\right) \longrightarrow s_v\left(Q\right)$; hence $\contextp{ s_v\left(P\right) } \longrightarrow \contextp{ s_v\left(Q\right) }$. Thereby $\vdash s_v\left(\context{ P }\right)$ holds, as required.
%\smallskip
%
%Now by induction on the length of a proof, if $\vdash P$ holds then $\vdash s_v\left( P \right)$ holds. So, as a special case, if $\vdash \context{ \boldsymbol\forall x P }$, where the universal quantifier is bold, then $\vdash \context{ P\sub{x}{v} }$ holds.
%\end{proof}
%\end{comment}

A corollary of Lemma~\ref{lemma:universal} is: if $\vdash \forall x. P \cpar Q$ then $\vdash P\sub{x}{y} \cpar Q$, where $\nfv{y}{\left(\forall x. P \cpar Q\right)}$.
This corollary is in the form of a \textit{splitting} lemma,
where we have a principal connective $\forall$ at the root of a formula inside a context of the form $\left\{\ \cdot\ \right\} \cpar Q$. This corollary of the above lemma should remind the reader of the (invertible) sequent calculus rule for universal quantifiers:
\[
\begin{prooftree}
\vdash P\sub{x}{y}, \Gamma
\justifies
\vdash \forall x. P, \Gamma
\using
\mbox{where $y$ is fresh for $\forall x. P$ and all formulae in $\Gamma$}
\end{prooftree}
\]
We discuss, the significance of splitting lemmas after some preliminary lemmas required for the main splitting result.

\subsection{Killing contexts and technical lemmas required for splitting}

We require a restricted form of context called a killing context (terminology is from~\cite{Chaudhuri2011}). 
A killing context is a context with one or more holes, defined as follows.
\begin{definition}
%\bogdan{Why it is called a killing context? according to this definition appears to be a simple context. However it seems that in the following proofs is used to remove some terms, but this attribute does not appear in the current definition}
A \textit{killing context} is a context defined by the following grammar.
\[
 \tcontext{} \Coloneqq \{\ \cdot\ \} \mid \tcontext{} \wwith \tcontext{} \mid \forall x. \tcontext{}
\mid \new x.\tcontext{}
\]
In the above, $\{\ \cdot\ \}$ is a hole into which any formula can be plugged. An $n$-ary killing context is a killing context in which $n$ holes appear.
\end{definition}

For readability of large formulae involving an $n$-ary killing context, for $n > 1$, we represent the holes using a comma-separated list, so for example, 
instead of writing $\tcontextsym \{\cdot\} \{ \cdot\}$, we write $\tcontext{\cdot, \cdot}$ for a binary context. Given an $n$-ary killing context $\tcontext{\ldots}$, we write $\tcontext{Q_1,\ldots,Q_n}$ to denote the formula obtained by filling the holes in the context with formulas $Q_1,\ldots, Q_n.$ 
%$\tcontext{}$ and family of $n$ predicates $Q_1$, $Q_2$, \ldots, $Q_n$,
We also introduce the notation $\tcontext{Q_i \colon 1 \leq i \leq n}$ as shorthand for $\tcontext{Q_1, Q_2, \hdots, Q_n}$; and $\tcontext{ Q_i \colon i \in I }$ for a family of formulae indexed by finite subset of natural numbers $I$. 
%

A killing context represents a context that cannot in general be removed until all other rules in a proof have been applied, hence the corresponding \textit{tidy} rules are suspended until the end of a proof.  A killing context has properties that are applied frequently in proofs, characterised by the following lemma. %Firstly, if $\cunit$ is plugged into all holes, then the context is provable. Secondly, the context distributes over parallel composition. 
\begin{lemma}\label{lemma:kill-distribute}
For any killing context $\tcontext{}$, $\vdash \tcontext{\cunit, \hdots, \cunit}$ holds; and, assuming the free variables of $P$ are not bound by $\tcontext{}$,
%(which is always possible by $\alpha$-conversion),
%, the following derivation holds.
we have derivation
\[
\infer[.]{
P \cpar \tcontext{Q_1, Q_2, \hdots Q_n}
}{
\tcontext{P \cpar Q_1, P \cpar Q_2, \hdots P \cpar Q_n}
}
\]
\end{lemma}
\fullproof{
% --- begin distribute-proof.tex ---
\begin{proof}
The proofs follow by straightforward inductions over the structure of a killing context.

%There are two base cases. When the killing context is the top only $P \cpar \top \longrightarrow \top$ and $\top \longrightarrow \cunit$, as required.
When the killing context is one hole only $P \cpar \ehole{Q} = \ehole{P \cpar Q}$ and $\ehole{\cunit} = \cunit$, as required.

Now assume that by the induction hypothesis the following hold for killing contexts $\tcontextn{1}{}$ and $\tcontextn{2}{}$, and also $\vdash \tcontextn{1}{\cunit, \hdots, \cunit}$ and $\vdash \tcontextn{2}{\cunit, \hdots, \cunit}$.
\[
\begin{array}{rl}
 P \cpar \tcontextn{1}{Q_1, \hdots, Q_m} &\longrightarrow \tcontextn{1}{P \cpar Q_1, \hdots, P \cpar Q_m}
 \\
 P \cpar \tcontextn{2}{Q_{m+1}, \hdots, Q_{m+n}} &\longrightarrow \tcontextn{2}{P \cpar Q_{m+1}, \hdots, P \cpar Q_{m+n}}
\end{array}
\]
Hence, by distributivity the following derivation can be constructed.
\[
\begin{array}{l}
P \cpar \left( \tcontextn{1}{Q_1, \hdots, Q_m} \wwith \tcontextn{2}{Q_{m+1}, \hdots, Q_{m+n}} \right)
\\\qquad
 \begin{array}{l}
 \longrightarrow P \cpar \tcontextn{1}{Q_1, \hdots, Q_m} \wwith P \cpar \tcontextn{2}{Q_{m+1}, \hdots, Q_{m+n}}
 \\
 \longrightarrow \tcontextn{1}{P \cpar Q_1, \hdots, P \cpar Q_m} \wwith \tcontextn{2}{P \cpar Q_{m+1}, \hdots, P \cpar Q_{m+n}}
 \end{array}
\end{array}
\]
Furthermore, $\tcontextn{1}{\cunit, \hdots, \cunit} \wwith \tcontextn{2}{\cunit, \hdots, \cunit} \longrightarrow \cunit \wwith \cunit \longrightarrow \cunit$ holds.

As the induction hypothesis, assume $P \cpar \tcontext{ Q_1, \hdots, Q_n } \longrightarrow \tcontext{ P \cpar Q_1, \hdots, P \cpar Q_n }$ and also that $\vdash \tcontext{ \cunit, \hdots, \cunit }$ . Hence, given $\nfv{x}{P}$, the following derivation holds.
\[
P \cpar \new x \tcontext{ Q_1, \hdots, Q_n } \longrightarrow \new x \left( P \cpar \tcontext{ Q_1, \hdots, Q_n }\right) \longrightarrow \new x \tcontext{ P \cpar Q_1, \hdots, P \cpar Q_n }
\]
Furthermore, $\new x \tcontext{ \cunit, \hdots, \cunit} \longrightarrow \new x \cunit \longrightarrow \cunit$.

A similar argument holds for a killing context of the form $\forall x \tcontext{}$.
\end{proof}
% --- end distribute-proof.tex ---
}

%A special case of the above distributivity property is that any nullary killing context (a killing context with no holes) annihilates over parallel composition; hence, for nullary killing context $\tcontext{}$, $P \cpar \tcontext{} \longrightarrow P$.
Killing contexts also satisfy the following property that is necessary for handling the \textit{seq} operator, which interacts subtly with killing contexts.
%The case analysis considers carefully the restriction on the \textit{medial} and \textit{medial1} rules, by using the observation that $\left(P \cseq \cunit\right) \wwith \left(Q \cseq \cunit\right) \equiv \left(P \wwith Q\right) \cseq \cunit$ and $\new x \left(P \cseq \cunit\right) \equiv \new x P \cseq \cunit$; hence the \textit{medial} rules are not required in such cases.
\begin{lemma}~\label{lemma:medial}
Assume that $I$ is a finite subset of natural numbers, $P_i$ and $Q_i$ are formulae, for $i \in I$, and $\tcontext{}$ is a killing context. There exist killing contexts $\tcontextn{0}{}$ and $\tcontextn{1}{}$ and sets of natural numbers $J \subseteq I$ and $K \subseteq I$ such that the following derivation holds:
\[
\infer[.]{
\tcontext{ P_i \cseq Q_i \colon i \in I }
}{
\tcontextn{0}{ P_j \colon j \in J } \cseq
\tcontextn{1}{ Q_k \colon k \in K }
}
\]
\end{lemma}

%\begin{comment}
The following lemma checks that \textit{wen} quantifiers can propagate to the front of a killing context.
Similarly, to the proof of the lemma above, the proof is by induction on the structure of a killing context, applying the \textit{all name}, $\textit{new wen}$, $\textit{with name}$, $\textit{left name}$ or $\textit{right name}$ rule, as appropriate.
\begin{lemma}\label{lemma:commute}
Consider an $n$-ary killing context $\tcontext{}$ and formulae such that $\nfv{x}{P_i}$ and either $P_i = \wen x. Q_i$ or $P_i = Q_i$,
for $1 \leq i \leq n$. If for some $i$ such that $1 \leq i \leq n$, $P_i = \wen x. Q_i$, then we have derivation 
$\vcenter{\infer{\tcontext{ P_1, P_2, \hdots P_n}}{\wen x. \tcontext{ Q_1, Q_2, \hdots, Q_n }}}$.
\end{lemma}
%\fullproof{
%\input{killing-medial-proof}
%}


To handle certain cases in splitting the following definitions and property is helpful.
Assume $\vec{y}$ defines a possibly empty list of variables $y_1, y_2, \hdots, y_n$ and $\quantifier \vec{y}. P$ abbreviates $\quantifier y_1. \quantifier y_2. \hdots \quantifier y_n. P$. Let $\nfv{\vec{y}}{P}$ hold only if $\nfv{y}{P}$ for every $y \in \vec{y}$.
By induction over the length of $\vec{z}$ we can establish the following lemma, by repeatedly applying the \textit{close}, \textit{fresh} and \textit{extrude new} rules.
\begin{lemma}\label{lemma:close}
If $\vec{y} \subseteq \vec{z}$ and $\nfv{\vec{z}}{\wen \vec{y}. P}$, then we have derivations 
$\vcenter{\infer{\wen \vec{y}. P \cpar \new \vec{z}. Q}{\mathopen{\new \vec{z}.} \left( P \cpar Q \right)}}$ 
and $\vcenter{\infer{\new \vec{y}. P \cpar \wen \vec{z}. Q}{\mathopen{\new \vec{z}.} \left( P \cpar Q \right)}}$.
\end{lemma}
%\ross{This paragraph can form a lemma.}

% --- end lemmas.tex ---
\tohide{
% --- begin multisets.tex ---
\subsection{An Affine Measure for the Size of a Proof.}

As an induction measure in the splitting lemmas, we employ a multiset-based measure~\cite{Dershowitz1979} of the size of a proof.
An \textit{occurrence count} is defined in terms of a multiset of multisets. To give weight to nominals, a \textit{wen} and \textit{new} count is employed.
The measure of the size of a proof, Definition~\ref{definition:size}, is then given by the lexicographical order induced by the occurrence count, wen count and new count for the formula in the conclusion of a proof, and the derivation length of the proof itself.

%\begin{comment}
In the sub-system \textsf{BV}~\cite{Guglielmi2007}, the occurrence count is simply the number of atom and co-atom occurrences.
For the sub-system corresponding to \textsf{MALL} (multiplicative-additive linear logic)~\cite{Strassburger2003thesis}, i.e. without \textit{seq}, 
a multiset of atom occurrences such that $\occ{\left(P \wwith Q\right) \cpar R} = \occ{\left(P  \cpar R\right) \wwith \left(Q \cpar R\right)}$ is sufficient, to ensure that the \textit{external} rule does not increase the size of the measure. 
The reason why a multiset of multisets is employed for extensions of \textsf{MAV}~\cite{Horne2015} is to handle subtle interactions between the unit, \textit{seq} and \textit{with} operators.
In particular, by applying the structural rules for units, such that $ \context{ P \wwith Q }
 \equiv
 \context{ \left(P \cseq \cunit\right) \wwith \left(\cunit \cseq Q\right) }$
and the \textit{medial} rule, we obtain the following inference.
\[
\infer[\mbox{by the \textit{medial} rule}]{
 \context{ P \wwith Q }
}{
% \context{ \left(P \cseq \cunit\right) \wwith \left(\cunit \cseq Q\right) }
% \rightarrow
 \context{ \left(P \wwith \cunit\right) \cseq \left(\cunit \wwith Q\right) }
}
\]
In the above derivation, the units cannot in general be removed from the formula in the premise; hence extra care should be taken that these units do not increase the size of the formula. This observation leads us to the notion of multisets of multisets of natural numbers defined below.

\begin{definition}
%A multiset of natural numbers is a non-empty set of natural numbers where numbers may occur more than once.
We denote the standard multiset disjoint union operator as $\discup$, a multiset sum operator defined such that $M + N = \left\{ m + n \colon m \in M \mbox{ and } n \in N \right\}$.
We also define pointwise plus and pointwise union over multisets of multisets of natural numbers, where $\mathcal{M}$ and $\mathcal{N}$ are multisets of multisets.
$\mathcal{M} \mplus \mathcal{N} = \left\{ M + N, M \in \mathcal{M} \mbox{ and } N \in \mathcal{N} \right\}$
and
$\mathcal{M} \mcup \mathcal{N} = \left\{ M \discup N, M \in \mathcal{M} \mbox{ and } N \in \mathcal{N} \right\}$.
\end{definition}

We employ two distinct multiset orderings over multisets and over multisets of multisets.
\begin{definition}
For multisets of natural numbers $M$ and $N$, define a multiset ordering $M \leq N$ if and only if there exists an injective multiset function $f \colon M \rightarrow N$ such that, for all $m \in M$, $m \leq f(m)$. Strict multiset ordering $M < N$ is defined such that $M \leq N$ but $M \not= N$.
\end{definition}

\begin{definition}
Given two multisets of multisets of natural numbers $\mathcal{M}$ and $\mathcal{N}$, $\mathcal{M} \sqsubseteq \mathcal{N}$ holds if and only if $\mathcal{M}$ can be obtained from $\mathcal{N}$ by repeatedly removing a multiset $N$ from $\mathcal{N}$ and replacing $N$ with zero or more multisets $M_i$ such that $M_i < N$. $\mathcal{M} \sqsubset \mathcal{N}$ is defined when $\mathcal{M} \sqsubseteq \mathcal{N}$ but $\mathcal{M} \not= \mathcal{N}$.
\end{definition}

\begin{definition}
The occurrence count is the following function from formulae to multiset of multisets of natural numbers.
\begin{gather*}
\begin{array}{l}
\occ{\cunit} = \left\{\left\{ 0 \right\}\right\}
\qquad
%\occ{0} = \occ{\top} =
\occ{\alpha} = \occ{\co{\alpha}} = \left\{\left\{ 1 \right\}\right\}
\\[6pt]
\occ{P \wwith Q} = \occ{P \ooplus Q} = \occ{P} \sqcup \occ{Q}
\\[6pt]
\occ{P \cpar Q} = \occ{P} \boxplus \occ{Q}
\\[6pt]
\occ{ \forall x. P } = \occ{ \exists x. P } = \left\{\left\{ 0 \right\}\right\} \mcup \occ{ P }
\end{array}
\!\!\!\!\!\!\!\!\!\!
\begin{array}{r}
\occ{ \new x. P } = \occ{ \wen x. P } = 
\left\{
\begin{array}{lr}
\left\{\left\{ 0, 0 \right\}\right\} & \mbox{if $P \equiv \cunit$} \\
\occ{P} & \mbox{otherwise}
\end{array}
\right.
\\[6pt]
\occ{P \tensor Q} = \occ{P \cseq Q} = 
\left\{
\begin{array}{lr}
\occ{P} & \mbox{if $Q \equiv \cunit$} \\
\occ{Q} & \mbox{if $P \equiv \cunit$} \\
\occ{P} \discup \occ{Q} & \mbox{otherwise}
\end{array}
\right.
\end{array}
\end{gather*}
\end{definition}

\begin{definition}
The wen count is the following function from formulae to natural numbers.
\begin{gather*}
\wensize{ \wen x. P } = 1 + \wensize{ P }
\qquad
\wensize{ \exists x. P } = \wensize{ \forall x. P } = \wensize{ \new x. P } = \wensize{ P }
\qquad
\wensize\alpha = \wensize{\co\alpha} = \wensize\cunit = 1
\\[4pt]
\wensize{ P \cseq Q } = \wensize{ P \tensor Q } = \wensize{ P \cpar Q } = \wensize{ P }\wensize{ Q }
\qquad
\wensize{ P \ooplus Q } = \wensize{ P \wwith Q } = \wensize P + \wensize Q
\end{gather*}
\begin{comment}
\begin{gather*}
\wensize{ \wen x P } = \left\{ 1 \right\} \discup \wensize{ P }
\qquad
\wensize{ \exists x P } = \wensize{ \forall x P } = \wensize{ \new x P } = \wensize{ P }
\qquad
\wensize\alpha = \wensize{\co\alpha} = \wensize\cunit = \left\{ 0 \right\}
\\[4pt]
\wensize{ P \cseq Q } = \wensize{ P \tensor Q } = \wensize{ P \cpar Q } = \wensize{ P } + \wensize{ Q }
\qquad
\wensize{ P \ooplus Q } = \wensize{ P \wwith Q } = \wensize P \discup \wensize Q
\end{gather*}
\end{comment}
\end{definition}


\begin{definition}
The new count is the following function from formulae to natural numbers.
\begin{gather*}
\newsize{ \new x. P } = 1 + \newsize{ P }
\qquad
\newsize{ \exists x. P } = \newsize{ \forall x. P } = \newsize{ \wen x. P } = \newsize{ P }
\qquad
\newsize\alpha = \newsize{\co\alpha} = \newsize\cunit = 1
\\[4pt]
\newsize{ P \cpar Q } = \newsize{ P }\newsize{ Q }
\quad
\newsize{ P \ooplus Q } = \newsize{ P \wwith Q } = \newsize P + \newsize Q
\quad
\newsize{ P \cseq Q } = \newsize{ P \tensor Q } = \mmax{\newsize{P}, \newsize{Q}}
\end{gather*}
\begin{comment}
\begin{gather*}
\newsize{ \new x P } = \left\{ 1 \right\} + \newsize{ P }
\qquad
\newsize{ \exists x P } = \newsize{ \forall x P } = \newsize{ \wen x P } = \newsize{ P }
\qquad
\newsize\alpha = \newsize{\co\alpha} = \newsize\cunit = \left\{ 0 \right\}
\\[4pt]
\newsize{ P \cpar Q } = \newsize{ P } + \newsize{ Q }
\qquad\qquad
\newsize{ P \ooplus Q } = \newsize{ P \wwith Q } = \newsize P \discup \wensize Q
\\[4pt]
\newsize{ P \cseq Q } = \newsize{ P \tensor Q } = \mmax{\newsize{P}, \newsize{Q}}
\end{gather*}
\end{comment}
\end{definition}
%\ross{In BV only $\wensize{P \cseq Q} = max \wensize{P} \wensize{Q}$.}

\begin{definition}
The size of a formula $\size{P}$ is defined as the triple $\left(\occ{P}, \wensize{P}, \newsize{P} \right)$ lexicographically ordered by $\prec$.
$\phi \preceq \psi$ is defined such that $\phi \prec \psi$ or $\phi = \psi$ pointwise.
%defined such that $\left( \mathcal{M}, M, N \right) \prec \left( \mathcal{N}, K, L \right)$ if and only if $\mathcal{M} \mstrict \mathcal{N}$, or $\mathcal{M} = \mathcal{N}$ and $M < K$, or $\mathcal{M} = \mathcal{N}$ and $M = K$ and $N < L$.
\end{definition}


%\paragraph{The size of a proof.}
%The \textit{size of a proof} $\vdash P$ is defined using lexicographically ordered pairs consisting of the size of the predicate, $\size{P}$, and the number of rules applied in the proof of the predicate. %The pairs representing the size of a proof are ordered.
\begin{definition}\label{definition:size}
%The length of a proof $P \longrightarrow \cunit$ is the length, say $n$, of the derivation from $P$ to $\cunit$.
The size of a proof of $P$ with derivation of length $n$ is given by the tuple of the form $\left(\size{P}, n\right)$, subject to lexicographical ordering.
%Hence $\left(\size{P}, m\right) \prec \left(\size{Q}, n\right)$ if and only if $\size{P} \mstrict \size{Q}$ or $\size{P} = \size{Q}$ and $m < n$.
\end{definition}
% such that $\left(u, j\right) \prec \left(v, k\right)$ only if $u < v$ or $u = v$ and $j < k$.

%
%\input{properties}
% --- end multisets.tex ---
% --- begin bound.tex ---
\begin{lemma}\label{lemma:size-sub}
For any formula $P$ and term $t$, $\size{ P } = \size{ P\sub{x}{t} }$.
\end{lemma}

\begin{lemma}\label{lemma:multiset-congruence}
If $P \equiv Q$ then $\size{P} = \size{Q}$.
\end{lemma}
\fullproof{
\begin{proof}
For commutativity the following arguments hold for \textit{par} and \textit{times}.
\begin{gather*}
\begin{array}{rl}
\occ{P \cpar Q} = \occ{P} \mplus \occ{Q} =&\!\!\left\{ M + N \colon M \in \occ{P}, N \in \occ{Q} \right\}
\\
=&\!\!\left\{ N + M \colon M \in \occ{P}, N \in \occ{Q} \right\} = \occ{Q} \mplus \occ{P} = \occ{ Q \cpar P }
\end{array}
\\
\occ{P \tensor Q} = \occ{P} \discup \occ{Q} = \occ{Q} \discup \occ{P} = \occ{Q \tensor P}
\\
\wensize{ P \tensor Q } = \wensize{ P \cpar Q } = \wensize{ P }\wensize{ Q } = \wensize{ Q }\wensize{ P } = \wensize{ Q \cpar P } = \wensize{ Q \tensor P }
\\
\newsize{ P \cpar Q } = \newsize{ P }\newsize{ Q } = \newsize{ Q }\newsize{ P } = \newsize{ Q \cpar P }
\\
\newsize{ P \tensor Q } = \mmax{\newsize{ P } , \newsize{ Q }} = \mmax{\newsize{ Q }, \newsize{ P }} = \newsize{ Q \tensor P }
\qquad
\end{gather*}


The identity rules hold by the following reasoning, where $\odot \in \left\{ \tensor, \cseq \right\}$, since $\newsize{P} \geq 1$ always $\newsize{P \odot \cunit} = \mmax{\newsize{ P}, 1} = \newsize{P}$ and the following hold. 
\begin{gather*}
\occ{P \cpar \cunit} = \occ{P} \mplus \left\{\left\{0\right\}\right\} = \occ{P}
\quad
\occ{\cunit \odot P} = \occ{P}
\quad
\wensize{P \odot \cunit} 
= \wensize{P}
= \wensize{P \cpar \cunit}
\quad
\newsize{P \cpar \cunit} = \newsize{ P }
\end{gather*}

Associativity properties hold by extending associativity of multisets to multisets of multisets.
\[
\begin{array}{rl}
\size{\left(P \cpar Q\right) \cpar R}
 =
 \left(\size{P} \mplus \size{Q}\right) \mplus \size{R}
 =&
 \left\{ \left(M + N\right) + K \colon M \in \size{P}, N \in \size{Q}, K \in \size{R} \right\}
 \\
 =&
 \left\{ M + \left(N + K\right) \colon M \in \size{P}, N \in \size{Q}, K \in \size{R} \right\}
 =  \size{P} \mplus \left(\size{Q} \mplus \size{R}\right) 
 = \size{P \cpar \left(Q \cpar R\right)}
\end{array}
\]
If any one of $P \equiv \cunit$, $Q \equiv \cunit$ or $R \equiv \cunit$ hold, then $\size{\left(P \cseq Q\right) \cseq R} = \size{P \cseq \left(Q \cseq R\right)}$ by definition. If $P \not\equiv \cunit$ and $Q \not\equiv \cunit$ and $R \not\equiv \cunit$, then the following equalities hold:
$
\size{\left(P \cseq Q\right) \cseq R} = \left(\size{P} \discup \size{Q}\right) \discup \size{R} = \size{P} \discup \left(\size{Q} \discup \size{R}\right) = \size{P \cseq \left(Q \cseq R\right)}
$.
The same associativity argument works for the \textit{times} operator. Furthermore, the following wen count and new count preserve associativity as follows, where $\odot \in \left\{ \tensor, \cseq \right\}$.
\[
\wensize{ \left( P \cpar Q\right) \cpar R } =
\wensize{ \left( P \odot Q\right) \odot R } = \left(\wensize{ P }\wensize{ Q }\right)\wensize{ R } = \wensize{ P }\left(\wensize{ Q }\wensize{R}\right) = \wensize{ P \odot \left( Q \odot R\right) }
= \wensize{ P \cpar \left(Q \cpar R\right) }
\]
\[
\newsize{ \left( P \cpar Q\right) \cpar R } = \left(\newsize{ P }\newsize{ Q }\right)\newsize{ R } = \newsize{ P }\left(\newsize{ Q }\newsize{R}\right) = \newsize{ P \cpar \left( Q \cpar R\right) }
\]
\[
\newsize{ \left( P \odot Q\right) \odot R } = \mmax{\mmax{\newsize{ P },\newsize{ Q }},\newsize{ R }} = \mmax{\newsize{ P },\mmax{\newsize{ Q },\newsize{R}}} = \newsize{ P \odot \left( Q \odot R\right) }
\]
The lemma then follows by induction over the structure of the derivation of $P \equiv Q$.
\end{proof}
}

The following lemma we will appeal to regularly in the splitting proofs in subsequent sections to bound the size of a derivation.
\begin{lemma}[affine]\label{lemma:bound}
Any derivation $\vcenter{\infer{Q}{P}}$, is bound such that $\size{P} \preceq \size{Q}$.
\end{lemma}
%\begin{proof}
%The proof proceeds by checking that each rule preserves the bound on the size of the formula, from which the result follows by induction on the length of a derivation.
%
%
%Consider the case of the \textit{close} rule. $\occ{\new x. P \cpar \wen x. Q} = \occ{ P } \mplus \occ{ Q } = \occ{\mathopen{\new x.} \left( P \cpar Q \right)}$,
%since $P \not\equiv \cunit$ and $Q \not\equiv \cunit$,
% and 
%%for the wen count 
%%strictly decreases as follows.
%$
%\wensize{\new x. P \cpar \wen x. Q}
%=
%\wensize{P} + \left( 1  +  \wensize{Q}\right)
%>
%\wensize{P} + \wensize{Q}
%=
%\wensize{\mathopen{\new x.} \left(P \cpar Q\right)}
%$.
%
%Consider the case of the \textit{fresh} rule.
%For the occurrence count, $\occ{ \wen x. P } = \occ{ \new x. P }$ and the wen count strictly decreases as follows:
%$
%\wensize{ \wen x. P }
%=
% 1 + \wensize{ P }
%>
%\wensize{ P }
%=
%\wensize{ \new x. P }
%$.
%
%
%Consider the case of the \textit{extrude new} rule, where $Q \not\equiv \cunit$.
%If $P \equiv \cunit$, then the occurence count is such that $\occ{ \new x. P \cpar Q } = \left\{\left\{ 0,0 \right\}\right\} \mplus \occ{ Q } > \occ{ Q } = \occ{ \mathopen{\new x.} \left( P \cpar Q \right)}$. If however $P \not\equiv \cunit$, then $\occ{ \new x. P \cpar Q } = \occ{ P } \mplus \occ{ Q } = \occ{ \mathopen{\new x.}\left( P \cpar Q \right) }$. Furthermore, for the new count, the following inequality holds:
%$
%\newsize{ \new x. P \cpar Q } = \left(1 + \newsize{P}\right)\wensize{Q} \geq 1 + \newsize{P}\newsize{Q} = \newsize{ \mathopen{\new x.} \left( P \cpar Q \right) }
%$.
%%The new count is also preserved by the same argument.
%
%
%
%Consider the case of the \textit{external} rule, where $R \not\equiv \cunit$.
%For the occurrence count, by distributivity of $\mcup$ over $\mplus$, 
%the following multiset equality holds:
%\[
%%\begin{array}{rl}
%\occ{\left( P \wwith Q \right) \cpar R }
%= \left( \occ{P} \mcup \occ{Q} \right) \mplus \occ{R}
%= \left( \occ{P}  \mplus \occ{R}\right) \mcup \left(\occ{Q}  \mplus \occ{R} \right)
%=
%\occ{\left( P \cpar R\right) \wwith \left(Q  \cpar R \right)}
%%\end{array}
%\]
%For the wen count 
%%the following equality holds:
%$
%\wensize{ \left( P \wwith Q \right) \cpar R }
%=
%\left(\wensize{ P } + \wensize{Q}\right) \wensize{R}
%=
%\wensize{ P }\wensize{R} + \wensize{Q}\wensize{R}
%=
%\wensize{ \left( P \cpar R\right) \wwith \left( Q \cpar R \right)  }
%$;
%and similarly for the new count.
%
%
%
%
%Consider the case of the \textit{suspend} rule, where $P \not\equiv \cunit$ and $Q \not\equiv \cunit$.
%For the occurrence count, 
%% and \textit{with}.
%$\occ{ \wen x. P \cseq \wen x. Q } = \occ{ P } \discup \occ{ Q } = \occ{ \mathopen{\wen x.} \left( P \cseq Q \right) }$
%and
%$\occ{ \wen x. P \cpar \wen x. Q } = \occ{ P } \mplus \occ{ Q } = \occ{ \mathopen{\wen x.} \left( P \cpar Q \right) }$
%for \textit{par} and \textit{seq} respectively.
%%\item $\occ{ \wen x. P \wwith \wen x. Q } = \occ{ P } \mcup \occ{ Q } = \occ{ \mathopen{\wen x.} \left( P \wwith Q \right) }$
%%\end{itemize}
%For the wen count for either operator, $\odot \in \left\{\cpar, \cseq \right\}$, the following strict inequality holds, noting $\wensize{P} \geq 1$ for any formula:
%\[
%\wensize{ \wen x. P \odot \wen x. Q }
%=
%\left( 1 + \wensize{ P }\right)\left(1 + \wensize{ Q }\right)
%=
%\wensize{ P } + \wensize{ P }\wensize{ Q } + \wensize{ Q }
%>
%1 + \wensize{P}\wensize{Q}
%=
%\wensize{ \mathopen{\wen x.} \left( P \cseq Q \right) }
%\]
%
%
%Consider the case of the \textit{left wen} rules, where $\nfv{x}{Q}$ and $Q \not\equiv \cunit$.
%For the occurrence count, there are four cases covering the operators \textit{seq} and \textit{par}. % and \textit{with}.
%\begin{itemize}
%\item
%If $P \equiv \cunit$ then, for \textit{seq}:
%$
% \occ{ \mathopen{\wen x.} \left( P \cseq Q \right) } = \occ{ Q }  \mstrict  \left\{\left\{ 0,0 \right\}\right\} \discup \occ{ Q } 
% = \occ{ \wen x. P \cseq Q }
%$.
%
%\item
%If $P \not\equiv \cunit$ then, for \textit{seq}:
%$
%\occ{ \wen x. P \cseq Q } = \occ{ P } \discup \occ{ Q } = \occ{ \mathopen{\wen x.} \left( P \cseq Q \right) }
%$.
%
%\item
%If $P \equiv \cunit$ then for \textit{par}:
%$
%\occ{ \mathopen{\wen x.} \left( P \cpar Q \right) }
%= \occ{Q}
%\mstrict
%\left\{\left\{ 0,0 \right\}\right\} \mplus \occ{ Q } 
%=
%\occ{ \wen x. P \cpar Q }
%$.
%
%\item
%If $P \not\equiv \cunit$ then for \textit{par}:
%$
%\occ{ \wen x. P \cpar Q } = \occ{ P } \mplus \occ{ Q } = \occ{ \mathopen{\wen x.} \left( P \cpar Q \right) }
%$.
%
%%\item %For the \textit{with} operator and $P \equiv \cunit$ the following holds:
%%$
%%\occ{ \wen x. P \wwith Q } = \left\{\left\{ 0,0 \right\}\right\} \mcup \occ{ Q }
%% = \left\{\left\{ 0,0 \right\}\right\} \mcup \occ{ Q }
%% = \occ{ \mathopen{\wen x.} \left( P \wwith Q \right) }
%%$
%
%%\item
%%For the \textit{with} operator and $P \not\equiv \cunit$ the following holds:
%%$
%%\occ{ \wen x. P \wwith Q } = \occ{ P } \mcup \occ{ Q }
%%  = \occ{ \mathopen{\wen x.} \left( P \wwith Q \right) }
%%$.
%\end{itemize}
%For the wen count 
%%multiset inequality
%$
%\begin{array}{rl}
%\wensize{ \wen x. P \odot Q }
%=
%\left(1 + \wensize{ P }\right)\wensize{ Q }
%=
%\wensize{Q} + \wensize{P}\wensize{ Q }
%\geq
%1 + \wensize{ P }\wensize{ Q }
%=
%\wensize{ \mathopen{\wen x.} \left( P \odot Q \right) }
%\end{array}
%$ holds, for $\odot \in \left\{\cpar, \cseq \right\}$.
%Also, for the new count 
%%the following equalites hold:
%$\newsize{ \wen x. P \cseq Q } = \mmax{\newsize{ P }, \newsize{ Q}} = \newsize{ \mathopen{\wen x.} \left( P \cseq Q \right) }$ and
%$\newsize{ \wen x. P \cpar Q } = \newsize{ P }\newsize{ Q } = \newsize{ \mathopen{\wen x.} \left( P \cpar Q \right) }$.
%The case \textit{right wen} follows a symmetric argument.
%
%
%
%Consider the case for the \textit{extrude} rule, where $Q \not\equiv \cunit$.
%$\occ{ \mathopen{\forall x.} \left( P \cpar Q \right) } \mstrict \occ{ \forall x. P \cpar Q }$ by the following:
%$
%%\begin{array}{rl}
% \left\{\left\{0\right\}\right\} \mcup \left( \occ{ P } \mplus \occ{ Q } \right)
%%\\
%\mstrict
% \left(\left\{\left\{0\right\}\right\} \mplus \occ{ Q } \right) \mcup \left( \occ{ P } \mplus \occ{ Q } \right)
%%\\
%=
% \left( \left\{\left\{0\right\}\right\} \mcup \occ{ P } \right) \mplus \occ{ Q }
%%\\=&
%%\end{array}
%$.
%
%
%
%Consider the case for the \textit{medial1} rule, where $P \not\equiv \cunit$ and $Q \not\equiv\cunit$.
%By distributivity of $\discup$ over $\mcup$,
%%the following multiset equality holds:
%$
%%\begin{array}{rl}
%\occ{ \mathopen{\forall x.} \left( P \cseq Q \right) }
%=
%\left\{\left\{ 0 \right\}\right\} \mcup \left( \occ{ P } \discup \occ{ Q } \right)
%=
%\left( \left\{\left\{ 0 \right\}\right\} \mcup \occ{ P } \right) \discup \left( \left\{\left\{ 0 \right\}\right\} \mcup \occ{ Q } \right)
%=
%\occ{ \forall x. P \cseq \forall x. Q }
%%\end{array}
%$.
%Also $\wensize{ \mathopen{\forall x.} \left( P \cseq Q \right) } = \wensize{ \forall x. P \cseq \forall x. Q }$
%and $\newsize{ \mathopen{\forall x.} \left( P \cseq Q \right) } = \newsize{ \forall x. P \cseq \forall x. Q }$.
%
%
%
%For the \textit{select} rule,
%$
% \occ{ \exists x. P }
%   = \left\{\left\{ 0 \right\}\right\} \mcup \occ{ P }
%   \mstrict \occ{ P }
%   = \occ{ P \sub{x}{t} }
%$,
%by Lemma~\ref{lemma:size-sub}.
%
%
%
%Consider the case for the \textit{switch} rule, where $P \not\equiv \cunit$ and $R \not\equiv \cunit$.
%If $Q \not\equiv \cunit$, then, since $R \not\equiv \cunit$ we have $\left\{\left\{ 0 \right\}\right\} \mstrict \occ{R}$ and hence $\occ{P} = \occ{P} \mplus \left\{\left\{ 0 \right\}\right\} \mstrict \occ{P} \mplus \occ{R}$; and therefore the following holds since $\discup$ distributes over $\mplus$.
%\[
%\begin{array}{rl}
%\occ{P \tensor \left(Q \cpar R\right)}
%=& \occ{P} \discup \left( \occ{Q} \mplus \occ{R} \right)
%\\
%\mstrict& \left(\occ{P} \mplus \occ{R}\right) \discup \left( \occ{Q} \mplus \occ{R} \right)
%\\
%=& \left(\occ{P} \discup \occ{Q}\right) \mplus \occ{R}
%= \occ{\left(P \tensor Q\right) \cpar R}
%\end{array}
%\]
%If $Q \equiv \cunit$ then, since $\left\{\left\{0\right\}\right\} \mstrict \occ{P}$ and $\left\{\left\{0\right\}\right\} \mstrict \occ{R}$, the following hold.
%\[
%\occ{P \tensor \left( \cunit \cpar R\right)}
%= \occ{P} \discup \occ{R}
%\mstrict \occ{P} \mplus \occ{R}
%= \occ{\left(P \tensor \cunit\right) \cpar R}
%\]
%
%
%
%Consider the case of the \textit{sequence} rule, where $P \not\equiv \cunit$ and $S \not\equiv \cunit$.
%If $Q \not\equiv \cunit$ and $R \not\equiv \cunit$, then the following holds since $\discup$ distributes over $\mplus$.
%\[
%\!\!
%\begin{array}{rl}
%\occ{\left(P \cpar R\right) \cseq \left(Q \cpar S\right)}
%=&\!\!\left( \occ{P} \mplus \occ{R} \right) \discup \left( \occ{Q} \mplus \occ{S} \right)
%\\
%\mstrict&\!\!\left( \occ{P} \mplus \occ{R} \right) \discup \left( \occ{Q} \mplus \occ{S} \right) \discup \left( \occ{P} \mplus \occ{S} \right) \discup \left( \occ{Q} \mplus \occ{R} \right)
%\\
%=&\!\!\left( \occ{P} \discup \occ{Q} \right) \mplus \left( \occ{R} \discup \occ{S} \right)
%= \occ{\left(P \cseq Q\right) \cpar \left(R \cpar S\right)}
%\end{array}
%\]
%If $Q \equiv \cunit$ and $R \not\equiv \cunit$, then, since $\left\{\left\{0\right\}\right\} \mstrict \occ{R}$, and hence $\occ{S} = \occ{S} \mplus \left\{\left\{0\right\}\right\} \mstrict \occ{S} \mplus \occ{R}$, therefore since $\discup$ distributes over $\mplus$.
%\[
%\begin{array}{rl}
%\!\!
%\occ{\left(P \cpar R\right) \cseq \left(\cunit \cpar S\right)}
%=\left( \occ{P} \mplus \occ{R} \right) \discup \occ{S}
%\mstrict&\!\!\left( \occ{P} \mplus \occ{R} \right) \discup \left(\occ{P} \mplus \occ{S}\right)
%\\
%=&\!\! \occ{P} \mplus \left( \occ{R} \discup \occ{S} \right)
%=
%\occ{\left(P \cseq \cunit \right) \cpar \left(R \cpar S\right)}
%\end{array}
%\]
%A symmetric argument holds when $Q \not\equiv \cunit$ and $R \equiv \cunit$.
%
%If $Q \equiv \cunit$ and $R \equiv \cunit$, then $\left\{\left\{0\right\}\right\} \mstrict \occ{P}$ and $\left\{\left\{0\right\}\right\} \mstrict \occ{S}$; hence the following strict inequality holds:
%$
%%\begin{array}{rl}
%\occ{\left(P \cpar \cunit\right) \cseq \left(\cunit \cpar S\right)}
%= \occ{P} \discup \occ{S}
%\mstrict \occ{P} \mplus \occ{S}
%= \occ{\left(P \cseq \cunit \right) \cpar \left(\cunit \cseq S\right)}
%%\end{array}
%$.
%
%
%
%
%Consider the case of the \textit{medial new} rule where $P \not\equiv \cunit$ and $Q \not\equiv \cunit$. For the occurrence count the equality $\occ{ \mathopen{\new x.} \left( P \cseq Q \right) } = \occ{ P } \mplus \occ{ Q } = \occ{ \new x. P \cseq \new x.Q }$ holds.
%For the wen count,
%$\wensize{ \mathopen{\new x.} \left( P \cseq Q \right) } = 
%\wensize{ P } \wensize{ Q } = 
%\wensize{ \new x. P \cseq \new x. Q }$.
%For the new count the following equality holds:
%$
%\newsize{ \mathopen{\new x.} \left( P \cseq Q \right) }
%=
%1 + \mmax{\newsize{ P }, \newsize{ Q }}
%=
%\mmax{ 1 + \newsize{ P }, 1 + \newsize{ Q }}
%=
%\newsize{ \new x. P \cseq \new x. Q }
%$.
%
%
%
%Consider the case for the \textit{medial} rule, where either $P \not\equiv \cunit$ or $R \not\equiv \cunit$ and also either $Q \not\equiv \cunit$ or $S \not\equiv \cunit$.
%When all of $P$, $Q$, $R$ and $S$ are not equivalent to the unit, we have the following.
%\[
%\begin{array}{rl}
%\!\!\!
%\occ{
% \left(P \wwith R\right) \cseq \left(Q \wwith S\right)
%}
%=&\!\!\!
%\left( \occ{P} \mcup \occ{R} \right) \discup \left( \occ{Q} \mcup \occ{S} \right)
%\\
%\mstrict
%&\!\!\!
%\left( \occ{P} \mcup \occ{R} \right) \discup \left( \occ{Q} \mcup \occ{S} \right) \discup \left( \occ{P} \mcup \occ{S} \right) \discup \left( \occ{Q} \mcup \occ{R} \right)
%\\
%=
%&\!\!\!
%\left( \occ{P} \discup \occ{Q} \right) \mcup \left( \occ{R} \discup \occ{S} \right)
%=
%\occ{
% \left(P \cseq Q\right) \wwith \left(R \cseq S\right)
%}
%\end{array}
%\]
%For when exactly one of $P$, $Q$, $R$ and $S$ is equivalent to the unit, all cases are symmetric. Without loss of generality suppose that $S \equiv \cunit$ (and possibly also $Q \equiv \cunit$, but $R \not\equiv \cunit$). By distributivity of $\discup$ over $\mcup$ the following holds.
%\[
%\begin{array}{rl}
%\occ{
% \left(P \wwith R\right) \cseq \left(Q \wwith \cunit\right)
%}
%=&
%\left( \occ{P} \mcup \occ{R} \right) \discup \left( \occ{Q} \mcup \left\{\left\{ 0 \right\}\right\} \right)
%\\
%\mstrict&
%\left( \occ{P} \mcup \occ{R} \right) \discup \left( \occ{Q} \mcup \occ{R} \right)
%\\
%=&
%\left( \occ{P} \discup \occ{Q} \right) \mcup \occ{R}
%=
%\occ{
% \left(P \cseq Q\right) \wwith \left(R \cseq \cunit \right)
%}
%\end{array}
%\]
%There is one more form of case to consider for the medial: either $P \not\equiv \cunit$, $Q \equiv \cunit$, $R \equiv \cunit$ and $S \not\equiv \cunit$; or $P \equiv \cunit$, $Q \not\equiv \cunit$, $R \not\equiv \cunit$ and $S \equiv \cunit$. We consider only the former case. The later case, can be treated symmetrically.
%Since $P \not\equiv \cunit$ and $S \not\equiv \cunit$, $\left\{\left\{ 0 \right\}\right\} \mstrict \occ{P}$ and $\left\{\left\{ 0 \right\}\right\} \mstrict \occ{S}$. Therefore, $\occ{P} \mcup \left\{\left\{ 0 \right\}\right\} \mstrict \occ{P} \mcup \occ{S}$ and $\occ{Q} \mcup \left\{\left\{ 0 \right\}\right\} \mstrict \occ{P} \mcup \occ{S}$. Hence, we have established that $\left(\occ{P} \mcup \left\{\left\{ 0 \right\}\right\}\right) \discup \left( \occ{Q} \mcup \left\{\left\{ 0 \right\}\right\} \right) \mstrict \occ{P} \mcup \occ{S}$.
%%
%Note that the restriction on the \textit{medial} rule, either $P \not\equiv \cunit$ or $R \not\equiv \cunit$ and also either $Q \not\equiv \cunit$ or $S \not\equiv \cunit$, excludes any further cases. Hence we have established that $
%\occ{\left(P \wwith R\right) \cseq \left(Q \wwith S\right)}
%\mstrict
%\occ{\left(P \cseq Q\right) \wwith \left(R \cseq S\right)}$. %Wen new counts need not be considered due to lexicographical ordering.
%%\ross{This is why we could simplify the measures!}
%
%For the \textit{with name} rule $\occ{\quantifier x.P \wwith \quantifier x.Q} = \occ{P} \mcup \occ{Q} = \occ{\mathopen{\quantifier x.}\left(P \wwith Q\right)}$,
%where $\quantifier \in \left\{ \new, \wen \right\}$. For the new count $\newsize{\new x.P \wwith \new x.Q} = 2 + \newsize{P} + \newsize{Q} > 1 + \newsize{P} + \newsize{Q} = \newsize{\mathopen{\new x.}\left(P \wwith Q\right)}$
%and $\newsize{\wen x.P \wwith \wen x.Q} = \newsize{\mathopen{\wen x.}\left(P \wwith Q\right)}$.
%Similarly, $\wensize{\wen x.P \wwith \wen x.Q} > \wensize{\mathopen{\wen x.}\left(P \wwith Q\right)}$.
%For \textit{left name}, \textit{right name} and \textit{all name}, the size of formulae are invariant.
%
%The cases for the rules \textit{tidy}, \textit{tidy name}, \textit{left}, \textit{right}, \textit{atomic interact} are established by the following inequalities:
%$\occ{ \cunit } \mstrict \occ{ \cunit \wwith \cunit }$,
%$\occ{ \cunit } \mstrict \occ{ \new x. \cunit }$,
%$\occ{ \cunit } \mstrict \occ{ \co{a} \cpar a }$,
%$\occ{ P } \mstrict \occ{ P \ooplus Q }$
%and
%$\occ{ Q } \mstrict \occ{ P \ooplus Q }$.
%
%Hence the lemma holds by induction on the length of the derivation.
%\end{proof}
%





% --- end bound.tex ---
}{}
% --- begin splitting.tex ---
\subsection{The splitting technique for simulating sequent-like rules}

The technique called splitting~\cite{Guglielmi2007,Guglielmi2011} generalises the application of rules in the sequent calculus.
In the sequent calculus, any root connective in a sequent can be selected and some rule for that connective can be applied.
For example, consider the following rules in linear logic forming part of a proof in the sequent calculus, where $\nfv{x}{P,Q,U,V,W}$.
\[
\begin{prooftree}
\begin{prooftree}
   \vdash P, U
   \quad
   \vdash Q, R
\justifies
   \vdash P \tensor Q, R, U
\end{prooftree}
  \quad
\begin{prooftree}
\begin{prooftree}
    \vdash P, R, V
    \quad
    \vdash Q, W
\justifies
    \vdash P \tensor Q, R, V, W
\end{prooftree}
\justifies
   \vdash P \tensor Q, R, V \cpar W
\end{prooftree}
\justifies
\begin{prooftree}
\vdash P \tensor Q, R, U\wwith \left(V \cpar W\right)
\justifies
 \vdash P \tensor Q, \forall x.R, U\wwith \left(V \cpar W\right)
\end{prooftree}
\end{prooftree}
\]
In the setting of the calculus of structures, the sequent at the conclusion of the above proof corresponds to a \textit{shallow context} of the form $\left\{\ \cdot\ \right\} \cpar \forall x.R \cpar \left(U\wwith \left(V \cpar W\right)\right)$
where the \textit{times} operator at the root of $P \tensor Q$ is a \textit{principal formula} that is plugged into the shallow context.
%The principal connective is like a formula at the root of a sequent that we would like to apply a rule to.
%In the setting of the calculus of structures, a sequent corresponds to a \textit{shallow context} of the form $\left\{\ \cdot\ \right\} \cpar Q$.
Splitting proves that there is always a derivation reorganising a shallow context 
%a distinguished \textit{principal predicate} $P$ in a shallow context $\left\{\ P\ \right\} \cpar Q$ can
into a form such that a rule for the root connective of the principal formula may be applied.
In the above example, this would correspond to the following derivation over contexts:
\[
\infer[\mbox{by the \textit{extrude1} rule}]{
\left\{\ \cdot\ \right\} \cpar \forall x.R \cpar \left(U\wwith \left(V \cpar W\right)\right)
}{
\infer[\mbox{by the \textit{external} rule}]{
\left\{\ \cdot\ \right\} \cpar \mathopen{\forall x.}\left(R \cpar \left(U\wwith \left(V \cpar W\right)\right)\right)
}{
\left\{\ \cdot\ \right\} \cpar \mathopen{\forall x.}\left(\left(R \cpar U\right) \wwith \left(R \cpar V \cpar W\right)\right)
}}
\]
By plugging in the principal formula, $P \tensor Q$, into the hole in the premise of the above derivation and applying distributivity properties of a killing context (Lemma~\ref{lemma:kill-distribute}),
the \textit{switch} rule involving the principal connective can be applied as follows.
%this leads to the following 
\[
\infer[\mbox{by Lemma~\ref{lemma:kill-distribute}}]{
\left(P \tensor Q\right) \cpar \mathopen{\forall x.}\left(\left(R \cpar U\right) \wwith \left(R \cpar V \cpar W\right)\right)
}{
\infer[\mbox{by the \textit{switch} rule}]{
\mathopen{\forall x.}\left(\left(\left(P \tensor Q\right) \cpar R \cpar U\right) \wwith \left(\left(P \tensor Q\right) \cpar R \cpar V \cpar W\right)\right)
}{
\mathopen{\forall x.}\left(\left(\left(P \cpar U\right)\tensor \left(Q \cpar R\right)\right) \wwith \left(\left(P \cpar R \cpar V\right)\tensor \left(Q \cpar W\right)\right)\right)
}}
\]
Notice that the final formula above holds when all of the following hold: $\vdash P \cpar U$, $\vdash Q \cpar R$, $\vdash P \cpar R \cpar V$ and $\vdash Q \cpar W$.
Notice that these correspond to the leaves of the example sequent above.

Splitting is sufficiently general that the technique can be applied to operators such as \textit{seq} that have no sequent calculus presentation~\cite{Tiu2006}.
The technique also extends to the pair of nominals \textit{new} and \textit{wen}, for which a sequent calculus presentation is an open problem.

The operators \textit{times}, \textit{seq}, \textit{new} and \textit{wen} are treated together in Lemma~\ref{lemma:split-times}.
These operators give rise to \textit{commutative cases}, where rules for these operators can permute with any principal formula, swapping the order of rules in a proof.
%The statement of the ``core'' splitting lemma is stronger than other splitting lemmas, since it also requires that the size of proofs are bound.
\textit{Principal cases} are where the root connective of the principal formula is directly involved in the bottommost rule of a proof.
As with \textsf{MAV}~\cite{Horne2015}, the \textit{principal cases} for \textit{seq} are challenging, demanding Lemma~\ref{lemma:medial}.
The principal case induced by \textit{medial new} demands Lemma~\ref{lemma:commute}. The cases where two nominal quantifiers commute are also interesting, particularly where the case arrises due to \textit{equivariance}.
%\tohide{}{
%~\cite{Dershowitz1979,Strassburger2003thesis} %adapted from~\cite{Horne2015}
%to give weight to quantifiers.
%Details of the multiset-based measure appear in the supporting technical report~\cite{Horne2016}.
%\footnote{For ICALP review purposes only: details of the induction measure are presented in the appendix.}
%Multisets of multisets handle subtle interactions between the unit, \textit{seq} and \textit{with} operators, explained in detail for \textsf{MAV}~\cite{Horne2015}.}
%All cases employ Lemma~\ref{lemma:bound} to ensure termination of the induction.
\begin{lemma}[Core Splitting]
\label{lemma:split-times}
The following statements hold.
\begin{enumerate}
%\item If $\vdash 0 \cpar U$ holds, then there exists nullary killing context $\tcontext{}$ such that $U \longrightarrow \tcontext{}$.
\item If $\vdash \left( P \tensor  Q \right) \cpar R$, then there exist formulae $V_i$ and $W_i$ such that $\vdash P \cpar V_i$ and $\vdash Q \cpar W_i$, where $1 \leq i \leq n$, and $n$-ary killing context $\tcontext{}$ such that $\vcenter{\infer{R}{\tcontext{V_1 \cpar W_1, V_2 \cpar W_2, \hdots, V_n \cpar W_n}}}$
and if $\tcontext{}$ binds $x$ then $\nfv{x}{\left(P \tensor Q\right)}$.

\item If $\vdash \left( P \cseq Q \right) \cpar R$, then there exist formulae $V_i$ and $W_i$ such that $\vdash P \cpar V_i$ and $\vdash Q \cpar W_i$, where $1 \leq i \leq n$, and $n$-ary killing context $\tcontext{}$ such that $\vcenter{\infer{R}{\tcontext{V_1 \andthen W_1, V_2 \andthen W_2, \hdots, V_n \andthen W_n}}}$
and if $\tcontext{}$ binds $x$ then $\nfv{x}{\left(P \cseq Q\right)}$.

\item If $\vdash \new x. P \cpar Q$, then there exist formulae $V$ and $W$
where $\nfv{x}{V}$ and $\vdash P \cpar W$ and either $V = W$ or $V = \wen x .W$,
such that there is a derivation $\vcenter{\infer{Q}{V}}$.

\item If $\vdash \wen x. P \cpar Q$, then there exist formulae $V$ and $W$ where $\nfv{x}{V}$ and $\vdash P \cpar W$ and either $V = W$ or $V = \new x. W$, such that there is a derivation $\vcenter{\infer{Q}{V}}$.

\item If $\vdash \left( P \wwith Q \right) \cpar R$, then $\vdash P \cpar R$ and $\vdash Q \cpar R$.
\end{enumerate}
%
Furthermore, for all $1 \leq i \leq n$, in the first two cases the size of the proofs
%\footnote{For the multiset-based measure of the size of a proof see technical report: \url{http://arxiv.org/abs/1602.06043}.}
of ${P \cpar V_i}$ and ${Q \cpar W_i}$ are strictly bounded above by the size of the proofs of ${ \left( P \tensor Q \right) \cpar R}$ and ${ \left( P \cseq Q \right) \cpar R }$.
In the third and fourth cases, the size of the proof ${P \cpar W}$ is strictly bounded above by the size of the proofs of $\new x .P \cpar Q$ and $\wen x. P \cpar Q$.
The size of a proof is measured according to Definition~\ref{definition:size}.
%For the fifth case bounds are immediate.
%\ross{Tighten size of proof statement.}
\end{lemma}
%\ross{Check freshness conditions on trivial contexts. N.B.\ I can't see possible problems.}
\tohide{
% --- begin split-times.tex ---
\begin{proof}
The proof proceeds by induction on the size of the proof, as in Defn.~\ref{definition:size}. 
In each of the following base cases, the conditions for splitting are immediately satisfied.
For the base case for the \textit{tidy name} rule, the bottommost rule of a proof is of the form 
$\vcenter{\infer[]{\new x.\new \vec{y}. \cunit \cpar P}{\new \vec{y}.\cunit \cpar P}}$, where $\nfv{\vec{y}}{P}$.
For the base case for the \textit{tidy} rule, the bottommost rule is of the form 
$\vcenter{\infer[]{\left(\cunit \wwith \cunit\right) \cpar P}{\cunit \cpar P}}$, such that $\vdash \cunit \cpar P$.
For the base case for $\textit{times}$ and $\textit{seq}$, $\vdash \left(\cunit \tensor \cunit\right) \cpar \cunit$ and $\vdash \left(\cunit \cseq \cunit\right) \cpar \cunit$ hold.

\begin{enumerate}[label=\textbf{\Alph*},ref=\Alph*,leftmargin=*]
%\textbf{(A) Principal cases for wen.}
\item \textbf{Principal cases for wen.}
There are principal cases for \textit{wen} where the rules \textit{close}, \textit{suspend}, \textit{left wen}, \textit{right wen} and \textit{fresh} interfere directly with \textit{wen} at the root of a principal formula. Three representative cases are presented.

\begin{enumerate}[label*=\textbf{.\arabic*}]
\item The first principal case for $\textit{wen}$ is when the bottommost rule of a proof is an instance of the \textit{close} rule of the form 
$
\vcenter{\infer[]{\wen x. P \cpar \new x. Q \cpar R
}{
\mathopen{\new x.} \left( P \cpar Q \right) \cpar R}}
$,
where $\vdash \mathopen{\new x.} \left( P \cpar Q \right) \cpar R$ and $\nfv{x}{R}$.
By the induction hypothesis,
there exist $S$ and $T$ such that $\vdash P \cpar Q \cpar T$ and $\nfv{x}{S}$ and either $S = T$ or $S = \wen x. T$, and also we have derivation $\infer{R}{S}$.
Since $\nfv{x}{S}$, if $S = T$ then $\vcenter{\infer[]{\new x. Q \cpar S}{\mathopen{\new x.}\left( Q \cpar T \right)}}$.
Furthermore, the size of the proof of $P \cpar Q \cpar T$ is no larger than the size of the proof of $\mathopen{\new x.} \left( P \cpar Q \right) \cpar R$; hence strictly bounded by the size of the proof of $\wen x. P \cpar \new x. Q \cpar R$.
If $S = \wen x. T$ then by the \textit{close} rule $\vcenter{\infer[]{\new x. Q \cpar \wen x. T}{\mathopen{\new x.}\left( Q \cpar T \right)}}$.
If $S = T$ then, since $\nfv{x}{S}$, by the \textit{extrude new} rule,
$\vcenter{\infer[]{\new x. Q \cpar T}{\mathopen{\new x.}\left( Q \cpar T \right)}}$.
Hence in either case $\vcenter{\infer[]{\new x. Q \cpar S}{\mathopen{\new x.}\left( Q \cpar T \right)}}$
and thereby the derivation 
$
\vcenter{
\infer[]{
\new x. Q \cpar R
}{
\infer[]{
\new x. Q \cpar S
}{
\mathopen{\new x.} \left(Q \cpar T \right)
}}}
$
can be constructed,
meeting the conditions for splitting for \textit{wen}.
%\begin{comment}
%The second principal case for $\textit{wen}$ is when the bottommost rule of a proof is of the form
%$
%\wen x P \cpar \new x Q \cpar R
%\longrightarrow
%\new x \left( P \cpar Q \right) \cpar R
%$, where $\vdash \new x \left( P \cpar Q \right) \cpar R$ and $\nfv{x}{R}$.
%By induction,
%there exist $S_i$ and $T_i$ such that and $\vdash P \cpar Q \cpar T_i$ and $\nfv{x}{S_i}$ either $S_i = T_i$ or $S_i = \wen x T_i$, for $1 \leq i \leq n$, and some $n$-ary killing context $\tcontext{}$ such that the follow derivation can be constructed.
%$R \longrightarrow
%\tcontext{ S_1, S_2, \hdots, S_n }
%$.
%Since $\nfv{x}{S_i}$, if $S_i = T_i$ then $\new x Q \cpar S_i \longrightarrow \new x\left( Q \cpar T_i \right)$.
%If $S_i = \wen x T_i$ then $\new x Q \cpar \wen x T_i \longrightarrow \new x\left( Q \cpar T_i \right)$.
%Hence the following derivation can be constructed, as required.
%\[
%\begin{array}{rl}
%\new x Q \cpar R
%\longrightarrow
%\new x Q \cpar \tcontext{ S_i \colon 1 \leq i \leq n }
%\longrightarrow&
%\tcontext{ \new x Q \cpar S_i \colon 1 \leq i \leq n }
%\\
%\longrightarrow&
%\tcontext{ \new x \left(Q \cpar T_i \right) \colon 1 \leq i \leq n }
%\end{array}
%\]
%\end{comment}
%%\smallskip

\item Consider the second principal case for \textit{wen} where the bottommost rule of a proof is an instance of the \textit{suspend} rule of the form
$
\vcenter{
\infer[]{
\wen x. P \cpar \wen x. Q \cpar R
}{
\mathopen{\wen x.} \left( P \cpar Q \right) \cpar R
}}
$, where $\vdash \mathopen{\wen x.} \left( P \cpar Q \right) \cpar R$ and $\nfv{x}{R}$.
By the induction hypothesis,
there exist $S$ and $T$ such that and $\vdash P \cpar Q \cpar T$ and $\nfv{x}{S}$ and either $S = T$ or $S = \new x. T$, and also $\vcenter{\infer[]{R}{S}}$.
Furthermore, the size of the proof of $P \cpar Q \cpar T$ is no larger than the size of the proof of $\mathopen{\wen x.} \left( P \cpar Q \right) \cpar R$; hence strictly bounded by the size of the proof of $\wen x. P \cpar \wen x. Q \cpar R$.
Since $\nfv{x}{S}$, if $S = T$ then, by the \textit{new wen} and \textit{extrude new} rules, 
$\vcenter{\infer[]{\wen x. Q \cpar T}{\infer[]{\new x. Q \cpar T}{\mathopen{\new x.} \left( Q \cpar T \right)}}}$.
If $S = \new x. T$ then, by the \textit{close} rule, $\vcenter{\infer[]{\wen x. Q \cpar \new x. T}{\mathopen{\new x.}\left( Q \cpar T \right)}}$.
So in either case, $\vcenter{\infer[]{\wen x. Q \cpar S}{\mathopen{\new x.} \left( Q \cpar T \right)}}$,
and hence the derivation
$
\vcenter{
\infer[]{
\wen x. Q \cpar R
}{
\infer[]{
\wen x. Q \cpar S
}{
\mathopen{\new x.}\left(Q \cpar T \right)
}}}
$
can be constructed, as required.
The principal cases for \textit{left wen} and \textit{right wen} are similar.
%\begin{comment}
%Consider the case where the bottommost rule of a proof is of the form
%$
%\wen x P \cpar \wen x Q \cpar R
%\longrightarrow
%\wen x \left( P \cpar Q \right) \cpar R
%$, where $\vdash \wen x \left( P \cpar Q \right) \cpar R$ and $\nfv{x}{R}$.
%By induction,
%there exist $S_i$ and $T_i$ such that and $\vdash P \cpar Q \cpar T_i$ and $\nfv{x}{S_i}$ and either $S_i = T_i$ or $S_i = \new x T_i$, for $1 \leq i \leq n$, and some $n$-ary killing context $\tcontext{}$ such that the follow derivation can be constructed.
%$R \longrightarrow
%\tcontext{ S_1, S_2, \hdots, S_n }
%$.
%Since $\nfv{x}{S_i}$, if $S_i = T_i$ then $\wen x Q \cpar S_i \longrightarrow \wen x \left( Q \cpar T_i \right) \longrightarrow \new x \left( Q \cpar T_i \right)$.
%If $S_i = \new x T_i$ then $\wen x Q \cpar \new x T_i \longrightarrow \new x\left( Q \cpar T_i \right)$.
%Hence the following derivation can be constructed, as required.
%\[
%\begin{array}{rl}
%\wen x Q \cpar R
%\longrightarrow
%\wen x Q \cpar \tcontext{ S_i \colon 1 \leq i \leq n }
%\longrightarrow&
%\tcontext{ \wen x Q \cpar S_i \colon 1 \leq i \leq n }
%\\
%\longrightarrow&
%\tcontext{ \new x \left(Q \cpar T_i \right) \colon 1 \leq i \leq n }
%\end{array}
%\]
%\end{comment}
%\smallskip

% For tech. rep. only.
\fullproof{
Another similar principal cases for $\textit{wen}$ is where the first the bottommost rule of a proof is of the form
$
\infer[]{
\wen x. P \cpar Q \cpar R
}{
\mathopen{\wen x.} \left( P \cpar Q \right) \cpar R
}
$, where $\vdash \mathopen{\wen x.} \left( P \cpar Q \right) \cpar R$ and $\nfv{x}{Q \cpar R}$.
By induction,
there exist $S$ and $T$ such that and $\vdash P \cpar Q \cpar T$ and either $S = T$ or $S = \new x. T$, and also $\vcenter{\infer[]{R}{S}}$.
Furthermore, the size of the proof of $P \cpar Q \cpar T$ is no larger than the size of the proof of $\mathopen{\wen x.} \left( P \cpar Q \right) \cpar R$; hence strictly bounded by the size of the proof of $\wen x. P \cpar Q \cpar R$.
If $S = T$ define $U = Q \cpar T$, and if $S = \new x. T$ define $U = \mathopen{\new x.} \left( Q \cpar T \right)$.
In the case $S = \new x. T$, since $\nfv{x}{Q}$, $\infer[]{Q \cpar \new x. T}{ \mathopen{\new x.} \left( Q \cpar T \right) }$.
Hence the following derivation
$
\infer[]{
Q \cpar R
}{
\infer[]{
Q \cpar S
}{
U
}}
$ can be constructed, as required.
}
%\begin{comment}
%Another similar principal cases for $\textit{wen}$ is where the first the bottommost rule of a proof is of the form
%$
%\wen x P \cpar Q \cpar R
%\longrightarrow
%\wen x \left( P \cpar Q \right) \cpar R
%$, where $\vdash \wen x \left( P \cpar Q \right) \cpar R$ and $\nfv{x}{Q \cpar R}$.
%By induction,
%there exist $S_i$ and $T_i$ such that and $\vdash P \cpar Q \cpar T_i$ and either $S_i = T_i$ or $S_i = \new x T_i$, for $1 \leq i \leq n$, and some $n$-ary killing context $\tcontext{}$ such that the follow derivation can be constructed.
%$R \longrightarrow
%\tcontext{ S_1, S_2, \hdots, S_n }
%$.
%If $S_i = T_i$ define $U_i = Q \cpar T_i$, and if $S_i = \new x T_i$ define $U_i = \new x \left( Q \cpar T_i \right)$.
%In the case $S_i = \new x T_i$, since $\nfv{x}{Q}$, $Q \cpar \new x T_i \longrightarrow \new x \left( Q \cpar T_i \right)$.
%Hence the following derivation can be constructed, as required.
%\[
%%\begin{array}{rl}
%Q \cpar R
%\longrightarrow
%Q \cpar \tcontext{ S_i \colon 1 \leq i \leq n }
%\longrightarrow
%\tcontext{ Q \cpar S_i \colon 1 \leq i \leq n }
%%\\
%\longrightarrow
%\tcontext{ U_i \colon 1 \leq i \leq n }
%%\end{array}
%\]
%\end{comment}
%

\item Consider the principal case for $\textit{wen}$ when the bottommost rule of a proof is an instance of the \textit{fresh} rule of the form 
$\vcenter{\infer[]{\wen x. \wen \vec{y}. P \cpar Q}{ \wen \vec{y}. \new x. P \cpar Q }}$, 
where $\vdash \wen \vec{y}. \new x. P \cpar Q$. Notice that $\vec{y}$ is required to handle the effect of \textit{equivariance}.
By applying the induction hypothesis inductively on the length of $\vec{y}$, there exist $\vec{z}$ and $\hat{Q}$ such that $\vec{z} \subseteq \vec{y}$ and $\nfv{\vec{y}}{\new \vec{z} \hat{Q}}$ and $\vdash \new x. P \cpar \hat{Q}$, and also $\vcenter{\infer[]{ Q }{ \new \vec{z}. \hat{Q} }}$. Furthermore, the size of the proof of $\new x. P \cpar \hat{Q}$ is bounded above by the size of the proof of $\wen \vec{y}. \new x. P \cpar Q$.
%
By the induction hypothesis, there exist $R$ and $S$ such that $\nfv{x}{R}$, $\vdash P \cpar S$ and either $R = S$ or $R = \wen x. S$, and also
$
\vcenter{ \infer[]{\hat{Q}}{R}}
$.
%Furthermore, the size of the proof of $P \cpar S$ is no larger than the size of the proof of $\new x P \cpar Q$; hence strictly bounded by the size of the proof of $\wen x P \cpar Q$.
There are two cases to consider. If $R = S$ then let $T = \new \vec{z}. S$; and if $R = \wen x. S$ then let $T = \new x. \new \vec{z}. S$, in which case, since $\new \vec{z}. \new x. S \equiv \new x. \new \vec{z}. S$ we have 
$\vcenter{\infer[]{\new \vec{z}. R}{T}}$. 
In either case $\nfv{x}{T}$. Thereby we can construct the derivation
$
\vcenter{
\infer[]{
Q
}{
\infer[]{
\new \vec{z}. \hat{Q}
}{
\infer[]{
\new \vec{z}. R
}{
T
}}}}
$.
Furthermore, appealing to Lemma~\ref{lemma:close}, the proof
$
\vcenter{
\infer[]{
\wen \vec{y}. P \cpar \new \vec{z}. S
}{
\infer[]{
\mathopen{\new \vec{y}.} \left( P \cpar S \right)
}{
\infer[]{
\new \vec{y}. \cunit
}{
\cunit
}}}}
$
can be constructed and, furthermore, $\size{ \wen \vec{y}. P \cpar \new \vec{z}. S } \prec \size{ \wen x. \wen \vec{y}. P \cpar Q }$, since by Lemma~\ref{lemma:bound} $\size{ \new \vec{z}. S } \preceq \size{ Q }$ and the \textit{wen} count strictly decreases.
\end{enumerate}

\begin{comment}
Consider the principal case for $\textit{wen}$ when $\wen x P \cpar Q \longrightarrow \new x P \cpar Q$ is the form of the bottommost rule of a proof, where $\vdash \new x P \cpar Q$ and $\nfv{x}{Q}$.
By induction,
there exist $R_i$ and $S_i$ such that $\vdash P \cpar S_i$ and either $R_i = S_i$ or $R_i = \wen x S_i$, for $1 \leq i \leq n$, and $n$-ary killing context $\tcontext{}$ such that 
$
 Q \longrightarrow \tcontext{ R_1, R_2, \hdots, R_n }
$.
There are two cases to consider. If $R_i = S_i$ then we are done. 
Otherwise, consider $R_i = \wen x S_i$ in which case $\wen x S_i \longrightarrow \new x S_i$, as required.
\end{comment}
%\smallskip


\item \textbf{Principal cases for new.}
The principal cases for \textit{new} are where the rules \textit{close}, \textit{extrude new}, \textit{medial new} and \textit{new wen} rules interfere directly with the $\textit{new}$ quantifier at the root of the principal formula.
Three cases are presented.

\begin{enumerate}[label*=\textbf{.\arabic*}]
\item The first principal case for $\textit{new}$ is when the bottommost rule of a proof is an instance of the \textit{close} rules of the form
$
\vcenter{
\infer[]{
\new x. P \cpar \wen x. Q \cpar R
}{
\mathopen{\new x.} \left( P \cpar Q \right) \cpar R
}}
$,
where $\vdash \mathopen{\new x.} \left( P \cpar Q \right) \cpar R$.
By the induction hypothesis, there exist formulae $U$ and $V$ such that $\vdash P \cpar Q \cpar V$ and $\nfv{x}{U}$ and either $U = V$ or $U = \wen x. V$, and also we have derivation
$\vcenter{\infer[]{R}{U}}$.
Furthermore, the size of the proof of $P \cpar Q \cpar V$ is no larger than the size of the proof of $\mathopen{\new x.} \left( P \cpar Q \right) \cpar R$; hence strictly bounded by the size of the proof of $\mathopen{\new x.} P \cpar \wen x. Q \cpar R$.
%
In the case $U = V$, we have $
\vcenter{\infer[]{
\wen x. Q \cpar V
}{
\mathopen{\wen x.} \left( Q \cpar V \right)
}}
$, since $\nfv{x}{U}$.
%
In the case $U = \wen x. V$, we have
$
\vcenter{\infer[]{
\wen x. Q \cpar \wen x. V
}{
\mathopen{\wen x.} \left( Q \cpar V \right)
}}
$.
%
Hence, by applying one of the above cases the following derivation
$
\vcenter{
\infer[]{
\wen x. Q \cpar R
}{
\infer[]{
\wen x. Q \cpar U
}{
\mathopen{\wen x.} \left( Q \cpar V \right)
}}}
$ can be constructed as required.
The principal case where the bottommost rule in a proof is the \textit{extrude new} rule follows a similar pattern.
\begin{comment}
There are four principal cases for $\textit{new}$.
This first is when the bottommost rule of a proof is
$
\new x P \cpar \wen x Q \cpar R
\longrightarrow
\new x \left( P \cpar Q \right) \cpar R
$,
where $\vdash \new x \left( P \cpar Q \right) \cpar R$ and $\nfv{x}{R}$. Observe that the proof of $\new x \left( P \cpar Q \right) \cpar R$ has proof one step shorter than $\new x P \cpar \wen x Q \cpar R$ and,
by Lemma~\ref{lemma:bound}, $\size{ \new x \left( P \cpar Q \right) \cpar R } \lsim \size{\new x P \cpar \wen x Q \cpar R}$
thereby strictly bounding the size of the proof.
N.B.\ this argument is used implicitly in every case for splitting.

By induction, there exist formulae $U_i$ and $V_i$ and $\vdash P \cpar Q \cpar V_i$ and $\nfv{x}{U_i}$ and either $U_i = V_i$ or $U_i = \wen x V_i$, for $1 \leq i \leq n$, and $n$-ary killing context such that the following holds.
$R \longrightarrow \tcontext{ U_1, U_1, \hdots, U_n }$.
Furthermore, the size of the proof of $P \cpar Q \cpar V_i$ is no larger than the size of the proof of $\new x \left( P \cpar Q \right) \cpar R$; hence strictly bounded by the size of the proof of $\new x P \cpar \wen x Q \cpar R$.

In the case that $U_i = V_i$, the following derivation can be constructed, since $\nfv{x}{U_i}$,
$
\wen x Q \cpar V_i
\longrightarrow
\wen x \left( Q \cpar V_i \right)
$.

In the case that $U_i = \wen x V_i$, we have derivation 
$
\wen x Q \cpar \wen x V_i
\longrightarrow 
\wen x \left( Q \cpar V_i \right)
$.

Hence, by applying one of the above cases for each $i$, the following derivation can be constructed as required.
\[
\begin{array}{rl}
\wen x Q \cpar R
\longrightarrow
\wen x Q \cpar \tcontext{ U_1, U_2, \hdots, U_n }
\longrightarrow&
\tcontext{ \wen x Q \cpar U_i \colon 1 \leq i \leq n }
\\
\longrightarrow&
\tcontext{ \wen x \left( Q \cpar V_i \right) \colon 1 \leq i \leq n }
\end{array}
\]
\end{comment}


\fullproof{
The second principal case for \textit{new} is when the bottommost rule of a proof is an instance of the \textit{extrude new} rule as follows, where $\nfv{x}{Q}$.
$
\vcenter{
\infer[]{
\new x. P \cpar Q \cpar R
}{
\mathopen{\new x.} \left( P \cpar Q \right) \cpar R
}}
$
where $\vdash \mathopen{\new x.} \left( P \cpar Q \right) \cpar R$ is provable.
By, the induction hypothesis, there exist formulae $U$ and $V$ where $\vdash P \cpar Q \cpar V$ and either $U = V$ or $U = \wen x. V$, and also
$\vcenter{\infer[]{R}{U}}$.
Furthermore, the size of the proof of $P \cpar Q \cpar V$ is bounded above by the size of the proof of $\mathopen{\new x.} \left( P \cpar Q \right) \cpar R$; thereby strictly bounded by the size of the proof of $\new x. P \cpar Q \cpar R$.
%
If $U = \wen x. V$, define $W = \mathopen{\wen x.} \left( Q \cpar V \right)$, and if $U = V$ define $W = Q \cpar V$.
In the cases where $U = \wen x. V$, since $\nfv{x}{Q}$, 
$\vcenter{\infer[]{Q \cpar \wen x. V}{\mathopen{\wen x.} \left( Q \cpar V \right)}}$, where the premise equals $W$.
In the cases where $U = V$, $Q \cpar U = W$.
Hence the derivation 
$
\vcenter{
\infer[]{
Q \cpar R
}{
\infer[]{
Q \cpar U
}{
W
}}}
$
can be constructed, as required.
}
\begin{comment}
The second principal case for \textit{new} is when the bottommost rule of a proof begins as follows, where $\nfv{x}{Q \cpar R}$.
$
\new x P \cpar Q \cpar R
\longrightarrow
\new x \left( P \cpar Q \right) \cpar R
$
where $\vdash \new x \left( P \cpar Q \right) \cpar R$ is provable.
By, the induction hypothesis, there exist formulae $U_i$ and $V_i$ where $\vdash P \cpar Q \cpar V_i$ and either $U_i = V_i$ or $U_i = \wen x V_i$, for $1 \leq i \leq n$, and $n$-ary killing context such that the following holds:
$R \longrightarrow \tcontext{ U_1, U_2, \hdots, U_n }$.
Furthermore, the size of the proof of $P \cpar Q \cpar V_i$ is bounded above by the size of the proof of $\new x \left( P \cpar Q \right) \cpar R$; thereby strictly bounded by the size of the proof of $\new x P \cpar Q \cpar R$.

If $U_i = \wen x V_i$, define $W_i = \wen x \left( Q \cpar V_i \right)$ and if $U_i = V_i$ define $W_i = Q \cpar V_i$ .
In the cases where $U_i = \wen x V_i$, since $\nfv{x}{Q}$, $Q \cpar \wen x V_i \longrightarrow \wen x \left( Q \cpar V_i \right) = W_i$.
In the cases where $U_i = V_i$, $Q \cpar U_i = W_i$.
Hence the following derivation can be constructed, as required.
\[
%\begin{array}{rl}
Q \cpar R
\longrightarrow
\tcontext{ U_1, U_2, \hdots, U_n } 
%\\
\longrightarrow
\tcontext{ Q \cpar U_i \colon 1 \leq i \leq n } 
\longrightarrow
\tcontext{ W_i \colon 1 \leq i \leq n } 
%\end{array}
\]
\end{comment}


\item Consider the second principal case for \textit{new} where the \textit{medial new} rule is the bottommost rule of a proof of the form
\[
\infer[\mbox{such that $\vdash \mathopen{\new \vec{y}.} \left(\new x. P \cseq \new x. Q\right) \cpar R$.}]{
\mathopen{\new x. \new \vec{y}.}\left( P \cseq Q \right) \cpar R
}{
\mathopen{\new \vec{y}.} \left( \new x. P \cseq \new x. Q \right) \cpar R
}
\]
The $\vec{y}$ is required to handle cases induced by equivariance.
By applying the induction hypothesis repeatedly, 
there exists $\vec{z}$ and $\hat{R}$ such that $\vec{z} \subseteq \vec{y}$ and
$\nfv{\vec{y}}{\wen \vec{z}. \hat{R} }$ and $\vdash \left( \new x. P \cseq \new x. Q \right) \cpar \hat{R}$, and also $\vcenter{\infer[]{R}{\hat{R}}}$.
Furthermore, the size of the proof of $\left( \new x. P \cseq \new x. Q \right) \cpar \hat{R}$ is bounded above by the size of the proof of $\mathopen{\new \vec{y}.} \left( \new x. P \cseq \new x. Q \right) \cpar R$.
%
By the induction hypothesis, there exist $S_i$ and $T_i$ such that $\vdash \new x. P \cpar S_i$ and $\vdash \new x. Q \cpar T_i$, for $1 \leq i \leq n$, and $n$-ary killing context such that 
$\vcenter{
\infer[]{\hat{R}}{\tcontext{ S_1 \cseq T_1, S_2 \cseq T_2, \hdots, S_n \cseq T_n }}}
$.
Furthermore, the size of the proofs of $\new x. P \cpar S_i$ and $\new x. Q \cpar T_i$ are bounded above by the size of the proof of $\left(\new x. P \cseq \new x. Q\right) \cpar R$.
%Observe that by $\alpha$-conversion, $\vdash \new y \left(P\sub{x}{y}\right) \cpar S_i$ and $\vdash \new y \left(Q\sub{x}{y}\right) \cpar T_i$ also hold, where $\nfv{y}{\new x P \cpar S_i}$ and $\nfv{y}{\new x Q \cpar T_i}$.
%
By the induction hypothesis again, there exist $U^i$ and $\hat{U}^i$ such that $\vdash P \cpar \hat{U}^i$ and $\nfv{x}{U^i}$ and either $U^i = \hat{U}^i$ or $U^i = \wen x. \hat{U}^i$,
and also
$
\vcenter{\infer[]{S_i}{U^i}}
$.
Also by the induction hypothesis, there exist $V^i$ and $\hat{V}^i$ such that $\vdash Q \cpar \hat{V}^i$ and $\nfv{x}{V^i}$ and either $V^i = \hat{V}^i$ or $V^i = \wen x. \hat{V}^i$, 
and also
$
\vcenter{\infer[]{T_i}{V^i}}
$.
%
Now define $W$ and $\hat{W}$ such that
$
\hat{W} = \wen \vec{z}. \tcontext{ \hat{U}^i \cseq \hat{V}^i \colon 1 \leq i \leq n }
$
and, 
if for all $1 \leq i \leq n$, $U^i = \hat{U}^i$ and $V^i = \hat{V}^i$, then $W = \hat{W}$;
otherwise $W = \wen x. \hat{W}$.
%
%\ross{Is lemma:commute required now, or is it built into splitting?}
%In cases where $U^i = \wen x \hat{U}^i$, by Lemma~\ref{lemma:commute}
%$U^i \longrightarrow \wen x \tcontextmn{i}{0}{ \hat{U}^i_1, \hat{U}^i_2, \hdots, \hat{U}^i_{m_i} }$;
%In cases where for some $1 \leq k \leq \ell_i$, $V^i_k = \wen x \hat{V}^i_k$, by Lemma~\ref{lemma:commute} $\tcontextmn{i}{1}{ V^i_1, V^i_2, \hdots, V^i_{\ell_i} } \longrightarrow \tcontextmn{i}{1}{ \hat{V}^i_1, \hat{V}^i_2, \hdots, \hat{V}^i_{\ell_i} }$.
Hence for each $i$, one of the following derivations holds.
\begin{itemize}
\item 
$U^i = \hat{U}^i$ and $V^i = \hat{V}^i$ hence
$U^i \cseq V^i = \hat{U}^i \cseq \hat{V}^i$.

\item
If $U^i = \wen x. \hat{U}^i$ and $V^i = \hat{V}^i$, hence $\nfv{x}{V^i}$, by the \textit{left wen} rule
$
\vcenter{
\infer[]{
\wen x. \hat{U}^i \cseq \hat{V}^i
}{
\mathopen{\wen x.} \left( \hat{U}^i \cseq \hat{V}^i \right)
}}
$.

\item
If $U^i = \hat{U}^i$, hence $\nfv{x}{\hat{U}^i}$, and $V^i = \wen x. \hat{V}^i$, by the \textit{right wen} rule
$
\vcenter{
\infer[]{
U^i \cseq \wen x. \hat{V}^i
}{
\mathopen{\wen x.} \left( \hat{U}^i \cseq \hat{V}^i \right)
}}
$.

\item
Otherwise by the \textit{suspend} rule $
\vcenter{\infer[]{
\wen x. \hat{U}^i \cseq \wen x. \hat{V}^i
}{
\mathopen{\wen x.} \left( \hat{U}^i \cseq \hat{V}^i \right)
}}
$
\end{itemize}
If for all $i$ such that $1 \leq i \leq n$, $U^i = \hat{U}^i$ and $V^i = \hat{V}^i$ then $W = \hat{W}$. Otherwise, by Lemma~\ref{lemma:commute},
$
\vcenter{
\infer[]{
\wen \vec{z}. \tcontext{
 U^i \cseq V^i \colon 1 \leq i \leq n
}
}{
\wen \vec{z}. \wen x. \tcontext{
 \hat{U}^i \cseq \hat{V}^i \colon 1 \leq i \leq n
}
}}$, where the premise is equialent to $W$.
Thereby the  derivation below left can be constructed, and furthermore, using Lemma~\ref{lemma:close}, 
the proof below right can also be constructed.
\[
\infer[]{
R
}{
\infer[]{
\wen \vec{z}. \hat{R}
}{
\infer[]{
\wen \vec{z}. \tcontext{ S_i \cseq T_i \colon 1 \leq i \leq n }
}{
\infer[]{
\wen \vec{z}. \tcontext{ U^i \cseq V^i \colon 1 \leq i \leq n }
}{
W
}}}}
\qquad
\infer[]{
\mathopen{\new \vec{y}.} \left( P \cseq Q \right) \cpar \hat{W}
}{
\infer[]{
\mathopen{\new \vec{y}.} \left(\left( P \cseq Q \right) \cpar
\tcontext{
 \hat{U}^i \cseq \hat{V}^i \colon 1 \leq i \leq n
}
\right)
}{
\infer[]{
\mathopen{\new \vec{y}.} 
\tcontext{
 \left( P \cseq Q \right) \cpar \left(\hat{U}^i \cseq \hat{V}^i\right) \colon 1 \leq i \leq n
}
}{
\infer[]{
\mathopen{\new \vec{y}.} 
\tcontext{
 \left( P \cpar \hat{U}^i \right) \cseq \left( Q \cpar \hat{V}^i\right) \colon 1 \leq i \leq n
}
}{
\infer[]{
\mathopen{\new \vec{y}.} 
\tcontext{
 \cunit \colon 1 \leq i \leq n
}
}{
\cunit
}}}}}
\]



By Lemma~\ref{lemma:bound},
$\size{\hat{W}} \preceq \size{R}$; hence $\size{ \mathopen{\new \vec{y}.} \left( P \cseq Q \right) \cpar \hat{W} } \prec \size{ \mathopen{\new x. \new \vec{y}.} \left( P \cseq Q \right) \cpar R }$ since the \textit{new count} strictly decreases, as required.
\begin{comment}
Consider the third principal case for \textit{new} where the \textit{medial new} rule is the bottommost rule of a proof of the form
$
\new x \left( P \cseq Q \right) \cpar R
\longrightarrow 
\left( \new x P \cseq \new x Q \right) \cpar R
$
such that $\vdash \left(\new x P \cpar \new x Q\right) \cpar R$ and $\nfv{x}{R}$.

By induction, there exist $S_i$ and $T_i$ such that $\vdash \new x P \cpar S_i$ and $\vdash \new x Q \cpar T_i$, for $1 \leq i \leq n$, and $n$-ary killing context such that 
$
R \longrightarrow \tcontext{ S_1 \cseq T_1, S_2 \cseq T_2, \hdots, S_n \cseq T_n }
$.
Furthermore, the size of the proofs of $\new x P \cpar S_i$ and $\new x Q \cpar T_i$ are bounded above by the size of the proof of $\left(\new x P \cseq \new x Q\right) \cpar R$.
Observe that by $\alpha$-conversion, $\vdash \new y \left(P\sub{x}{y}\right) \cpar S_i$ and $\vdash \new y \left(Q\sub{x}{y}\right) \cpar T_i$ also hold, where $\nfv{y}{\new x P \cpar S_i}$ and $\nfv{y}{\new x Q \cpar T_i}$.

By induction, there exist $U^i_j$ and $\hat{U}^i_j$ such that $\vdash P\sub{x}{y} \cpar \hat{U}^i_j$ and $\nfv{y}{U^i_j}$ and either $U^i_j = \hat{U}^i_j$ or $U^i_j = \wen y \hat{U}^i_j$, for $1 \leq j \leq m_i$, and $m_i$-ary killing context $\tcontextmn{i}{0}{}$ such that 
$
S_i \longrightarrow \tcontextmn{i}{0}{ U^i_1, U^i_2, \hdots, U^i_{m_i} }
$.
Also by induction, there exist $V^i_k$ and $\hat{V}^i_k$ such that $\vdash Q\sub{x}{y} \cpar \hat{V}^i_k$ and $\nfv{y}{V^i_k}$ and either $V^i_k = \hat{V}^i_k$ or $V^i_k = \wen y \hat{V}^i_k$, for $1 \leq k \leq \ell_i$, and $\ell_i$-ary killing context $\tcontextmn{i}{1}{}$ such that 
$
T_i \longrightarrow \tcontextmn{i}{1}{ V^i_1, V^i_2, \hdots, V^i_{\ell_i} }
$.

Now define $W_i$ and $\hat{W}_i$ such that
$
\hat{W}_i = \tcontextmn{i}{0}{ \hat{U}^i_j \colon 1 \leq j \leq m_i } \cseq \tcontextmn{i}{1}{ \hat{V}^i_k \colon 1 \leq k \leq \ell_i }
$
and, if $U^i_j = \hat{U}^i_j$ and $V^i_k = \hat{V}^i_k$ for all $j$ and $k$, then $W_i = \hat{W}_i$;
otherwise $W_i = \wen x \hat{W}_i$. 
In cases where for some $1 \leq j \leq m_i$, $U^i_j = \wen x \hat{U}^i_j$, by Lemma~\ref{lemma:commute} $\tcontextmn{i}{0}{ U^i_1, U^i_2, \hdots, U^i_{m_i} } \longrightarrow \wen x \tcontextmn{i}{0}{ \hat{U}^i_1, \hat{U}^i_2, \hdots, \hat{U}^i_{m_i} }$;
In cases where for some $1 \leq k \leq \ell_i$, $V^i_k = \wen x \hat{V}^i_k$, by Lemma~\ref{lemma:commute} $\tcontextmn{i}{1}{ V^i_1, V^i_2, \hdots, V^i_{\ell_i} } \longrightarrow \tcontextmn{i}{1}{ \hat{V}^i_1, \hat{V}^i_2, \hdots, \hat{V}^i_{\ell_i} }$.
Hence for each $i$, one of the of the following derivations holds.
\begin{itemize}
\item 
$U^i_j = \hat{U}^i_j$ and $V^i_k = \hat{V}^i_k$ for all $j$ and $k$ hence
$\tcontextmn{i}{0}{ U^i_j \colon 1 \leq j \leq m_i } \cseq \tcontextmn{i}{1}{ V^i_k \colon 1 \leq k \leq \ell_i } = W_i$.

\item
If $U^i_j = \hat{U}^i_j$ for all $j$, hence $\nfv{y}{U^i_j}$ for all $j$, and hence $\nfv{y}{Q}$ thereby $\nfv{y}{\tcontextmn{i}{0}{ \hat{U}^i_j \colon 1 \leq j \leq m_i }}$, hence by using the \textit{left wen} rule
$
\tcontextmn{i}{0}{ U^i_j \colon 1 \leq j \leq m_i } \cseq \tcontextmn{i}{1}{ V^i_k \colon 1 \leq k \leq \ell_i }
\longrightarrow
\wen x \tcontextmn{i}{0}{ \hat{U}^i_j \colon 1 \leq j \leq m_i } \cseq \tcontextmn{i}{1}{ \hat{V}^i_k \colon 1 \leq k \leq \ell_i }
\longrightarrow
W_i
$.

\item
If $V^i_k = \hat{V}^i_k$ for all $k$, hence $\nfv{y}{V^i_j}$ for all $k$, thereby $\nfv{y}{\tcontextmn{i}{1}{ \hat{V}^i_k \colon 1 \leq k \leq \ell_i }}$, hence by using the \textit{right wen} rule
$
\tcontextmn{i}{1}{ U^i_j \colon 1 \leq j \leq m_i } \cseq \tcontextmn{i}{1}{ V^i_k \colon 1 \leq k \leq \ell_i }
\longrightarrow
\tcontextmn{i}{1}{ \hat{U}^i_j \colon 1 \leq j \leq m_i } \cseq \wen x \tcontextmn{i}{1}{ \hat{V}^i_k \colon 1 \leq k \leq \ell_i }
\longrightarrow
W_i
$.

\item
Otherwise by using the \textit{suspend} rule the following derivation holds.
$
\tcontextmn{i}{0}{ U^i_j \colon 1 \leq j \leq m_i } \cseq \tcontextmn{i}{1}{ V^i_k \colon 1 \leq k \leq \ell_i }
\longrightarrow
\wen x \tcontextmn{i}{0}{ \hat{U}^i_j \colon 1 \leq j \leq m_i } \cseq \wen x \tcontextmn{i}{1}{ \hat{V}^i_k \colon 1 \leq k \leq \ell_i }
\longrightarrow
W_i
$
\end{itemize}
Thereby the following derivation can be constructed.
\[
\begin{array}{rl}
R
\longrightarrow&
\tcontext{ S_i \cseq T_i \colon 1 \leq i \leq n }
\\
\longrightarrow&
\tcontext{ \tcontextmn{i}{0}{ U^i_j \colon 1 \leq j \leq m_i } \cseq \tcontextmn{i}{1}{ V^i_k \colon 1 \leq k \leq \ell_i } \colon 1 \leq i \leq n }
\longrightarrow
\tcontext{ W_i }
\end{array}
\]
Furthermore, the following proof can be constructed.
\[
\begin{array}{rl}
\left( P \cseq Q \right)\sub{x}{y} \cpar \hat{W}_i
\longrightarrow&
\left( P\sub{x}{y} \cpar \tcontextmn{i}{0}{ \hat{U}^i_j \colon 1 \leq j \leq m_i } \right) \cseq
\left( Q\sub{x}{y} \cpar \tcontextmn{i}{1}{ \hat{V}^i_k \colon 1 \leq k \leq \ell_i } \right)
\\
\longrightarrow&
\tcontextmn{i}{0}{ P\sub{x}{y} \cpar \hat{U}^i_j \colon 1 \leq j \leq m_i } \cseq
\tcontextmn{i}{1}{ Q\sub{x}{y} \cpar \hat{V}^i_k \colon 1 \leq k \leq \ell_i }
\\
\longrightarrow&
\tcontextmn{i}{0}{ \cunit \colon 1 \leq j \leq m_i } \cseq
\tcontextmn{i}{1}{ \cunit \colon 1 \leq k \leq \ell_i }
\longrightarrow \cunit
\end{array}
\]
By Lemma~\ref{lemma:bound} and since the \textit{new count} strictly decreases,
$\size{\hat{W}_i} \preceq \size{R}$; thereby $\size{ \left( P \cseq Q \right)\sub{x}{y} \cpar \hat{W}_i } \prec \size{ \new x \left( P \cseq Q \right) \cpar R }$ as required.
\end{comment}
\smallskip

\item Consider the third principal  case for \textit{new} where the bottommost rule of a proof is the \textit{new wen} rule of the form
\[
\infer[\mbox{, where $\vdash \new \vec{z}. \wen y. \new x. P \cpar Q$.} ]{
\new x. \new \vec{z}. \wen y. P \cpar Q
}{
\new \vec{z}. \wen y. \new x. P \cpar Q
}
\]
%, $\nfv{x}{Q}$ and $\nfv{y}{Q}$.
By applying the induction hypothesis repeatedly, 
there exist $\vec{w}$ and $\hat{Q}$ such that $\vec{w} \subseteq \vec{z}$
and $\nfv{\vec{z}}{\wen \vec{w} .\hat{Q}}$ and $\vdash \wen y. \new x. P \cpar \hat{Q}$, and also $\vcenter{\infer[]{Q}{\wen \vec{w}. \hat{Q}}}$. Furthermore, the size of the proof of $\wen y. \new x. P \cpar \hat{Q}$ is bounded above by the size of the proof of $\new \vec{z}. \wen y.\new x. P \cpar Q$.
%
By the induction hypothesis, there exist $R$ and $S$ such that $\nfv{x}{R}$ and $\vdash \new x. P \cpar S$ and either $R = S$ or $R = \new y. S$, and also 
$
\vcenter{
\infer[]{\hat{Q}}{R}}
$.
Furthermore, the size of the proof of $\new x. P \cpar S$ is bounded above by the size of the proof of $\wen y. \new x. P \cpar Q$, hence strictly bounded above by the size of the proof of $\new x. \wen y. P \cpar Q$ enabling the induction hypothesis.
By the induction hypothesis again, there exist $U$ and $V$ such that $\nfv{x}{U}$ and $\vdash P \cpar V$ and either $U = V$ or $U = \wen x. V$, and also
$
\vcenter{
\infer[]{S}{U}}
$.

Let $W$ and $\hat{W}$ be defined such that, if $R = \new y. S$, then $\hat{W} = \new y. V$;
or, if $R = S$, then $\hat{W} = V$.
If $V = U$ then define $W = \wen \vec{w}. \hat{W}$.
If $U = \wen x. V$, then define $W = \wen x. \wen \vec{w}. \hat{W}$.
There are four scenarios for constructing a derivation with premise $W$ and conclusion $\wen \vec{w}. R$.
\begin{itemize}
\item In the case $V = U$ and $R = \new y. S$ then $\wen \vec{w}. \new y. U = W$.

\item If $V = U$ and $R = S$ then $\wen \vec{w}. U = W$.

% \new y V^i_j =  $\nfv{x}{R}$, hence, by Lemma~\ref{lemma:nfv}, $\nfv{x}{W_i}$.
\item If both $U = \wen x. V$ and $R = \new y. S$ hold, then 
we have 
\[
\infer[\mbox{, where the premise is $W$.}]{
\wen \vec{w}.R
% = \wen \vec{w}. \new y. S
}{
\infer[]{
%\wen \vec{w}. \new y. U %\\
%}{
 \wen \vec{w}. \new y. \wen x. V %\\
}{
\wen x. \wen \vec{w}. \new y. V
}}
\]

\item If both $U = \wen x. V$ and $R = S$ then 
%$\wen \vec{w}.R = \wen \vec{w}. S$
%and so
$
\vcenter{
\infer[]{
\wen \vec{w}.R
}{
 \wen \vec{w}. U
}}
$,
where the premise is equivalent to 
%= \wen \vec{w}. \wen x. V \equiv 
$W$.
\end{itemize}
Thereby, by applying one of the above cases,  we have %the following derivation can be constructed.
$\vcenter{
\infer[.]{
\hat{Q}
}{
\infer[]{
\wen \vec{w}. Q
}{
\infer[]{
\wen \vec{w}. R
}{
W
}}}}
$

In the case that $\hat{W} = \new y. V$,  the left most derivation below holds. 
In the case, $\hat{W} = V$ and $\nfv{y}{V}$ the middle derivation below holds.  
Hence in either case, appealing to Lemma~\ref{lemma:close}, the proof below right can be constructed:
\[
\infer[]{
\wen y. P \cpar \new y. V
}{
\mathopen{\new y.} \left( P \cpar V \right)
}
\qquad
\infer[]{
\wen y. P \cpar \hat{W}
}{
\infer[]{
\mathopen{\wen y.} \left( P \cpar V \right)
}{
\mathopen{\new y.} \left( P \cpar V \right)
}}
\qquad
\infer[]{
\new \vec{z}. \wen y. P \cpar \wen \vec{w}. \hat{W}
}{
\infer[]{
\mathopen{\new \vec{z}.} \left( \wen y. P \cpar \hat{W} \right)
}{
\infer[]{
\mathopen{\new \vec{z}. \new y.} \left( P \cpar V \right)
}{
\infer[]{
\new \vec{z}. \new y. \cunit
}{
\cunit
}}}}
\]

\noindent Furthermore, by Lemma~\ref{lemma:bound}, $\size{ \wen \vec{w}. \hat{W} } \preceq \size{ Q }$.
Hence $\size{\wen y. P \cpar \wen \vec{w}. \hat{W}} \prec \size{ \new x. \new \vec{z}. \wen y. P \cpar Q }$ since the \textit{new} count strictly decreases.
\begin{comment}
Consider the fourth principal  case for \textit{new} where the bottommost rule of a proof is the \textit{new wen} rule of the form
$
\new x \wen y P \cpar Q
\longrightarrow
\wen y \new x P \cpar Q
$
where $\vdash \wen y \new x P \cpar Q$, $\nfv{x}{Q}$ and $\nfv{y}{Q}$.

By induction, there exist $R_i$ and $S_i$ such that $\vdash \new x P \cpar S_i$ and either $R_i = S_i$ or $R_i = \new y S_i$, for $1 \leq i \leq n$,
and $n$-ary killing context $\tcontext{}$ such that
$
Q \longrightarrow \tcontext{ R_1, R_2, \hdots, R_n }
$.
Furthermore, the size of the proof of $\new x P \cpar S_i$ is bounded above by the size of the proof of $\wen y \new x P \cpar Q$ enabling the induction hypothesis.
By induction again, there exist $U^i_j$ and $V^i_j$ such that $\vdash P \cpar V^i_k$ and either $U^i_j = V^i_j$ or $U^i_j = \wen x V^i_j$, for $1 \leq j \leq m_i$, and $n$-ary killing context $\tcontextn{i}{}$ such that 
$
S_i \longrightarrow
\tcontextn{i}{ U^i_1, U^i_2, \hdots, U^i_{m_i} }
$.

Let $W_i$ and $\hat{W}_i$ be defined such that, if $R_i = \new y S_i$, then $\hat{W}_i = \new y \tcontextn{i}{ V^i_j \colon 1 \leq j \leq m_i }$;
or, if $R_i = S_i$, then $\hat{W}_i = \tcontextn{i}{ V^i_j \colon 1 \leq j \leq m_i }$.
If $V^i_j = U^i_j$ for all $1 \leq j \leq m_i$, then define $W_i = \hat{W}_i$.
If for some $1 \leq j \leq m_i$ $U^i_j = \wen x V^i_j$, then define $W_i = \wen x \hat{W}_i$.

In the following fix some $i$ such that $1 \leq i \leq n$. There are four scenarios for $R_i \longrightarrow W_i$.
In the case $V^i_j = U^i_j$ for all $1 \leq j \leq m_i$, if $R_i = \new y S_i$ then $\new y \tcontextn{i}{ U^i_j \colon 1 \leq j \leq m_i } = W_i$ and if $R_i = S_i$ then $\tcontextn{i}{ U^i_j \colon 1 \leq j \leq m_i } = W_i$.
% \new y V^i_j =  $\nfv{x}{R}$, hence, by Lemma~\ref{lemma:nfv}, $\nfv{x}{W_i}$.

If, for some $1 \leq j \leq m_i$, $U^i_j = \wen x V^i_j$, then $\tcontextn{i}{ U^i_j \colon 1 \leq j \leq m_i } \longrightarrow \wen x \tcontextn{i}{V^i_j \colon 1 \leq j \leq m_i}$, by Lemma~\ref{lemma:commute}, and one of the following hold.
\begin{itemize}
\item If, further to $U^i_j = \wen x V^i_j$ for some $1 \leq j \leq m_i$, $R_i = \new y S_i$ then 
$
%\begin{array}{rl}
R_i \longrightarrow \new y \tcontextn{i}{ U^i_j \colon 1 \leq j \leq m_i } %\\
    \longrightarrow \new y \wen x \tcontextn{i}{ V^i_j \colon 1 \leq j \leq m_i } %\\
    \longrightarrow \wen x \new y \tcontextn{i}{ V^i_j \colon 1 \leq j \leq m_i } =  W_i
%\end{array}
$.
\item
If, further to $U^i_j = \wen x V^i_j$ for some $1 \leq j \leq m_i$,  $R_i = S_i$ then 
$
%\begin{array}{rl}
R_i \longrightarrow \tcontextn{i}{ U^i_1, U^i_2, \hdots, U^i_{m_i} }
%\\
    \longrightarrow \wen x \tcontextn{i}{ V^i_1, V^i_2, \hdots, V^i_{m_i} } = W_i
%\end{array}
$.
\end{itemize}
Thereby, by applying one of the above cases for each $i$, the following derivation can be constructed.
$
%\begin{array}{rl}
Q
\longrightarrow
 \tcontext{ R_1, R_2, \hdots, R_n }
%\\
\longrightarrow
 \tcontext{ W_1, W_2, \hdots, W_n }
%\end{array}
$.

Now observe that the following proof can always be constructed for every $i$ such that $1 \leq i \leq n$.
\[
\begin{array}{rl}
\new y \left( P \cpar \tcontextn{i}{ V^i_j \colon 1 \leq j \leq m_i } \right)
\longrightarrow
\new y \tcontextn{i}{ P \cpar V^i_j \colon 1 \leq j \leq m_i }
%\\
\longrightarrow
\new y \tcontextn{i}{ \cunit \colon 1 \leq j \leq m_i }
%\\
\longrightarrow
\cunit
\end{array}
\]
By using the proof above, a proof can be constructed for each $i$ in the following two cases.

In the first case where $\hat{W}_i = \new y \tcontextn{i}{ V^i_j \colon 1 \leq j \leq m_i }$, a proof can be constructed such that
$
%\begin{array}{rl}
\wen y P \cpar \hat{W}_i
\longrightarrow
\new y \left( P \cpar \tcontextn{i}{ V^i_j \colon 1 \leq j \leq m_i } \right)
\longrightarrow
\cunit
%\end{array}
$
Furthermore, by Lemma~\ref{lemma:bound}, $\size{ \new y \tcontextn{i}{ V^i_j \colon 1 \leq j \leq m_i } } \preceq \size{ Q }$. Hence 
Hence $\size{ \wen y P \cpar \new y \tcontextn{i}{ V^i_j \colon 1 \leq j \leq m_i }} \prec \size{ \new x \wen y P \cpar Q }$.

In the second case, $\hat{W}_i = \tcontextn{i}{ V^i_j \colon 1 \leq j \leq m_i }$ and, since $\nfv{y}{R}$, by Lemma~\ref{lemma:nfv} we have $\nfv{y}{\hat{W}_i}$.
Hence the following proof can be constructed, using the above proof:
$
%\begin{array}{rl}
\wen y P \cpar \hat{W}_i
\longrightarrow
\new y \left( P \cpar \tcontextn{i}{ V^i_j \colon 1 \leq j \leq m_i } \right)
%\\
\longrightarrow
\cunit
%\end{array}
$.
Furthermore, by Lemma~\ref{lemma:bound}, $\size{ \tcontextn{i}{ V^i_j \colon 1 \leq j \leq m_i } } \preceq \size{ Q }$.
Hence $\size{\wen y P \cpar \tcontextn{i}{ V^i_j \colon 1 \leq j \leq m_i }} \prec \size{\new x \wen y P \cpar Q}$.
\end{comment}

\end{enumerate}



\item \textbf{Principal cases for seq.} There are two forms of principal cases for \textit{seq}.
The first case, induced by the \textit{sequence} rule, is the case that forces the \textit{medial}, \textit{medial1} and \textit{medial new} rules.
The other cases are induced by the \textit{suspend}, \textit{left wen} and \textit{right wen} rules (which are forced as a knock on effect of the \textit{medial new} rule).

\begin{enumerate}[label*=\textbf{.\arabic*}]
\item Consider the first principal case for \textit{seq}.
The difficulty in this case is that, due to associativity of \textit{seq}, the \rseq rule may be applied in several ways when there are multiple occurrences of \textit{seq}.
Consider a principal formula of the form $\left(T_0 \cseq T_1\right) \cseq T_2$, where we aim to split the formula around the second \textit{seq} operator. The difficulty is that the bottommost rule may be an instance of the \rseq rule applied between $T_0$ and $T_1 \cseq T_2$. Symmetrically, the principal formula may be of the form $T_0 \cseq \left(T_1 \cseq T_2\right)$ but the bottommost rule may be an instance of the \rseq rule applied between $T_0 \cseq T_1$ and $T_2$.
In the following analysis, only the former case is considered; the symmetric case follows a similar pattern.
The principal formula is $\left(T_0 \cseq T_1\right) \cseq T_2$ and the bottommost rule is an instance of the \textit{sequence} rule of the form
\[
\infer[]{
\left( T_0 \cseq T_1 \cseq T_2 \right) \cpar \left( U \cseq V \right) \cpar W
}{
\left(\left( T_0 \cpar U \right) \cseq \left( \left(T_1 \cseq T_2\right) \cpar V \right)\right) \cpar W
}
\] 
where $T_0 \not\equiv \cunit$, $T_2 \not\equiv \cunit$ (otherwise splitting is trivial), and either $U \not\equiv \cunit$ or $V \not\equiv \cunit$ (otherwise the \rseq rule cannot be applied);
%\]
and also $\vdash \left(\left( T_0 \cpar U \right) \cseq \left( \left(T_1 \cseq T_2\right) \cpar V \right)\right) \cpar W$.
By the induction hypothesis, %either $W \longrightarrow \top$ in which case $\left(U \cseq V\right) \cpar W \longrightarrow \left(U \cseq V\right) \cpar \top \longrightarrow \top$ hence we are done, or
there exist $P_i$ and $Q_i$ such that $\vdash T_0 \cpar U \cpar P_i$ and $\vdash \left(T_1 \cseq T_2\right) \cpar V \cpar Q_i$ hold, for $1 \leq i \leq n$, and an $n$-ary killing context $\tcontext{}$ such that 
\[
\infer[.]{W}{
\tcontext{ P_1 \cseq Q_1, \hdots, P_n \cseq Q_n }
}
\]
Furthermore,
%$\size{T_0 \cpar U \cpar P_i} \mleq \size{ \left(\left( T_0 \cpar U \right) \cseq \left( \left(T_1 \cseq T_2\right) \cpar V \right)\right) \cpar W}$ and also 
the size of the proof of formula
$
\left(T_1 \cseq T_2\right) \cpar V \cpar Q_i
$
is bounded above by the size of the proof of
$\left(\left( T_0 \cpar U \right) \cseq \left( \left(T_1 \cseq T_2\right) \cpar V \right)\right) \cpar W$, hence the induction hypothesis is enabled.
%
By the induction hypothesis, %either $V \cpar Q_i \longrightarrow \top$, or
there exists $R^i_j$ and $S^i_j$ such that $\vdash T_1 \cpar R^i_j$ and $\vdash T_2 \cpar S^i_j$, for $1 \leq j \leq m_i$, and $m_i$-ary killing context $\tcontextn{i}{}$ such that
\[
\infer[.]{
V \cpar Q_i
}{
\tcontextn{i}{R^i_1 \cseq S^i_1, \hdots, R^i_{m_i} \cseq S^i_{m_i}}
}
\]
Furthermore, by Lemma~\ref{lemma:medial} there exist killing contexts $\tcontextmn{i}{0}{}$ and $\tcontextmn{i}{1}{}$ and sets of integers $J^i \subseteq \left\{1, \hdots, n\right\}$, $K^i \subseteq \left\{1, \hdots, n\right\}$ such that
\[
\infer[.]{
\tcontextn{i}{R^i_1 \cseq S^i_1, \hdots, R^i_{m_i} \cseq S^i_{m_i}}
}{
\tcontextmn{i}{0}{R^i_j \colon j \in J^i}
 \cseq
\tcontextmn{i}{1}{S^i_k \colon k \in K^i}
}
\]
Thereby, the following derivation can be constructed.
\[
\infer[]{
\left( U \cseq V \right) \cpar W
}{
\infer[]{
\left( U \cseq V \right) \cpar \tcontext{ P_1 \cseq Q_1, \hdots, P_n \cseq Q_n } 
}{
\infer[]{
 \tcontext{ \left( U \cseq V \right) \cpar \left( P_1 \cseq Q_1 \right), \hdots, \left( U \cseq V \right) \cpar \left( P_n \cseq Q_n \right) } 
}{
\infer[]{
\tcontext{ \left( U \cpar P_1 \right) \cseq \left( V \cpar Q_1 \right), \hdots, \left( U \cpar P_n \right) \cseq \left( V \cpar Q_n \right) } 
}{
\infer[]{
 \tcontext{ \left( U \cpar P_i \right) \cseq \tcontextn{i}{R^i_j \cseq S^i_j \colon 1 \leq j \leq m_i } \colon 1 \leq i \leq n } 
}{
\tcontext{
 \left( U \cpar P_i \right) \cseq \tcontextmn{i}{0}{R^i_j \colon j \in J^i} \cseq \tcontextmn{i}{1}{S^i_k \colon k \in K^i}
 \colon 1 \leq i \leq n 
} 
}}}}}
\]
Furthermore, the following two proofs can be constructed.
\[
\infer[]{
T_2 \cpar \tcontextn{i}{ S^i_j \colon 1 \leq j \leq m_i }
}{
\infer[]{
\tcontextn{i}{ T_2 \cpar S^i_j \colon 1 \leq j \leq m_i }
}{
\infer[]{
\tcontextn{i}{ \cunit \colon 1 \leq j \leq m_i }
}{
\cunit
}}}
\qquad\qquad
\infer[]{
\left(T_0 \cseq T_1\right) \cpar \left(\left( U \cpar P_i \right) \cseq \tcontextn{i}{R^i_j \colon 1 \leq j \leq m_i }\right)
%\\\qquad\qquad\qquad\quad
%\begin{array}{rl}
}{
\infer[]{
\left(T_0 \cpar U \cpar P_i\right) \cseq \left(T_1 \cpar \tcontextn{i}{R^i_j \colon 1 \leq j \leq m_i }\right)
}{
\infer[]{
T_1 \cpar \tcontextn{i}{R^i_j \colon 1 \leq j \leq m_i }
}{
\infer[]{
\tcontextn{i}{T_1 \cpar R^i_j \colon 1 \leq j \leq m_i }
}{
\infer[]{
\tcontextn{i}{\cunit \colon 1 \leq j \leq m_i }
}{
\cunit
}}}}}
%\end{array}
\]
By Lemma~\ref{lemma:bound}, 
\[\size{ \tcontext{
 \left( U \cpar P_1 \right) \cseq \tcontextmn{i}{0}{R^i_j \colon j \in J^i} \cseq \tcontextmn{i}{1}{S^i_k \colon k \in K^i}
 \colon 1 \leq i \leq n 
} }  \preceq \size{ \left( U \cseq V \right) \cpar W }
\]
which are also upper bounds for
$
\size{ \tcontextmn{i}{0}{R^i_j \colon j \in J^i} 
}$
and
$\size{
\tcontextmn{i}{1}{S^i_k \colon k \in K^i}
}$.
Furthermore, $T_0 \not\equiv \cunit$ and $T_2 \not\equiv \cunit$ both $\occ{ T_0 } \mstrict \occ{ T_0 \cseq T_1 \cseq T_2 }$ and  $\occ{ T_2 } \mstrict \occ {T_0 \cseq T_1 \cseq T_2 }$
Hence the sizes of the above proofs of 
$T_2 \cpar \tcontextn{i}{ S^i_j \colon 1 \leq j \leq m_i }$ 
and 
\[
\left(T_0 \cseq T_1\right) \cpar \left(\left( U \cpar P_i \right) \cseq \tcontextn{i}{R^i_j \colon 1 \leq j \leq m_i }\right)
\]
 are strictly less than the size of the proof of $\left( T_0 \cseq T_1 \cseq T_2 \right) \cpar \left( U \cseq V \right) \cpar W$.
%\smallskip


% --- begin seq-medial.tex ---
%%%%%%%%%%%%%%%%%%%%% PRINCIPAL FOR SEQ and MEDIAL %%%%%%%%%
\item Consider the principal case for \textit{seq} where the bottommost rule of a proof is an instance of the \textit{suspend} rule of the form
\[
\infer[\mbox{, where $\vdash \left(P_0\cseq \mathopen{\wen x.} \left( P_1\cseq P_2  \right) \cseq P_3 \right) \cpar Q$ holds.}]{
\left(P_0\cseq \wen x. P_1\cseq \wen x. P_2  \cseq P_3 \right) \cpar Q
}{
\left(P_0\cseq \mathopen{\wen x.} \left( P_1\cseq P_2  \right) \cseq P_3 \right) \cpar Q
}
\] %and also $\nfv{x}{P}$, $\nfv{x}{S}$ and $\nfv{x}{Q}$.
% has a proof one rule shorter and, by Lemma~\ref{lemma:bound}, the size of the formula is bound.
By induction, there exist $U^0_i$ and $U^1_i$ such that $\vdash P_0\cpar U^0_i$ and $\vdash \left( \mathopen{\wen x.} \left( P_1\cseq P_2  \right) \cseq P_3 \right) \cpar U^1_i$ hold, for $1 \leq i \leq n$,
and $n$-ary killing context $\tcontext{}$ such that 
$
\vcenter{
\infer[]{Q}{
\tcontext{ U^0_i \cseq U^1_i \colon 1 \leq i \leq n }
}}
$.
Furthermore the size of the proof of $\left( \mathopen{\wen x.} \left( P_1\cseq P_2  \right) \cseq P_3 \right) \cpar U^1_i$ is bounded above by the size of the proof of $\left(P_0\cseq \wen x. P_1\cseq \wen x. P_2  \cseq P_3 \right) \cpar Q$.
%
By induction again, there exist $V^i_j$ and $W^i_j$ such that $\vdash \mathopen{\wen x.} \left( P_1\cseq P_2  \right) \cpar V^i_j$ and $\vdash P_3 \cpar W^i_j$, for $1 \leq j \leq m_i$, and $m_i$-ary killing context $\tcontextn{i}{}$ such that the following derivation holds.
$
\vcenter{
\infer[]{
U^1_i
}{
\tcontextn{i}{ V^i_j \cseq W^i_j \colon 1 \leq j \leq m_i }
}}
$.
Furthermore, the size of the proof of $\mathopen{\wen x.} \left( P_1\cseq P_2  \right) \cpar V^i_j$ is bounded by the size of the proof of $\left( \mathopen{\wen x.} \left( P_1\cseq P_2  \right) \cseq P_3 \right) \cpar U^1_i$. 
%
By applying the induction hypothesis again, there exist $R^i_j$ and $\hat{R}^i_j$ such that $\nfv{x}{R^i_j}$ and $\vdash \left(P_1\cseq P_2 \right) \cpar \hat{R}^i_j$ and either $R^i_j = \hat{R}^i_j$ or $R^i_j = \new x. \hat{R}^i_j$, and also
$
\vcenter{
\infer[]{
V^i_j
}{
 R^i_j
}}
$.
Furthermore, the size of the proof of $\left(P_1\cseq P_2 \right) \cpar \hat{R}^i_j$ is bounded above by the size of the proof of $\left( \mathopen{\wen x.} \left( P_1\cseq P_2  \right) \cseq P_3 \right) \cpar U^1_i$.
%
By a fourth induction, there exist $S^{i,j}_k$ and $T^{i,j}_k$ such that both $\vdash P_1\cpar S^{i,j}_k$ and $\vdash P_2  \cpar T^{i,j}_k$ hold, for $1 \leq k \leq \ell^{i,j}$, and $\ell^{i,j}$-ary killing context $\tcontextn{i,j}{}$ such that the following derivation holds:
\[
\infer[]{
\hat{R}^i_j
}{
\tcontextn{i,j}{ S^{i,j}_1 \cseq T^{i,j}_1, S^{i,j}_2 \cseq T^{i,j}_2, \hdots, S^{i,j}_{\ell^{i,j}} \cseq T^{i,j}_{\ell^{i,j}} }
}.
\]
By Lemma~\ref{lemma:medial}, there exists some $I^i_j \subseteq \left\{ 1 \hdots \ell^{i,j} \right\}$ and $J^i_j \subseteq \left\{ 1 \hdots \ell^{i,j} \right\}$ and killing contexts $\tcontextmn{i,j}{0}{}$ and $\tcontextmn{i,j}{1}{}$ such that
\[
\infer[]{
\hat{R}^i_j
}{
\infer[]{
\tcontextn{i,j}{ S^{i,j}_k \cseq T^{i,j}_k \colon 1 \leq k \leq \ell^{i,j} }
}{
\tcontextmn{i,j}{0}{S^{i,j}_k \colon k \in I^i_j } \cseq
\tcontextmn{i,j}{1}{T^{i,j}_k \colon k \in J^i_j }
}}.
\]
%
Define $\hat{S}^i_j$ and $\hat{T}^i_j$ as follows.
%\begin{itemize}
%\item 
If $R^i_j = \hat{R}^i_j$, then 
\[
\hat{S}^i_j = \tcontextmn{i,j}{0}{S^{i,j}_k \colon k \in I^i_j }
\mbox{ and }
\hat{T}^i_j = \tcontextmn{i,j}{1}{T^{i,j}_k \colon k \in J^i_j };
\]
and hence, we can construct the derivation 
\[
\infer[]{
R^i_j
}{
\tcontextmn{i,j}{0}{S^{i,j}_k \colon k \in I^i_j } \cseq
\tcontextmn{i,j}{1}{T^{i,j}_k \colon k \in J^i_j }
}
\] where the premise equals $\hat{S}^i_j  \cseq \hat{T}^i_j$.
%\item
If however $R^i_j = \new x. \hat{R}^i_j$, then define
\[
\hat{S}^i_j = \new x. \tcontextmn{i,j}{0}{S^{i,j}_k \colon k \in I^i_j }
\mbox{ and }
\hat{T}^i_j = \new x. \tcontextmn{i,j}{1}{T^{i,j}_k \colon k \in J^i_j };
\]
and hence, the derivation 
\[
\infer[]{
R^i_j
}{
\infer[]{
\mathopen{\new x.} \left(
 \tcontextmn{i,j}{0}{S^{i,j}_k \colon k \in I^i_j } \cseq
 \tcontextmn{i,j}{1}{T^{i,j}_k \colon k \in J^i_j }
\right)
}{
\hat{S}^i_j \cseq
\hat{T}^i_j
}}
\]
can be constructed.
%\end{itemize}
By Lemma~\ref{lemma:medial}, for some $K^i \subseteq \left\{ 1\hdots m_i \right\}$, $L^i \subseteq \left\{ 1\hdots m_i \right\}$ and killing contexts $\tcontextmn{i}{0}{}$ and $\tcontextmn{i}{1}{}$, we obtain the following derivation:
\[
\infer[]{
\tcontextn{i}{
    \hat{S}^i_j \cseq 
    \hat{T}^i_j \cseq 
    W^i_j \colon 1 \leq j \leq m_i
}
}{
\tcontextmn{i}{0}{
    \hat{S}^i_j
    \colon j \in K^i
} \cseq
\tcontextmn{i}{1}{
    \hat{T}^i_j \cseq 
    W^i_j \colon j \in L^i
}
}
\]
%
By using the above derivations we can construct the following derivation:
\[
\infer[]{
Q
}{
\infer[]{
\tcontext{ U^0_i \cseq U^1_i \colon 1 \leq i \leq n }
}{
\infer[]{
\tcontext{ U^0_i \cseq \tcontextn{i}{ V^i_j \cseq W^i_j \colon 1 \leq j \leq m_i } \colon 1 \leq i \leq n }
}{
\infer[]{
\tcontext{ U^0_i \cseq \tcontextn{i}{ %\tcontextmn{i}{j}{
                                    R^i_j %\colon 1 \leq k \leq g^i_j } 
                                                     \cseq W^i_j \colon 1 \leq j \leq m_i
                                    } \colon 1 \leq i \leq n }
}{
\infer[]{
\tcontext{
           U^0_i \cseq \tcontextn{i}{ %\tcontextmn{i}{j}{ 
                                                                        \hat{S}^i_j \cseq 
                                                                        \hat{T}^i_j 
     %\colon 1 \leq k \leq g^i_j } 
                                                     \cseq W^i_j \colon 1 \leq j \leq m_i
                                    }  \colon 1 \leq i \leq n }
}{
\tcontext{ 
           U^0_i \cseq
             \tcontextmn{i}{0}{
    \hat{S}^i_j 
    \colon j \in K^i
} \cseq
\tcontextmn{i}{1}{
    \hat{T}^i_j \cseq 
    W^i_j \colon j \in L^i
}
           \colon 1 \leq i \leq n  }
}}}}}
\]
%%%%%%%% Derivation done now proofs %%%%%%%%%%%%
\begin{comment}
We aim now to establish $
\vdash \left( P_0\cpar \wen x Q\right) \cpar
U^0_i \cseq
             \tcontextmn{i}{0}{
    \hat{S}^i_j 
    \colon j \in K^i
}
$ and
$\vdash \left( \wen x P_2  \cpar P_3 \right) \cpar
\tcontextmn{i}{1}{
    \hat{T}^i_j \cseq 
    W^i_j \colon j \in L^i
}
$.
\end{comment}

\noindent Consider whether the judgement $\vdash \wen x. P_1 \cpar \hat{S}^i_j$ holds.
We have two cases: in the first, $\hat{S}^i_j = \tcontextmn{i,j}{0}{S^{i,j}_k \colon k \in I^i_j }$ and $\nfv{x}{\hat{S}^i_j}$; in the second $\hat{S}^i_j = \new x. \tcontextmn{i,j}{0}{S^{i,j}_k \colon k \in I^i_j }$. In each case, one of the following derivations can be respectively constructed.
\[
\infer[]{
\wen x. P_1\cpar \tcontextmn{i,j}{0}{S^{i,j}_k \colon k \in I^i_j }
}{
\infer[]{
\new x. P_1\cpar \tcontextmn{i,j}{0}{S^{i,j}_k \colon k \in I^i_j }
}{
\mathopen{\new x.}\left( P_1\cpar \tcontextmn{i,j}{0}{S^{i,j}_k \colon k \in I^i_j } \right)
}}
\qquad\qquad
\infer[]{
\wen x. P_1\cpar \new x. \tcontextmn{i,j}{0}{S^{i,j}_k \colon k \in I^i_j }
}{
\mathopen{\new x.} \left( P_1\cpar \tcontextmn{i,j}{0}{S^{i,j}_k \colon k \in I^i_j } \right)
}
\]
Similarly, consider whether judgement $\vdash \wen x. P_2 \cpar \hat{T}^i_j$ holds.
Either we have 
\[
\hat{T}^i_j = \tcontextmn{i,j}{1}{T^{i,j}_k \colon k \in J^i_j } \mbox{ and } \nfv{x}{\hat{T}^i_j}; 
\]
or we have $\hat{T}^i_j = \new x. \tcontextmn{i,j}{1}{T^{i,j}_k \colon k \in J^i_j }$. In each case, one of the following derivations holds, respectively.
%we have that
\[
\infer[]{
\wen x. P_2\cpar \tcontextmn{i,j}{1}{T^{i,j}_k \colon k \in J^i_j }
}{
\infer[]{
\new x. P_2\cpar \tcontextmn{i,j}{1}{T^{i,j}_k \colon k \in J^i_j }
}{
\mathopen{\new x.} \left( P_2\cpar \tcontextmn{i,j}{1}{T^{i,j}_k \colon k \in J^i_j } \right)
}}
\qquad\qquad
\infer[]{
\wen x. P_2\cpar 
\hat{T}^i_j
%\new x. \tcontextmn{i,j}{1}{T^{i,j}_k \colon k \in J^i_j }
}{
\mathopen{\new x.} \left( P_2\cpar \tcontextmn{i,j}{1}{T^{i,j}_k \colon k \in J^i_j } \right)
}
\]
Thereby, by applying one of the above cases for each $i$ and $j$, 
the following two proofs exist.
\begin{comment}
\[
\begin{array}{rl}
\wen x. P_1\cpar \hat{S}^i_j
\longrightarrow&
\mathopen{\new x.} \left( P_1\cpar \tcontextmn{i,j}{0}{S^{i,j}_k \colon k \in I^i_j } \right)
\\
\longrightarrow&
\new x. \tcontextmn{i,j}{0}{ P_1\cpar S^{i,j}_k  \colon k \in I^i_j }
\longrightarrow
\new x. \tcontextmn{i,j}{0}{ \cunit \colon k \in I^i_j }
\longrightarrow
\cunit
\end{array}
\]
\end{comment}
\[
\infer[]{
\left( P_0 \cseq \wen x. P_1\right) \cpar
\left(
    U^0_i \cseq
    \tcontextmn{i}{0}{
         \hat{S}^i_j 
         \colon j \in K^i
    }
\right)
}{
\infer[]{
\left( P_0\cpar U^0_i \right) \cseq
\left(
    \wen x. P_1 \cpar
    \tcontextmn{i}{0}{
             \hat{S}^i_j 
      \colon j \in K^i
    }
\right)
}{
\infer[]{
    \wen x. P_1 \cpar
    \tcontextmn{i}{0}{
             \hat{S}^i_j 
      \colon j \in K^i
    }
}{
\infer[]{
    \tcontextmn{i}{0}{
         \wen x. P_1 \cpar
         \hat{S}^i_j 
      \colon j \in K^i
    }
}{
\infer[]{
    \tcontextmn{i}{0}{
         \mathopen{\new x.} \left( P_1\cpar \tcontextmn{i,j}{0}{S^{i,j}_k \colon k \in I^i_j } \right)
      \colon j \in K^i
    }
}{
\infer[]{
    \tcontextmn{i}{0}{
         \mathopen{\new x.} \tcontextmn{i,j}{0}{ P_1\cpar S^{i,j}_k  \colon k \in I^i_j }
      \colon j \in K^i
    }
}{
\infer[]{
    \tcontextmn{i}{0}{
         \new x. \tcontextmn{i,j}{0}{ \cunit \colon k \in I^i_j }
      \colon j \in K^i
    }
}{
\cunit
}}}}}}}
%Thereby, by applying one of the above derivations for each $i$ and $j$, the following proof can be constructed.
\qquad
\infer[]{
\left( \wen x. P_2  \cseq P_3 \right) \cpar
\left(
\tcontextmn{i}{1}{
    \hat{T}^i_j \cseq 
    W^i_j \colon j \in L^i
}\right)
}{
\infer[]{
\tcontextmn{i}{1}{
    \left( \wen x. P_2  \cseq P_3 \right) \cpar
    \left( \hat{T}^i_j \cseq W^i_j \right)
    \colon j \in L^i
}
}{
\infer[]{
\tcontextmn{i}{1}{
    \left( \wen x. P_2  \cpar \hat{T}^i_j\right) \cseq
    \left( P_3 \cpar W^i_j \right)
    \colon j \in L^i
}
}{
\infer[]{
\tcontextmn{i}{1}{
    \wen x. P_2  \cpar \hat{T}^i_j
    \colon j \in L^i
}
}{
\infer[]{
\tcontextmn{i}{1}{
    \mathopen{\new x.} \left( P_2\cpar \tcontextmn{i,j}{1}{T^{i,j}_k \colon k \in J^i_j } \right)
    \colon j \in L^i
}
}{
\infer[]{
\tcontextmn{i}{1}{
    \new x. \tcontextmn{i,j}{1}{ P_2\cpar T^{i,j}_k  \colon k \in J^i_j }
    \colon j \in L^i
}
}{
\infer[]{
\tcontextmn{i}{1}{
    \new x. \tcontextmn{i,j}{1}{ \cunit \colon k \in J^i_j }
    \colon j \in L^i
}
}{
\cunit
}}}}}}}
\]
Furthermore, by Lemma~\ref{lemma:bound}, 
\[
\size{ 
    U^0_i \cseq
    \tcontextmn{i}{0}{
         \hat{S}^i_j 
         \colon j \in K^i
    }
} \preceq \size{ Q }
\mbox{ and }
\size{ 
 \tcontextmn{i}{1}{
    \hat{T}^i_j \cseq 
    W^i_j \colon j \in L^i
} } \preceq \size{Q}.
\]
Hence, sizes
\[
\size{
\left( P_0 \cseq \wen x. P_1\right) \cpar
\left(
    U^0_i \cseq
    \tcontextmn{i}{0}{
         \hat{S}^i_j 
         \colon j \in K^i
    }
\right)
}
\mbox{ and }
\size {
\left( \wen x. P_2  \cseq P_3 \right) \cpar
\left(
\tcontextmn{i}{1}{
    \hat{T}^i_j \cseq 
    W^i_j \colon j \in L^i
}\right)
}
\]
are strictly bounded above by 
$\size{ \left(P_0\cseq \wen x. P_1\cseq \wen x. P_2  \cseq P_3 \right) \cpar Q }$, as required.
Cases for \textit{left wen} and \textit{right wen} rules are similar.
\begin{comment}
There are three similar principal case for \textit{seq} that do not correspond to cases in \textsf{MAV}~\cite{Horne2015}, induced by the \textit{suspend}, \textit{left wen} and \textit{right wen} rules.
We present here only the case induced by \textit{suspend}.
Consider the case where the bottommost rule of a proof is a \textit{suspend} rule that interferes with the \textit{seq} connective to which splitting is applied as follows.
\[
\left(P_0\cseq \wen x P_1\cseq \wen x P_2  \cseq P_3 \right) \cpar T
\longrightarrow
\left(P_0\cseq \wen x \left( P_1\cseq P_2  \right) \cseq P_3 \right) \cpar T
\]
In the above, it is assumed that $\vdash \left(P_0\cseq \wen x \left( P_1\cseq P_2  \right) \cseq P_3 \right) \cpar T$ and also $\nfv{x}{P}$, $\nfv{x}{S}$ and $\nfv{x}{T}$.
% has a proof one rule shorter and, by Lemma~\ref{lemma:bound}, the size of the formula is bound.
Hence, by induction, there exist $U^0_i$ and $U^1_i$ such that $\vdash P_0\cpar U^0_i$ and $\vdash \left( \wen x \left( P_1\cseq P_2  \right) \cseq P_3 \right) \cpar U^1_i$, for $1 \leq i \leq n$,
and $n$-ary killing context $\tcontext{}$ such that the following derivation holds.
$
T \longrightarrow \tcontext{ U^0_i \cseq U^1_i \colon 1 \leq i \leq n }
$.
Furthermore the size of the proof of $\left( \wen x \left( P_1\cseq P_2  \right) \cseq P_3 \right) \cpar U^1_i$ is bounded by the size of the proof of $\left(P_0\cseq \wen x P_1\cseq \wen x P_2  \cseq P_3 \right) \cpar T$.

By induction again, there exist $V^i_j$ and $W^i_j$ such that $\vdash \wen x \left( P_1\cseq P_2  \right) \cpar V^i_k$ and $\vdash P_3 \cpar W^i_j$, for $1 \leq j \leq m_i$, and $m_i$-ary killing context $\tcontextn{i}{}$ such that the following derivation holds.
$
U^1_i
\longrightarrow
\tcontextn{i}{ V^i_j \cseq W^i_j \colon 1 \leq j \leq m_i }
$.
Furthermore, the size of the proof of $\wen x \left( P_1\cseq P_2  \right) \cpar V^i_k$ is bounded by the size of the proof of $\left( \wen x \left( P_1\cseq P_2  \right) \cseq P_3 \right) \cpar U^1_i$. 

Since $\nfv{x}{T}$, by Lemma~\ref{lemma:nfv}, $\nfv{x}{V^i_k}$. By appling the induction hypothesis again, there exist $R^i_j_k$ and $\hat{R}^i_j_k$ such that $\vdash \left(P_1\cseq P_2 \right) \cpar \hat{R}^i_j_k$ and either $R^i_j_k = \hat{R}^i_j_k$ or $R^i_j_k = \new x \hat{R}^i_j_k$, for $1 \leq k \leq g^i_j$, and $g^i_j$-ary killing context $\tcontextmn{i}{j}{}$ such that 
$
V^i_j \longrightarrow \tcontextmn{i}{j}{ R^i_j_1, R^i_j_2, \hdots, R^i_j_{g^i_j} }
$.
Furthermore, the size of the proof of $\left(P_1\cseq P_2 \right) \cpar \hat{R}^i_j_k$ is bounded above by the size of the proof of $\left( \wen x \left( P_1\cseq P_2  \right) \cseq P_3 \right) \cpar U^1_i$.

By a fourth induction, there exist $S^{i,j,k}_k$ and $T^{i,j,k}_k$ such that $\vdash P_1\cpar S^{i,j,k}_k$ and $\vdash P_2  \cpar T^{i,j,k}_k$, for $1 \leq k \leq \ell^{i,j}_k$, and $\ell^{i,j}_k$-ary killing context $\tcontextmn{i,j}{k}{}$ such that the following derivation holds.
\[
\hat{R}^i_j_k
\longrightarrow
\tcontextmn{i,j}{k}{ S^{i,j,k}_1 \cseq T^{i,j,k}_1, S^{i,j,k}_2 \cseq T^{i,j,k}_2, \hdots, S^{i,j,k}_{\ell^{i,j}_k} \cseq T^{i,j,k}_{\ell^{i,j}_k} }
\]
In the case $R^i_j_k = \hat{R}^i_j_k$, by Lemma~\ref{lemma:medial}, the following derivation holds, where $I^i_j_k \subseteq \left\{ 1 \hdots \ell^{i,j}_k \right\}$ and $J^i_j_k \subseteq \left\{ 1 \hdots \ell^{i,j}_k \right\}$ and killing contexts $\tcontextmn{i,j}{k,0}{}$ and $\tcontextmn{i,j}{k,1}{}$.
\[
\hat{R}^i_j_k
\longrightarrow
\tcontextmn{i,j}{k}{ S^{i,j,k}_k \cseq T^{i,j,k}_k \colon 1 \leq k \leq \ell^{i,j}_k }
\longrightarrow
\tcontextmn{i,j}{k,0}{S^{i,j,k}_k \colon k \in I^i_j_k } \cseq
\tcontextmn{i,j}{k,1}{T^{i,j,k}_k \colon k \in J^i_j_k }
\]
In the case $R^i_j_k = \new x \hat{R}^i_j_k$, by Lemma~\ref{lemma:medial}, the following derivation holds, where $I^i_j_k \subseteq \left\{ 1 \hdots \ell^{i,j}_k \right\}$ and $J^i_j_k \subseteq \left\{ 1 \hdots \ell^{i,j}_k \right\}$ and killing contexts $\tcontextmn{i,j}{k,0}{}$ and $\tcontextmn{i,j}{k,1}{}$.
\[
\begin{array}{rl}
\new x \hat{R}^i_j_k
 \longrightarrow&
\new x \tcontextmn{i,j}{k}{ S^{i,j,k}_k \cseq T^{i,j,k}_k \colon 1 \leq k \leq \ell^{i,j}_k }
\\
\longrightarrow&
\new x
\left(
\tcontextmn{i,j}{k}{S^{i,j,k,0}_k \colon k \in I^i_j_k } \cseq
\tcontextmn{i,j}{k}{T^{i,j,k,1}_k \colon k \in J^i_j_k }
\right)
\\
\longrightarrow&
\new x
\tcontextmn{i,j}{k}{S^{i,j,k,0}_k \colon k \in I^i_j_k } \cseq
\new x
\tcontextmn{i,j}{k}{T^{i,j,k,1}_k \colon k \in J^i_j_k }
\end{array}
\]
Define the formulae
$
\hat{S}^i_j_k = \tcontextmn{i,j}{k,0}{S^{i,j,k}_k \colon k \in I^i_j_k }
$ and 
$
\hat{T}^i_j_k = \tcontextmn{i,j}{k,1}{T^{i,j,k}_k \colon k \in J^i_j_k }
$.
If $R^i_j_k = \hat{R}^i_j_k$, then define $\hat{S}^i_j_k = \hat{S}^i_j_k $
and $\hat{T}^i_j_k = \hat{T}^i_j_k$.
	If $R^i_j_k = \new x \hat{R}^i_j_k$, then define $\hat{S}^i_j_k = \new x \hat{S}^i_j_k$
and $\hat{T}^i_j_k = \new x \hat{T}^i_j_k$.
Hence we construct the following derivation, using Lemma~\ref{lemma:medial} again, for some $K^i_j \subseteq \left\{ 1\hdots g^i_j \right\}$ and $L^i_j \subseteq \left\{ 1\hdots g^i_j \right\}$ and killing contexts $\tcontextmn{i}{j,0}{}$ and $\tcontextmn{i}{j,1}{}$.
\[
\begin{array}{rl}
T \longrightarrow&
\tcontext{ U^0_i \cseq U^1_i \colon 1 \leq i \leq n }
\\
\longrightarrow&
\tcontext{ U^0_i \cseq \tcontextn{i}{ V^i_j \cseq W^i_j \colon 1 \leq j \leq m_i } \colon 1 \leq i \leq n }
\\
\longrightarrow&
\tcontext{ U^0_i \cseq \tcontextn{i}{ \tcontextmn{i}{j}{ R^i_j_k \colon 1 \leq k \leq g^i_j } 
                                                     \cseq W^i_j \colon 1 \leq j \leq m_i
                                    } \colon 1 \leq i \leq n }
\\
\longrightarrow&
\tcontext{
           U^0_i \cseq \tcontextn{i}{ \tcontextmn{i}{j}{ 
                                                                        \hat{S}^i_j_k \cseq 
                                                                        \hat{T}^i_j_k 
                                                                        \colon 1 \leq k \leq g^i_j } 
                                                     \cseq W^i_j \colon 1 \leq j \leq m_i
                                    }  \colon 1 \leq i \leq n }
\\
\longrightarrow&
\tcontext{ \begin{array}{l}
           U^0_i \cseq \tcontextn{i}{ \tcontextmn{i}{j,0}{ \hat{S}^i_j_k \colon k \in K^i_j } \colon 1 \leq j \leq m_i } \cseq \\
           \tcontextn{i}{ \tcontextmn{i}{j,1}{ \hat{T}^i_j_k \colon k \in L^i_j } \cseq W^i_j \colon 1 \leq j \leq m_i }
           \end{array}
           \colon 1 \leq i \leq n  }
\end{array}
\]
We aim now to establish $\vdash \left( P_0\cpar \wen x Q\right) \cpar U^0_i \cseq \tcontextn{i}{ \tcontextmn{i}{j,0}{ \hat{S}^i_j_k \colon k \in K^i_j } \colon 1 \leq j \leq m_i }$ and
$\vdash \left( \wen x P_2  \cpar P_3 \right) \cpar \tcontextn{i}{ \tcontextmn{i}{j,1}{ \hat{T}^i_j_k \colon k \in L^i_j } \cseq W^i_j \colon 1 \leq j \leq m_i }$.
Firstly, observe that $\hat{S}^i_j_k = \hat{S}^i_j_k$ then
$
\wen x P_1\cpar \hat{S}^i_j_k
\longrightarrow 
\new x P_1\cpar \hat{S}^i_j_k
\longrightarrow 
\new x \left( P_1\cpar \hat{S}^i_j_k \right)
$, since $\nfv{x}{T}$, by Lemma~\ref{lemma:nfv}, $\nfv{x}{U^{i,j}_k}$.
Also, if $\hat{S}^i_j_k = \new x \hat{S}^i_j_k$ then 
$
\wen x P_1\cpar \new x \hat{S}^i_j_k
\longrightarrow
\new x \left( P_1\cpar \hat{S}^i_j_k \right)
$.
Furthermore, the following proof can be constructed.
\[
%\begin{array}{rl}
P_1\cpar \tcontextmn{i,j}{k}{S^{i,j,k}_k \colon 1 \leq k \leq \ell^{i,j}_k }
\longrightarrow
\tcontextmn{i,j}{k}{ P_1\cpar S^{i,j,k}_k \colon 1 \leq k \leq \ell^{i,j}_k }
%\\
\longrightarrow
\tcontextmn{i,j}{k}{ \cunit \colon 1 \leq k \leq \ell^{i,j}_k }
%\\
\longrightarrow
\cunit
%\end{array}
\]
Thereby, in either case $\vdash \wen x P_1\cpar \hat{S}^i_j_k$.
A similar argument establishes 
$
\vdash \wen x P_2  \cpar \hat{T}^i_j_k
$.

By applying the above proofs appropriately, the following two proofs can be constructed.
\[
\begin{array}{l}
\left( P_0\cseq \wen x P_1\right)  \cpar \left( U^0_i \cseq \tcontextn{i}{ \tcontextmn{i}{j,0}{ \hat{S}^i_j_k \colon k \in K^i_j } \colon 1 \leq j \leq m_i } \right)
\\\qquad\qquad\qquad\qquad\qquad
\begin{array}{rl}
\longrightarrow&
\left( P_0\cpar U^0_i  \right)  \cseq 
\left(  \wen x P_1\cpar \tcontextn{i}{ \tcontextmn{i}{j}{ \hat{S}^i_j_k \colon  k \in K^i_j } \colon 1 \leq j \leq m_i } \right)
\\
\longrightarrow&
\left(  \wen x P_1\cpar \tcontextn{i}{ \tcontextmn{i}{j}{ \hat{S}^i_j_k \colon  k \in K^i_j } \colon 1 \leq j \leq m_i } \right)
\\
\longrightarrow&
\left(  \tcontextn{i}{ \tcontextmn{i}{j}{ \wen x P_1\cpar \hat{S}^i_j_k \colon  k \in K^i_j } \colon 1 \leq j \leq m_i } \right)
\\
\longrightarrow&
\left(  \tcontextn{i}{ \tcontextmn{i}{j}{ \cunit \colon  k \in K^i_j } \colon 1 \leq j \leq m_i } \right)
\longrightarrow
\cunit
\end{array}
\end{array}
\]
\[
\begin{array}{l}
\left( \wen x P_2  \cseq P_3 \right) \cpar \tcontextn{i}{ \tcontextmn{i}{j}{ \hat{T}^i_j_k \colon k \in L^i_j } \cseq W^i_j \colon 1 \leq j \leq m_i }
\\\qquad\qquad\qquad\qquad\qquad
\begin{array}{rl}
\longrightarrow&
\tcontextn{i}{ \left( \wen x P_2  \cseq P_3 \right) \cpar \tcontextmn{i}{j}{ \hat{T}^i_j_k \colon k \in L^i_j } \cseq W^i_j \colon 1 \leq j \leq m_i }
\\
\longrightarrow&
\tcontextn{i}{ 
               \left( \wen x P_2  \cpar \tcontextmn{i}{j}{ \hat{T}^i_j_k \colon k \in L^i_j } \right)
                \cseq \left(P_3 \cpar W^i_j\right) \colon 1 \leq j \leq m_i }
\\
\longrightarrow&
\tcontextn{i}{ \left( \wen x P_2  \cpar \tcontextmn{i}{j}{ \hat{T}^i_j_k \colon k \in L^i_j } \right) \colon 1 \leq j \leq m_i }
\\
\longrightarrow&
\tcontextn{i}{ \tcontextmn{i}{j}{ \wen x P_2  \cpar \hat{T}^i_j_k \colon k \in L^i_j } \colon 1 \leq j \leq m_i }
\\
\longrightarrow&
\tcontextn{i}{ \tcontextmn{i}{j}{ \cunit \colon k \in L^i_j } \colon 1 \leq j \leq m_i }
\longrightarrow
\cunit
\end{array}
\end{array}
\]
Furthermore, by Lemma~\ref{lemma:bound}, the following inequalities holds.
\[
\begin{array}{rl}
\size{ \left( U^0_i \cseq \tcontextn{i}{ \tcontextmn{i}{j,0}{ \hat{S}^i_j_k \colon k \in K^i_j } \colon 1 \leq j \leq m_i } \right) } &\preceq \size{ T }
\\
\size{ \tcontextn{i}{ \tcontextmn{i}{j}{ \hat{T}^i_j_k \colon k \in L^i_j } \cseq W^i_j \colon 1 \leq j \leq m_i } } &\preceq \size{T}
\end{array}
\]
Hence, by Lemma~\ref{lemma:multisets}, the size of the formulae
$\left( \wen x P_2  \cseq P_3 \right) \cpar \tcontextn{i}{ \tcontextmn{i}{j}{ \hat{T}^i_j_k \colon k \in L^i_j } \cseq W^i_j \colon 1 \leq j \leq m_i }$ and
$\left( P_0\cseq \wen x P_1\right)  \cpar \left( U^0_i \cseq \tcontextn{i}{ \tcontextmn{i}{j,0}{ \hat{S}^i_j_k \colon k \in K^i_j } \colon 1 \leq j \leq m_i } \right)$
  are strictly bounded above by the size of 
$\left(P_0\cseq \wen x P_1\cseq \wen x P_2  \cseq P_3 \right) \cpar T$, as required.
\end{comment}


% --- end seq-medial.tex ---

\end{enumerate}


\item \textbf{Principal case for times.}
There is only one principal case for \textit{times}, which does not differ significantly from the corresponding case in \textsf{BV} and its extensions. A proof may begin with an instance of the \textit{switch} rule of the form
\[
\infer[\mbox{where $\vdash \left(T_0 \tensor U_0 \tensor \left( \left(T_1 \tensor U_1\right) \cpar V \right)\right) \cpar W$,}]{
\left( T_0 \tensor T_1 \tensor U_0 \tensor U_1 \right) \cpar V \cpar W
}{
\left(T_0 \tensor U_0 \tensor \left( \left(T_1 \tensor U_1\right) \cpar V \right)\right) \cpar W
}
\]
such that $T_0 \tensor U_0 \not\equiv \cunit$ and $V \not\equiv \cunit$ (otherwise the \textit{switch} rule cannot be applied), and also $T_0 \tensor T_1 \not\equiv \cunit$ and $U_0 \tensor U_1 \not\equiv \cunit$ (otherwise splitting holds trivially).
%
By the induction hypothesis, there exist $R_i$ and $S_i$ such that $\vdash \left(T_0 \tensor U_0\right) \cpar R_i$ and $\vdash \left(T_1 \tensor U_1\right) \cpar V \cpar S_i$ hold, for $1 \leq i \leq n$, and an $n$-ary killing context $\tcontext{}$ such that 
derivation $\vcenter{\infer{W}{\tcontext{ R_1 \cpar S_1, \hdots, R_n \cpar S_n }}}$ holds.
Furthermore $\size{\left(T_0 \tensor U_0\right) \cpar R_i}$ and $\size{\left(T_1 \tensor U_1\right) \cpar V \cpar S_i}$ are bounded above by $\size{\left(T_0 \tensor U_0 \tensor \left( \left(T_1 \tensor U_1\right) \cpar V \right)\right) \cpar W}$.
%
Hence, by the induction hypothesis twice there exist formulae $P^{i,0}_j$, $Q^{i,0}_j$, $P^{i,1}_k$ and $Q^{i,1}_k$ such that $\vdash T_0 \cpar P^{i,0}_j$, $\vdash U_0 \cpar Q^{i,0}_j$, $\vdash T_1 \cpar P^{i,1}_k$ and $\vdash U_1 \cpar Q^{i,1}_k$, for $1 \leq j \leq m^0_i$ and $1 \leq k \leq m^1_i$, and $m^0_i$-ary killing context $\tcontextmn{0}{i}{}$ and $m^1_i$-ary killing context $\tcontextmn{1}{i}{}$ such that derivations
\[
\vcenter{
\infer[]{
R_i}{
\tcontextmn{0}{i}{P^{i,0}_j \cpar Q^{i,0}_j \colon 1 \leq j \leq m^0_i}
}
}
\quad
\mbox{ and }
\quad
\vcenter{
\infer[]{
V \cpar S_i
}{
\tcontextmn{1}{i}{P^{i,1}_k \cpar Q^{i,1}_k \colon 1 \leq k \leq m^1_i}
}}
\] 
can be constructed.
%Let $I \subset \left\{1, \hdots, n\right\}$ be such that if $R_i \cpar S_i \cpar V \longrightarrow \top$ then $i \not\in I$ and $\tcontext{V \cpar R_i \cpar S_i \colon 1 \leq i \leq n} \longrightarrow \tcontextp{ V \cpar R_i \cpar S_i \colon i \in I}$
Thereby the following derivation can be constructed.
\[
\infer[]{
V \cpar W
}{
\infer[]{
V \cpar \tcontext{ R_i \cpar S_i \colon 1 \leq i \leq n } 
}{
\infer[]{
\tcontext{ R_i \cpar V \cpar S_i \colon 1 \leq i \leq n }
}{
\infer[]{
\tcontext{
%  \begin{array}{rl}
  \tcontextmn{0}{i}{P^{i,0}_j \cpar Q^{i,0}_j \colon 1 \leq j \leq m^0_i}
   \cpar
  \tcontextmn{1}{i}{P^{i,1}_k \cpar Q^{i,1}_k \colon 1 \leq k \leq m^1_i}
% \end{array}
\colon 1 \leq i \leq n }
}{
\infer[]{
 \tcontext{ \tcontextmn{1}{i}{
%  \begin{array}{l}
  \tcontextmn{0}{i}{P^{i,0}_j \cpar Q^{i,0}_j \colon 1 \leq j \leq m^0_i}
  \cpar
  P^{i,1}_k \cpar Q^{i,1}_k \colon 1 \leq k \leq m^1_i
%  \end{array}
} \colon 1 \leq i \leq n }
}{
 \tcontext{ \tcontextmn{1}{i}{
% \begin{array}{l}
  \tcontextmn{0}{i}{
%  \begin{array}{l}
  P^{i,0}_j \cpar P^{i,1}_k \cpar Q^{i,0}_j \cpar Q^{i,1}_k
  \colon 1 \leq j \leq m^0_i
%  \end{array}
  }
 \colon 1 \leq k \leq m^1_i
% \end{array}
 } \colon 1 \leq i \leq n }
}}}}}
\]
Now observe that the following two proofs can be constructed.
\[
\infer[]{
\left(T_0 \tensor T_1\right)
\cpar
P^{i,0}_j \cpar P^{i,1}_k
}{
\infer[]{
\left(T_0 \cpar P^{i,0}_j\right) \tensor \left(T_1 \cpar P^{i,1}_k\right)
}{
\cunit
}}
\qquad\qquad\qquad\qquad
\infer[]{
\left(U_0 \tensor U_1\right)
\cpar
Q^{i,0}_j \cpar Q^{i,1}_k
}{
\infer[]{
\left(U_0 \cpar Q^{i,0}_j\right) \tensor \left(U_1 \cpar Q^{i,1}_k\right)
}{
 \cunit
}}
\]
Furthermore,
%we have that 
$\occ{ T_0 \tensor T_1 } \mstrict \occ{ T_0 \tensor T_1 \tensor U_0 \tensor U_1 }$
and 
$\occ{ U_0 \tensor U_1 } \mstrict \occ{ T_0 \tensor T_1 \tensor U_0 \tensor U_1 }$,
since $T_0 \tensor T_1 \not\equiv \cunit$ and $U_0 \tensor U_1 \not\equiv \cunit$.
Also, by Lemma~\ref{lemma:bound}, the following inequality holds.
\[
\size{
 \tcontext{ \tcontextmn{1}{i}{
% \begin{array}{l}
  \tcontextmn{0}{i}{
%  \begin{array}{l}
  P^{i,0}_j \cpar P^{i,1}_k \cpar Q^{i,0}_j \cpar Q^{i,1}_k
  \colon 1 \leq j \leq m^0_i
%  \end{array}
  }
 \colon 1 \leq k \leq m^1_i
% \end{array}
 } \colon 1 \leq i \leq n }
} \preceq \size{ V \cpar W }
\]
Hence both
$
\size{ P^{i,0}_j \cpar P^{i,1}_k }
\preceq \size{ V \cpar W }$
and $
\size{ Q^{i,0}_j \cpar Q^{i,1}_k }
\preceq \size{ V \cpar W }$ hold.
Thereby the size of each of the above proofs is strictly bounded above by the size of the proof of $\left(T_0 \tensor T_1 \tensor U_0 \tensor U_1 \right) \cpar V \cpar W$.
%\smallskip



\item \textbf{Principal cases for with.}
There are three forms of principal case where the \textit{with} operator is directly involved in the bottommost rules.
Note that in \textsf{MAV} the \textit{with} operator is separated from the core splitting lemma, much like universal quantification in this paper.
However, in the case of \textsf{MAV1} the \textit{left name} and \textit{right name} rules introduce inter-dependencies between nominals and \textit{with}, forcing cases for \textit{with} to be checked in this lemma.

\begin{enumerate}[label*=\textbf{.\arabic*}]

\item Consider the principal case involving the \textit{extrude} rule.
In this case, the bottommost rule is of the form 
\[
\infer[\mbox{where $\vdash \left(P  \cpar R\right)\wwith\left( Q \cpar R\right) \cpar S$ holds.}]{
\left(P \wwith Q\right) \cpar R \cpar S
}{
\left(P  \cpar R\right)\wwith\left( Q \cpar R\right) \cpar S
}
\]
Now, by the induction hypothesis, 
since $\vdash \left(P  \cpar R\right)\wwith\left( Q \cpar R\right) \cpar S$ holds, we have that 
$\vdash P  \cpar R \cpar S$ and $\vdash Q \cpar R \cpar S$ hold, as required.
%with proof strictly smaller than the proof of $\left(P  \cpar R\right)\wwith\left( Q \cpar R\right) \cpar S$, which, in turn, has a one step shorter proof than $\left(P \wwith Q\right) \cpar R \cpar S$, as required.


\item Consider the principal case involving the \textit{left name} rule.
In this case, the bottommost rule is of the form 
\[
\infer[\mbox{, where $\nfv{x}{Q}$, such that $\vdash \mathopen{\wen x.}\left( P \wwith Q \right) \cpar R$.}]{
\left( \wen x.P \wwith Q \right) \cpar R
}{
\mathopen{\wen x.}\left( P \wwith Q \right) \cpar R
}
\] 
By the induction hypothesis, there exist $S$ and $\hat{S}$ such that
$\vcenter{\infer[]{R}{S}}$ and $\nfv{x}{S}$ and $\vdash \left(P \wwith Q\right) \cpar \hat{S}$ 
and either $S = \hat{S}$ or $S = \mathopen{\new x.}\hat{S}$. Furthermore, the size of the proof of $\left(P \wwith Q\right) \cpar \hat{S}$ is strictly less than the size of the proof of $\left( \wen x.P \wwith Q \right) \cpar R$, since the \textit{wen} count strictly decreases, and by Lemma~\ref{lemma:bound}, $\size{\hat{S}} \leq \size{R}$.
By the induction hypothesis again, $\vdash P \cpar \hat{S}$ and $\vdash Q \cpar \hat{S}$ hold.
% with strictly smaller proofs than $\left(P \wwith Q\right) \cpar \hat{S}$.

Now if $S = \hat{S}$ then $\nfv{x}{\hat{S}}$ and $\vdash Q \cpar S$ holds immediately, whereas 
$\vdash \wen x.P \cpar R$ is proved as below left. Otherwise, $S = \new x.\hat{S}$ and
$\vdash \wen x.P \cpar R$ is proved in the middle derivation below, whereas $\vdash  Q \cpar S$ is proved in the right derivation below.
\[
\infer[]{
\wen x.P \cpar R 
}{
\infer[]{
\wen x.P \cpar \hat S
}{
\infer[]
{\wen x.\left(P \cpar \hat S \right)}
{
\infer[]{
\mathopen{\new x.}\left(P \cpar \hat{S} \right)
}{
\infer[]{
\mathopen{\new x.}\cunit
}{
\cunit
}}}}}
\qquad \qquad
\infer[]{
\wen x.P \cpar R 
}{
\infer[]{
\wen x.P \cpar \new x. \hat S
}{
\infer[]{
\mathopen{\new x.}\left(P \cpar \hat{S} \right)
}{
\infer[]{
\mathopen{\new x.}\cunit
}{
\cunit
}}}}
\qquad
\qquad
\infer[.]{
Q \cpar \wen x.\hat{S}
}{
\infer[]{
\mathopen{\wen x.}\left( Q \cpar \hat{S} \right)
}{
\infer[]{
 \mathopen{\new x.}\left( Q \cpar \hat{S} \right)
}{
\infer[]{
 \new x.\cunit 
}{
 \cunit
}}}}
\]

%so 
%\[
%\infer[.]{
%\wen x. P \cpar \hat{S}}{
%\infer[]{
%\mathopen{\wen x.}\left( P \cpar \hat{S} \right)
%}{
%\mathopen{\new x.}\left( P \cpar \hat{S} \right)
%}}
%\]
%Otherwise $S = \new x.\hat{S}$ so $
%\vcenter{\infer[]{
%\wen x. P \cpar \new x.\hat{S}
%}{
%\mathopen{\new x.} \left( P \cpar \hat{S} \right)
%}}$.
%Hence $
%\vcenter{\infer[]{
%\wen x. P \cpar S
%}{
%\mathopen{\new x.} \left( P \cpar \hat{S} \right)
%}}$,
%using which we construct: % can construct proof:
%\[
%\infer[]{
%\wen x.P \cpar R 
%}{
%\infer[.]{
%\wen x.P \cpar S
%}{
%\infer[]{
%\mathopen{\new x.}\left(P \cpar \hat{S} \right)
%}{
%\infer[]{
%\mathopen{\new x.}\cunit
%}{
%\cunit
%}}}}
%\]

%If $S = \hat{S}$ then $\vdash Q \cpar S$ immediately.
%Otherwise $S = \wen x.\hat{S}$, in which case, since $\nfv{x}{Q}$, by the \textit{left wen} rule, we have the proof
%\[
%\infer[.]{
%Q \cpar \wen x.\hat{S}
%}{
%\infer[]{
%\mathopen{\wen x.}\left( Q \cpar \hat{S} \right)
%}{
%\infer[]{
% \mathopen{\new x.}\left( Q \cpar \hat{S} \right)
%}{
%\infer[]{
% \new x.\cunit 
%}{
% \cunit
%}}}}
%\]
Hence, in either case, $\vdash Q \cpar S$ 
and since $
\vcenter{
\infer[]{
Q \cpar R
}{
Q \cpar S
}}$,
we have that $\vdash Q \cpar R$ holds.
Thereby $\vdash \wen x.P \cpar R$ and $\vdash Q \cpar R$ hold, as required.
% and furthermore $\occ{\wen x.P \cpar R} \prec \occ{\left(P \wwith Q\right) \cpar R \cpar S}$ and $\occ{Q \cpar R} \prec \occ{\left(P \wwith Q\right) \cpar R \cpar S}$, as required.
The case for the \textit{left name} rule, where $\new$ replaces $\wen$ is similar; as are the cases for the \textit{right name} and \textit{with name} rules.


\begin{comment} Case for \textit{with name}, omitted for readability.
Consider the principal case involving the \textit{with name} rule.
In this case, the bottommost rule is of the form $\left( \wen x.P \wwith \wen x.Q \right) \cpar R \longrightarrow \mathopen{\wen x.}\left( P \wwith Q \right) \cpar R$ such that $\vdash \mathopen{\wen x.}\left( P \wwith Q \right) \cpar R$.
By the induction hypothesis, there exist $S$ and $\hat{S}$ such that $R \longrightarrow S$ and $\nfv{x}{S}$ and $\vdash \left(P \wwith Q\right) \cpar \hat{S}$ 
and either $S = \hat{S}$ or $S = \mathopen{\new x.}\hat{S}$. Furthermore, the size of the proof of $\left(P \wwith Q\right) \cpar \hat{S}$ is strictly less than the size of the proof of $\left(P \wwith Q\right) \cpar R \cpar S$.
By the induction hypothesis again, $\vdash P \cpar \hat{S}$ and $\vdash Q \cpar \hat{S}$ hold.
% with strictly smaller proofs than $\left(P \wwith Q\right) \cpar \hat{S}$.
Now if $S = \hat{S}$ then $\nfv{x}{\hat{S}}$ and so $\wen x. P \cpar \hat{S} \longrightarrow \mathopen{\wen x.}\left( P \cpar \hat{S} \right) \longrightarrow \mathopen{\new x.}\left( P \cpar \hat{S} \right)$.
Otherwise $S = \new x.\hat{S}$ so $\wen x. P \cpar \new x.\hat{S} \longrightarrow \mathopen{\new x.} \left( P \cpar \hat{S} \right)$.
Hence $\wen x. P \cpar S \longrightarrow \mathopen{\new x.} \left( P \cpar \hat{S} \right)$.
By a similar argument, $\wen x. Q \cpar S \longrightarrow \mathopen{\new x.} \left( Q \cpar \hat{S} \right)$.
Thereby we can construct the following two proofs as follows.
\[
\begin{array}{rl}
\wen x.P \cpar R 
\longrightarrow&
\wen x.P \cpar S
\\
\longrightarrow&
\mathopen{\new x.}\left(P \cpar \hat{S} \right)
\longrightarrow
\mathopen{\new x.}\cunit
\longrightarrow 
\cunit
\end{array}
\begin{array}{rl}
\wen x.Q \cpar R 
\longrightarrow&
\wen x.Q \cpar S
\\
\longrightarrow&
\mathopen{\new x.}\left(Q \cpar \hat{S} \right)
\longrightarrow
\mathopen{\new x.}\cunit
\longrightarrow
\cunit
\end{array}
\]
Furthermore $\occ{\wen x.P \cpar R} \prec \occ{\left(P \wwith Q\right) \cpar R \cpar S}$
and $\occ{Q \cpar R} \prec \occ{\left(P \wwith Q\right) \cpar R \cpar S}$, hence the size of proofs strictly decrease as requried.
\smallskip
\end{comment}


\item Consider the principal case involving the \textit{medial} rule.
In this case, the bottommost rule of a proof is of the form 
\[
\infer[\mbox{such that $\vdash \left(\left( P \wwith R \right) \cseq \left( Q \wwith S \right)\right) \cpar W$ holds.}]{
\left(\left( P \cseq Q \right) \wwith \left( R \cseq S \right)\right) \cpar W
}{
\left(\left( P \wwith R \right) \cseq \left( Q \wwith S \right)\right) \cpar W
}
\]
By the induction hypothesis, for $1 \leq i \leq n$ there exists $U_i$ and $V_i$ such that $\vdash \left( P \wwith R \right) \cpar U_i$ and $\vdash \left( Q \wwith S \right) \cpar V_i$ hold, and $n$-ary killing context $\tcontext{}$ such that 
$
\vcenter{\infer[]{W}{\tcontext{U_i \cseq V_i \colon 1 \leq i \leq n}}}$.
Furthermore, the size of the proofs of $\left( P \wwith R \right) \cpar U_i$ and $\left( Q \wwith S \right) \cpar V_i$ are strictly less than the size of the proof of $\left(\left( P \wwith R \right) \cseq \left( Q \wwith S \right)\right) \cpar W$. Hence by the induction hypothesis again, $\vdash P \cpar U_i$, $\vdash R \cpar U_i$, $\vdash Q \cpar V_i$ and $\vdash S \cpar V_i$.
Hence we can construct the following two proofs, as required.
\[
\infer[]{
\left( P \cseq Q \right) \cpar W
}{
\infer[]{
\left( P \cseq Q \right) \cpar \tcontext{U_i \cseq V_i \colon 1 \leq i \leq n}
}{
\infer[]{
\tcontext{\left( P \cseq Q \right) \cpar \left(U_i \cseq V_i\right) \colon 1 \leq i \leq n}
}{
\infer[]{
\tcontext{\left( P \cpar U_i \right) \cseq \left(Q \cpar V_i\right) \colon 1 \leq i \leq n}
}{
\infer[]{
\tcontext{\cunit  \colon 1 \leq i \leq n}
}{
\cunit
}}}}}
\qquad
\infer[]{
\left( R \cseq S \right) \cpar W
}{
\infer[]{
\left( R \cseq S \right) \cpar \tcontext{U_i \cseq V_i \colon 1 \leq i \leq n}
}{
\infer[]{
\tcontext{\left( R \cseq S \right) \cpar \left(U_i \cseq V_i\right) \colon 1 \leq i \leq n}
}{
\infer[]{
\tcontext{\left( R \cpar U_i \right) \cseq \left(S \cpar V_i\right) \colon 1 \leq i \leq n}
}{
\infer[]{
\tcontext{\cunit \colon 1 \leq i \leq n}
}{
\cunit
}}}}}
\]
%Furthermore, $\occ{\left( P \cseq Q \right) \cpar W} \prec \occ{\left(\left( P \cseq Q \right) \wwith \left( R \cseq S \right)\right) \cpar W}$ and $\occ{\left( R \cseq S \right) \cpar W} \prec \occ{\left(\left( P \cseq Q \right) \wwith \left( R \cseq S \right)\right) \cpar W}$; hence the size of the proofs strictly decrease as required.

\end{enumerate}


% --- begin equivariance.tex ---
\item \textbf{Commutative cases induced by equivariance.}
%
%two rules that do not involving the same connectives swap with each other.
There are certain commutative cases induced by the \textit{equivariance} rule for nominal quantifiers.
% that are unlike cases that appear in splitting lemmas in the literature.
%\begin{remark}
These are the cases that force the rules \textit{all name}, \textit{with name}, \textit{left name} and \textit{right name} to be included.
Notice also that \textit{equivariance} for \textit{new} is required when handling the case induced by \textit{equivariance} for \textit{wen}; hence \textit{equivariance} for both nominal quantifiers must be explicit structural rules rather than properties derived from each other.
%\end{remark}

\begin{enumerate}[label*=\textbf{.\arabic*}]
\item Consider the commutative case for \textit{wen} where the bottommost rule of a proof is an instance of the \textit{close} rule of following form
\[
\infer[\mbox{,where $\vdash \mathopen{\new y.}\left( \wen x. P \cpar Q \right) \cpar R$, $\nfv{y}{R}$ and $\nfv{x}{R}$.} ]{
\wen x. \wen y. P \cpar \new y. Q \cpar R
}{
\mathopen{\new y.}\left( \wen x. P \cpar Q \right) \cpar R
}
\]
Notice that $\wen x$ is the principal connective but the \textit{close} rule is applied to $\wen y$ behind the principal connective.
Thus we desire some formula $R'$ such that $
\vcenter{
\infer[]{
\new y. Q \cpar R
}{R'}}$
and $\nfv{x}{R'}$ and either $\vdash \wen y. P \cpar R'$ or there exists $Q'$ such that $R' = \new x.Q'$ and $\vdash \wen y. P \cpar Q'$, and the size of $ \wen y. P \cpar R'$ is strictly smaller than $\wen x. \wen y. P \cpar \new y. Q \cpar R$.
%
By the induction hypothesis, there exist $S$ and $T$ such that $\nfv{y}{S}$ and $\vdash \wen x. P \cpar Q \cpar T$ and either $S = T$ or $S = \wen y. T$ and the derivation
$
\vcenter{
\infer{R}{S}}
$ holds.
Furthermore the size of the proof of $\wen x. P \cpar Q \cpar T$ is bounded above by the size of the proof of $\mathopen{\new y.}\left( \wen x.P \cpar Q \right) \cpar R$; hence strictly bounded by the size of the proof of $\wen x. \wen y. P \cpar \new y.Q \cpar R$.
%Notice that $\size{\wen x \wen y P \cpar \new y Q \cpar R} = \size{ \new y \left( \wen x P \cpar Q \right) \cpar R }$.
Hence, by induction, there exist $U$ and $V$ such that $\vdash P \cpar V$ and $\nfv{x}{U}$ and either $U = V$ or $U = \new x. V$ the derivation 
$
\vcenter{
\infer[]{
Q \cpar T
}{
U
}}
$ holds.
Observe that if $S = T$, then 
$
\vcenter{
\infer[]{
\new y. Q \cpar S
}{
\mathopen{\new y.} \left( Q \cpar T \right)
}}$, since $\nfv{y}{S}$.
If $S = \wen y. T$ then 
$
\vcenter{
\infer[]{
\new y. Q \cpar \wen y. T
}{
\mathopen{\new y.} \left( Q \cpar T \right)
}}$.
Thereby the following derivation can be constructed, where if $U = V$ then $W = \new y. V$ and if $U = \new x. V$ then $W = \new x. \new y. V$, and also the premise is equivalent to $W$ by \textit{equivariance} for \textit{new}:
$
\vcenter{
\infer[]{
\new y. Q \cpar R
}{
\infer[]{
\new y. Q \cpar S
}{
\infer[]{
\mathopen{\new y.} \left( Q \cpar T \right)
}{
\new y. U
}}}}
$.
Furthermore, the following proof can be constructed 
$
\vcenter{
\infer[]{
\wen y. P \cpar \new y. V
}{
\infer[]{
\mathopen{\new y.} \left(P \cpar V\right)
}{
\infer[]{
 \new y. \cunit
}{
\cunit
}}}}$
and, by Lemma~\ref{lemma:bound}, $\size{\new y. V} \preceq \size{\new y. Q \cpar R}$ hence 
$\size{ \wen y. P \cpar \new y. V } \prec \size{ \wen x. \wen y. P \cpar \new y. Q \cpar R }$, as required.


\item Consider a commutative case for \textit{new} induced by \textit{equivariance} for \textit{new}, where the bottommost rule is an instance of \textit{extrude new} of the form
\[
\infer[, \mbox{where $\nfv{y}{Q}$ and $\vdash \mathopen{\new y.} \left( \new x. P \cpar Q\right) \cpar R$.}]{
\new x. \new y. P \cpar Q \cpar R
}{
\mathopen{\new y.} \left( \new x. P \cpar Q\right) \cpar R
}
\]
By the induction hypothesis, there exist $S$ and $T$ such that $\nfv{y}{S}$ and $\vdash \new x. P \cpar Q \cpar T$ and either $S = T$ or $S = \wen y. T$, where
$\vcenter{\infer[]{R}{S}}$.
Furthermore, the size of the proof of $\new x. P \cpar Q \cpar T$ is bound above by the size of the proof of $\mathopen{\new y.} \left( \new x. P \cpar Q\right) \cpar R$, hence strictly bound above by the size of the proof of $\new x. \new y. P \cpar Q \cpar R$.
Hence, by induction again, there exist $U$ and $V$ such that $\nfv{x}{U}$ and $\vdash P \cpar V$ and either $U = V$ or $U = \wen x. V$, and also $\infer[]{Q \cpar T}{U}$.
%
Now define $\hat{W}$ and $W$ as follows.
If $S = T$ then let $\hat{W} = V$.
If $S = \wen y. T$ then let $\hat{W} = \wen y. V$.
If $U = V$ then let $W = \hat{W}$.
If $U = \wen x. V$ then let $W = \wen x. \hat{W}$.
Now observe if $S = T$ then
$
\vcenter{
\infer[]{
Q \cpar R
}{
\infer[]{
Q \cpar T
}{
U
}}}$ and $U = W$.
For $S = \wen y. T$ observe
$
\vcenter{
\infer[]{
Q \cpar R
}{
\infer[]{
Q \cpar \wen y. T
}{
\infer[]{
\mathopen{\wen y.} \left( Q \cpar T \right)
}{
\mathopen{\wen y.} U
}}}}
$, since $\nfv{y}{Q}$,
and if $U = V$ then $\mathopen{\wen y.} U = \hat{W}$,
while if $U = \wen x. V$ then $\wen y.U \equiv \wen x.\hat{W}$, by \textit{equivariance} for \textit{wen}.
Hence in all cases $\vcenter{\infer[]{Q \cpar R}{W}}$
and, since $\nfv{y}{Q}$ and $\nfv{y}{T}$, we can arrange that $\nfv{y}{W}$.
Now, for the cases where $\hat{W} = V$, we have $\nfv{y}{V}$,
and hence $\vcenter{\infer[]{\new y. P \cpar V}{\mathopen{\new y.}\left( P \cpar V \right)}}$.
Also if $\hat{W} = \wen y.V$, then $
\vcenter{
\infer[]{
\new y. P \cpar \wen y.V
}{
\mathopen{\new y.}\left( P \cpar V \right)
}}$. Hence in either case we can construct the proof $
\vcenter{
\infer[]{
\new y. P \cpar \hat{W}
}{
\infer[]{
\mathopen{\new y.}\left( P \cpar V \right)
}{
\infer[]{
 \new y.\cunit
}{ \cunit
}}}}$.
Furthermore, $\size{ \new y. P \cpar \hat{W} } \prec \size{ \new x. \new y. P \cpar Q \cpar R }$,
since by Lemma~\ref{lemma:bound} $\size{ \hat{W} } \preceq \size{ Q \cpar R }$.
%\smallskip


\item Similar commutative cases for \textit{wen} and \textit{new} as principal formulae are induced by \textit{equivariance} where the bottommost rule in a proof is an instance of the \textit{close}, \textit{right wen} or \textit{suspend} rules.
In each case, the quantifier involved in the bottommost rule appears behind the principal connective and is propagated in front of the principal connective using \textit{equivariance}.
%\smallskip
\end{enumerate}

\fullproof{
A similar commutative case for \textit{wen} is induced where the bottommost rule in a proof is an instance of the \textit{left wen} rule of the form 
$
\infer[]{
\wen x. \wen y. P \cpar Q \cpar R
}{
\mathopen{\wen y.} \left( \wen x. P \cpar Q \right) \cpar R
}
$
where $\nfv{y}{Q}$ and $\vdash \mathopen{\wen y.} \left( \wen x. P \cpar Q \right) \cpar R$.
By the induction hypothesis, there exist $S$ and $T$ such that $\nfv{y}{S}$ and $\vdash \wen x. P \cpar Q \cpar T$ and either $S = T$ or $S = \new y. T$ such that $
\infer[]{R}{S}$.
Furthermore, the size of the proof of $\wen x. P \cpar Q \cpar T$ is bounded by the size of the proof of $\mathopen{\wen y.} \left( \wen x. P \cpar Q \right) \cpar R$ hence strictly bounded by the size of the proof of $\wen x. \wen y. P \cpar Q \cpar R$.
Hence, by the induction hypothesis again, there exist $U$ and $V$ such that $\nfv{x}{U}$ and $\vdash P \cpar V$ and either $U = V$ or $U = \new x. V$, and derivation $
\infer[]{Q \cpar T}{U}$ holds.
%
If $S = \new y. T$ let $\hat{W} = \new y. V$ otherwise $\hat{W} = V$. If $U = \new x. V$ let $W = \new x. \hat{W}$, otherwise $W = \hat{W}$.
Thereby the following derivations hold:
\begin{itemize}
\item If $S = \new y. T$ and $U = \new x. V$ then
$
\infer[]{
Q \cpar S
}{
\infer[]{
\mathopen{\new y.} \left( Q \cpar T \right)
}{
 \new y. U \equiv W
}}$.

\item If $S = \new y. T$ and $U = V$ then $
\infer[]{Q \cpar S
}{
\infer[]{
\mathopen{\new y.} \left( Q \cpar T \right) 
}{
 \new y. U
}}$, where the premise is equivalent to $W$.

\item If $S = T$ then $Q \cpar S = Q \cpar T$
and $
\infer[]{
Q \cpar T
}{
 U \equiv W
}$.
\end{itemize}
Hence in any of the above cases, $
\infer[]{
Q \cpar R
}{
\infer[]{
 Q \cpar S
}{ W
}}$.
Now, if $\hat{W} = \new y. V$, then $
\infer[]{
\wen y. P \cpar \hat{W}
}{
 \new y. \left(P \cpar V\right)
}$;
and if $\hat{W} = V$ then $S = T$, hence $
\infer[]{
\wen y. P \cpar \hat{W}
}{
\infer[]{
 \new y. P \cpar \hat{W}
}{
 \mathopen{\new y.} \left( P \cpar V\right)
}}$, since $\nfv{y}{V}$.
Clearly $
\infer[]{
\mathopen{\new y.} \left( P \cpar V\right)
}{
\infer[]{
\new y. \cunit
}{
 \cunit
}}
$, hence $\vdash \wen y.P \cpar \hat{W}$ holds.  Furthermore, $\size{\wen y. P \cpar \hat{W}} \prec \size{\wen x. \wen y. P \cpar Q \cpar R}$, since by Lemma~\ref{lemma:bound} $\size{\hat{W}} \preceq \size{ Q \cpar R }$.


The third commutative case for \textit{wen} induced by \textit{equivariance} is where the bottommost rule is an instance of the \textit{suspend} rule of the form 
$
\infer[]{
\wen x. \wen y. P \cpar \wen y. Q \cpar R
}{
\mathopen{\wen y.} \left( \wen x. P \cpar Q \right) \cpar R
}
$, where $\vdash \mathopen{\wen y.} \left( \wen x. P \cpar Q \right) \cpar R$.
By the induction hypothesis, there exist $S$ and $T$ where $\nfv{y}{S}$ and $\vdash \wen x. P \cpar Q \cpar T$ and either $S = T$ or $S = \new y. T$ such that $\infer{R}{S}$.
Furthermore, the size of the proof of $\wen x. P \cpar Q \cpar T$ is bounded above by the size of the proof of $\mathopen{\wen y.} \left( \wen x. P \cpar Q \right) \cpar R$, hence strictly bounded above by the size of the proof of $\wen x. \wen y. P \cpar \wen y. Q \cpar R$.
Hence by the induction hypothesis again, there exist $U$ and $V$ where $\nfv{x}{U}$ and $\vdash P \cpar V$ and either $U = V$ or $U = \new x. V$, such that $\infer{Q \cpar T}{U}$.
%
If $U = V$ then let $W = \new y. V$, and if $U = \new x. V$ then  let $W = \new x. \new y. V$.
Also observe that whether we have $S = T$ or $S = \new y. T$, we have $
\infer[]{
\wen y. Q \cpar S
}{
\mathopen{\wen y.} \left( Q \cpar T \right)
}$.
Thereby the following derivation can be constructed:
$\infer[]{
\wen y. Q \cpar R
}{
\infer[]{
\mathopen{\wen y.} \left( Q \cpar T \right)
}{
\wen y. U \equiv W
}}$.
Furthermore, the following proof can be constructed $
\infer[]{
\wen y. P \cpar \new y. V
}{
\infer[]{
\mathopen{\new y.} \left( P \cpar V \right)
}{
\infer[]{
\new y. \cunit
}{
\cunit
}}}$;
and the size of the proof of $\size{ \wen y. P \cpar \new y. V } \prec \size{ \wen x. \wen y. P \cpar \wen y. Q \cpar R}$, since $\size{ \new y. V } \preceq \size{ \wen y. Q \cpar R }$ by Lemma~\ref{lemma:bound}.
%\smallskip

Consider the first commutative case for \textit{new} induced by equivariance, where the bottommost rule is of the form
$
\infer[]{
\new x. \new y. P \cpar \wen y. Q \cpar R
}{
\mathopen{\new y.} \left( \new x. P \cpar Q\right) \cpar R
}
$,
where $\vdash \mathopen{\new y.} \left( \new x. P \cpar Q\right) \cpar R$.
By the induction hypothesis, there exist $S$ and $T$ such that $\nfv{y}{S}$ and $\vdash \wen x. P \cpar Q \cpar T$ and either $S = T$ or $S = \wen y .T$, where
$\infer[]{R}{S}$.
Furthermore, the size of the proof of $\wen x. P \cpar Q \cpar T$ is bound above by the size of the proof of $\mathopen{\new y.} \left( \new x. P \cpar Q\right) \cpar R$, hence strictly bound above by the size of the proof of $\new x. \new y. P \cpar \wen y. Q \cpar R$.
Hence, by induction again, there exist $U$ and $V$ such that $\nfv{x}{U}$ and $\vdash P \cpar V$ and either $U = V$ or $U = \wen x. V$, and also $\infer[]{Q \cpar T}{U}$.
If $U = V$ then let $W = \wen y. V$, and if $U = \wen x. V$ then let $W = \wen x. \wen y. V$. 
Also, regardless of whether $S = T$ or $S = \wen y. T$, $\infer[]{\wen y. Q \cpar S}{\mathopen{\wen y.} \left( Q \cpar S \right)}$. Hence derivation %the following derivation can be constructed:
$
\infer[]{
\wen y. Q \cpar R
}{
\infer[]{
\wen y. Q \cpar S
}{
\infer[]{
\mathopen{\wen y.} \left( Q \cpar T \right)
}{
\wen y. U
}}}$ can be constructed, where the premise is equivalent to $W$.
Furthermore, $
\infer[]{
\new y. P \cpar \wen y. V
}{
\infer[]{
\mathopen{\new y.} \left( P \cpar V \right)
}{
\infer[]{
\new y. \cunit 
}{ \cunit
}}}$,
and $\size{ \new y. P \cpar \wen y. V } \prec \size{ \new x. \new y. P \cpar \wen y. Q \cpar R }$,
since by Lemma~\ref{lemma:bound} $\size{ \wen y. V } \preceq \size{ \wen y. Q \cpar R }$.

}




% --- end equivariance.tex ---

% --- begin commutative.tex ---


\item \textbf{Regular commutative cases.} % involving nominals.}
As in every splitting lemma, there are numerous \textit{commutative} cases where the bottommost rule in a proof does not directly involve the principal connective.
For each principal formula handled by this splitting lemma (\textit{new}, \textit{wen}, \textit{with}, \textit{seq} and \textit{times}) there are commutative cases induced by \textit{new}, \textit{wen}, \textit{all}, \textit{with} and \textit{times} and also two commutative cases induced by \textit{seq}.
Thus there are 35 similar commutative cases to check, that all follow a pattern, hence only a representative selection of four cases are presented that make special use of $\alpha$-conversion and the rules \textit{new wen},
\textit{all name}, \textit{with name}, \textit{left name} and \textit{right name}.
Further, representative cases appear in the proof for existential quantifiers.
%that involve the \textit{new wen}, \textit{all name}, \textit{with name}.

\begin{enumerate}[label*=\textbf{.\arabic*}]

\item Consider the commutative case where the principal formula is $\new x. P$ and the bottommost rule is an instance of \textit{extrude new} but applied to a distinct \textit{new} quantifier $\new y. Q$, as in the following rule instance 
\[
\infer[\mbox{, where $\nfv{y}{\new x. P \cpar R}$.}]{
\new x. P \cpar \new y. Q \cpar R \cpar S
}{
\mathopen{\new y.} \left( \new x. P \cpar Q \cpar R\right) \cpar S
}
\] 
Also assume, by $\alpha$-conversion, that $x \not= y$.
By induction, there exist $T$ and $U$ such that $\vdash \new x. P \cpar Q \cpar R \cpar U$, $\nfv{y}{T}$ and either $T = U$ or $T = \wen y. U$, and also
$\vcenter{\infer{S}{T}}$.
%
Furthermore, the size of the proof of $\new x. P \cpar Q \cpar R \cpar U$ is bounded above by the size of the proof of $\mathopen{\new y.} \left( \new x. P \cpar Q \cpar R\right) \cpar S$ and hence strictly bounded above by the size of the proof of $\new x. P \cpar \new y. Q \cpar R \cpar S$, enabling the induction hypothesis.
Hence, by the induction hypothesis, there exist formulae $V$ and $\hat{V}$ such that $\vdash P \cpar \hat{V}$ and $\nfv{x}{V}$ and either $V = \hat{V}$ or $V = \wen x. \hat{V}$, and also $
\vcenter{
\infer[]{
Q \cpar R \cpar U
}{V}}$.
%
Define $W$ such that if $V = \hat{V}$ then $W = \new y. \hat{V}$ and if $V = \wen x. \hat{V}$ then $W = \wen x. \new y. \hat{V}$.
Hence if $V = \wen x. \hat{V}$ then $
\vcenter{
\infer[]{
\new y. V
}{
% = \new y. \wen x.\hat{V}
 \wen x. \new y. \hat{V}
}}$ by applying the \textit{new wen} rule, where the premise equals $W$.
%\ross{Case using new-wen rule. Include in body?}. 
If $V = \hat{V}$ then $\new y. V = W$. In both cases, $\nfv{x}{W}$.
Now observe that either $T = U$ and $\nfv{y}{U}$, hence the derivation $(a)$ below holds;
or $T = \wen y. U$, hence the derivation $(b)$ below holds. Given these, the derivation $(c)$ can be constructed:
\[
\begin{array}{cccc}
\infer[]{
\new y. Q \cpar R \cpar T
}{
\mathopen{\new y.} \left(Q \cpar R \cpar U\right)
}
&
\infer[]{
\new y. Q \cpar R \cpar \wen y. U
}{
\infer[]{
 \mathopen{\new y.} \left(Q \cpar R\right) \cpar \wen y. U
}{
 \mathopen{\new y.} \left(Q \cpar R \cpar U\right)
}}
& 
\infer[]{
\new y. Q \cpar R \cpar S
}{
\infer[]{
\new y. Q \cpar R \cpar T
}{
\infer[]{
\mathopen{\new y.} \left( Q \cpar R \cpar U \right)
}{
\infer[]{
\new y. V
}{
W
}}}} &
\infer[]{
P \cpar \new y. \hat{V}
}{
\infer[]{
 \mathopen{\new y.} \left( P \cpar \hat{V} \right) 
}{
\infer[]{
 \new y. \cunit
}{
 \cunit
}}}\\
(a) & (b) & (c) & (d)
\end{array}
\]
%\[
%\infer[]{
%\new y. Q \cpar R \cpar S
%}{
%\infer[]{
%\new y. Q \cpar R \cpar T
%}{
%\infer[]{
%\mathopen{\new y.} \left( Q \cpar R \cpar U \right)
%}{
%\infer[]{
%\new y. V
%}{
%W
%}}}}
%\]
Since $\nfv{y}{\new x. P \cpar R}$ and $x \not= y$, we have $\nfv{y}{P}$;
thereby the proof $(d)$ above can be constructed.
%$
%\vcenter{
%\infer[]{
%P \cpar \new y. \hat{V}
%}{
%\infer[]{
% \mathopen{\new y.} \left( P \cpar \hat{V} \right) 
%}{
%\infer[]{
% \new y. \cunit
%}{
% \cunit
%}}}}$.
Furthermore, $\size{ P \cpar \new y. \hat{V} } \prec \size{ \wen x. P \cpar \new y. Q \cpar R \cpar S }$ since by Lemma~\ref{lemma:bound} $\size{ \new y. \hat{V} } \preceq \size{\new y. Q \cpar R \cpar S}$
and the \textit{wen count} strictly decreases.
%\smallskip


%%%%%%%%%%%%%%%% Standard commutative cases for wen %%%%%%%%%%%%%%%%%%%%%%%
\item Consider the commutative case for principal formula $\wen x. T$ where the bottommost rule is \textit{external}: 
%for \textit{with} where . The bottommost rule is the form
\[
\infer[]{
\wen x. T \cpar \left( U \wwith V \right) \cpar W \cpar P
}{
\left(\left(\wen x. T \cpar U \cpar W\right) \wwith
\left(\wen x. T \cpar V \cpar W\right)\right)
\cpar P
}
\]
where $\vdash \left(\left(\wen x. T \cpar U \cpar W\right) \wwith \left(\wen x. T \cpar V \cpar W\right)\right) \cpar P$ holds.
By the induction hypothesis, we have that both $\vdash \wen x. T \cpar U \cpar W \cpar P$ and $\vdash \wen x. T \cpar V \cpar W \cpar P$ hold;
and furthermore the multiset inequalities 
\[
\begin{array}{rcl}
\occ{\wen x. T \cpar U \cpar W \cpar P} & \mstrict & \occ{ \wen x. T \cpar \left( U \wwith V \right) \cpar W \cpar P } \mbox{ and }\\
\occ{\wen x. T \cpar V \cpar W \cpar P} & \mstrict & \occ{ \wen x. T \cpar \left( U \wwith V \right) \cpar W \cpar P }
\end{array}
\]
hold.
%
Hence, by the induction hypothesis, there exist $Q$ and $\hat{Q}$ such that $\vdash T \cpar \hat{Q}$, $\nfv{x}{Q}$ and either $Q = \hat{Q}$ or $Q = \new x. \hat{Q}$.
Also, by the induction hypothesis, there exist $R$ and $\hat{R}$ such that
$\vdash T \cpar \hat{R}$, $\nfv{x}{R}$ and either $R = \hat{R}$ or $R = \new x. \hat{R}$.
Furthermore the two derivations 
$
\vcenter{
\infer[]{
U \cpar W \cpar P
}{ Q }}
$
and
$
\vcenter{
\infer[]{
V \cpar W \cpar P
}{ R }
}
$ hold.
%
Now define $S$ such that if $Q = \hat{Q}$ and $R = \hat{R}$ then $S = \hat{Q} \wwith \hat{R}$, and 
$S = \mathopen{\wen x.} \left( \hat{Q} \wwith \hat{R} \right)$ otherwise, observing that in either case $\nfv{x}{S}$.
In the case $Q = \wen x. \hat{Q}$ and $R = \wen x. \hat{R}$, by the \textit{with name} rule,
$
\vcenter{
\infer[]{
\wen x. \hat{Q} \wwith \wen x. \hat{R}
}{
\mathopen{\wen x.} \left( \hat{Q} \wwith \hat{R} \right)
}}$.
In the case $Q = \wen x. \hat{Q}$ and $R = \hat{R}$, by the \textit{left name} rule,
$
\vcenter{
\infer[]{
\wen x. \hat{Q} \wwith \hat{R}
}{
\mathopen{\wen x.} \left( \hat{Q} \wwith \hat{R} \right)
}}$.
In the case that $Q = \hat{Q}$ and $R = \wen x. \hat{R}$, by the \textit{right name} rule,
$
\vcenter{
\infer[]{
\hat{Q} \wwith \wen x. \hat{R}
}{
\mathopen{\wen x.} \left( \hat{Q} \wwith \hat{R} \right)
}}$.
%
Thereby the following derivation and proof can be constructed:
\[
\infer[]{
\left( U \wwith V \right) \cpar W \cpar P
}{
\infer[]{
\left(U \cpar W \cpar P\right) \wwith \left(V \cpar W \cpar P\right) 
}{
\infer[]{
Q \wwith R
}{
S
}}}
\qquad
\infer[.]{
T \cpar \left(\hat{Q} \wwith \hat{R}\right)
}{
\infer[]{
\left(T \cpar \hat{Q}\right) \wwith \left( T \cpar \hat{R}\right)
}{
\infer[]{
\cunit \wwith \cunit
}{
\cunit
}}}
\]
Furthermore, by Lemma~\ref{lemma:bound}, $\size{ S } \preceq \size{ \left( U \wwith V \right) \cpar W \cpar P }$; and,
since the \textit{wen} count strictly decreases,
$\size{ T \cpar \hat{Q} \wwith \hat{R} } \prec \size{\wen x. T \cpar \left( U \wwith V \right) \cpar W \cpar P}$.
%A similar argument holds for \textit{new} as the principal formula.



\item Consider the commutative case where the principal formula is $\wen x. T$ and the bottommost rule is an instance of the
\textit{extrude1} rule of the form
\[
\infer[]{
\wen x. T \cpar \forall y. U \cpar V \cpar W
}{
\mathopen{\forall y.}\left( \wen x. T \cpar U \cpar V \right) \cpar W
}\]
 assuming $\nfv{y}{(\wen x. T \cpar V)}$
and $\vdash \mathopen{\forall y.}\left( \wen x. T \cpar U \cpar V \right) \cpar W$ holds.
By Lemma~\ref{lemma:universal}, for every variable $z$, $\vdash \left( \wen x. T \cpar U \cpar V \right)\sub{y}{z} \cpar W$ holds.
Furthermore, since $\nfv{y}{(\wen x. T \cpar V)}$, we have equivalence $\left( \wen x. T \cpar U \cpar V \right)\sub{y}{z} \cpar W \equiv \wen x. T \cpar U\sub{y}{z} \cpar V \cpar W$.
%
The strict multiset inequality $
\occ{ \wen x. T \cpar U\sub{y}{z} \cpar V \cpar W } \mstrict \occ{ \wen x. T \cpar \forall y. U \cpar V \cpar W }
$ holds.
Hence, by the induction hypothesis, for every variable $z$, there exist formulae $P^z$ and $Q^z$ such that $\vdash T \cpar Q^z$ and $\nfv{x}{P^z}$ and either $P^z = Q^z$ or $P^z = \new x. Q^z$, and also
$
\vcenter{
\infer[]{
U\sub{y}{z} \cpar V \cpar W
}{
P^z
}}
$.
%
Define $W^z$ such that if $P^z = Q^z$ then $W^z = \forall z. Q^z$ and if $P^z = \new x. Q^z$ then $W^z = \new x. \forall z. Q^z$. Hence if $P^z = \new x. Q^z$ then, since $\forall$ permutes with any quantifier using the \textit{all name} rule, $
\vcenter{
\infer[]{
\forall z. \new x. Q^z
}{
 \new x. \forall z. Q^z
}}$.
Hence, for a fresh $z$ such that $\nfv{z}{(\forall y. U \cpar V \cpar W)}$ and $\nfv{z}{T}$, 
the following derivations can be constructed:
\[
\infer[]{
\forall y. U \cpar V \cpar W
}{
\infer[]{
\mathopen{\forall z.} \left( U\sub{y}{z} \cpar V \cpar W \right) 
}{
\infer[]{
\forall z. P^z
}{
W^z
}}}
\qquad
\infer[]{
T \cpar \forall z. Q^z
}{
\infer[]{
 \mathopen{\forall z.} \left( T \cpar  Q^z\right) 
}{
\infer[]{
 \mathopen{\forall z.} \cunit 
}{
 \cunit
}}}
\]
%Notice that $\forall z \left( \tcontext{} \right)$ is a $n$-ary killing context as required.
Furthermore, $\size{ W^z } \preceq \size{ \forall y. U \cpar V \cpar W }$ by Lemma~\ref{lemma:bound}; hence 
\[
\size{ T \cpar \forall z. Q^z } \prec \size{ \wen x. T \cpar \forall y. U \cpar V \cpar W }
\] 
since the wen count strictly decreases.


%\ross{Include in body since uses $\alpha$-conversion to avoid $x = y$.}
\item Consider the commutative case where the principal connective is $\textit{wen}$ and the bottommost rule is an instance of the extrude new rule of the form 
%when the \textit{new} quantifier commutes with a \textit{wen} quantifier.
%In this case, the bottommost rule of a proof is of the form
\[
\infer[,]{
\wen x. P \cpar \new y. Q \cpar R \cpar S
}{
\mathopen{\new y.} \left( \wen x. P \cpar Q \cpar R\right) \cpar S
}
\]
 where $\nfv{y}{\wen x. P \cpar R}$ and also $x \not= y$, where the second condition can be achieved by $\alpha$-conversion.
%
By the induction hypothesis, there exist $T$ and $U$ such that $\vdash \wen x. P \cpar Q \cpar R \cpar U$, $\nfv{y}{T}$ and either $T = U$ or $T = \wen y. U$, and also
$\vcenter{\infer{S}{T}}$.
Furthermore, the size of the proof of $\wen x. P \cpar Q \cpar R \cpar U$ is bounded above by the size of the proof of $\mathopen{\new y.} \left( \wen x. P \cpar Q \cpar R\right) \cpar S$ and hence strictly bounded above by the size of the proof of $\wen x. P \cpar \new y. Q \cpar R \cpar S$, enabling the induction hypothesis.
Hence, by the induction hypothesis, there exist formulae $V$ and $\hat{V}$ such that $\vdash P \cpar \hat{V}$ and $\nfv{x}{V}$ and either $V = \hat{V}$ or $V = \new x. \hat{V}$, and also $
\vcenter{
\infer[]{
Q \cpar R \cpar U}{
 V
}}$.
%
Define $W$ such that if $V = \hat{V}$ then $W = \new y. \hat{V}$ and if $V = \new x. \hat{V}$ then $W = \new x. \new y. \hat{V}$.
Now observe that either we have that $T = U$ and $\nfv{y}{U}$ and hence
the derivation $(a)$ below left holds;  or we have that $T = \wen y. U$ and hence the derivation $(b)$ belw  holds. 
Hence,
% the following derivation can be constructed, 
by applying one of these cases, we have the derivation $(c)$ below, where the premise is equivalent to $W$.
\[
\begin{array}{cccc}
\infer[]{
\new y. Q \cpar R \cpar T
}{
 \mathopen{\new y.} \left(Q \cpar R \cpar U\right)
}
&
\infer[]{
\new y. Q \cpar R \cpar \wen y. U 
}{
\infer[]{
\mathopen{\new y.} \left(Q \cpar R\right) \cpar \wen y. U
}{
 \mathopen{\new y.} \left(Q \cpar R \cpar U\right)
}}
&
\infer[]{
\new y. Q \cpar R \cpar S
}{
\infer[]{
\new y. Q \cpar R \cpar T
}{
\infer[]{
\mathopen{\new y.} \left( Q \cpar R \cpar U \right)
}{
\new y. V
}}} &
\infer[.]{
P \cpar \new y. \hat{V}
}{
\infer[]{
 \mathopen{\new y.} \left( P \cpar \hat{V} \right)
}{
\infer[]{
 \new y. \cunit
}{
 \cunit
}}}
\\
(a)  & (b) & (c) & (d)
\end{array}
\]
%$
%%\begin{array}{rl}
%\infer[]{
%\new y. Q \cpar R \cpar S
%}{
%\infer[]{
%\new y. Q \cpar R \cpar T
%}{
%\infer[]{
%\mathopen{\new y.} \left( Q \cpar R \cpar U \right)
%}{
%\new y. V
%}}}
%%\end{array}
%$ 
Since $\nfv{y}{\wen x. P}$ and $x \not= y$, we have $\nfv{y}{P}$;
thereby the proof $(d)$ above can be constructed.
%\[
%\infer[.]{
%P \cpar \new y. \hat{V}
%}{
%\infer[]{
% \mathopen{\new y.} \left( P \cpar \hat{V} \right)
%}{
%\infer[]{
% \new y. \cunit
%}{
% \cunit
%}}}
%\]
Furthermore, $\size{ P \cpar \new y.\hat{V} } \prec \size{ \wen x. P \cpar \new y. Q \cpar R \cpar S }$ since by Lemma~\ref{lemma:bound} 
\[
\size{ \new y. \hat{V} } \preceq \size{\new y. Q \cpar R \cpar S}
\]
and the \textit{wen count} strictly decreases.
%\smallskip
\end{enumerate}

\fullproof{
Consider the commutative case when the \textit{wen} quantifier commutes with another \textit{wen} quantifier at the root of the principal formula.
In this case, the bottommost rule of a proof is of the form $\wen x. P \cpar \wen y. Q \cpar R \cpar \infer[]{
S}{
 \mathopen{\wen y.} \left( \wen x. P \cpar Q \cpar R\right) \cpar S
}$, where $\nfv{y}{\wen x. P \cpar R}$ and $x \not= y$, where the second condition can be achieved by $\alpha$-conversion.
By the induction hypothesis, there exist $T$ and $U$ such that $\vdash \wen x. P \cpar Q \cpar R \cpar U$, $\nfv{y}{T}$ and either $T = U$ or $T = \new y. U$, and also
$\infer{S}{T}$.
Furthermore, the size of the proof of $\wen x. P \cpar Q \cpar R \cpar U$ is bounded above by the size of the proof of $\mathopen{\wen y.} \left( \wen x. P \cpar Q \cpar R\right) \cpar S$ and hence strictly bounded above by the size of the proof of $\wen x. P \cpar \wen y. Q \cpar R \cpar S$, enabling the induction hypothesis.
Hence, by the induction hypothesis, there exist formulae $V$ and $\hat{V}$ such that $\vdash P \cpar \hat{V}$ and $\nfv{x}{V}$ and either $V = \hat{V}$ or $V = \new x. \hat{V}$, and also $
\infer[]{
Q \cpar R \cpar U 
}{ V }$.

Define $W$ such that if $V = \hat{V}$ then $W = \new y. \hat{V}$ and if $V = \new x. \hat{V}$ then $W = \new x. \new y. \hat{V}$. In either case, $\nfv{x}{W}$.
Now observe that either $T = U$ and $\nfv{y}{U}$ hence
$\infer[]{
\wen y. Q \cpar R \cpar T
}{
\infer[]{
 \mathopen{\wen y.} \left(Q \cpar R \cpar U\right) 
}{
 \mathopen{\new y.} \left(Q \cpar R \cpar U\right)
}}$;
or $T = \new y. U$ hence $
\infer[]{
\wen y. Q \cpar R \cpar \new y. U
}{
\infer[]{
 \mathopen{\wen y.} \left(Q \cpar R\right) \cpar \mathopen{\new y.} U 
}{
\mathopen{\new y.} \left(Q \cpar R \cpar U\right)
}}$.
Hence the following derivation can be constructed, by applying one of these cases:
$
\infer[]{
\wen y. Q \cpar R \cpar S
}{
\infer[]{
 \wen y. Q \cpar R \cpar T
}{
\infer[]{
\mathopen{\new y.} \left( Q \cpar R \cpar U \right)
}{
\new y. V
}}}
$, where the premise is equivalent to $W$.
Since $\nfv{y}{\wen x. P}$ and $x \not= y$ we have $\nfv{y}{P}$,
hence the following proof can be constructed:
$
\infer[]{
P \cpar \new y. \hat{V}
}{
\infer[]{
\mathopen{\new y.}\left( P \cpar \hat{V} \right)
}{
\infer[]{
\new y. \cunit 
}{
\cunit
}}}$.
Furthermore, $\size{ P \cpar \new y. \hat{V} } \prec \size{ \wen x. P \cpar \new y. Q \cpar R \cpar S }$ since by Lemma~\ref{lemma:bound} $\size{ \new y. \hat{V} } \preceq \size{\new y. Q \cpar R \cpar S}$
and the \textit{wen count} strictly decreases.
\smallskip











Consider the commutative case when the \textit{new} quantifier commutes with $P_0 \tensor P_1$ as the principal formula.
In this case the bottommost rule of a proof is of the form $
\infer[]{
\left( P_0 \tensor P_1 \right) \cpar \new y. Q \cpar R \cpar S
}{
\mathopen{\new y.} \left( \left( P_0 \tensor P_1 \right) \cpar Q \cpar R\right) \cpar S
}$,
assuming that $\nfv{y}{\left( P_0 \tensor P_1 \right) \cpar R}$.
By induction, there exist $T$ and $U$ such that $\vdash \left( P_0 \tensor P_1 \right) \cpar Q \cpar R \cpar U$, $\nfv{y}{T}$ and either $T = U$ or $T = \wen y. U$, and also
$\infer{S}{T}$.
%
Furthermore, the size of the proof of $\left( P_0 \tensor P_1 \right) \cpar Q \cpar R \cpar U$ is bounded above by the size of the proof of $\mathopen{\new y.} \left( \left( P_0 \tensor P_1 \right) \cpar Q \cpar R\right) \cpar S$ and hence strictly bounded above by the size of the proof of $\left( P_0 \tensor P_1 \right) \cpar \new y. Q \cpar R \cpar S$, enabling the induction hypothesis.
Hence, by the induction hypothesis, there exist formulae $V_i$ and $W_i$ such that $\vdash P_0 \cpar V_i$ and $\vdash P_1 \cpar W_i$, for $1 \leq i \leq n$, and also $n$-ary killing context $\tcontext{}$ such that
$
\infer[]{
Q \cpar R \cpar U 
}{
\tcontext{ V_i \cpar W_i \colon 1 \leq i \leq n }
}$.
Furthermore, the size of the proofs of $P_0 \cpar V_i$ and $P_1 \cpar W_i$ are bounded above by the size of the proof of $\left( P_0 \tensor P_1 \right) \cpar \new y. Q \cpar R \cpar S$.
%
Now observe that either $T = U$ and $\nfv{y}{U}$ hence
$\infer[]{
\new y. Q \cpar R \cpar T
}{
 \mathopen{\new y.} \left(Q \cpar R \cpar U\right)
}$;
or $T = \wen y. U$ hence $
\infer[]{
\new y. Q \cpar R \cpar \wen y. U
}{
\infer[]{
 \mathopen{\new y.} \left(Q \cpar R\right) \cpar \wen y. U
}{
 \mathopen{\new y.} \left(Q \cpar R \cpar U\right)
}}$.
Hence the following derivation can be constructed, by the above observations:
$
\infer[]{
\new y. Q \cpar R \cpar S
}{
\infer[]{
\new y. Q \cpar R \cpar T
}{
\infer[]{
\mathopen{\new y.} \left( Q \cpar R \cpar U \right)
}{
\new y. \tcontext{ V_i \cpar W_i \colon 1 \leq i \leq n }
}}}
$.
Observe that $\new y. \tcontext{ }$ is a $n$-ary killing context as required.
\smallskip


Consider the commutative case when the \textit{wen} quantifier commutes with $P_0 \tensor P_1$ as the principal formula.
In this case the bottommost rule of a proof is of the form $
\infer[]{
\left( P_0 \tensor P_1 \right) \cpar \wen y. Q \cpar R \cpar S
}{
 \mathopen{\wen y.} \left( \left( P_0 \tensor P_1 \right) \cpar Q \cpar R\right) \cpar S
}$, where $\nfv{y}{\left( P_0 \tensor P_1 \right) \cpar R}$.
By induction, there exist $T$ and $U$ such that $\vdash \left( P_0 \tensor P_1 \right) \cpar Q \cpar R \cpar U$, $\nfv{y}{T}$ and either $T = U$ or $T = \new y. U$, and also
$\infer{S}{T}$.
%
Furthermore, the size of the proof of $\left( P_0 \tensor P_1 \right) \cpar Q \cpar R \cpar U$ is bounded above by the size of the proof of $\mathopen{\wen y.} \left( \left( P_0 \tensor P_1 \right) \cpar Q \cpar R\right) \cpar S$ and hence strictly bounded above by the size of the proof of $\left( P_0 \tensor P_1 \right) \cpar \wen y. Q \cpar R \cpar S$, enabling the induction hypothesis.
Hence, by the induction hypothesis, there exist formulae $V_i$ and $W_i$ such that $\vdash P_0 \cpar V_i$ and $\vdash P_1 \cpar W_i$, for $1 \leq i \leq n$, and also $n$-ary killing context $\tcontext{}$ such that
$\infer[]{
Q \cpar R \cpar U
}{
\tcontext{ V_i \cpar W_i \colon 1 \leq i \leq n }
}$.
Furthermore, the size of the proofs of $P_0 \cpar V_i$ and $P_1 \cpar W_i$ are bounded above by the size of the proof of $\left( P_0 \tensor P_1 \right) \cpar \wen y. Q \cpar R \cpar S$.
%
Now observe that either $T = U$ and $\nfv{y}{U}$ hence
$\infer[]{
\wen y. Q \cpar R \cpar T
}{
\infer[]{
\mathopen{\wen y.} \left(Q \cpar R \cpar U\right)
}{
 \mathopen{\new y.} \left(Q \cpar R \cpar U\right)
}}$;
or $T = \new y. U$ hence $
\infer[]{
\wen y. Q \cpar R \cpar \new y. U
}{
\infer[]{
\mathopen{\wen y.} \left(Q \cpar R\right) \cpar \new y. U
}{
 \mathopen{\new y.} \left(Q \cpar R \cpar U\right)
}}$.
Hence the following derivation can be constructed, by the above observations:
$
%\begin{array}{rl}
\infer[]{
\wen y. Q \cpar R \cpar S
}{
\infer[]{
\wen y. Q \cpar R \cpar T
}{
\infer[]{
\mathopen{\new y.} \left( Q \cpar R \cpar U \right)
}{
\mathopen{\new y.} \tcontext{ V_i \cpar W_i \colon 1 \leq i \leq n }
}}}
%\end{array}
$.
Observe that $\new y. \tcontext{ }$ is a $n$-ary killing context as required.
\smallskip



Consider the commutative case where the \textit{new} quantifier commutes with $P_0 \cseq P_1$ as the principal formula.
In this case the bottommost rule of a proof is of the form $
\infer[]{
\left( P_0 \cseq P_1 \right) \cpar \new y. Q \cpar R \cpar S 
}{
\mathopen{\new y.} \left( \left( P_0 \cseq P_1 \right) \cpar Q \cpar R\right) \cpar S
}$,
where $\nfv{y}{\left( P_0 \cseq P_1 \right) \cpar R}$.
By induction, there exist $T$ and $U$ such that $\vdash \left( P_0 \cseq P_1 \right) \cpar Q \cpar R \cpar U$, $\nfv{y}{T}$ and either $T = U$ or $T = \wen y. U$, and also
$\infer{S}{T}$.
%
Furthermore, the size of the proof of $\left( P_0 \cseq P_1 \right) \cpar Q \cpar R \cpar U$ is bounded above by the size of the proof of $\mathopen{\new y.} \left( \left( P_0 \cseq P_1 \right) \cpar Q \cpar R\right) \cpar S$ and hence strictly bounded above by the size of the proof of $\left( P_0 \cseq P_1 \right) \cpar \new y. Q \cpar R \cpar S$, enabling the induction hypothesis.
Hence, by the induction hypothesis, there exist formulae $V_i$ and $W_i$ such that $\vdash P_0 \cpar V_i$ and $\vdash P_1 \cpar W_i$, for $1 \leq i \leq n$, and also $n$-ary killing context $\tcontext{}$ such that
$\infer[]{
Q \cpar R \cpar U
}{
\tcontext{ V_i \cseq W_i \colon 1 \leq i \leq n }
}$.
Furthermore, the size of the proofs of $P_0 \cpar V_i$ and $P_1 \cpar W_i$ are bounded above by the size of the proof of $\left( P_0 \cseq P_1 \right) \cpar \new y. Q \cpar R \cpar S$.
%
Now observe that either $T = U$ and $\nfv{y}{U}$ hence
$
\infer[]{
\new y. Q \cpar R \cpar T
}{
\mathopen{\new y.} \left(Q \cpar R \cpar U\right)
}$;
or $T = \wen y. U$ hence $
\infer[]{
\new y. Q \cpar R \cpar \wen y. U
}{
\infer[]{
 \mathopen{\new y.} \left(Q \cpar R\right) \cpar \wen y. U 
}{ \mathopen{\new y.} \left(Q \cpar R \cpar U\right)
}}$.
Hence the following derivation can be constructed:
$
%\begin{array}{rl}
\infer[]{
\new y. Q \cpar R \cpar S
}{
\infer[]{
\new y. Q \cpar R \cpar T
}{
\infer[]{
\mathopen{\new y.} \left( Q \cpar R \cpar U \right)
}{
\new y. \tcontext{ V_i \cseq W_i \colon 1 \leq i \leq n }
}}}
%\end{array}
$.
Observe that $\new y. \tcontext{ }$ is a $n$-ary killing context as required.
\smallskip


Consider the commutative case when the \textit{wen} quantifier commutes with $P_0 \cseq P_1$ as the principal formula.
In this case the bottommost rule of a proof is of the form $
\infer[]{
\left( P_0 \cseq P_1 \right) \cpar \wen y. Q \cpar R \cpar S
}{
\mathopen{\wen y.} \left( \left( P_0 \cseq P_1 \right) \cpar Q \cpar R\right) \cpar S
}$. where $\nfv{y}{\left( P_0 \cseq P_1 \right) \cpar R}$.
By induction, there exist $T$ and $U$ such that $\vdash \left( P_0 \cseq P_1 \right) \cpar Q \cpar R \cpar U$, $\nfv{y}{T}$ and either $T = U$ or $T = \new y. U$, and also
$\infer{S}{T}$.
%
Furthermore, the size of the proof of $\left( P_0 \cseq P_1 \right) \cpar Q \cpar R \cpar U$ is bounded above by the size of the proof of $\mathopen{\wen y.} \left( \left( P_0 \cseq P_1 \right) \cpar Q \cpar R\right) \cpar S$ and hence strictly bounded above by the size of the proof of $\left( P_0 \cseq P_1 \right) \cpar \wen y. Q \cpar R \cpar S$, enabling the induction hypothesis.
Hence, by the induction hypothesis, there exist formulae $V_i$ and $W_i$ such that $\vdash P_0 \cpar V_i$ and $\vdash P_1 \cpar W_i$, for $1 \leq i \leq n$, and also $n$-ary killing context $\tcontext{}$ such that
$
\infer[]{
Q \cpar R \cpar U 
}{
 \tcontext{ V_i \cseq W_i \colon 1 \leq i \leq n }
}$.
Furthermore, the size of the proofs of $P_0 \cpar V_i$ and $P_1 \cpar W_i$ are bounded above by the size of the proof of $\left( P_0 \cseq P_1 \right) \cpar \wen y. Q \cpar R \cpar S$.
%
Now observe that either $T = U$ and $\nfv{y}{U}$ hence
$
\infer[]{
\wen y. Q \cpar R \cpar T
}{
\infer[]{
 \wen y. \left(Q \cpar R \cpar U\right) 
}{ \mathopen{\new y.} \left(Q \cpar R \cpar U\right)
}}$;
or $T = \new y. U$ hence $
\infer[]{
\wen y. Q \cpar R \cpar \new y. U
}{
\infer[]{
 \wen y. \left(Q \cpar R\right) \cpar \new y. U 
}{
 \mathopen{\new y.} \left(Q \cpar R \cpar U\right)
}}$.
Hence the following derivation can be constructed:
$
%\begin{array}{rl}
\infer[]{
\wen y. Q \cpar R \cpar S
}{
\infer[]{
\wen y. Q \cpar R \cpar T
}{
\infer[]{
\mathopen{\new y.} \left( Q \cpar R \cpar U \right)
}{
\mathopen{\new y.} \tcontext{ V_i \cseq W_i \colon 1 \leq i \leq n }
}}}
%\end{array}
$.
Observe that $\new y. \tcontext{ }$ is a $n$-ary killing context as required.
\smallskip




Consider commutative cases where the principal formula moves entirely to the left hand side of a \textit{seq} operator. 
For principal formula $\wen x. U$, the bottommost rule in a proof is of the form
$
\infer[]{
\wen x. U \cpar \left( V \cseq P \right) \cpar W \cpar Q
}{
\left( \left( \wen x. U \cpar V \cpar W\right) \cseq P \right) \cpar Q
}
$
such that $\vdash \left( \left( \wen x. U \cpar V \cpar W\right) \cseq P \right) \cpar Q$ holds.
%
By the induction hypothesis, there exist $R_i$ and $S_i$ such that $\vdash \wen x. U \cpar V \cpar W \cpar R_i$ and $\vdash P \cpar S_i$, for $1 \leq i \leq n$, and $n$-ary killing context $\tcontext{}$ such that the derivation
$
\infer[]{
 Q
}{
 \tcontext{ R_1 \cseq S_1, \hdots, R_n \cseq S_n }
}
$ holds, 
and furthermore the size of the proof of $\wen x. U \cpar V \cpar W \cpar R_i$ is bounded above by the size of the proof of $\left( \left( \wen x. U \cpar V \cpar W\right) \cseq P \right) \cpar Q$ hence strictly bounded above by the size of the proof of $\wen x. U \cpar \left( V \cseq P \right) \cpar W \cpar Q$ enabling the induction hypothesis.
By the induction hypothesis again, there exist formulae $T_i$ and $\hat{T}_i$ such that $\vdash U \cpar \hat{T}_i$ and $\nfv{x}{T_i}$ and either $T_i = U_i$ or $T_i = \new x. U_i$, and also the derivation $
\infer[]{
V \cpar W \cpar R_i
}{
 T_i
}$ holds.
%
Hence the following derivation can be constructed.
\[
\infer[]{
\left( V \cseq P \right) \cpar W \cpar Q
}{
\infer[]{
\left( V \cseq P \right) \cpar W \cpar \tcontext{ R_1 \cseq S_1, \hdots, R_n \cseq S_n } \\
}{
\infer[]{
\tcontext{ \left( V \cseq P \right) \cpar W \cpar R_i \cseq S_i \colon 1 \leq i \leq n } \\
}{
\infer[]{
\tcontext{ \left( V \cpar W \cpar R_i \right) \cseq \left( P \cpar S_i \right) \colon 1 \leq i \leq n } 
}{
\infer[]{
\tcontext{ V \cpar W \cpar R_i \colon 1 \leq i \leq n } 
}{
\tcontext{ T_i \colon 1 \leq i \leq n }
}}}}}
\]
%Notice that $\tcontext{ \tcontextn{1}{}, \hdots, \tcontextn{n}{} }$ is a $\sum_{i=1}^n{m_i}$-ary killing context as required.
Furthermore, by Lemma~\ref{lemma:bound},
$\size{T_i} \preceq \size{\left( V \cseq P \right) \cpar W \cpar Q}$
and hence $\size{ U \cpar \hat{T}_i } \prec \size{\wen x. U \cpar \left( V \cseq P \right) \cpar W \cpar Q}$, since the \textit{wen} count strictly decreases.
The cases where the principal formula moves entirely to the right hand side of the \textit{seq} operator, and the analogous case for \textit{times}, are similar to the above case.
%\smallskip





\item \textbf{Commutative cases involving all and with.}
%%%%%%%%%%%%%%%% Standard commutative cases for times and seq %%%%%%%%%%%%%%%%%%%%%%%
Consider the commutative case for \textit{with} where $T_0 \tensor T_1$ is the principal formula. In this case the bottommost rule is the following form, such that $\vdash \left(\left(\left(T_0 \tensor T_1\right) \cpar U \cpar W\right) \wwith \left(\left(T_0 \tensor T_1\right) \cpar V \cpar W\right)\right) \cpar P$ holds:
$
\infer[]{
\left(T_0 \tensor T_1\right) \cpar \left( U \wwith V \right) \cpar W \cpar P
}{
\left(\left(T_0 \tensor T_1\right) \cpar U \cpar W \wwith
\left(T_0 \tensor T_1\right) \cpar V \cpar W\right)
\cpar P
}
$.
By the induction hypothesis, $\vdash \left(T_0 \tensor T_1\right) \cpar U \cpar W \cpar P$ and $\vdash \left(T_0 \tensor T_1\right) \cpar V \cpar W \cpar P$; and furthermore strict multiset inequalities 
$\occ{\left(T_0 \tensor T_1\right) \cpar U \cpar W \cpar P} \mstrict \occ{ \left(T_0 \tensor T_1\right) \cpar \left( U \wwith V \right) \cpar W \cpar P }$
and
$\occ{\left(T_0 \tensor T_1\right) \cpar V \cpar W \cpar P} \mstrict \occ{ \left(T_0 \tensor T_1\right) \cpar \left( U \wwith V \right) \cpar W \cpar P }$
hold.
%
Hence, by the induction hypothesis, there exist $Q^0_i$ and $Q^1_i$ such that $\vdash T_0 \cpar Q^0_i$ and $\vdash T_1 \cpar Q^1_i$, for $1 \leq i \leq m$; and $R_j^0$ and $R_j^1$ such that
$\vdash T_0 \cpar R^0_j$ and $\vdash T_1 \cpar R^1_j$, for $1 \leq j \leq n$;
and also $m$-ary killing context $\tcontextn{0}{}$ and $n$-ary killing context $\tcontextn{1}{}$
such that the two derivations 
$
%\begin{array}{rl}
\infer[]{
U \cpar W \cpar P 
}{
 \tcontextn{0}{Q^0_i \cpar Q^1_i \colon 1 \leq i \leq m }
}
$
and
$
\infer[]{
V \cpar W \cpar P
}{
 \tcontextn{1}{R^0_j \cpar R^1_j \colon 1 \leq j \leq n }
}
%\end{array}
$ hold.
Furthermore, the size of the proofs of $T_0 \cpar R^0_j$, $T_1 \cpar R^1_j$, $T_0 \cpar Q^0_i$ and $T_1 \cpar Q^1_i$ are bounded above by the size of the proof of $\left(T_0 \tensor T_1\right) \cpar \left( U \wwith V \right) \cpar W \cpar P$.
%
Thereby the following derivation can be constructed.
\[
\infer[]{
\left( U \wwith V \right) \cpar W \cpar P
}{
\infer[]{
\left(U \cpar W \cpar P\right) \wwith \left(V \cpar W \cpar P\right) 
}{
\tcontextn{0}{Q^0_i \cpar Q^1_i \colon 1 \leq i \leq m }
\wwith
\tcontextn{1}{R^0_j \cpar R^1_j \colon 1 \leq j \leq n }
}}
\]
%Furthermore, by Lemma~\ref{lemma:bound}, $\size{ S } \preceq \size{ \left( U \wwith V \right) \cpar W \cpar P }$; hence, since the \textit{wen} count strictly decreases,
%$\size{ T \cpar \hat{Q} \wwith \hat{R} } \prec \wen x T \cpar \left( U \wwith V \right) \cpar W \cpar P$.
%\smallskip


Consider the commutative case where universal quantification commutes with $T_0 \tensor T_1$  as the principal formula. Suppose the bottommost rule is of the form
$
\infer[]{
\left(T_0 \tensor T_1\right) \cpar \forall y. U \cpar V \cpar W
}{
\mathopen{\forall y.}\left( \left(T_0 \tensor T_1\right) \cpar U \cpar V \right) \cpar W
}$, 
assuming $\nfv{y}{(\left(T_0 \tensor T_1\right) \cpar V)}$
where $\vdash \mathopen{\forall y.}\left( \left(T_0 \tensor T_1\right) \cpar U \cpar V \right) \cpar W$ holds.
By Lemma~\ref{lemma:universal}, for every variable $z$, $\vdash \left( \left(T_0 \tensor T_1\right) \cpar U \cpar V \right)\sub{y}{z} \cpar W$ holds.
Furthermore, since $\nfv{y}{(\left(T_0 \tensor T_1\right) \cpar V)}$, $\left( \left(T_0 \tensor T_1\right) \cpar U \cpar V \right)\sub{y}{z} \cpar W \equiv \left(T_0 \tensor T_1\right) \cpar U\sub{y}{z} \cpar V \cpar W$.
Since $\forall$ is removed, $
\occ{ \left(T_0 \tensor T_1\right) \cpar U\sub{y}{z} \cpar V \cpar W } \mstrict \occ{ \left(T_0 \tensor T_1\right) \cpar \forall y. U \cpar V \cpar W }
$ holds.
Pick a fresh $z$ such that $\nfv{z}{(\forall y. U \cpar V \cpar W)}$.
Hence, by the induction hypothesis, there exist formulae $P_i$ and $Q_i$ such that $\vdash T_0 \cpar P_i$ and $\vdash T_1 \cpar Q_i$, for $1 \leq i \leq n$; and also $n$-ary killing context $\tcontext{}$ such that
$
\infer[]{
U\sub{y}{z} \cpar V \cpar W
}{
\tcontext{ P_i \cpar Q_i \colon 1 \leq i \leq n }
}
$.
Furthermore, the size of the proof of $T_0 \cpar P_i$ and $T_1 \cpar Q_i$ are bounded above by the size of the proof of $\left(T_0 \tensor T_1\right) \cpar \forall y. U \cpar V \cpar W$.
Since $z$ was chosen such that $\nfv{z}{ \forall y. U \cpar V \cpar W }$ the following derivation can be constructed, as required.
$
\infer[]{
\forall y. U \cpar V \cpar W
}{
\infer[]{
\mathopen{\forall z.} \left( U\sub{y}{z} \cpar V \cpar W \right) 
}{
\mathopen{\forall z.} \tcontext{ P_i \cpar Q_i \colon 1 \leq i \leq n }
}}
$.
\smallskip


Consider the commutative case for \textit{with} where $T_0 \cseq T_1$ is the principal formula. The bottommost rule is the form
$
\infer[]{
\left(T_0 \cseq T_1\right) \cpar \left( U \wwith V \right) \cpar W \cpar P
}{
\left(\left(T_0 \cseq T_1\right) \cpar U \cpar W \wwith
\left(T_0 \cseq T_1\right) \cpar V \cpar W\right)
\cpar P
}
$
where $\vdash \left(\left(\left(T_0 \cseq T_1\right) \cpar U \cpar W\right) \wwith \left(\left(T_0 \cseq T_1\right) \cpar V \cpar W\right)\right) \cpar P$ holds.
By the induction hypothesis, $\vdash \left(T_0 \cseq T_1\right) \cpar U \cpar W \cpar P$ and $\vdash \left(T_0 \cseq T_1\right) \cpar V \cpar W \cpar P$; and furthermore the strict multiset inequalities
$\occ{\left(T_0 \cseq T_1\right) \cpar U \cpar W \cpar P} \mstrict \occ{ \wen x. T \cpar \left( U \wwith V \right) \cpar W \cpar P }$
and
$\occ{\left(T_0 \cseq T_1\right) \cpar V \cpar W \cpar P} \mstrict \occ{ \wen x. T \cpar \left( U \wwith V \right) \cpar W \cpar P }$ hold.

Hence, by the induction hypothesis, there exist $Q^0_i$ and $Q^1_i$ such that $\vdash T_0 \cpar Q^0_i$ and $\vdash T_1 \cpar Q^1_i$, for $1 \leq i \leq m$; and $R_j^0$ and $R_j^1$ such that
$\vdash T_0 \cpar R^0_j$ and $\vdash T_1 \cpar R^1_j$, for $1 \leq j \leq n$;
and also $m$-ary killing context $\tcontextn{0}{}$ and $n$-ary killing context $\tcontextn{1}{}$
such that the two derivations 
$
%\begin{array}{rl}
\infer[]{
U \cpar W \cpar P
}{
 \tcontextn{0}{Q^0_i \cseq Q^1_i \colon 1 \leq i \leq m }
}
$
and
$
\infer[]{
V \cpar W \cpar P 
}{
  \tcontextn{1}{R^0_j \cseq R^1_j \colon 1 \leq j \leq n }
}
%\end{array}
$ hold.
Furthermore, the size of the proofs of $T_0 \cpar R^0_j$, $T_1 \cpar R^1_j$, $T_0 \cpar Q^0_i$ and $T_1 \cpar Q^1_i$ are bounded above by the size of the proof of $\left(T_0 \cseq T_1\right) \cpar \left( U \wwith V \right) \cpar W \cpar P$.
%
Thereby the following derivation can be constructed.
\[
\infer[]{
\left( U \wwith V \right) \cpar W \cpar P
}{
}{
\infer[]{
\left(U \cpar W \cpar P\right) \wwith \left(V \cpar W \cpar P\right) 
}{
\tcontextn{0}{Q^0_i \cseq Q^1_i \colon 1 \leq i \leq m }
\wwith
\tcontextn{1}{R^0_j \cseq R^1_j \colon 1 \leq j \leq n }
}}
\]
%Furthermore, by Lemma~\ref{lemma:bound}, $\size{ S } \preceq \size{ \left( U \wwith V \right) \cpar W \cpar P }$; hence, since the \textit{wen} count strictly decreases,
%$\size{ T \cpar \hat{Q} \wwith \hat{R} } \prec \wen x T \cpar \left( U \wwith V \right) \cpar W \cpar P$.
\smallskip



Consider the commutative case where universal quantification commutes with $T_0 \cseq T_1$  as the principal formula. Suppose the bottommost rule is of the form
$
\infer[]{
\left(T_0 \cseq T_1\right) \cpar \forall y. U \cpar V \cpar W
}{
\mathopen{\forall y.}\left( \left(T_0 \cseq T_1\right) \cpar U \cpar V \right) \cpar W
}$,
 assuming $\nfv{y}{(\wen x. T \cpar V)}$
where $\vdash \mathopen{\forall y.}\left( \left(T_0 \cseq T_1\right) \cpar U \cpar V \right) \cpar W$ holds.
By Lemma~\ref{lemma:universal}, for every variable $z$, $\vdash \left( \left(T_0 \cseq T_1\right) \cpar U \cpar V \right)\sub{y}{z} \cpar W$ holds.
Furthermore, since $\nfv{y}{(\left(T_0 \cseq T_1\right) \cpar V)}$, $\left( \left(T_0 \cseq T_1\right) \cpar U \cpar V \right)\sub{y}{z} \cpar W \equiv \left(T_0 \cseq T_1\right) \cpar U\sub{y}{z} \cpar V \cpar W$.

The strict multiset inequality $
\occ{ \left(T_0 \cseq T_1\right) \cpar U\sub{y}{z} \cpar V \cpar W } \mstrict \occ{ \left(T_0 \tensor T_1\right) \cpar \forall y. U \cpar V \cpar W }
$ holds.
Pick a fresh $z$ such that $\nfv{z}{(\forall y. U \cpar V \cpar W)}$.
Hence, by the induction hypothesis, there exist formulae $P_i$ and $Q_i$ such that $\vdash T_0 \cpar P_i$ and $\vdash T_1 \cpar Q_i$, for $1 \leq i \leq n$; and also $n$-ary killing context $\tcontext{}$ such that
$
\infer[]{
U\sub{y}{z} \cpar V \cpar W
}{
\tcontext{ P_i \cseq Q_i \colon 1 \leq i \leq n }
}
$.
Furthermore, the size of the proof of $T_0 \cpar P_i$ and $T_1 \cpar Q_i$ are bounded above by the size of the proof of $\left(T_0 \cseq T_1\right) \cpar \forall y. U \cpar V \cpar W$.

Since $z$ was chosen such that $\nfv{z}{ \forall y. U \cpar V \cpar W }$ the following derivation can be constructed, as required.
$
\infer[]{
\forall y. U \cpar V \cpar W
}{
\infer[]{
\mathopen{\forall z.} \left( U\sub{y}{z} \cpar V \cpar W \right) 
}{
\forall z. \tcontext{ P_i \cseq Q_i \colon 1 \leq i \leq n }
}}
$.
\smallskip


Consider the commutative case for \textit{sequence} rule in the presence of principal formula $T \cseq U$, where the \textit{seq} connective in the principal formula is not active on the \textit{sequence} rule. In this case, the bottommost rule in a proof is an instance of the \textit{sequence} rule of the form
$
\infer[]{
\left( T \cseq U \right) \cpar \left( V \cseq P \right) \cpar W \cpar Q
}{
\left( \left( \left( T \cseq U\right) \cpar V \cpar W\right) \cseq P \right) \cpar Q
}
$,
where $T \cseq U \not\equiv \cunit$ and $P \not\equiv \cunit$ and $\vdash \left( \left( \left( T \cseq U\right) \cpar V \cpar W\right) \cseq P \right) \cpar Q$ holds.
By the induction hypothesis, there exists $R_i$, $S_i$ such that $\vdash  \left( T \cseq U\right) \cpar V \cpar W \cpar R_i$ and $\vdash P \cpar S_i$, for $1 \leq i \leq n$, and $n$-ary killing context $\tcontext{}$ such that the following derivation holds:
$
\infer[]{
 Q
}{
 \tcontext{ R_1 \cseq S_1, \hdots, R_n \cseq S_n }
}
$.
Furthermore, $\size{\left( T \cseq U\right) \cpar V \cpar W \cpar R_i} \mleq \size{\left( \left( \left( T \cseq U\right) \cpar V \cpar W\right) \cseq P \right) \cpar Q}$ hence the induction hypothesis is enabled again.
%
By the induction hypothesis, for $1 \leq i \leq n$, there exist formulae $P^i_j$ and $Q^i_j$ such that $\vdash T \cpar P^i_j$ and $\vdash U \cpar Q^i_j$ hold, for $1 \leq j \leq m_i$, and killing contexts $\tcontextn{i}{}$ such that the following derivation holds.
\[
\infer[]{
V \cpar W \cpar R_i
}{
 \tcontextn{i}{P^i_1 \cseq Q^i_1, \hdots, P^i_{m_i} \cseq Q^i_{m_i}}
}
\]
Furthermore the following strict multiset inequalities hold.
\[
\occ{T \cpar P^i_j} \mstrict \occ{\left( T \cseq U \right) \cpar \left( V \cseq P \right) \cpar W \cpar Q
}
\qquad
\mbox{and}
\qquad
\occ{U \cpar Q^i_j} \mstrict \occ{\left( T \cseq U \right) \cpar \left( V \cseq P \right) \cpar W \cpar Q
}
\]
Hence the following derivation can be constructed, as required.
\[
\infer[]{
\left( V \cseq P \right) \cpar W \cpar Q
}{
\infer[]{
\left( V \cseq P \right) \cpar W \cpar \tcontext{ R_1 \cseq S_1, \hdots, R_n \cseq S_n }
}{
\infer[]{
\tcontext{ \left( V \cseq P \right) \cpar W \cpar \left(R_1 \cseq S_1\right), \hdots, \left( V \cseq P \right) \cpar W \cpar \left(R_n \cseq S_n\right) }
}{
\infer[]{
\tcontext{ \left( V \cpar W \cpar R_1 \right) \cseq \left( P \cpar S_1 \right), \hdots, \left( V \cpar W \cpar R_n \right) \cseq \left( P \cpar S_n \right) }
}{
\infer[]{
\tcontext{ V \cpar W \cpar R_1, \hdots, V \cpar W \cpar R_n }
}{
\tcontext{ \tcontextn{i}{P^i_j \cseq Q^i_j \colon 1 \leq j \leq m_i } \colon 1 \leq i \leq n }
}}}}}
\]
The case for the \rseq rule commuting with the principal formula $T \tensor U$ is similar to the above case. Also the cases for the \textit{switch} rule commuting with seq and times as the principal formula, follow a similar pattern.
\smallskip

}  %END fullproof


%%%%%%%%%%% DEEP CASES %%%%%%%%%%%%%%%%%%%%%%%%%%%%%%%%5

\item \textbf{Commutative cases deep in contexts.}
In many commutative cases, the bottommost rule does not interfere with the principal formula either directly or indirectly. Two such cases are presented for \textit{wen} as the principal connective. Other such cases use almost identical reasoning.

\begin{enumerate}[label*=\textbf{.\arabic*}]

\item Consider when a rule is applied outside the scope of the principal formula.
%
In this case, the bottommost rule in a proof is of the form
\[
\infer[\mbox{, such that $\vdash \wen x. U \cpar \context{ W }$.}]{
\wen x. U \cpar \context{ V }
}{
 \wen x. U \cpar \context{ W }
}
\]
% has a proof of length $k$, and also $\size{ \exists x U \cpar \context{ W } } \mleq \size{ \exists x U \cpar \context{ V } }$ by Lemma~\ref{lemma:bound}; thereby the size of the proof strictly decreases.
By the induction hypothesis, there exist formulae $P$ and $Q$ such that $\vdash U \cpar Q$ and $\nfv{x}{P}$ and either $P = Q$ or $P = \new x. Q$, and also $
\vcenter{\infer[]{
\context{ W }
}{
 P
}}
$.
%Furthermore, by induction, the size of the proof of $U \cpar Q$ is bounded above by the size of the proof of $\wen x U \cpar \context{ W }$.
Hence clearly derivation $
\vcenter{\infer[]{
\context{ V }
}{
\infer[]{
 \context{ W } 
}{
  P
}}}
$
 holds.
Furthermore, by Lemma~\ref{lemma:bound}, $\size{ \wen x. U \cpar \context{ W }} \prec \size{ U \cpar \context{ W }}$ 
and $\size{ U \cpar \context{ W }} \preceq \size{\wen x. U \cpar \context{ V }}$.
%A similar proof holds for $\new$.
%\smallskip

\fullproof{
Assume that the following application of any rule
$
\infer[]{
\left( T \cseq U \right) \cpar \context{ V }
}{
 \left( T \cseq U \right) \cpar \context{ W }
}$
is the bottommost rule in a proof,
such that $\vdash \left( T \cseq U \right) \cpar \context{ W }$.
By the induction hypothesis, there exist $n$-ary killing context $\tcontext{}$ and formulae $Q_i$ and $R_i$ such that $\vdash T \cpar Q_i$ and $\vdash U \cpar R_i$, for $1 \leq i \leq n$, such that $
\infer[]{
\context{ W } 
}{
\tcontext{ Q_1 \cseq R_1, \ldots, Q_n \cseq R_n }
}
$.
Hence, the derivation 
$
\infer[]{
\context{ V }
}{
\infer[]{
 \context{ W } 
 \tcontext{ Q_1 \cseq R_1, \ldots, Q_n \cseq R_n }
}}
$
holds, satisfying the induction invariant.
%A similar proof holds for \textit{times}.
%\smallskip


Alternatively, the bottommost rule may appear inside the context of principal formula without affecting the root connective of the principal formula. 
%We provide one illustrative case where 
%
Consider the case where \textit{seq} is the principal formula.
Assume that the following application of any rule is the bottommost rule in a proof
$
\infer[]{
\left( \context{T} \cseq V \right) \cpar W 
}{
 \left( \context{U} \cseq V \right) \cpar W
}
$
such that $\vdash \left( \context{U} \cseq V \right) \cpar W$ has a proof of length $n$. 
Hence by induction, there exist $n$-ary killing context $\tcontext{}$ and formulae $P_i$ and $Q_i$ such that $\vdash \context{U} \cpar P_i$ and $\vdash V \cpar Q_i$ hold and have a proof no longer than $n$, for $1 \leq i \leq n$, and also 
$
\infer[]{
W
}{
 \tcontext{ P_1 \cseq Q_1, \ldots, P_n \cseq Q_n }
}
$.
Hence we can construct the following proof of length no longer than $n+1$, for all $i$, as required:
$
\infer[]{
\context{T} \cpar P_i
}{
\infer[]{
   \context{U} \cpar P_i 
}{
 \cunit
}}
$.
%A similar proof holds for \textit{times}.
%\smallskip
}

\item Consider the case where the following application of any rule in a derivation of the form 
\[
\infer[]{
\wen x. \context{T} \cpar W
}{
 \wen x. \context{U} \cpar W
}
\]
is the bottommost rule is a proof of length $k+1$, where $\vdash \wen x. \context{U} \cpar W$ has a proof of length $k$.
%Furthermore, by Lemma~\ref{lemma:bound}, $\size{ \wen x \context{U} \cpar W } \mleq \size{ \wen x \context{T} \cpar W }$; hence the size of the proof strictly decreases.
Hence, by induction, there exist formulae $P$ and $Q$ such that $\vdash \context{U} \cpar Q$ and $\nfv{x}{P}$ and either $P = Q$ or $P = \new x. Q$, and also
$
\vcenter{
\infer[]{
W
}{
P
}}
$.
Furthermore, the size of the proof of $\context{U} \cpar Q$ is bounded above by the size of the proof of $\wen x. \context{U} \cpar W$; hence either $\size{\context{U} \cpar Q} \prec \size{\wen x. \context{U} \cpar W}$ or $\size{\context{U} \cpar Q} = \size{\wen x. \context{U} \cpar W }$ and the length of the proof of $U \cpar Q$ is bound by $k$.
%
The proof 
$
\vcenter{
\infer[]{
\context{T} \cpar Q
}{
\infer[]{
 \context{U} \cpar Q 
}{
   \cunit
}}}
$ can be constructed as required.
Furthermore, if $\size{\context{U} \cpar Q} \prec \wen x. \size{\context{U} \cpar W}$ then $\size{\context{U} \cpar Q} \prec \size{\wen x. \context{U} \cpar \context{ V }}$, by Lemma~\ref{lemma:bound}. Otherwise, $\size{\context{U} \cpar Q} = \size{\wen x. \context{U} \cpar W }$
hence $\size{U \cpar Q} \preceq \size{\wen x. U \cpar \context{ V }}$ by Lemma~\ref{lemma:bound} and the length of the proof of $\vdash \context{T} \cpar Q$ is $k+1$.
Thereby in either case, the size of the proof of $\context{T} \cpar Q$ is bounded above by the size of the proof of $\wen x. \context{T} \cpar W$.
\end{enumerate}
%A similar proof holds for $\new$.
%\smallskip
% --- end commutative.tex ---

\end{enumerate}

This covers all scenarios for the bottommost rule, hence splitting follows by induction over the size of the proof.
\end{proof}
% --- end split-times.tex ---
}{}
The final three splitting lemmas mainly involve checking commutative cases.
The commutative cases follow a similar pattern to the commutative cases in Lemma~\ref{lemma:split-times}.
% cases for \textit{new}, \textit{wen} and \textit{all} and \textit{with}.
% as corresponding commutative cases in other lemmas, hence proofs are omitted.
\begin{lemma}\label{lemma:split-exists}
If $\vdash \exists x. P \cpar Q$, then there exist formulae $V_i$ and values $v_i$ such that $\vdash P\sub{x}{v_i} \cpar V_i$, where $1 \leq i \leq n$, and $n$-ary killing context $\tcontext{}$ such that $\vcenter{\infer{Q}{ \tcontext{V_1, V_2, \hdots, V_n}}}$
and if $\tcontext{}$ binds $y$ then $\nfv{y}{\left(\exists x . P\right)}$.
\end{lemma}
%\input{split-exists}

%The splitting proof for atoms is generic for both channels and tests.
%Recall that the preorder $\leq$ over atoms could be any preorder including classical implication ($\Sigma \vdash \phi \Rightarrow \psi$) and equivalence ($\Sigma \vdash t = u$), with respect to signature $\Sigma$.
%
%
%
%
The proofs of the splitting lemmas for \textit{plus} and atoms offer no new insight or difficulties compared to their treatment in \textsf{MAV}~\cite{Horne2015}.
Similarly, to the above lemma for existential quantifiers, the proofs mainly involve commutative cases of a standard form.
\begin{lemma}\label{lemma:split-plus}
If $\vdash \left( P \ooplus Q \right) \cpar R$, then there exist formulae $W_i$ such that either $\vdash P \cpar W_i$ or $\vdash Q \cpar W_i$ where $1 \leq i \leq n$, and $n$-ary killing context $\tcontext{}$ such that $\vcenter{\infer{R}{\tcontext{W_1, W_2, \hdots, W_n}}}$
and if $\tcontext{}$ binds $x$ then $\nfv{x}{\left(P \ooplus Q\right)}$.
\end{lemma}


\begin{lemma}\label{lemma:split-atoms}
The following statements hold, for any atom $\alpha$, where if $\tcontext{}$ binds $x$ then $\nfv{x}{\alpha}$.
\begin{itemize}
\item If $\vdash \co{\alpha} \cpar Q$, then there exist $n$-ary killing context $\tcontext{}$ such that 
$\vcenter{\infer{Q}{\tcontext{\alpha, \alpha, \hdots, \alpha}}}$.

\item If $\vdash \alpha \cpar Q$, then there exist $n$-ary killing context $\tcontext{}$ such that 
$\vcenter{\infer{Q}{\tcontext{\co{\alpha}, \co{\alpha}, \hdots, \co{\alpha}}}}$.
\end{itemize}
\end{lemma}

% --- end splitting.tex ---
% --- begin context.tex ---
\section{Context Reduction and the Admissibility of Co-rules}
\label{section:context}

The splitting lemmas  in the previous section are formulated for sequent-like \textit{shallow contexts}.
By applying splitting repeatedly, \textit{context reduction} (Lemma~\ref{lemma:context}) is established, which can be used to extends normalisation properties to an arbitrary (deep) context.
In particular, we extend a series of proof normalisation properties called \textit{co-rule elimination} properties to any context, by first establishing the normalisation property in a shallow context, then applying context reduction to extend to any context.
Together, these \textit{co-rule elimination} properties establish cut elimination, by eliminating each connective directly involved in a cut one-by-one.

%The splitting lemmas in the previous section only work for a 
%

\subsection{Extending from a sequent-like context to a deep context}

Context reduction extends rules simulated by splitting to any context.
This appears to be the first context reduction lemma in the literature to handle first-order quantifiers.
Of particular note is the use of substitutions to account for the effect of existential quantifiers in the context.
The trick is to first establish the following stronger invariant.
% (this can be considered to be the induction invariant for the context lemma).
\begin{lemma}\label{lemma:invariant}
If $\vdash \context{T}$, then there exist formulae $U_i$ and substitutions $\sigma_i$, for $1 \leq i \leq n$, and $n$-ary killing context $\tcontext{}$ such that $\vdash T\sigma_i \cpar U_i$; and, for any formula $V$ there exist $W_i$ such that either $W_i = V\sigma_i \cpar U_i$ or $W_i = \cunit$ and the following holds:
$
\vcenter{
\infer[]{
 \context{V}
}{
 \tcontext{ W_1, W_2, \hdots, W_n }
}
}$.
\end{lemma}

%\input{context-invariant}

Having established the above stronger invariant, the context lemma follows directly.
\begin{lemma}[Context reduction]
\label{lemma:context}
If $\vdash P\mathclose\sigma \cpar R$ yields that $\vdash Q\mathclose\sigma \cpar R$, for any formula $R$ and substitution of terms for variables $\sigma$, then $\vdash \context{ P }$ yields $\vdash \context{ Q }$, for any context $\context{}$.
\end{lemma}
\tohide{
% --- begin contextproof.tex ---
\begin{proof}
Assume that for any formula $U$, $\vdash S \cpar U$ yields $\vdash T \cpar U$, and fix any context $\context{}$ such that $\vdash \context{S}$ holds. By Lemma~\ref{lemma:invariant}, there exist $n$-ary killing context $\tcontext{}$ and,
for $1 \leq i \leq n$, $P_i$ such that either $P_i = \cunit$ or there exists $W_i$ where $P_i = T \cpar W_i$ and $\vdash S \cpar W_i$, and furthermore
$
\vcenter{
\infer[]{
\context{T}
}{
\tcontext{P_1, \hdots, P_n}
}}$.
Since, by our assumption, also $\vdash T \cpar W_i$ holds for $1 \leq i \leq n$, the proof 
$
\vcenter{
\infer[]{
\context{T}
}{
\infer[]{
 \tcontext{P_1, \hdots, P_n} 
}{
\infer[]{
 \tcontext{\cunit, \hdots, \cunit} 
}{
 \cunit
}}}}
$ can be constructed.
Therefore $\vdash \context{T}$ holds.
\end{proof}
% --- end contextproof.tex ---
Note that the case for existential quantifiers will not work for second-order quantifiers, since termination of the induction is reliant on the size of the term-free part of the formula being reduced. Thus the techniques in the above proof apply to first-order quantifiers only.
}{
%%%%%%%%%%%%%%%%%%%%%%%%%%%%%%%%%% SHORTER VERSION %%%%%%%%%%%%%%%%%%%%%%%%%%%%%%%%%%%%
The subtelty of context reduction is to accommodate \textit{plus} and \textit{some} by the following stronger induction invariant:
If $\vdash \context{T}$, then there exist formulae $U_i$ and substitutions $\sigma_i$ such that $\vdash T\mathclose{\sigma_i} \cpar U_i$, for $1 \leq i \leq n$; and $n$-ary killing context $\tcontext{}$  such that for any formula $V$ there exist $W_i$ such that either $W_i = V\mathclose{\sigma_i} \cpar U_i$ or $W_i = \cunit$, for $1 \leq i \leq n$, and the following holds:
$
\vcenter{
\infer[]{
 \context{V}
}{
 \tcontext{ W_1, W_2, \hdots, W_n }
}
}$.
%%%%%%%%%%%%%%%%%%%%%%%%%%%%%%%%%% SHORTER VERSION %%%%%%%%%%%%%%%%%%%%%%%%%%%%%%%%%%%%
}
\begin{comment}
The proof first establishes the following stronger property.
If $\vdash \context{T}$, then there exist predicates $U_i$ and substitutions $\sigma_i$, for $1 \leq i \leq n$, and $n$-ary killing context $\tcontext{}$ such that $\vdash T\sigma_i \cpar U_i$; and, for any local predicate $V$ these exists $W_i$ such that either $W_i = V\sigma_i \cpar U_1$ or $W_i = \cunit$ and the following holds:
$
 \context{V} \longrightarrow \tcontext{ W_1, W_2, \hdots, W_n }
$.
\end{comment}

\begin{comment}
\proof{
Firstly we establish, by induction on the size of the context, the following stronger property.
If $\vdash \context{T}$, then there exist predicates $U_i$ and substitutions $\sigma_i$, for $1 \leq i \leq n$, and $n$-ary killing context $\tcontext{}$ such that $\vdash T\sigma_i \cpar U_i$; and, for any local predicate $V$ these exists $W_i$ such that either $W_i = V\sigma_i \cpar U_1$ or $W_i = \cunit$ and the following holds:
\[
 \context{V} \longrightarrow \tcontext{ W_1, W_2, \hdots, W_n }
\]
\medskip

The base case is when the context is of the from $\chole \cpar P$, where the hole appears directly inside a parallel composition, and $\vdash T \cpar P$ holds, in which case we are done.
\medskip


Consider the case for a context of the form $\exists x \context{} \cpar P$, where $\vdash \exists x \context{ T } \cpar P$.
By Lemma~\ref{lemma:splitting},
there exist predicates $Q_i$ and values $v_i$ such that $\vdash \context{ T } \sub{x}{v_i} \cpar Q_i$, for $1 \leq i \leq n$;
and $n$-ary killing context $\tcontext{}$ such that the following derivation holds.
\[
P \longrightarrow \tcontext{ Q_1, Q_2, \hdots, Q_n }
\]

For context $\context{}$ and any predicate $U$, let $\contextn{i}{}$ and $\sigma_{i}$ be such that $\context{U}\sub{x}{v_i} \equiv \contextn{i}{U\sigma_{i}}$.
Notice that for first-order quantifiers, the substitutions does not increase the size of the predicate part of the context. It can only increases the size of values in atoms, which are not counted in this induction.

Since $\vdash \context{ T } \sub{x}{v_i} \cpar Q_i$ holds, then $\vdash \contextn{i}{ T\sigma_i } \sub{x}{v_i} \cpar Q_i$ holds.
Therefore, by induction on the size of the predicate part of the context (n.b.\ not counting the size of values), there exists predicate $V^i_j$ such that either $V^i_j = \cunit$ or $V^i_j = \left(U\sigma_i\right) \sigma^i_j \cpar W^i_j$, where $\vdash \left(T \sigma_i\right) \sigma^i_j \cpar W^i_j$, for $1 \leq j \leq m_i$;
and $m_i$-ary killing context $\tcontextn{i}{}$
%\bogdan{shouldn't be $\tcontextn{i}{}$ here?}
such that the following derivation holds.
\[
\context{ U } \sub{x}{v_i} \cpar Q_i
\equiv
\contextn{i}{ U\sigma_i } \cpar Q_i
\longrightarrow
\tcontextn{i}{ V^i_1, V^i_2, \hdots, V^i_{m_i} }
\]
Hence the following derivation can be constructed for all predicates $U$.
\[
\begin{array}{rl}
\exists x \context{ U } \cpar P
\longrightarrow&
\exists x \context{ U } \cpar \tcontext{ Q_1, \hdots, Q_n }
\\
\longrightarrow&
\exists x \context{ U } \cpar \tcontext{ Q_i \colon 1 \leq i \leq n }
\\
\longrightarrow&
\tcontext{ \exists x \context{ U } \cpar Q_i \colon 1 \leq i \leq n }
\\
\longrightarrow&
\tcontext{ \context{ U } \sub{x}{v_i} \cpar Q_i \colon 1 \leq i \leq n }
\\
\longrightarrow&
\tcontext{ \tcontextn{i}{ V^i_j \colon 1 \leq j \leq m_i } \colon 1 \leq i \leq n }
\end{array}
\]
where $V^i_j = \cunit$ or $V^i_j = U \left(\sigma_i \cdot \sigma^{i}_j\right) \cpar W^i_j$, such that $\vdash T \left(\sigma_i \cdot \sigma^{i}_j\right) \cpar W^i_j$, for all $i$ and $j$. Thereby the induction invariant is satisfied.
\medskip



Consider the case for a context of the form $\left( \context{ } \tensor U \right) \cpar P$ such that $\vdash \left( \context{T} \tensor U \right) \cpar P$.
By Lemma~\ref{lemma:splitting}, there exist predicates $Q_i$, $R_i$ such that $\vdash \context{T} \cpar Q_i$ and $\vdash U \cpar R_i$, for $1 \leq i \leq n$, and $n$-ary killing context $\tcontext{}$ such that:
\[
 P \longrightarrow \tcontext{ Q_1 \cpar R_1, \hdots, Q_n \cpar R_n }
\]

By the induction hypothesis, for every $i$ such that $1 \leq i \leq n$, and for any predicate $V$, there exist predicates $S^i_j$ such that either $S^i_j = \cunit$ or $S^i_j = V\sigma^i_j \cpar W^i_j$ for some predicate $W^i_j$ and substitution $\sigma^i_j$ where $\vdash T\sigma^i_j \cpar W^i_j$ holds, for $1 \leq j \leq m_i$; and, $m_i$-ary killing context $\tcontextn{i}{}$
such that the following derivations holds:
\[
\context{V} \cpar Q_i \longrightarrow \tcontextn{i}{ S^i_1, \hdots, S^i_{m_i} }
\]
Hence we can construct the following derivation for any predicate $V$.
\[
\begin{array}{rl}
\left( \context{V} \tensor U \right) \cpar P
  \longrightarrow& \left( \context{V} \tensor U \right) \cpar \tcontext{ Q_1 \cpar R_1, \hdots, Q_n \cpar R_n } \\
  \longrightarrow& \tcontext{ \left( \context{V} \tensor U \right) \cpar Q_i \cpar R_i \colon 1 \leq i \leq n } \\
  \longrightarrow& \tcontext{ \left( \context{V} \cpar Q_i \right) \tensor \left( U \cpar R_i \right) \colon 1 \leq i \leq n } \\
  \longrightarrow& \tcontext{ \context{V} \cpar Q_i \colon 1 \leq i \leq n } \\
  \longrightarrow& \tcontext{ \tcontextn{i}{ S^i_1, \hdots, S^i_{m_i} } \colon 1 \leq i \leq n } \\
\end{array}
\]
The final predicate above forms a $\left(\sum_i m_i\right)$-ary killing context such that either $S^i_j = \cunit$ or $S^i_j = V\sigma^i_j \cpar W^i_j$, such that $\vdash T\sigma^i_j \cpar W^i_j$, for all $i$, $j$. Thereby, the induction invariant is satisfied.
The two cases for when the hole appears on the left or right of a sequential composition are similar to the above case for tensor.
\medskip


Consider the case for a context of the form $\left( \context{ } \wwith U \right) \cpar P$ such that $\vdash \left( \context{T} \wwith U \right) \cpar P$.
By Lemma~\ref{lemma:branching}, $\vdash \context{T} \cpar P$ and $\vdash U \cpar P$ hold.

Hence, by the induction hypothesis, there exist killing context $\tcontext{}$ and predicates $V_i$ for $1 \leq i \leq n$ such that $\vdash T \cpar V_i$; and, for all predicates $W$, for $1 \leq i \leq n$, there exists $S_i$ such that either $S_i = \cunit$ or $S_i = W \cpar V_i$ and the following derivation holds:
\[ \context{W} \cpar P
    \longrightarrow \tcontext{ S_1, \hdots , S_n }
\]
Hence we can construct a derivation as follows for all predicates $W$.
\[
\begin{array}{rl}
\left( \context{W} \wwith U \right) \cpar P
  \longrightarrow& \context{W} \cpar P \wwith  U \cpar P \\
  \longrightarrow& \context{W} \cpar P \wwith \cunit \\
  \longrightarrow& \tcontext{ S_1, \hdots , S_n } \wwith \cunit
\end{array}
\]
where $\tcontext{} \wwith \left\{\ \cdot\ \right\}$ is a $\left(n+1\right)$-ary killing context
and $S_{n+1} = \cunit$, thereby satisfying the induction invariant.
The case when the hole appears on the right of an external choice is similar.



Consider the case for a context of the form $\left( \context{ } \ooplus U \right) \cpar P$ where $\vdash \left( \context{T} \ooplus U \right) \cpar P$.
By Lemma~\ref{lemma:splitting}, there exist killing context $\tcontext{}$ and $V_i$ for $1 \leq i \leq n$ such that either $\vdash \context{T} \cpar V_i$ or $\vdash U \cpar V_i$ and the following derivation holds:
\[
 P \longrightarrow \tcontext{ V_1, \hdots, V_n }
\]
Now, by the induction hypothesis, if $\vdash \context{T} \cpar V_i$ holds, then there exist $\tcontextn{i}{}$ and predicates $W^i_j$ for $1 \leq j \leq m_i$ such that $\vdash T \cpar W^i_j$; and, for any predicate $Q$ for $1 \leq j \leq m_i$, there exists $S^i_j$ where wither $S^i_j = \cunit$ or $S^i_j = Q \cpar W^i_j$ and the following derivation holds:
\[
 \context{Q} \cpar V_i 
   \longrightarrow \tcontextn{i}{ S^i_1, \hdots, S^i_{m_i} }
\]
Hence we can construct the following derivation for any predicate $Q$.
\[
\begin{array}{rl}
\left( \context{Q} \ooplus U \right) \cpar P
  \longrightarrow& \left( \context{Q} \ooplus U \right) \cpar \tcontext{ V_1, \hdots, V_n } \\
  \longrightarrow& \tcontext{ \left( \context{Q} \ooplus U \right) \cpar V_1, \hdots, \left( \context{Q} \ooplus U \right) \cpar V_n } \\
  \longrightarrow& \tcontext{ R_1, \hdots, R_n }
\end{array}
\]
where $R_i$ is defined as follows.
\[
R_i = 
\left\{
\begin{array}{lr}
\cunit                                                     & \mbox{if $\vdash U \cpar V_i$} \\
\tcontextn{i}{ S^i_1, \hdots, S^i_{m_i} }  & \mbox{otherwise}
\end{array}
\right.
\] 
The above is well-defined since if $\vdash U \cpar V_i$ holds, then 
\[
\left( \context{Q} \ooplus U \right) \cpar V_i
 \longrightarrow U \cpar V_i
 \longrightarrow \cunit = R_i
\]
and, if $\vdash U \cpar V_i$ does not hold, then $\vdash \context{T} \cpar V_i$ must hold;
hence the following derivation can be applied:
\[
\begin{array}{rl}
\left( \context{Q} \ooplus U \right) \cpar V_i
  \longrightarrow& \context{Q} \cpar V_i 
  \longrightarrow \tcontextn{i}{ S^i_1, \hdots, S^i_{m_i} } = R_i
\end{array}
\]
Hence the induction invariant is satisfied.

%\hfill\qed
}%\end{proof}
\end{comment}
%In the second order case, second order variables may appear in the context hence the substitution may cause the size of the predicate to blow up.

% --- end context.tex ---
% --- begin corule.tex ---
\subsection{Cut elimination as co-rule elimination}

\begin{figure}
\begin{gather*}
\infer[\mbox{(atomic co-interaction)}]{
\context{ \cunit }
}{
\context{ \alpha \tensor \co{\alpha} }
}
\qquad
\infer[\mbox{(co-select1)}]{
\context{ P \sub{x}{v} }
}{
\context{ \forall x. P }
}
\\[9pt]
\infer[\mbox{(co-sequence)}]{
\context{(P \tensor U) \cseq (Q \tensor V)}
}{
\context{(P \cseq Q) \tensor (U \cseq V)}
}
\qquad\qquad
\infer[\mbox{(co-external)}]{
\context{ (P \cpar R) \ooplus (Q \cpar S) }
}{
\context{ (P \ooplus Q) \cpar S }
}
\\[9pt]
\infer[\mbox{(co-tidy)}]{
\context{ \cunit }
}{
\context{ \cunit \ooplus \cunit }
}
\qquad\qquad
\infer[\mbox{(co-left)}]{
\context{ P}
}{
\context{ P \wwith Q }
}
\qquad\qquad
\infer[\mbox{(co-right)}]{
\context{ Q}
}{
\context{ P \wwith Q }
}
\\[9pt]
\infer[\mbox{(co-extrude1)}]{
\context{ \mathopen{\exists x.}\left( P \tensor R \right) }
}{
\context{ \exists x. P \tensor R }
}
%\mbox{ only if $\nfv{x}{R}$}
\qquad\qquad\qquad
\infer[\mbox{(co-tidy1)}]{
\context{ \cunit }
}{
\context{ \exists x. \cunit }
}
\\[9pt]
\infer[\mbox{(co-close)}]{
\context{ \mathopen{\wen x .}\left( P \tensor Q \right) }
}{
\context{ \new x. P \tensor \wen x. Q }
}
\qquad\qquad\qquad
\infer[\mbox{(co-tidy name)}]{
\context{ \cunit }
}{
\context{ \wen x. \cunit }
}
%\mbox{ only if $\nfv{x}{R}$}
%\mbox{ for $\quantifier \in \left\{ \new, \wen \right\}$}
%~~\mbox{(co right name)}
\end{gather*}
\caption[Co-rules extending the system \textsf{MAV1} to \textsf{SMAV1}.]
{Co-rules extending the system \textsf{MAV1} to \textsf{SMAV1}, where 
%$\protect{\quantifier \in \left\{ \new, \wen \right\}}$ and
 $\nfv{x}{R}$.}
\label{figure:co-rules}
\end{figure}
\begin{comment}
\\[9pt]
\hline
\hfill
\textit{co-medials}
\\
\context{ \left(P \ooplus U\right) \cseq \left(Q \ooplus V\right) }
\longrightarrow
\context{ \left(P \cseq Q\right) \ooplus \left(U \cseq V\right) }
%~~\mbox{(co-medial)}
\qquad
\context{ \exists x .P \cseq \exists x .S }
\longrightarrow
\context{ \mathopen{\exists x .}\left(P \cseq S \right) }
%~~\mbox{(co-medial1)}
\\[9pt]
\context{ \mathopen{\wen x .}\left( P \tensor R \right) }
\longrightarrow
\context{ \wen x .P \tensor R }
%\mbox{ for $\quantifier \in \left\{ \new, \wen \right\}$}
%\mbox{ only if $\nfv{x}{R}$}
%~~\mbox{(co extrude name)}
\qquad%\\[9pt]
\context{ \mathopen{\new x .}\left( P \tensor S \right) }
\longrightarrow
\context{ \new x. P \tensor \new x. S }
%~~\mbox{(co-medial name)}
\\[9pt]
\context{ \mathopen{\new x .}\left( P \cseq R \right) }
\longrightarrow
\context{ \new x .P \cseq R }
%\mbox{ only if $\nfv{x}{R}$}
%~~\mbox{(left new)}
\qquad
\context{ \mathopen{\new x .}\left( R \cseq Q \right) }
\longrightarrow
\context{ R \cseq \new x. Q }
%\mbox{ only if $\nfv{x}{R}$}
%~~\mbox{(right new)}
\\[9pt]
\context{ \quantifier y. \exists x .P }
\longrightarrow
\context{ \exists x .\quantifier y. P }
%\mbox{ for $\quantifier \in \left\{\new, \wen\right\}$}
%~~\mbox{(some name)}
\qquad
\context{ \mathopen{\quantifier x .}\left( P \ooplus S \right) }
\longrightarrow
\context{ \quantifier x .P \ooplus \quantifier x. S }
%\mbox{ for $\quantifier \in \left\{ \new, \wen \right\}$}
%~~\mbox{(plus name)}
\\[9pt]
\context{ \mathopen{\quantifier x .}\left( P \ooplus R \right) }
\longrightarrow
\context{ \quantifier x .P \ooplus R }
%\mbox{ only if $\nfv{x}{R}$}
%\mbox{ for $\quantifier \in \left\{ \new, \wen \right\}$}
%~~\mbox{(co left name)}
\qquad
\context{ \mathopen{\quantifier x .}\left( R \ooplus Q \right) }
\longrightarrow
\context{ R \ooplus \quantifier x. Q }
\end{comment}

%For every rule there is a \textit{co-rule}.
For a rule of the form $\vcenter{\infer{P}{Q}}$, there is a  corresponding \textit{co-rule} of the form $\vcenter{\infer{\co{Q}}{\co{P}}}$, where premise and conclusion are interchanged and each formula is dualised using negation. 
The rules \textit{switch}, \textit{fresh} and \textit{new wen} are their own co-rules. 
Also the co-rule of the \textit{medial new} rule is an instance of the \textit{suspend} rule.
All other rules give rise to distinct co-rules, presented in Figure~\ref{figure:co-rules}.
Note co-rules with no role in cut elimination are ommitted from the figure.


%For example \textit{close} has co-rule $\context{ \mathopen{\wen x.} \left( P \tensor Q \right) } \longrightarrow \context{ \wen x. P \tensor \new x. Q }$
%and \textit{extrude1} has co-rule if $\nfv{x}{Q}$ then $\context{ \mathopen{\exists x.} \left( P \tensor Q \right) } \longrightarrow \context{ \exists x. P \tensor Q }$.

The following nine lemmas each establish that a co-rule is admissible in \textsf{MAV1}.
Only the following co-rules need be handled directly in order to establish cut elimination:
\textit{co-close}, \textit{co-tidy name}, \textit{co-extrude1}, \textit{co-select1}, \textit{co-tidy1}, \textit{co-left}, \textit{co-right}, \textit{co-external}, \textit{co-tidy}, \textit{co-sequence} and \textit{atomic co-interaction}.
In each case, the proof proceeds by applying splitting in a shallow context, forming a new proof, and finally applying Lemma~\ref{lemma:context}. 
Each co-rule can be treated independently, hence are established as separate lemmas.
%
\begin{lemma}[co-close]
\label{lemma:co-close}
If $\vdash \context{ \wen x. P \tensor \new x. Q }$ holds
then $\vdash \context{ \mathopen{\wen x.} \left( P \tensor Q \right) }$ holds.
\end{lemma}
% --- begin coclose.tex ---
\begin{proof}
Assume that $\vdash \left(\wen x. P \tensor \new x. Q \right)\mathclose\sigma \cpar R$ for some substitution of terms for variables $\sigma$.
By Lemma~\ref{lemma:split-times}, there exist $S_i$ and $T_i$ such that $\vdash \left(\wen x. P\right)\mathclose\sigma \cpar S_i$ and $\vdash \left(\new x. Q\right)\mathclose\sigma \cpar T_i$, for $1 \leq i \leq n$,
and $n$-ary killing context such that the derivation
\[
\infer{
R
}{
 \tcontext{ S_i \cpar T_i \colon 1 \leq i \leq n }
}
\] holds.
%
Also for some $y$ such that $\nfv{y}{\wen x. P}$, $\nfv{y}{\new x. Q}$ and $\nfv{y}{\sigma}$, $\left(\wen x. P\right)\mathclose\sigma \equiv \mathopen{\wen y.} \left( P\sub{x}{y}\mathclose\sigma \right)$ and $\left(\new x. Q\right)\mathclose\sigma \equiv \mathopen{\new y.}\left(Q\sub{x}{y}\mathclose\sigma\right)$, where $\nfv{y}{\sigma}$ is defined such that $y$ does not appear in the domain of $\sigma$ nor free in any term in the range of $\sigma$.
Hence both $\vdash \mathopen{\wen y.}\left( P\sub{x}{y}\mathclose\sigma \right) \cpar S_i$ and $\vdash \mathopen{\new y.} \left(Q\sub{x}{y}\mathclose\sigma\right) \cpar T_i$ hold.

Hence, by Lemma~\ref{lemma:split-times}, there exist $U_i$ and $\hat{U}_i$ such that $\vdash P\sub{x}{y}\mathclose\sigma \cpar \hat{U}_i$ and either $U_i = \hat{U}_i$ or $U_i = \new y. \hat{U}_i$, and also the derivation
$
\vcenter{
\infer{
S_i
}{
 U_i
}}
$ holds.

Similarly, by Lemma~\ref{lemma:split-times}, there exist $W_i$ and $\hat{W}_i$ such that $\vdash Q\sub{x}{y}\mathclose\sigma \cpar \hat{W}_i$ and either $W_i = \hat{W}_i$ or $W_i = \wen y. \hat{W}_i$, and also the derivation 
$
\vcenter{
\infer{
T_i
}{
 W_i
}}
$ holds.

There are four cases to consider for each $i$. Three of the cases are as follows.
\begin{itemize}
\item
If $U_i = \new y. \hat{U}_i$ and $W_i = \wen y. \hat{W}_i$ then 
\[
\infer[.]{
\new y. \hat{U}_i \cpar \wen y. \hat{W}_i
}{
\new y. \left( \hat{U}_i \cpar \hat{W}_i\right)
}
\]

\item 
If $U_i = \hat{U}_i$, $\nfv{y}{\hat{U}_i}$, and $W_i = \wen y. \hat{W}_i$, 
then %the derivation
\[
\infer[.]{
U_i \cpar \wen y. \hat{W}_i
}{
\infer[]{
\mathopen{\wen y.}\left( \hat{U}_i \cpar \hat{W}_i\right)
}{
\mathopen{\new y.} \left( \hat{U}_i \cpar \hat{W}_i\right)
}}
\]
%can be constructed.

\item
If $U_i = \new y. \hat{U}_i$ and $W_i = \hat{W}_i$, such that $\nfv{y}{\hat{W}_i}$ then %the derivation
\[
\infer[.]{
\new x. U_i \cpar \hat{W}_i
}{
\mathopen{\new y.} \left( \hat{U}_i \cpar \hat{W}_i\right)
}
\]
%can be constructed.
\end{itemize}
Thereby in any of the above three cases the following derivation can be constructed.
\[
%\begin{array}{rl}
\infer[]{
\left(\mathopen{\wen x.} \left( P \tensor Q \right)\right)\mathclose{\sigma} \cpar U_i \cpar W_i 
}{
\infer[]{
\left(\mathopen{\wen x.} \left( P \tensor Q \right)\right)\mathclose{\sigma} \cpar \new y. \left( \hat{U}_i \cpar \hat{W}_i \right)
}{
\mathopen{\new y.} \left(\left( P \tensor Q \right)\sub{x}{y}\mathclose\sigma \cpar \hat{U}_i \cpar
   \hat{W}_i \right)
}}
%\end{array}
\]
In the fourth case $U_i = \hat{U}_i$ and $W_i = \hat{W}_i$, such that $\nfv{y}{\hat{W}_i}$ and  $\nfv{y}{\hat{U}_i}$ yielding the following.%, in which case the following derivation can be constructed.
\[
%\begin{array}{rl}
\infer[]{
\left(\mathopen{\wen x.} \left( P \tensor Q \right)\right)\mathclose\sigma \cpar \hat{U}_i \cpar
   \hat{W}_i
}{
\infer[]{
\mathopen{\new y.} \left(\left( P \tensor Q\right)\sub{x}{y}\mathclose\sigma\right) \cpar \hat{U}_i \cpar
   \hat{W}_i
}{
\mathopen{\new y.} \left(\left( P \tensor Q \right)\sub{x}{y}\mathclose\sigma \cpar \hat{U}_i \cpar
   \hat{W}_i \right)
}}
%\end{array}
\]
By applying one of the above possible derivations for every $i$, the following proof can be constructed.
\[
\infer[]{
\left(\mathopen{\wen x.} \left( P \tensor Q \right)\right)\mathclose\sigma \cpar R
}{
\infer[]{
\left(\mathopen{\wen x.} \left( P \tensor Q \right)\right)\mathclose\sigma \cpar \tcontext{S_i \cpar T_i \colon 1 \leq i \leq n }
}{
\infer[]{
\left(\mathopen{\wen x.} \left( P \tensor Q \right)\right)\mathclose\sigma \cpar
\tcontext{
 U_i
 \cpar 
 W_i
 \colon 1 \leq i \leq n
}
}{
\infer[]{
\tcontext{
 \left(\mathopen{\wen x.} \left( P \tensor Q \right)\right)\mathclose\sigma
 \cpar
 U_i
 \cpar 
 W_i
 \colon 1 \leq i \leq n
}
}{
\infer[]{
\tcontext{
 \mathopen{\new y.} \left(\left( P \tensor Q \right)\sub{x}{y}\mathclose\sigma \cpar \hat{U}_i \cpar
   \hat{W}_i \right)
 \colon 1 \leq i \leq n
}
}{
\infer[]{
\tcontext{
   \mathopen{\new y.} \left(\left( P\sub{x}{y}\mathclose\sigma \cpar \hat{U}_i \right) \tensor
   \left( Q\sub{x}{y}\mathclose\sigma \cpar \hat{W}_i \right)\right)
 \colon 1 \leq i \leq n
}
}{
\infer[]{
\tcontext{
   \new y. \cunit
 \colon 1 \leq i \leq n
}
}{
 \cunit
}}}}}}}
\]
%Hence $\vdash \left(\wen x \left( P \tensor Q \right)\right)\sigma \cpar R$.
Therefore, by Lemma~\ref{lemma:context}, for all contexts $\context{}$,
if $\vdash \context{ \wen x. P \tensor \new x. Q }$ then
$\vdash \context{ \mathopen{\new x.} \left( P \tensor Q \right) }$.
\end{proof}
\begin{comment}
\begin{proof}
Assume that $\vdash \left(\wen x P \tensor \new x Q \right)\sigma \cpar R$ for some substitution $\sigma$.
By Lemma~\ref{lemma:split-times}, there exist $S_i$ and $T_i$ such that $\vdash \left(\wen x P\right)\sigma \cpar S_i$ and $\vdash \left(\new x Q\right)\sigma \cpar T_i$, for $1 \leq i \leq n$,
and $n$-ary killing context such that the derivation
$
R \longrightarrow \tcontext{ S_i \cpar T_i \colon 1 \leq i \leq n }
$ holds.

Also observe that for some $y$ such that $\nfv{y}{\wen x P}$, $\nfv{y}{\new x Q}$ and $\nfv{y}{\sigma}$, $\left(\wen x P\right)\sigma \equiv \wen y \left( P\sub{x}{y}\sigma \right)$ and $\left(\new x Q\right)\sigma \equiv \new y \left(Q\sub{x}{y}\sigma\right)$, where $\nfv{y}{\sigma}$ is defined such that $y$ does not appear in the domain of $\sigma$ nor free in any terms in the range of $\sigma$.
Hence both $\vdash \wen y \left( P\sub{x}{y}\sigma \right) \cpar S_i$ and $\vdash \new y \left(Q\sub{x}{y}\sigma\right) \cpar T_i$ hold.

Hence, by Lemma~\ref{lemma:split-times}, there exist $U^i_j$ and $\hat{U}^i_j$ such that $\vdash P\sub{x}{y}\sigma \cpar \hat{U}^i_j$ and either $U^i_j = \hat{U}^i_j$ or $U^i_j \equiv \new y \hat{U}^i_j$ and $\vdash P\sub{x}{y}\sigma \cpar \hat{U}^i_j$, for $1 \leq j \leq m_i$, and $m_i$-ary killing context $\tcontextmn{i}{0}{}$ such that the derivation
$
S_i \longrightarrow \tcontextmn{i}{0}{
 U^i_j \colon 1 \leq j \leq m_i
}
$ holds.

Also, by Lemma~\ref{lemma:split-times}, there exist $W^i_h$ and $\hat{W}^i_h$ such that $\vdash Q\sub{x}{y}\sigma \cpar \hat{W}^i_h$ and either $W^i_h = \hat{W}^i_h$ or $W^i_h = \wen y \hat{W}^i_h$, for $1 \leq h \leq k_i$, and $k_i$-ary killing context $\tcontextmn{i}{1}{}$ such that the derivation 
$
T_i \longrightarrow
\tcontextmn{i}{1}{
 W^i_h
 \colon
 1 \leq h \leq k_i
}
$ holds.

There are four cases to consider for each $i$, $j$ and $h$. Three of the cases are as follows.
\begin{itemize}
\item
If $U^i_j = \new y \hat{U}^i_j$ and $W^i_h = \wen y \hat{W}^i_h$ then the derivation
$
\new y \hat{U}^i_j \cpar \wen y \hat{W}^i_h
\longrightarrow
\new y \left( \hat{U}^i_j \cpar \hat{W}^i_h\right)
$
can be constructed.

\item 
If $U^i_j = \hat{U}^i_j$, such that $\nfv{y}{\hat{U}^i_j}$, and $W^i_h = \wen y \hat{W}^i_h$ then the derivation
$
U^i_j \cpar \wen y \hat{W}^i_h
\longrightarrow
\wen y \left( \hat{U}^i_j \cpar \hat{W}^i_h\right)
\longrightarrow
\new y \left( \hat{U}^i_j \cpar \hat{W}^i_h\right)
$
can be constructed.

\item
If $U^i_j = \new y \hat{U}^i_j$ and $W^i_h = \hat{W}^i_h$, such that $\nfv{y}{\hat{W}^i_j}$, then the derivation
$
\new x U^i_j \cpar \hat{W}^i_h
\longrightarrow
\new y \left( \hat{U}^i_j \cpar \hat{W}^i_h\right)
$
can be constructed.
\end{itemize}
Thereby in any of the above three cases the following derivation can be construced.
\[
%\begin{array}{rl}
\left(\wen x \left( P \tensor Q \right)\right)\sigma \cpar U^i_j \cpar W^i_h 
\longrightarrow
\left(\wen x \left( P \tensor Q \right)\right)\sigma \cpar \new y \left( \hat{U}^i_j \cpar \hat{W}^i_h \right)
\longrightarrow
\new y \left(\left( P \tensor Q \right)\sub{x}{y}\sigma \cpar \hat{U}^i_j \cpar
   \hat{W}^i_h \right)
%\end{array}
\]

In the fourth cases $U^i_j = \hat{U}^i_j$, such that $\nfv{y}{\hat{U}^i_j}$, and $W^i_h = \hat{W}^i_h$, such that $\nfv{y}{\hat{W}^i_j}$, in which case the following derivation can be constructed.
\[
%\begin{array}{rl}
\left(\wen x \left( P \tensor Q \right)\right)\sigma \cpar \hat{U}^i_j \cpar
   \hat{W}^i_h
\longrightarrow
\new y \left(\left( P \tensor Q\right)\sub{x}{y}\sigma\right) \cpar \hat{U}^i_j \cpar
   \hat{W}^i_h
\longrightarrow
\new y \left(\left( P \tensor Q \right)\sub{x}{y}\sigma \cpar \hat{U}^i_j \cpar
   \hat{W}^i_h \right)
%\end{array}
\]
By applying one of the above possible derivations for every $i$, $j$ and $h$, the following proof can be constructed.
\[
\begin{array}{l}
\left(\wen x \left( P \tensor Q \right)\right)\sigma \cpar R
\\
\begin{array}{rl}
\longrightarrow&
\left(\wen x \left( P \tensor Q \right)\right)\sigma \cpar \tcontext{S_i \cpar T_i \colon 1 \leq i \leq n }
\\
\longrightarrow&
%\begin{array}[t]{l}
\left(\wen x \left( P \tensor Q \right)\right)\sigma \cpar
%\\
\tcontext{
% \begin{array}{l}
 \tcontextmn{i}{0}{ U^i_j \colon 1 \leq j \leq m_i }
 \cpar 
 \tcontextmn{i}{1}{
  W^i_h
  \colon
  1 \leq h \leq k_i
 }
 \colon 1 \leq i \leq n
% \end{array}
}
%\end{array}
\\
\longrightarrow&
\tcontext{
 %\begin{array}{r}
 \tcontextmn{i}{0}{
  %\begin{array}{r}
  \tcontextmn{i}{1}{
   %\begin{array}{l}
   \left(\wen y \left( P \tensor Q \right)\right)\sigma \cpar U^i_j \cpar
   W^i_h
   \colon 1 \leq h \leq k_i
   %\end{array}
  }
  \colon 1 \leq j \leq m_i
  %\end{array}
 }
 \colon 1 \leq i \leq n
 %\end{array}
}
\\
\longrightarrow&
\tcontext{
 %\begin{array}{r}
 \tcontextmn{i}{0}{
  %\begin{array}{r}
  \tcontextmn{i}{1}{
   %\begin{array}{l}
\new y \left(\left( P \tensor Q \right)\sub{x}{y}\sigma \cpar \hat{U}^i_j \cpar
   \hat{W}^i_h \right)
   \colon 1 \leq h \leq k_i
   %\end{array}
  }
  \colon 1 \leq j \leq m_i
  %\end{array}
 }
 \colon 1 \leq i \leq n
 %\end{array}
}
\\
\longrightarrow&
\tcontext{
 %\begin{array}{r}
 \tcontextmn{i}{0}{
  %\begin{array}{r}
  \tcontextmn{i}{1}{
   %\begin{array}{l}
\new y \left(
   \left( P\sub{x}{y}\sigma  \cpar \hat{U}^i_j \right) \tensor
   \left(Q\sub{x}{y}\sigma \cpar \hat{W}^i_h \right)\right)
   \colon 1 \leq h \leq k_i
   %\end{array}
  }
  \colon 1 \leq j \leq m_i
  %\end{array}
 }
 \colon 1 \leq i \leq n
 %\end{array}
}
\\
\longrightarrow&
\tcontext{
 %\begin{array}{r}
 \tcontextmn{i}{0}{
  %\begin{array}{r}
  \tcontextmn{i}{1}{
   %\begin{array}{l}
\new y \cunit
   \colon 1 \leq h \leq k_i
   %\end{array}
  }
  \colon 1 \leq j \leq m_i
  %\end{array}
 }
 \colon 1 \leq i \leq n
 %\end{array}
}
\\
\longrightarrow&
\tcontext{
 \tcontextmn{i}{0}{
  \tcontextmn{i}{1}{
   \cunit
  \colon 1 \leq h \leq k_i }
 \colon 1 \leq j \leq m_i }
\colon 1 \leq i \leq n}
\longrightarrow \cunit
\end{array}
\end{array}
\]
Hence $\vdash \left(\wen x \left( P \tensor Q \right)\right)\sigma \cpar R$.
Therefore, by Lemma~\ref{lemma:context}, for all contexts $\context{}$,
if $\vdash \context{ \wen x P \tensor \new x Q }$ then
$\vdash \context{ \new x \left( P \tensor Q \right) }$.
\end{proof}
\end{comment}
% --- end coclose.tex ---

\begin{lemma}[co-tidy name]
\label{lemma:co-tidynew}
If $\vdash \context{ \wen x. \cunit }$ holds
then $\vdash \context{ \cunit }$ holds.
\end{lemma}
\begin{proof}
Assume that $\vdash \wen x. \cunit \cpar P$ holds.
By Lemma~\ref{lemma:split-times}, there exists $Q$ such that $\vdash Q$ and $\vcenter{\infer{P}{Q}}$.
Hence the following proof of $P$ can be constructed:
$
\vcenter{
\infer{
P
}{
\infer[]{
Q
}{
\cunit
}}}
$.
Therefore, by Lemma~\ref{lemma:context}, for any context $\context{}$, if $\vdash \context{\wen{x}. \cunit}$ then $\vdash \context{ \cunit }$, as required.
\end{proof}

\begin{lemma}[co-extrude1]
\label{lemma:co-extrude}
If  $\nfv{x}{Q}$ and $\vdash \context{ \exists x. P \tensor Q }$ holds then $\vdash \context{ \mathopen{\exists x.} \left( P \tensor Q \right) }$ holds.
\end{lemma}
% --- begin coextrude.tex ---
\begin{proof}
Assume that $\vdash \left( \exists x. P \tensor Q \right)\mathclose{\sigma} \cpar V$ holds, where $\nfv{x}{Q}$. Now, since
$\nfv{y}{(\exists x. P \tensor Q)}$ and $\nfv{y}{\sigma}$, we have
$\left( \exists x. P \tensor Q \right)\mathclose{\sigma} \cpar V
\equiv
\left( \mathopen{\exists y.} \left(P\sub{x}{y} \sigma\right) \tensor Q\sigma \right)  \cpar V$.
So, by Lemma~\ref{lemma:split-times}, there exist $T_i$ and $U_i$ such that $\vdash \mathopen{\exists y.} \left(P\sub{x}{y} \sigma\right) \cpar T_i$ and $\vdash Q\sigma \cpar U_i$, for $1 \leq i \leq n$, and $n$-ary killing context such that the derivation
\[
\infer[]{
V
}{ \tcontext{T_1 \cpar U_1, \hdots, T_n \cpar U_n}
}
\] holds.
%
By Lemma~\ref{lemma:split-exists}, there exist $R^i_j$ and $v^i_j$ such that $\vdash P\sub{x}{y} \sigma \sub{y}{v^i_j} \cpar R^i_j$, for $1 \leq j \leq m_i$, and $m_i$-ary killing context $\tcontextn{i}{}$ such that the derivation
\[
\infer[]{
T_i
}{
 \tcontextn{i}{ R^i_1, R^i_2, \hdots, R^i_{m_i} }
}
\] holds.
%
Hence the following proof can be constructed, where we appeal to $\alpha$-conversion in the conclusion.
\[
\infer[]{
%\left( \mathopen{\exists x.} \left( P \tensor Q \right) \right)\mathclose{\sigma} \cpar V \equiv
\mathopen{\exists y.} \left( P\sub{x}{y}\sigma \tensor Q\sigma \right) \cpar V
}{
\infer[]{
\mathopen{\exists y.} \left( P\sub{x}{y}\sigma \tensor Q\sigma \right)
\cpar \tcontext{T_i \cpar U_i \colon 1 \leq i \leq n }
}{
\infer[]{
\mathopen{\exists y.} \left( P\sub{x}{y}\sigma \tensor Q\sigma \right)
\cpar \tcontext{
%            \begin{array}{r}
            \tcontextn{i}{ R^i_j \colon 1 \leq j \leq m_i } \cpar U_i
            %\\
            \colon 1 \leq i \leq n 
%           \end{array}
           }
}{
\infer[]{
\tcontext{
% \begin{array}{r}
 \mathopen{\exists y.} \left( P\sub{x}{y}\sigma \tensor Q\sigma \right)
 \cpar \tcontextn{i}{ R^i_j \colon 1 \leq j \leq m_i } \cpar U_i
% \\
 \colon 1 \leq i \leq n 
% \end{array}
 }
}{
\infer[]{
\tcontext{
% \begin{array}{r}
 \tcontextn{i}{ \mathopen{\exists y.} \left( P\sub{x}{y}\sigma \tensor Q\sigma \right)
 \cpar R^i_j \cpar U_i \colon 1 \leq j \leq m_i }
% \\
 \colon 1 \leq i \leq n 
% \end{array}
 }
}{
\infer[]{
\tcontext{
% \begin{array}{r}
 \tcontextn{i}{
  \left( P\sub{x}{y}\sigma\sub{y}{v^i_j} \tensor Q\sigma \right)
  \cpar R^i_j \cpar U_i \colon 1 \leq j \leq m_i }
% \\
 \colon 1 \leq i \leq n
% \end{array}
 }
}{
\infer[]{
\tcontext{
% \begin{array}{r}
 \tcontextn{i}{
%  \begin{array}{l}
  \left( P\sub{x}{y}\sigma\sub{y}{v^i_j} \cpar R^i_j \right) 
  \tensor
  \left( Q\sigma \cpar U_i \right)
%  \end{array}
  \colon 1 \leq j \leq m_i }
% \\
 \colon 1 \leq i \leq n
% \end{array}
 }
}{
\infer[]{
\tcontext{
 \tcontextn{i}{
  \cunit
  \colon 1 \leq j \leq m_i }
 \colon 1 \leq i \leq n } 
}{
 \cunit
}}}}}}}}
\]
\noindent Hence, by Lemma~\ref{lemma:context}, if $\vdash \context{ \exists x. P \tensor Q }$, where $\nfv{x}{Q}$, then $\vdash \context{ \mathopen{\exists x.} \left( P \tensor Q \right) }$.
\end{proof}
% --- end coextrude.tex ---

\begin{lemma}[co-tidy1]
\label{lemma:co-tidy1}
If $\vdash \context{ \exists x. \cunit }$ holds
then $\vdash \context{ \cunit }$ holds.
\end{lemma}
\begin{proof}
Assume that $\vdash \exists x. \cunit \cpar T$ holds.
By Lemma~\ref{lemma:split-exists}, there exists $U_i$ such that $\vdash U_i$, for $1 \leq i \leq n$, and $n$-ary killing context $\tcontext{}$ such that 
$\vcenter{\infer{T}{\tcontext{U_1, \hdots, U_n}}}$.
Hence the following proof of $T$ can be constructed:
\[
\infer[.]{
\cunit \cpar T 
}{
\infer[]{
\tcontext{U_1, \hdots, U_n}
}{
\infer[]{
\tcontext{\cunit, \hdots, \cunit}
}{
\cunit
}}}
\]
Therefore, by Lemma~\ref{lemma:context}, if $\vdash \context{\exists{x} \cunit}$ then $\vdash \context{ \cunit }$, as required.
\end{proof}

The above four lemmas are particular to \textsf{MAV1}.
The following lemma is proven directly for \textsf{MAV}, similarly to Lemma~\ref{lemma:universal}; however, for \textsf{MAV1} the proof is more indirect due to interdependencies between $\wwith$ and nominals.
\begin{lemma}[co-left and co-right]\label{lemma:co-branching}
If $\vdash \context{ P \wwith Q }$ holds
then both $\vdash \context{ P }$ and $\vdash \context{ Q }$ hold.
\end{lemma}

The proofs for the four co-rule elimination lemmas below are similar to the corresponding cases in \textsf{MAV}~\cite{Horne2015}.
%
\begin{lemma}[co-external]
\label{lemma:co-external}
If $\vdash \context{ P \tensor  \left( Q \ooplus R \right) }$ holds
 then $\vdash \context{ \left( P \tensor Q \right) \ooplus \left( P \tensor R \right) }$ holds.
\end{lemma}

\begin{lemma}[co-sequence]
\label{lemma:co-sequence}
If $\vdash \context{ \left( P \cseq Q \right) \tensor \left( R \andthen S \right) }$ holds
 then $\vdash \context{ \left( P \tensor R \right) \cseq \left( Q \tensor S \right) }$ holds.
\end{lemma}

%
\begin{lemma}[co-tidy]
\label{lemma:co-tidy}
If $\vdash \context{ \cunit \ooplus \cunit }$ holds, then $\vdash \context{ \cunit }$ holds.
\end{lemma}

\begin{lemma}[atomic co-interaction]
\label{lemma:co-atoms}
If $\vdash \context{ \alpha \tensor \co{\alpha} }$ holds
then $\vdash \context{ \cunit }$ holds.
\end{lemma}


\subsection{The proof of cut elimination}
The main result of this paper, Theorem~\ref{theorem:cut}, follows by induction on the structure of $P$ in a formula of the form $\vdash \context{ P \tensor \co{P} }$, by applying the above eight co-rule elimination lemmas and also Lemma~\ref{lemma:universal} in the cases for \textit{all} and \textit{some}.

\tohide{
\begin{proof}
The base cases for any atom $\alpha$ follows since if $\vdash \context{\co{\alpha} \tensor \alpha}$ then $\vdash \context{\cunit}$ by Lemma~\ref{lemma:co-atoms}. The base case for the unit is immediate. %The base case for $\top$ and $0$ follow from Lemma~\ref{lemma:co-top} since if $\vdash \context{ \top \cpar 0 }$ then $\vdash \context{ \zero }$ and therefore, by Lemma~\ref{lemma:co-unit}, $\vdash \context{ \cunit }$ holds.
As the induction hypothesis in the following cases assume for any context $\context{}$, $\vdash \context{P \tensor \co{P}}$ yields $\context{ \cunit }$ and $\vdash \contextb{Q \tensor \co{Q}}$ yields $\contextb{ \cunit }$.

Consider the case for \textit{times}. 
Assume that $\vdash \context{ P \tensor Q \tensor \left(\co{P} \cpar \co{Q}\right) }$ holds.
By the \textit{switch} rule, $\vdash \context{ \left( P \tensor \co{P} \right) \cpar \left( Q \tensor \co{Q} \right) }$ holds.
Hence, by the induction hypothesis twice, $\vdash \context{ \cunit }$ holds.
%
The case for \textit{par} is symmetric to the case for \textit{times}.

Consider the case for \textit{seq}. 
Assuming that $\vdash \context{ \left( P \cseq Q \right) \tensor \left(\co{P} \cseq \co{Q}\right) }$ holds,
by Lemma~\ref{lemma:co-sequence}, it holds that $\vdash \context{ \left( P \tensor \co{P} \right) \cseq \left( Q \tensor \co{Q} \right) }$.
Hence, by the induction hypothesis twice, $\vdash \context{ \cunit }$ holds.


Consider the case for \textit{with}.
Assume that $\vdash \context{ \left(P \wwith Q\right) \tensor \left(\co{P} \ooplus \co{Q}\right) }$ holds.
By Lemma~\ref{lemma:co-external},
$\vdash \context{ \left(\left(P \wwith Q\right) \tensor \co{P}\right) \ooplus \left(\left(P \wwith Q\right) \tensor \co{Q}\right) }$
holds.
By Lemma~\ref{lemma:co-branching} twice, $\vdash \context{ \left(P \tensor \co{P}\right) \ooplus \left(Q \tensor \co{Q}\right) }$ holds.
Hence by the induction hypothesis twice, $\vdash \context{ \cunit \ooplus \cunit }$ holds.
Hence by Lemma~\ref{lemma:co-tidy}, $\vdash \context{ \cunit }$ holds, as required.
The case for \textit{plus} is symmetric to the case for \textit{with}.

Consider the case for universal quantification.
Assume that $\vdash \context{ \forall x. P \tensor \exists x. \co{P} }$ holds.
By Lemma~\ref{lemma:co-extrude}, it holds that $\vdash \context{ \mathopen{\exists x.} \left( \forall x. P \tensor \co{P} \right) }$, since $\nfv{x}{\exists x. P}$.
By Lemma~\ref{lemma:universal}, $\vdash \context{ \mathopen{\exists x.} \left( P \tensor \co{P} \right) }$ holds.
Hence by the induction hypothesis, $\vdash \context{ \exists x. \cunit }$ holds.
Hence by Lemma~\ref{lemma:co-tidy1}, $\vdash \context{ \cunit }$ holds, as required.
The case for existential quantification is symmetric to the case for universal quantification.

Consider the case for \textit{new}.
Assume that $\vdash \context{ \new x. P \tensor \wen x. \co{P} }$ holds.
By Lemma~\ref{lemma:co-close}, it holds that $\vdash \context{ \mathopen{\wen x.} \left( P \tensor \co{P} \right) }$.
Hence by the induction hypothesis, $\vdash \context{ \wen x. \cunit }$ holds.
Hence by Lemma~\ref{lemma:co-tidynew}, $\vdash \context{ \cunit }$ holds, as required.
The case for \textit{wen} is symmetric to the case for \textit{new}.

%This completes the case analysis.
Therefore, by induction on the structure of $P$, if $\vdash \context{ P \tensor \co{P} }$ holds, then $\vdash \context{ \cunit }$ holds.
\end{proof}
}{}

Notice that the structure of the above argument is similar to the structure of the argument for Proposition~\ref{proposition:reflexivity}.
The only difference is that the formulae are dualised and co-rule lemmas are applied instead of rules.
% --- end corule.tex ---
% --- begin symmetric.tex ---
\subsection{Discussion on alternative presentations of rules for \textsf{MAV1}}

Having established cut elimination (Theorem~\ref{theorem:cut}), an immediate corollary is that all co-rules in Fig.~\ref{figure:co-rules} are admissible. This can be formulated by demonstrating that linear implication coincides with the inverse of a derivation in the symmetric system \textit{SMAV1}.
\begin{corollary}
$\vdash P \multimap Q$ in \textsf{MAV1} if and only if $\infer{Q}{P}$ in \textsf{SMAV1}.
\end{corollary}
\begin{proof}
Firstly, assume $\vdash P \multimap Q$ in \textsf{MAV1},
in which case the following can be constructed in \textsf{SMAV1}: 
\[
\infer{
 Q
}{
\infer[.]{
 \left(P \tensor \co{P}\right) \cpar Q
}{
\infer[]{
 P \tensor \left(\co{P} \cpar Q\right)
}{
 P
}}}
\]
For the converse, assume $\vcenter{\infer{Q}{P}}$ in \textsf{SMAV1}; hence 
\[
\infer{
 \co{P} \cpar Q
}{\infer[]{
 \co{P} \cpar P
}{
 \cunit
}}
\] 
can be constructed.
%Also, observe that the remaining co-rules are provable as linear implications in $\textsf{MAV1}$.
%Hence, by using Theorem~\ref{theorem:cut}, an instance of any co-rule can be eliminated.
Thereby by Lemma~\ref{lemma:universal} and Lemmas~\ref{lemma:co-close}~to~\ref{lemma:co-sequence},
the above derivation in \textsf{SMAV1} can be transformed into a proof in \textsf{MAV1}.
\end{proof}
The advantage of the definition of linear implication using provability in \textsf{MAV} rather than derivations in \textsf{SMAV1}, is that \textsf{MAV1} is \textit{analytic}~\cite{Bruscoli2009}; hence, with some care taken for existential quantifiers~\cite{Lincoln1994,Brunnler2003}, each formula gives rise to finitely many derivations up-to congruence. In contrast, in $\textsf{SMAV1}$, many co-rules can be applied indefinitely.
Notice co-rules including \textit{atomic co-interaction}, \textit{co-left} and \textit{co-tidy} can infinitely increase the size of a formula during proof search.
%\end{comment}

\textbf{A small rule set.}
Alternatively, we could extend the structural congruence with the following. %for both nominal quantifiers.
\[
\wen x. P \equiv P 
~~\mbox{only if $\nfv{x}{P}$}
\qquad
\new x. P \equiv P 
~~\mbox{only if $\nfv{x}{P}$}
\qquad\mbox{(vacuous)}
\]
Vacuous allows nominals to be defined by the smaller set of rules \textit{close}, \textit{medial new}, \textit{suspend}, \textit{new wen}, \textit{with name}, and \textit{all wen}.
\begin{comment}
 following rules, where $\quantifier \in \left\{\new, \wen\right\}$ and $\odot \in \left\{ \cpar, \cseq \right\}$.
\begin{gather*}
\context{ \new x. P \cpar \wen x. Q }
\longrightarrow
\context{ \mathopen{\new x .}\left( P \cpar Q \right) }
~~\mbox{(close)}
\\[9pt]
\context{ \mathopen{\new x .}\left( P \cseq S \right) }
\longrightarrow
\context{ \new x. P \cseq \new x. S }
~~\mbox{(medial new)}
\\[9pt]
\context{ \wen x. P \odot \wen x. S }
\longrightarrow
\context{ \mathopen{\wen x .}\left( P \odot S \right) }
~~\mbox{(suspend)}
\\[9pt]
\context{ \quantifier x .P \wwith \quantifier x. S }
\longrightarrow
\context{ \mathopen{\quantifier x .}\left( P \wwith S \right) }
~~\mbox{(with name)}
\\[9pt]
\context{ \new x .\wen y .P }
\longrightarrow
\context{ \wen y. \new x .P }
~~\mbox{(new wen)}
\qquad
\context{ \forall x .\quantifier y. P }
\longrightarrow
\context{ \quantifier y. \forall x .P }
~~\mbox{(all name)}
\end{gather*}
\end{comment}
Any formula provable in this smaller system is also provable in \textsf{MAV1}, since all rules of \textsf{MAV1} can be simulated by the rules above.
Perhaps the least obvious case is the \textit{fresh} rule,
where since 
$
\vcenter{
\infer[]{\new x. \wen x. P}{\wen x. \new x. P}}
$, 
by the \textit{new wen} rule
and both 
$\wen x. \new x. P \equiv \new x. P$
and
$\wen x. P \equiv \new x. \wen x. P$
hold
using the \textit{vacuous} rule,
we have 
$
\vcenter{
\infer[]{\wen x. P}{\new x. P}}
$.

Conversely, \textit{vacuous} is a provable equivalence in \textsf{MAV1}; hence, by inductively applying cut elimination to eliminate each \textit{vacuous} rule in a proof using the smaller set of rules, we can obtain a proof with the same conclusion in \textsf{MAV1}.
The disadvantage of the above system is that the \textit{vacuous} rules can introduce an arbitrary number of nominal quantifiers at any stage in the proof leading to infinite paths in proof search,
i.e., the above system is not \textit{analytic}. 
Indeed the multiset-based measure used to guide splitting would not be respected, hence our cut elimination strategy would fail.
None the less, the smaller rule set above offers insight into design decisions.

\textbf{Alternative approaches to cut elimination.}
Further styles of proof system are possible. For example, again as a consequence of cut elimination, we can show the equivalence of \textsf{MAV1} and a system which reduces the implicit contraction in the \textit{external} rule to an atomic form $\vcenter{\infer{\alpha}{\alpha \ooplus \alpha}}$, in which additional medial rules play a central role for propagating contraction~\cite{Brunnler2001,Strassburger2002,Bruscoli2016}.
% (in a sense the medial rules are the source of proof size explosion in normalisation~\cite{Tubella}).
Similarly, the implicit vacuous existential quantifier introduction can be given an explicit atomic treatment~\cite{Strassburger09}.
The point is that, although the cut elimination result in this work is sufficient to establish the equivalent expressive power of systems mentioned in this subsection, further proof theoretic insight may be gained by attempting direct proofs of cut elimination in such alternative systems.
Indeed a different approach to cut elimination is required for tackling \textsf{MAV2} with second-order quantifiers.


\textbf{Note on probabilistic choice.} 
Insight from investigating the proof theory of \textsf{MAV1} led to the surprising observation that
probabilistic choice has similar proof theoretic properties to \textit{new}.
A proof theory of \textsf{MAV} extended with \textit{sub-additive} operators is explored in related work~\cite{Horne2019}.
The sub-additives, similarly to nominal quantifiers which lie between universal and existential quantifiers, lie between the traditional additives \textit{with} and \textit{plus}.
Sub-additives can either be self-dual, similarly to $\nabla$, or de Morgan dual, similarly to $\new$ and $\wen$ --- controlling distributivity properties which are undesirable when embedding probabilistic processes, much like the quantifiers in this work avoid undesirable distributivity properties when embedding processes with private names.

We remark that adapting recent work on splitting in \textit{subatomic logic}~\cite{Tubella2018} may help explain general patterns emerging, connecting the nominal quantifiers and sub-additives.
Subatomic logic may also be used to provide a more concise proof of splitting by exploiting the evident general patterns in the case analysis. Beside abstractly explaining general patterns, the study of \textsf{MAV1} in terms of subatomic logic would likely expose alternative formulations of the rules of \textsf{MAV1}.

% --- end symmetric.tex ---
% --- begin decidable.tex ---
\section{Decidability of Proof Search}
\label{section:complexity}

Here we identify complexity classes for proof search in fragments of \textsf{MAV1}. % and justify this claim.
The hardness results in this section are consequences of cut elimination (Theorem~\ref{theorem:cut}) and established complexity results for fragments of linear logic and extensions of \textsf{BV}.

NEXPTIME-hardness follows from the NEXPTIME-hardness of \textsf{MALL1}~\cite{Lincoln1994}; while membership in NEXPTIME follows a similar argument as for \textsf{MALL1}~\cite{Lincoln1994b} (in a proof there are at most exponentially many \textit{atomic interaction} rules, each involving quadratically bounded terms).
\begin{proposition}\label{proposition:NEXPTIME}
Deciding provability in \textsf{MAV1} is NEXPTIME-complete.
\end{proposition}
%
If we restrict terms to a nominal type, i.e.\ \textit{some} can only be instantiated with variables and constants, we obtain a tighter complexity bound.
PSPACE-hardness is a consequence of the PSPACE-hardness of \textsf{MAV}~\cite{Horne2015}, which in turn follows from the PSPACE-hardness of MALL~\cite{Lincoln1992}. Membership in PSPACE follows a similar argument as for \textsf{MALL1} without function symbols~\cite{Lincoln1994}.
\begin{proposition}\label{proposition:PSPACE}
Deciding provability in \textsf{MAV1} without function symbols is PSPACE-complete.
\end{proposition}
%
If we consider the sub-system without \textit{with} and \textit{plus}, named \textsf{BV1}, we obtain a tighter complexity bound again, even with function symbols in terms.
NP-hardness is a consequence of the NP-hardness of \textsf{BV}~\cite{Ozan2008a}; while membership in NP follows a similar argument as for \textsf{MLL1}~\cite{Lincoln1994b}
\begin{proposition}\label{proposition:NP}
Deciding provability in \textsf{BV1} is NP-complete.
\end{proposition}
For problems in the complexity class NEXPTIME, we can always check a proof in exponential time.
The high worst-case complexity means that proof search in general is considered to be infeasible.
Implementations of NEXPTIME-complete problems that regularly work efficiently, include reasoning in description logic $\mathcal{ALCI(D)}$~\cite{Lutz2004}.


\begin{comment}
The proof is a consequence of cut elimination, Theorem~\ref{theorem:cut}, and the established result that proof search in first-order multiplicative linear logic, \textsf{MALL1}~\cite{Lincoln1994}, is NEXPTIME-complete. The proof of NEXPTIME-hardness in \textsf{MALL1} relies on an operational model of Turing machines where an exponential number of steps can be simulated using a polynomial size predicate. 
%
This NEXPTIME-hardness proof tells us that the worst case problems are intricate encodings that are not required for many problems.
%and furthermore that those encodings employ arbitrary nesting of function symbols to form values.
Furthermore, if only constants and variables are permitted as terms (as with channel-only process calculi or other calculi restricted to a nominal type), proof search is reduced to a PSPACE-complete problem. This matches the complexity results for \textsf{MAV}, where deciding provability is also a PSPACE-complete problem~\cite{Ciobanu2015,Lincoln1994}.
\begin{proposition}\label{proposition:nominal}
Deciding provability in \textsf{MAV1} without function symbols, i.e.\ only variables and constants as terms (\textit{nominal} \textsf{MAV1}), is PSPACE-complete.
\end{proposition}
\end{comment}

%\begin{comment}
\begin{figure*}[t]
\begin{tabular}{l|l|l}
Complexity class &
Linear logic &
Calculus of structures
\\ \hline
NP-complete &
\textsf{MLL1}~with functions~\cite{Kanovich1994} &
\txt{
\textsf{BV1}~with functions
\\
(Proposition~\ref{proposition:NP})
}
\\ \hline
PSPACE-complete &
\textsf{MALL1}~without functions~\cite{Lincoln1992} &
\txt{
\textsf{MAV1}~without functions
\\
(Proposition~\ref{proposition:PSPACE})
}
% \\ & &
%nominal \textsf{MAV1}~(Proposition~\ref{proposition:nominal})
\\ \hline
NEXPTIME-complete &
\textsf{MALL1}~with functions~\cite{Lincoln1994,Lincoln1994b} &
\txt{
\textsf{MAV1}~with functions
\\
(Proposition~\ref{proposition:NEXPTIME})
}
\\ \hline
Undecidable &
%\textsf{MLL2}~\cite{Lafont1996,Lincoln1995} &
%\textsf{BV2}~(by embedding MLL2) \\
%&
\textsf{MAELL}~\cite{Lincoln1992}
and 
\textsf{MLL2}~\cite{Lincoln1995}
&
\textsf{NEL}~\cite{Strassburger2003}
\end{tabular}
\caption{Complexity results.}\label{figure:complexity}
\end{figure*}
%\end{comment}

%Note that each inclusion result above assumes that deciding the preorder over atoms is also in the respective complexity class.

%\begin{comment}
Figure~\ref{figure:complexity} summarises complexity results for related calculi.
Notice the pattern that each fragment of linear logic has the same complexity as the calculus that is a conservative extension of that fragment of linear logic (with mix), where the extra operator is the self-dual non-commutative operator \textit{seq}. The complexity classes match since the source of the NP-completeness in multiplicative-only linear logic (\textsf{MLL}) lies in the number of ways of partitioning resources (formulae), while the mix rule and sequence rule are also ways of partitioning the same resources.
%The same technique can be applied to obtain the undecidability of \textsf{BV} extended with second-order quantifiers (\textsf{BV2}), since second-order multiplicative linear logic (\textsf{MLL2}) is undecidable.

An exceptional case is that \textsf{BV} extended with exponentials (\textsf{NEL}) is undecidable, whereas the decidability of multiplicative linear logic with exponentials (\textsf{MELL}) is unknown.\footnote{\textsf{MELL} was claimed to be decidable in~\cite{Bimbo2015}, but this was later refuted~\cite{Strassburger19}. } 
However, by including additives to obtain full propositional linear logic (\textsf{MAELL} or simply \textsf{LL}) provability is known to be undecidable.
%\end{comment}

By the above observations, the complexity of deciding linear implication for embeddings of finite name passing processes, as in $\pi$-calculus, is in PSPACE.
However, extending to finite value passing processes where terms constructed using function symbols can be communicated, e.g.\ capturing tuples in the polyadic $\pi$-calculus~\cite{Milner1993}, the complexity class increases, but only for processes involving choice.
Further extensions to \textsf{MAV1}
introducing second-order quantifiers, exponentials or fixed points would lead to undecidable proof search~\cite{Lafont1996,Lincoln1995,Strassburger2003}.
%(without for example careful restricting existential quantifiers~\cite{Perrier1997}).

% --- end decidable.tex ---
% --- begin conclusion.tex ---
\section{Conclusion}

%From a proof theoretic perspective 
This paper makes two significant contributions to proof theory: the first cut elimination result for a novel de Morgan dual pair of nominal quantifiers;
and the first direct cut elimination result for first-order quantifiers in the calculus of structures. 
As a consequence of cut-elimination (Theorem~\ref{theorem:cut}), we obtain the first proof system that features both
non-commutative operator \textit{seq} and first-order quantifiers $\forall$ and $\exists$.
%
%The novel 
A novelty of the nominal quantifiers $\new$ and $\wen$
%\textit{new} and \textit{wen}
 compared to established self-dual nominal quantifiers
%are dual to each other and 
is in how they distribute over positive and negative operators. 
%The greater control of the bookkeeping of names compared to established self-dual nominal quantifiers.
This greater control of bookkeeping of names enables private names to be modelled in direct embeddings of 
processes as formulae in \textsf{MAV1}.
%Logical implication respects causality for embedded processes.
In Section~\ref{section:syntax}, every rule in \textsf{MAV1} is justified as necessary either: for soundly embedding processes;
or for ensuring cut elimination holds. Of particular note, some rules were introduced for ensuring cut elimination holds in the presence of \textit{equivariance}.

The cut elimination result is an essential prerequisite for recommending the system \textsf{MAV1} as a logical system.
This paper only hints about formal connections between \textsf{MAV1} and models of processes, which receives separate attention in a companion paper~\cite{MSCS}.
In particular, we know that linear implication defines a precongruence over processes embedded as formulae, that is sound with respect to both weak simulation and pomset traces.

Further to connections with process calculi, there are several problems exposed as future work.
%exposed by this work include the following.
%The pair of nominals, of 
%$\wen x.P \not\multimap \new x.P$ in general with respect to linear implication, so the operators do not collapse in extensions of linear logic; but it is not clear whether the nominals can be meaningfully distinguished in the setting of classical logic.
Regarding the sequent calculus, in the setting of linear logic (i.e.,\ without \textsf{seq}), it is an open problem to determine whether there is a sequent calculus presentation of \textit{new} and \textit{wen}.
Regarding model theory, a model theory or game semantics may help to explain the nature of the de Morgan dual pair of nominal quantifiers, 
%(perhaps involving uniform choices from a seemingly infinite set of names),
although note that it remains an open problem just to establish a sound and complete denotational model of \textsf{BV}.
Another open question is whether quantifiers \textit{new} and \textit{wen} are relevant in a classical or intuitionistic setting, or whether these operators are uniquely interesting in a linear setting.
Since \textit{new} must distribute over classical disjunction (recall, in contrast, \textit{new} does not distribute over multiplicative disjunction), nominal operators \textit{new} and \textit{wen} likely collapse to an established self-dual nominal operator in the classical setting; hence \textit{wen} is probably unrelated to the `generous' operator proposed in related work on stratifiable languages~\cite{Gabbay}.
% where perhaps nominal quantifiers are connected to names chosen uniformly from an set exponential in size with respect to a security parameter (similarly to the security intuition for private name binders in process calculi).
%Regarding consistency, although there are propositions such as $\left(\alpha \tensor \co{\alpha}\right) \ooplus \cunit$ that are both provable and their complement is provable, there may be a notion of consistency satisfied by a ``guarded'' fragment of \textsf{MAV1} where $\cunit$ is forbidden in the conclusion of a proof.
Regarding implementation, it is a challenge to reduce non-determinism in proof search~\cite{Chaudhuri2011,Ozan2014,Andreoli1992};
a problem that can perhaps be tackled by restricting to well-behaved fragments of \textsf{MAV1} or by exploiting complexity results to embed rules as constraints for a suitable solver. 
Regarding proof normalisation, systems including classical propositional logic~\cite{Tubella2017}, first-order logic~\cite{Tubella2017}, intuitionistic logic~\cite{Guenot2014} and \textsf{NEL} (\textsf{BV} with exponentials)~\cite{Strassburger2011} 
satisfy a proof normalisation property called \textit{decomposition} related to interpolation;
leading to the question of whether there is an alternative presentation of the rules of \textsf{MAV1}, for which a decomposition result can be established. 
%Decomposition can be used as part of an alternative proof for establishing cut elimination, but is an independent and rather strong property to investigate.
Finally, an expressivity problem, perhaps related to decomposition, is how to establish cut elimination for second-order extensions suitable for modelling infinite processes.



%% We provide the first direct
%% proof of cut elimination in the calculus of structures; and thereby
%% the first consistent proof system that features both the self-dual
%% non-commutative operator \textit{seq} and first-order quantifiers.

\begin{comment}
This paper continues a line of work on logical systems defined using
the calculus of structures with applications to modelling processes~\cite{Bruscoli2002,Ciobanu2015,Horne2015}.  Our approach is
distinct from related work on nominal logic~\cite{Alexandru2013,Gabbay2003,Tiu2010}
where processes are terms, rather than predicates, and an operational semantics is given either as an inductive definition or a logical theory.
The related approach is capable of encoding observational preorders and bisimulations, but has the drawback that implication cannot be directly used to compare processes.
%In contrast, in the approach presented in this paper, implication itself directly defines a preorder over processes.
%That apporoch enables a variety of operational equivalences to be specified, but in contrast to the this paper logical implication itself cannot characterise 
%A drawback of that distinct approach is that linear implication cannot be directly used to relate processes.
%% This approach is sometimes referred to as a {\em deep embedding}
%% approach in formal methods literature. 
%In that distinct approach, only the diagonalization property of the nominal quantifier is needed to prove the adequacy of the embeddings (see e.g., \cite{Tiu2010}).
%with respect to the operational semantics of the $\pi$-calculus 
%% For example, in work such as
%% \cite{Bengtson2007,Tiu2010}, process equivalence needs to be explicitly defined as 
%% greatest fixed points in a suitable expressive logic.
%% where
%% private names are modelled indirectly by exploiting the
%% `diagonalization' property of a self-dual nominal quantifier in
%% encodings of rules. 
Our approach is also distinct from the
proofs-as-processes Curry-Howard inspired approach to session
types~\cite{Caires2010,Wadler2014}.
%\ross{How do we phrase this to dispell the "but these guys don't even know the seminal work of Pfenning and Wadler'' reviews?}.
Instead, we adopt a processes-as-predicates approach, setting up a
discussion on the relationship between linear implication in
\textsf{MAV1} and observational preorders. %to be elaborated in future work. 
%Results such as the multiparty compatibility of coherent protocols for subsystem~\textsf{MAV}, can be extended to the more expressive \textsf{MAV1}, providing a decidable subtype system for dependent session types.
%The complexity results (Propositions~\ref{proposition:NEXPTIME},~\ref{proposition:PSPACE} and~\ref{proposition:NP}) will help guide tool support.
%Future work includes undecidable extensions.
\end{comment}

%In particular, \textsf{MAV1} extends the expressive power of \textsf{MAV} to enabling reasonning about values passed in protocols and private names. \ross{Extend this if we decide to include more results and less proofs in the body.}
\begin{comment}
\begin{itemize}
\item Sequential composition is used instead of prefixing, allowing DAGs of causality, instead of trees of causality to be expresses.

\item The semantics of parallel composition permits a non-interleaving interpretation, where events represented by atoms may be independent.

\item The operator $\tensor$ enables a high level approach to synchronisation or separation.

\item Internal and external choices are both captured, enabling the static analysis of multiple branches at compile time.

\item Existential quantification generalises input binders.

\item Universal %and renaming
quantifiers can usually fulfil the role of name binders, universal is not an exact match for ``new name'' in the applied $\pi$ calculus.
\end{itemize}
Besides the above differences in expressive power, predicates can be compared using an internal notion of linear implication. This contrasts to the (bi-)simulation or trace semantics for the applied $\pi$-calculus.
\end{comment}
%, hence labelled transitions and frames are not required.


\paragraph*{Acknowledgements.}
We thank the anonymous reviewers, whose thorough reading led to improvements in the presentation of \textsf{MAV1}.


%\newpage

% --- end conclusion.tex ---




%\newpage
\bibliographystyle{plain} %Lipics standard is plain.
\bibliography{WWV15}

% --- begin appendix.tex ---
\newpage

\appendix

\section{Electronic Appendix}


\begin{proposition}[Reflexivity: Proposition~\ref{proposition:reflexivity}]
For any formula $P$, $\vdash \co{P} \cpar P$ holds, i.e., $\vdash P \multimap P$.
\end{proposition}
\begin{proof}
% --- begin proof_reflexivity.tex ---
The proof proceeds by induction on the structure of a formula $P$.
The base cases for any atom $\alpha$ follows immediately from the \textit{atomic interaction} rule.% Since subsorting over atoms is reflexive, $\vdash \co{\alpha} \cpar \alpha$.
The base case for the unit is immediate by definition of a proof. %The cases for top and zero hold since $0 \cpar \top \longrightarrow \top \longrightarrow \cunit$.
For the following inductive cases assume that $\vdash \co{P} \cpar P$ and $\vdash \co{Q} \cpar Q$ hold.
%Thereby, the following cases hold.

Consider when the root connective in the formula is the $\tensor$ operator. 
Observe, by definition, $\co{\left(P \tensor Q\right)} \cpar \left(P \tensor Q\right) = \co{P} \cpar \co{Q} \cpar \left(P \tensor Q\right)$ and 
by applying the \textit{switch} rule and then the \textit{induction hypothesis}
we have the following proof:
\[
\infer[.]{
\co{P} \cpar \co{Q} \cpar \left(P \tensor Q\right) 
}{
\infer[%\mbox{by the }
]{
\left( \co{P} \cpar P \right) \tensor \left( \co{Q} \cpar Q\right)
}{
\cunit
}}
\]

The case when the root connective is the \textit{par} operator is symmetric to the case for \textit{times}.

Consider when the root connective in the formula is the \textit{seq} operator. We have, by definition,
$ \co{\left(P \cseq Q\right)} \cpar \left(P \cseq Q\right)
    = \left(\co{P} \cseq \co{Q}\right) \cpar \left(P \cseq Q\right)$
and, by applying the \textit{sequence} rule and then the \textit{induction hypothesis}, the following proof holds:
\[
\infer[.]{
 \left(\co{P} \cseq \co{Q}\right) \cpar \left(P \cseq Q\right) 
}{
\infer[%\mbox{by the \textit{induction hypothesis}}
]{
\left( \co{P} \cpar P \right) \cseq \left( \co{Q} \cpar Q\right)
}{
\cunit
}}
\]

Consider when the root connective in the formula is the \textit{with} operator.
By definition we have 
$\co{\left(P \wwith Q\right)} \cpar \left(P \wwith Q\right)
    =
\left(\co{P} \ooplus \co{Q}\right) \cpar \left(P \wwith Q\right)$
and the following proof holds.
\[
\infer[\mbox{by the \textit{external} rule}]{
\left(\co{P} \ooplus \co{Q}\right) \cpar \left(P \wwith Q\right) 
}{
\infer[\mbox{by the \textit{left} and \textit{right} rules}]{
\left(\left(\co{P} \ooplus \co{Q}\right) \cpar P\right)
\wwith
\left(\left(\co{P} \ooplus \co{Q}\right) \cpar Q\right) 
}{
\infer[\mbox{by the \textit{induction hypothesis}}]{
\left(\co{P} \cpar P\right) \wwith \left(\co{Q} \cpar Q\right) 
}{
\infer[\mbox{by \textit{tidy}}]{
\cunit \wwith \cunit
}{
\cunit
}}}
}
\]
The case for when \textit{plus}, $\ooplus$, is the root connective is symmetric to the case for \textit{with}.


Consider when the root connective in the formula is $\forall$.
By definition,
$ \co{\forall x. P} \cpar \forall x. P
    = \exists x. \co{P} \cpar \forall x. P $
and the following proof holds:
\[
\infer[\mbox{by the \textit{extrude1} rule}]{
\exists x. \co{P} \cpar \forall x. P 
}{
\infer[\mbox{by the \textit{select1} rule}]{
\mathopen{\forall x.} \left(\exists x. \co{P} \cpar P\right) 
}{
\infer[\mbox{by the \textit{induction hypothesis}}]{
\mathopen{\forall x.} \left(\co{P} \cpar P\right) 
}{
\infer[\mbox{by the \textit{tidy1} rule}]{
\forall x. \cunit
}{
 \cunit
}}}}
\]

 The case for when $\exists$ is the root connective is symmetric to the case for $\forall$.

Consider when the root connective in the formula is $\new$.
By definition $ \co{\new{x}. P} \cpar \new{x}. P
    = \wen{x}. \co{P} \cpar \new{x}. P $
and the following proof holds:
\[
\infer[\mbox{by the \textit{close} rule}]{
 \wen{x}. \co{P} \cpar \new{x}. P 
}{
\infer[\mbox{by the \textit{induction hypothesis}}]{
 \mathopen{\new{x}.} \left( \co{P} \cpar P \right) 
}{
\infer[\mbox{by the \textit{tidy name} rule}]{
 \mathopen{\new x.} \cunit
}{
  \cunit
}}}
\]

The case for when the root connective is $\wen$ is symmetric to the case for $\new$.

Hence, by induction on the number of connectives in the formula, reflexivity holds.
% --- end proof_reflexivity.tex ---
\end{proof}

\begin{comment}
\begin{corollary}[Corollary~\ref{corollary:precongruence}]
Linear implication defines a precongruence.
\end{corollary}
\begin{proof}
For transitivity, if $\vdash P \multimap Q$ and $\vdash Q \multimap R$ hold, we have the following.
%the following proof can be constructed.
\[
\infer[\mbox{by the \textit{switch} rule}]{
 \left(\co{P} \cpar \left( Q \tensor \co{Q} \right) \cpar R \right)
}{
\infer[\mbox{by the assumptions $\vdash \co{P} \cpar Q$ and $\vdash \co{Q} \cpar R$}]{
 \left(\co{P} \cpar Q \right) \tensor \left( \co{Q} \cpar R \right)
}{
\cunit
}}
\]
Hence, by Theorem~\ref{theorem:cut}, $\vdash P \multimap R$ as required.

For contextual closure, if $\vdash P \multimap Q$ holds, we have the following.
\[
\infer[\mbox{by the \textit{switch} rule}]{
\co{\context{P}} \cpar \context{ \left( P \tensor \co{P} \right) \cpar Q}
}{
\infer[\mbox{by the assumption $\vdash P \multimap Q$}]{
\co{\context{P}} \cpar \context{ P \tensor \left( \co{P}  \cpar Q \right) } 
}{
\infer[\mbox{by Proposition~\ref{proposition:reflexivity}}]{
\co{\context{P}} \cpar \context{ P } 
}{
 \cunit
}}}
\]
Hence by Theorem~\ref{theorem:cut}, $\vdash \context{P} \multimap \context{Q}$ as required.
Reflexivity holds by Proposition~\ref{proposition:reflexivity}.
\end{proof}
\end{comment}

%%%

\begin{lemma}[Universal: Lemma~\ref{lemma:universal}] 
If $\vdash \context{ \forall x. P }$ holds
 then, for all terms~$v$, $\vdash \context{ P\sub{x}{v} }$ holds.
\end{lemma}
% --- begin LemmaUniversal.tex ---
\begin{proof}
We require a function over formulae $s_v(T)$ that replaces a certain universal quantifier in $T$ with a substitution for a value $v$. 
The universal quantifiers to be replaced are highlighted in bold $\boldall$. Note that during a proof the bold operator may be duplicated by the \textit{external} rule and \textit{medial1} rule, hence there may be multiple bold occurrences in a formula.
The function is defined as follows, where $\odot \in \left\{ \cseq, \cpar, \tensor, \ooplus, \wwith \right\}$ is any binary connective, $\quantifier \in \left\{ \forall, \exists, \new, \wen \right\}$ is any quantifier except bold universal quantification and $\kappa \in \left\{ \alpha, \co{\alpha}, \cunit \right\}$ is any constant or atom.
\begin{gather*}
%s_v(\boldall x. T) = T \sub{x}{v}
% [AT] the bold quantifier may be nested
\mathopen{s_v}\left(\boldall x. T\right) = \mathopen{s_v}\left(T \sub{x}{v}\right)
\qquad
\mathopen{s_v}\left(\quantifier x. T\right) =
\quantifier x. \mathopen{s_v}\left(T\right)
\qquad
\mathopen{s_v}\left( T \odot U \right) = 
\mathopen{s_v}\left(T\right) \odot \mathopen{s_v}\left(U\right)
\qquad
\mathopen{s_v}\left(\kappa\right) = \kappa
\end{gather*}
In what follows we use that $\mathopen{s_v}\left(\context{ U }\right) = \contextp{ \mathopen{s_v}\left(U'\right) }$, 
for some context $\context{}$ and $U'$ 
such that $\contextp{}$ is obtained from $\context{}$ by applying the $s_v$ function and
$U'$ is obtained by substituting free variables in $U$, 
that are bound by $\boldall$ quantifiers in the context $\context{}$, with $v$.

We shall prove a stronger statement in the following: 
for every $R$, if $\vdash R$ holds then for all terms $v$, $\vdash \mathopen{s_v}\left(R\right)$ holds.

Without loss of generality, we can assume that the bound and the free variables in $R$ are pairwise
distinct and that the bound variables in $R$ are also distinct from the 
variables in $v$. This simplifies the proof below since substitutions of $\boldall$-quantified variables
commute with other connectives and quantifiers in $R.$

For the base case, $s_v(R) = R$, in which case trivially if $\vdash R$ then $\vdash s_v(R)$, for example where $R \equiv \cunit$.

Consider the case when the bottommost rule in a proof is an instance of the \textit{extrude1} rule involving a bold universal quantifier, as follows, 
$\vcenter{\infer[]{\context{ \boldall x. T \cpar U }}{\context{ \mathopen{\boldall x.} \left( T \cpar U \right) }}}$, 
where $\nfv{x}{U}$ and $\vdash \context{ \mathopen{\boldall x.} \left( T \cpar U \right) }$.

By the induction hypothesis, $\vdash \mathopen{s_v}\left(\context{ \mathopen{\boldall x.} \left( T \cpar U \right) }\right)$ holds.
Now the following equalities hold.
\[
\begin{array}{rl}
\mathopen{s_v}\left(\context{ \mathopen{\boldall x.} \left( T \cpar U \right) }\right)
= & \contextp{ \mathopen{s_v}\left(\left( T' \cpar U' \right) \sub{x}{v}\right)} \\
= & \contextp{ \mathopen{s_v}\left(T'\sub{x}{v}\right) \cpar \mathopen{s_v}\left(U'\right) } \\
= & \mathopen{s_v}\left(\context{ \boldall x. T \cpar U }\right)
\end{array}
\]
Hence $\vdash \mathopen{s_v}\left(\context{ \boldall x. T \cpar U }\right)$ holds as required.
%\smallskip

Consider the case where the bottommost rule of a proof is an instance of the \textit{tidy1} rule of the form 
$\vcenter{\infer[]{\context{ \boldall x. \cunit }}{\context{ \cunit }}}$, where $\vdash \context{\cunit}$ holds.
By the induction hypothesis, $\vdash \mathopen{s_v}\left( \context{ \cunit } \right)$ holds.
Since $\mathopen{s_v}\left( \context{ \boldall x. \cunit } \right) = \mathopen{s_v}\left( \context{ \cunit } \right)$, we have $\vdash \mathopen{s_v}\left( \context{ \boldall x. \cunit } \right)$ holds, as required.
%\smallskip

%\ross{Case induced by forall and wen.}
Consider the case where the bottommost rule of a proof is an instance of the \textit{all name} rule of the form 
$\vcenter{\infer[]{\context{ \boldall x. \wen y. P }}{\context{ \wen y. \boldall x. P }}}$, where $\vdash \context{ \wen y. \boldall x. P }$ holds.
By the induction hypothesis, $\vdash s_v\left( \context{ \wen y. \boldall x. P } \right)$ holds. Observe that the following equalities hold, by definition of function $s_v$.
%Thereby the following derivation holds.
\[
\begin{array}{rl}
\mathopen{s_v}\left( \context{  \boldall x. \wen y. P } \right)
=
\contextp{ \mathopen{s_v}\left(\left( \wen y. P' \right)\sub{x}{v}\right) }
=
\contextp{ \mathopen{\wen y. s_v}\left(P'\sub{x}{v} \right)}
=
\mathopen{s_v}\left( \context{ \wen y. \boldall x. P } \right)
\end{array}
\]
Hence $\vdash \mathopen{s_v}\left( \context{ \wen y. \boldall x. P } \right)$ holds, as required.
The case where \textit{all name} involves \textit{new} is similar.
%\smallskip


Consider the case when the bottommost rule does not involve a bold universal quantifier.
We show here one instance where the rule involved is \textit{extrude1}; the other cases are similar.
So suppose the bottommost rule instance is 
\[
\infer[.]{
\context{ \forall x. T \cpar U }
}{
\context{ \mathopen{\forall x.} \left( T \cpar U \right) }
}
\]
By the induction hypothesis, $\vdash \mathopen{s_v}\left(\context{ \mathopen{\forall x.} \left( T \cpar U \right) }\right)$.
So, since
\[
\mathopen{s_v}\left(\context{ \mathopen{\forall x.}\left(T \cpar U\right) }\right)
=
\contextp{ \mathopen{\forall x.}\left(\mathopen{s_v}\left(T'\right) \cpar \mathopen{s_v}\left(U'\right)\right) } 
\]
we have $\vdash \contextp{ \mathopen{\forall x.}\left(\mathopen{s_v}\left(T'\right) \cpar \mathopen{s_v}\left(U'\right)\right) }$ also holds.
Hence, since
\[
\mathopen{s_v}\left(\context{ \forall x. T \cpar U }\right) 
%=  \contextp{ \mathopen{s_v}\left(\forall x. T' \cpar U'\right) } 
= \contextp{ \forall x. \mathopen{s_v}\left(T'\right) \cpar \mathopen{s_v}\left(U'\right) }
\]
and 
\[
\infer[]{
\contextp{ \forall x. \mathopen{s_v}\left(T'\right) \cpar \mathopen{s_v}\left(U'\right) }
}{
\contextp{ \mathopen{\forall x.}\left(\mathopen{s_v}\left(T'\right) \cpar \mathopen{s_v} \left(U'\right)\right) } 
}
\]
we have $\vdash \mathopen{s_v}\left(\context{ \forall x. T \cpar U }\right)$ holds, as required.

The statement of the lemma is then a special case of the stronger statement established by induction. If $\vdash \context{ \boldall x. T }$, where no further bold universal quantifiers occur in the context, then $\vdash \context{ T\sub{x}{v} }$ holds, since in such a scenario $\mathopen{s_v}\left( \context{ \boldall x. T } \right) = \context{ T\sub{x}{v} }$.
\end{proof}


%\begin{proof}
%We require a function over predicates $s_v(T)$ that replaces a certain universal quantifier in $T$ with a substitution for a value $v$. 
%The universal quantifiers to be replaced are highlighted in bold $\boldall$. Note that during a proof the bold operator may be duplicated by the \textit{external} rule and \textit{medial1} rule, hence there may be multiple bold occurrences in a formula.
%The function is defined as follows, where $\odot \in \left\{ \cseq, \cpar, \tensor, \ooplus, \wwith \right\}$ is any binary connective, $\quantifier \in \left\{ \forall, \exists, \new, \wen \right\}$ is any quantifier except bold universal quantification and $\kappa \in \left\{ \alpha, \co{\alpha}, \cunit \right\}$ is any constant or atom.
%\begin{gather*}
%s_v(\boldall x. T) = T \sub{x}{v}
%\qquad
%s_v(\quantifier x. T) =
%\quantifier x. s_v(T)
%\qquad
%s_v\left( T \odot U \right) = 
%s_v(T) \odot s_v(U)
%\qquad
%s_v(\kappa) = \kappa
%\end{gather*}
%In what follows we use that $s_v\left(\context{ U }\right) = \contextp{ s_v\left(U\right) }$, for some context $\context{}$ and such that $\contextp{}$ is obtained from $\context{}$ by applying the $s_v$ function.
%
%Firstly, observe that if $P \equiv Q$, then $s_v(P) \equiv s_v(Q)$. Consider the case where $\alpha$-conversion is applied to a bold universal quantifier. In this case, $\mathopen{s_v}\left( \boldall x. P \right) = P\sub{x}{v} = P\sub{x}{y}\sub{y}{v} = \mathopen{s_v}\left( \boldall y. P\sub{x}{y} \right)$, where $\nfv{y}{\forall x. P}$.
%For the base case, $s_v(T) = T$, in which case trivially if $\vdash T$ then $\vdash s_v(T)$, for example where $T \equiv \cunit$.
%
%Consider the case when the bottommost rule in a proof is an instance of the \textit{extrude1} rule involving a bold universal quantifier, as follows, $\context{ \boldall x. T \cpar U } \longrightarrow \context{ \mathopen{\boldall x.} \left( T \cpar U \right) }$, where $\nfv{x}{U}$ and $\vdash \context{ \mathopen{\boldall x.} \left( T \cpar U \right) }$.
%
%By the induction hypothesis, $\vdash s_v\left(\context{ \mathopen{\boldall x.} \left( T \cpar U \right) }\right)$ holds.
%Now the following equalities hold.
%\[
%\begin{array}{rl}
%\mathopen{s_v}\left(\context{ \mathopen{\boldall x.} \left( T \cpar U \right) }\right)
%= \contextp{ \left( T \cpar U \right) \sub{x}{v}} 
%= \contextp{ T\sub{x}{v} \cpar U } 
%= \mathopen{s_v}\left(\context{ \boldall x. T \cpar U }\right)
%\end{array}
%\]
%Hence $\vdash \mathopen{s_v}\left(\context{ \boldall x. T \cpar U }\right)$ holds as required.
%%\smallskip
%
%Consider the case where the bottommost rule of a proof is an instance of the \textit{tidy1} rule of the form $\context{ \boldall x. \cunit } \longrightarrow \context{ \cunit }$, where $\vdash \context{\cunit}$ holds.
%By the induction hypothesis, $\vdash \mathopen{s_v}\left( \context{ \cunit } \right)$ holds.
%Since $\mathopen{s_v}\left( \context{ \boldall x. \cunit } \right) = \mathopen{s_v}\left( \context{ \cunit } \right)$, we have $\vdash \mathopen{s_v}\left( \context{ \boldall x. \cunit } \right)$ holds, as required.
%%\smallskip
%
%
%%\ross{Case induced by forall and wen.}
%Consider the case where the bottommost rule of a proof is an instance of the \textit{all name} rule of the form $\context{ \boldall x. \wen y. P } \longrightarrow \context{ \wen y. \boldall x. P }$, where $\vdash \context{ \wen y. \boldall x. P }$ holds.
%By the induction hypothesis, $\vdash s_v\left( \context{ \wen y. \boldall x. P } \right)$ holds. Thereby the following derivation holds for some fresh $z$ such that $z \not= x$, $\nfv{z}{v}$ and $\nfv{z}{\wen y. P}$.
%\[
%\begin{array}{rl}
%s_v\left( \context{  \boldall x. \wen y. P } \right)
%\equiv
%\context{ \left( \wen y. P \right)\sub{x}{v} }
%\equiv&
%\context{ \mathopen{\wen z.} \left(P\sub{y}{z} \sub{x}{v} \right) }
%%\\
%%\equiv&
%%\context{ \mathopen{\wen z.} \left(P\sub{y}{z} \sub{x}{v} \right) }
%\\
%\equiv&
%s_v\left( \context{ \mathopen{\wen z. \boldall x.} \left(P\sub{y}{z}\right) } \right)
%\equiv
%s_v\left( \context{ \wen y. \boldall x. P } \right)
%\end{array}
%\]
%Hence $\vdash s_v \left( \context{ \wen y. \boldall x. P } \right)$ holds, as required.
%The case where \textit{all name} involves \textit{new} is similar.
%%\smallskip
%
%\begin{comment}
%Consider the case where the bottommost rule of a proof is an instance of the \textit{all name} rule of the form $\context{ \boldall x. \new y. P } \longrightarrow \context{ \new y. \boldall x. P }$, where $\vdash \context{ \new y. \boldall x. P }$ holds.
%By the induction hypothesis, $\vdash s_v\left( \context{ \new y. \boldall x. P } \right)$ holds. Thereby the following derivation holds for some fresh $z$ such that $z \not= x$, $\nfv{z}{v}$ and $\nfv{z}{\new y. P}$.
%\[
%\begin{array}{rl}
%s_v\left( \context{  \boldall x. \new y. P } \right)
%\equiv
%\context{ \left( \new y. P \right)\sub{x}{v} }
%\equiv&
%\context{ \mathopen{\new z.} \left(P\sub{y}{z} \sub{x}{v} \right) }
%\\
%\equiv&
%\context{ \mathopen{\new z.} \left(P\sub{y}{z} \sub{x}{v} \right) }
%\\
%\equiv&
%s_v\left( \context{ \mathopen{\new z. \boldall x.} \left(P\sub{y}{z}\right) } \right)
%\equiv
%s_v\left( \context{ \new y. \boldall x. P } \right)
%\end{array}
%\]
%Hence $\vdash s_v \left( \context{ \new y. \boldall x. P } \right)$ holds, as required.
%\smallskip
%\end{comment}
%
%Consider when a rule appears within the scope of a bold universal quantifier, of the form $\context{ \boldall x. \contextb{ T } }
%\longrightarrow
%\context{ \boldall x. \contextb{ U } }$, where $\vdash \context{ \boldall x. \contextb{ U } }$ holds.
%By the induction hypothesis, $\vdash s_v\left( \context{ \boldall x. \contextb{ U } } \right)$ holds.
%Now both equivalences
%%of the following equivalences hold:
%$
%s_v\left( \context{ \boldall x. \contextb{ U } } \right)
%\equiv
%\contextp{ \contextb{ U }\sub{x}{v} }
%$
%and
%$
%s_v\left( \context{ \boldall x. \contextb{ T } } \right)
%\equiv
%\contextp{ \contextb{ T }\sub{x}{v} } 
%$ hold.
%Hence, by Lemma~\ref{lemma:substitution}, 
%$
%\vdash s_v\left( \context{ \boldall x. \contextb{ T } } \right)
%$
%holds.
%%\smallskip
%
%Otherwise, the bottommost rule does not directly involve a bold universal quantifier nor is it applied inside the scope of a bold universal quantifier. In this case, the bottommost rule is of the form $\context{ T } \longrightarrow \context{ U }$, where $\vdash \context{ U }$ holds.
%By the induction hypothesis, $\vdash s_v\left(\context{ U }\right)$ holds. 
%Now $s_v\left(\context{ U }\right) = \contextp{ s_v\left(U\right) }$ for some context and %\bogdan{Who is $V$? Shouldn't be $T$ here and in what follows (3 times in total)?}
%$s_v\left(\context{ T }\right) = \contextp{ s_v\left(T\right) }$ for the same context. By Lemma~\ref{lemma:substitution}, since $T \longrightarrow U$, $s_v\left(T\right) \longrightarrow s_v\left(U\right)$; hence $\contextp{ s_v\left(T\right) } \longrightarrow \contextp{ s_v\left(U\right) }$. Thereby $\vdash s_v\left(\context{ T }\right)$ holds, as required.
%%\smallskip
%
%Now by induction on the length of a proof, if $\vdash T$ then $\vdash s_v\left( T \right)$. So, as a special case, if $\vdash \context{ \boldall x T }$, where not further bold universal quantifiers occur in the context, then $\vdash \context{ T\sub{x}{v} }$ holds.
%\end{proof}
% --- end LemmaUniversal.tex ---


\begin{lemma}[Lemma~\ref{lemma:medial}]
Assume that $I$ is a finite subset of natural numbers, $P_i$ and $Q_i$ are formulae, for $i \in I$, and $\tcontext{}$ is a killing context. There exist killing contexts $\tcontextn{0}{}$ and $\tcontextn{1}{}$ and sets of natural numbers $J \subseteq I$ and $K \subseteq I$ such that the following derivation holds:
$
\vcenter{
\infer[]{
\tcontext{ P_i \cseq Q_i \colon i \in I }
}{
\tcontextn{0}{ P_j \colon j \in J } \cseq
\tcontextn{1}{ Q_k \colon k \in K }
}}
$.
\end{lemma}
% --- begin killing-proof.tex ---
\begin{proof}
Proceed by induction on the structure of the killing context. The base case is immediate.
%For the base case consider the killing context consisting of one hole, hence $\left\{ P \cseq Q \right\} = \left\{ P \right\} \cseq \left\{ Q \right\}$, as required.

Consider a predicate of the form $\new x. \tcontext{P_i \cseq Q_i \colon i \in I }$.
By the induction hypothesis, assume there exists $\tcontextn{0}{}$ and $\tcontextn{1}{}$ such that 
\[
\begin{prooftree}
\tcontextn{0}{P_j \colon j \in J} \cseq \tcontextn{1}{ Q_k \colon k \in K }
\justifies
\tcontext{P_i \cseq Q_i \colon i \in I }
\end{prooftree}
\]
where $J \subseteq I$ and $K \subseteq I$.
There are three cases to consider.

If $\tcontextn{0}{P_j \colon j \in J} \equiv \cunit$, then we have derivation 
\[
\infer[.]{
\new x. \tcontext{P_i \cseq Q_i \colon i \in I }
}{
\infer[\mbox{by using $\equiv$}]{
 \cunit \cseq \new x. \tcontextn{1}{ Q_k \colon k \in K }
}{
  \mathopen{\new x.} \left( \cunit \cseq \tcontextn{1}{Q_k \colon k \in K } \right)
}
}
\]

\begin{comment}
\begin{prooftree}
\begin{prooftree}
  \justifies
\using \mbox{by using $\equiv$}
\end{prooftree}
\justifies
\end{prooftree}
\]
\end{comment}

If $\tcontextn{1}{Q_k \colon k \in K} \equiv \cunit$, then we have derivation
\[
\infer[.]{
\new x. \tcontext{P_i \cseq Q_i \colon i \in I }
}{
\infer[\mbox{using $\equiv$}]{
\new x. \tcontextn{0}{P_j \colon j \in J } \cseq \cunit
}{
  \mathopen{\new x.} \left( \tcontextn{0}{P_j \colon j \in J } \cseq \cunit \right)
}
}
\]

Otherwise, $\tcontextn{0}{P_j \colon j \in J} \not\equiv \cunit$ and $\tcontextn{1}{Q_k \colon k \in K} \not\equiv \cunit$ in which case the \textit{medial new} rule can be applied as follows:
\[
\infer[.]{
 \new x. \tcontext{P_i \cseq Q_i \colon i \in I }
}{
 \infer[\mbox{by the \textit{medial new} rule}]{
  \mathopen{\new x.} \left( \tcontextn{0}{P_j \colon j \in J } \cseq \tcontextn{1}{Q_k \colon k \in K } \right) 
 }{
  \new x. \tcontextn{0}{P_j \colon j \in J } \cseq \new x. \tcontextn{1}{Q_k \colon k \in K }
 }
}
\]

In each of the three cases above, killing contexts of the correct form are obtained. The arguments in the cases of universal quantifiers and with follow a similar pattern.
\end{proof}
%\smallskip

\begin{comment}


%\ross{If include above the instead write: A similar argument holds for predicates of the form $\forall x \tcontext{P_i \cseq Q_i \colon i \in I }$. The remaining case for \textit{with} appear in the companion paper~\cite{Horne2015}.}
Now consider a predicate of the form
$
\tcontextn{0}{P_i \cseq Q_i \colon i \in I_0 } \wwith \tcontextn{1}{P_i \cseq Q_i \colon i \in I_1 }
$.
Assume as the induction hypothesis that the following derivations can be constructed, where $J_0 \subseteq I_0$, $K_0 \subseteq I_0$, $J_1 \subseteq I_1$ and $K_1 \subseteq I_1$:
$
\infer{
 \tcontextn{0}{ P_i \cseq Q_i \colon i \in I_0 }
}{
 \tcontextmn{0}{0}{ P_j \colon j \in J_0 } \cseq
 \tcontextmn{0}{1}{ Q_k \colon k \in K_0 }
}
$
and

\begin{array}{l}

%\\\qquad\qquad\qquad\qquad\qquad
\longrightarrow
\end{array}
\]
\[
\begin{array}{l}
\tcontextn{1}{ P_i \cseq Q_i \colon i \in I_1 }
%\\\qquad\qquad\qquad\qquad\qquad\qquad\qquad
\longrightarrow
\tcontextmn{1}{0}{ P_j \colon j \in J_1 } \cseq
\tcontextmn{1}{1}{ Q_k \colon k \in K_1 }
\end{array}
\]
There are three cases to consider. If $\tcontextmn{0}{0}{P_j \colon j \in J_0} \equiv \cunit$ and $\tcontextmn{1}{0}{P_j \colon j \in J_1} \equiv \cunit$, then the following derivation holds.
\[
\begin{array}{l}
\tcontextn{0}{P_i \cseq Q_i \colon i \in I_0 } \wwith \tcontextn{1}{P_i \cseq Q_i \colon i \in I_1 }
\\\quad
\begin{array}{rl}
\longrightarrow& \tcontextmn{0}{1}{ Q_k \colon k \in K_0 } \wwith \tcontextmn{1}{1}{ Q_k \colon k \in K_1 }
\equiv
\cunit \cseq \left(\tcontextn{0}{ Q_k \colon k \in K_0 } \wwith \tcontextn{1}{ Q_k \colon k \in K_1 }\right)
\end{array}
\end{array}
\]
Similarly, if $\tcontextmn{0}{1}{Q_k \colon k \in K_0} \equiv \cunit$ and $\tcontextmn{1}{1}{Q_k \colon k \in K_1} \equiv \cunit$, then the following derivation holds.
\[
\begin{array}{l}
\tcontextn{0}{P_i \cseq Q_i \colon i \in I_0 } \wwith \tcontextn{1}{P_i \cseq Q_i \colon i \in I_1 }
\\\qquad
\begin{array}{rl}
\longrightarrow& \tcontextmn{0}{0}{ P_j \colon j \in J_0 } \wwith \tcontextmn{1}{0}{ P_j \colon j \in J_1 }
\equiv
\left(\tcontextn{0}{ P_j \colon j \in J_0 } \wwith \tcontextn{1}{ P_j \colon j \in J_1 }\right) \cseq \cunit
\end{array}
\end{array}
\]
Otherwise, either $\tcontextmn{0}{0}{ P_j \colon j \in J_0 } \not\equiv \cunit$ or $\tcontextmn{1}{0}{ P_j \colon j \in J_1 } \not\equiv \cunit$ and also either $\tcontextmn{0}{1}{ Q_k \colon k \in K_0 } \not\equiv \cunit$ or $\tcontextmn{1}{1}{ Q_k \colon k \in K_1 } \not\equiv \cunit$. Hence by the \textit{medial} rule the following derivation can be constructed.
\[
\begin{array}{l}
\tcontextn{0}{ P_i \cseq Q_i \colon i \in I_0 }
\wwith
\tcontextn{1}{ P_i \cseq Q_i \colon i \in I_1 }
\\\qquad\qquad
\begin{array}{l}
\longrightarrow
 %\begin{array}[t]{l}
 \left(
 \tcontextmn{0}{0}{ P_j \colon j \in J_0 } \cseq
 \tcontextmn{0}{1}{ Q_k \colon k \in K_0 }
 \right) \wwith
 %\\
  \left(
 \tcontextmn{1}{0}{ P_j \colon j \in J_1 } \cseq
 \tcontextmn{1}{1}{ Q_k \colon k \in K_1 }
 \right)
 %\end{array}
\\
\longrightarrow
 %\begin{array}[t]{l}
 \left(
 \tcontextmn{0}{0}{ P_j \colon j \in J_0 } \wwith
 \tcontextmn{1}{0}{ P_j \colon j \in J_1 }
 \right) \cseq
 %\\
  \left(
 \tcontextmn{0}{1}{ Q_k \colon k \in K_0 } \wwith
 \tcontextmn{1}{1}{ Q_k \colon k \in K_1 }
 \right)
 %\end{array}
\end{array}
\end{array}
\]
Notice that 
$\tcontextmn{0}{0}{} \wwith
\tcontextmn{1}{0}{}$
and
$ \tcontextmn{0}{1}{} \wwith
\tcontextmn{1}{1}{}$
are two killing contexts and $J_0 \cup J_1 \subseteq I_0 \cup I_1$ and 
$K_0 \cup K_1 \subseteq I_0 \cup I_1$ as required.
\end{proof}

\end{comment}
% --- end killing-proof.tex ---


%%
\begin{lemma}[Affine: Lemma~\ref{lemma:bound}]
Any derivation $\vcenter{\infer{Q}{P}}$, is bound such that $\size{P} \preceq \size{Q}$.
\end{lemma}
\begin{proof}
The proof proceeds by checking that each rule preserves the bound on the size of the formula, from which the result follows by induction on the length of a derivation.


Consider the case of the \textit{close} rule. $\occ{\new x. P \cpar \wen x. Q} = \occ{ P } \mplus \occ{ Q } = \occ{\mathopen{\new x.} \left( P \cpar Q \right)}$,
since $P \not\equiv \cunit$ and $Q \not\equiv \cunit$,
 and 
%for the wen count 
%strictly decreases as follows.
$
\wensize{\new x. P \cpar \wen x. Q}
=
\wensize{P} + \left( 1  +  \wensize{Q}\right)
>
\wensize{P} + \wensize{Q}
=
\wensize{\mathopen{\new x.} \left(P \cpar Q\right)}
$.

Consider the case of the \textit{fresh} rule.
For the occurrence count, $\occ{ \wen x. P } = \occ{ \new x. P }$ and the wen count strictly decreases as follows:
$
\wensize{ \wen x. P }
=
 1 + \wensize{ P }
>
\wensize{ P }
=
\wensize{ \new x. P }
$.


Consider the case of the \textit{extrude new} rule, where $Q \not\equiv \cunit$.
If $P \equiv \cunit$, then the occurence count is such that $\occ{ \new x. P \cpar Q } = \left\{\left\{ 0,0 \right\}\right\} \mplus \occ{ Q } > \occ{ Q } = \occ{ \mathopen{\new x.} \left( P \cpar Q \right)}$. If however $P \not\equiv \cunit$, then $\occ{ \new x. P \cpar Q } = \occ{ P } \mplus \occ{ Q } = \occ{ \mathopen{\new x.}\left( P \cpar Q \right) }$. Furthermore, for the new count, the following inequality holds:
$
\newsize{ \new x. P \cpar Q } = \left(1 + \newsize{P}\right)\newsize{Q} \geq 1 + \newsize{P}\newsize{Q} = \newsize{ \mathopen{\new x.} \left( P \cpar Q \right) }
$.
%The new count is also preserved by the same argument.



Consider the case of the \textit{external} rule, where $R \not\equiv \cunit$.
For the occurrence count, by distributivity of $\mcup$ over $\mplus$, 
the following multiset equality holds:
\[
\begin{array}{rl}
\occ{\left( P \wwith Q \right) \cpar R }
& = \left( \occ{P} \mcup \occ{Q} \right) \mplus \occ{R} \\
& = \left( \occ{P}  \mplus \occ{R}\right) \mcup \left(\occ{Q}  \mplus \occ{R} \right) \\
& = \occ{\left( P \cpar R\right) \wwith \left(Q  \cpar R \right)}
\end{array}
\]
For the wen count 
%the following equality holds:
\[
\wensize{ \left( P \wwith Q \right) \cpar R }
=
\left(\wensize{ P } + \wensize{Q}\right) \wensize{R}
=
\wensize{ P }\wensize{R} + \wensize{Q}\wensize{R}
=
\wensize{ \left( P \cpar R\right) \wwith \left( Q \cpar R \right)  }
\]
and similarly for the new count.




Consider the case of the \textit{suspend} rule, where $P \not\equiv \cunit$ and $Q \not\equiv \cunit$.
For the occurrence count, 
% and \textit{with}.
$\occ{ \wen x. P \cseq \wen x. Q } = \occ{ P } \discup \occ{ Q } = \occ{ \mathopen{\wen x.} \left( P \cseq Q \right) }$
and
$\occ{ \wen x. P \cpar \wen x. Q } = \occ{ P } \mplus \occ{ Q } = \occ{ \mathopen{\wen x.} \left( P \cpar Q \right) }$
for \textit{par} and \textit{seq} respectively.
%\item $\occ{ \wen x. P \wwith \wen x. Q } = \occ{ P } \mcup \occ{ Q } = \occ{ \mathopen{\wen x.} \left( P \wwith Q \right) }$
%\end{itemize}
For the wen count for either operator, $\odot \in \left\{\cpar, \cseq \right\}$, the following strict inequality holds, noting $\wensize{P} \geq 1$ for any formula:
\[
\wensize{ \wen x. P \odot \wen x. Q }
=
\left( 1 + \wensize{ P }\right)\left(1 + \wensize{ Q }\right)
=
\wensize{ P } + \wensize{ P }\wensize{ Q } + \wensize{ Q }
>
1 + \wensize{P}\wensize{Q}
=
\wensize{ \mathopen{\wen x.} \left( P \cseq Q \right) }
\]


Consider the case of the \textit{left wen} rules, where $\nfv{x}{Q}$ and $Q \not\equiv \cunit$.
For the occurrence count, there are four cases covering the operators \textit{seq} and \textit{par}. % and \textit{with}.
\begin{itemize}
\item
If $P \equiv \cunit$ then, for \textit{seq}:
$
 \occ{ \mathopen{\wen x.} \left( P \cseq Q \right) } = \occ{ Q }  \mstrict  \left\{\left\{ 0,0 \right\}\right\} \discup \occ{ Q } 
 = \occ{ \wen x. P \cseq Q }
$.

\item
If $P \not\equiv \cunit$ then, for \textit{seq}:
$
\occ{ \wen x. P \cseq Q } = \occ{ P } \discup \occ{ Q } = \occ{ \mathopen{\wen x.} \left( P \cseq Q \right) }
$.

\item
If $P \equiv \cunit$ then for \textit{par}:
$
\occ{ \mathopen{\wen x.} \left( P \cpar Q \right) }
= \occ{Q}
\mstrict
\left\{\left\{ 0,0 \right\}\right\} \mplus \occ{ Q } 
=
\occ{ \wen x. P \cpar Q }
$.

\item
If $P \not\equiv \cunit$ then for \textit{par}:
$
\occ{ \wen x. P \cpar Q } = \occ{ P } \mplus \occ{ Q } = \occ{ \mathopen{\wen x.} \left( P \cpar Q \right) }
$.

%\item %For the \textit{with} operator and $P \equiv \cunit$ the following holds:
%$
%\occ{ \wen x. P \wwith Q } = \left\{\left\{ 0,0 \right\}\right\} \mcup \occ{ Q }
% = \left\{\left\{ 0,0 \right\}\right\} \mcup \occ{ Q }
% = \occ{ \mathopen{\wen x.} \left( P \wwith Q \right) }
%$

%\item
%For the \textit{with} operator and $P \not\equiv \cunit$ the following holds:
%$
%\occ{ \wen x. P \wwith Q } = \occ{ P } \mcup \occ{ Q }
%  = \occ{ \mathopen{\wen x.} \left( P \wwith Q \right) }
%$.
\end{itemize}
For the wen count 
%multiset inequality
$
\begin{array}{rl}
\wensize{ \wen x. P \odot Q }
=
\left(1 + \wensize{ P }\right)\wensize{ Q }
=
\wensize{Q} + \wensize{P}\wensize{ Q }
\geq
1 + \wensize{ P }\wensize{ Q }
=
\wensize{ \mathopen{\wen x.} \left( P \odot Q \right) }
\end{array}
$ holds, for $\odot \in \left\{\cpar, \cseq \right\}$.
Also, for the new count 
%the following equalites hold:
$\newsize{ \wen x. P \cseq Q } = \mmax{\newsize{ P }, \newsize{ Q}} = \newsize{ \mathopen{\wen x.} \left( P \cseq Q \right) }$ and
$\newsize{ \wen x. P \cpar Q } = \newsize{ P }\newsize{ Q } = \newsize{ \mathopen{\wen x.} \left( P \cpar Q \right) }$.
The case \textit{right wen} follows a symmetric argument.



Consider the case for the \textit{extrude} rule, where $Q \not\equiv \cunit$.
$\occ{ \mathopen{\forall x.} \left( P \cpar Q \right) } \mstrict \occ{ \forall x. P \cpar Q }$ by the following:
$
%\begin{array}{rl}
 \left\{\left\{0\right\}\right\} \mcup \left( \occ{ P } \mplus \occ{ Q } \right)
%\\
\mstrict
 \left(\left\{\left\{0\right\}\right\} \mplus \occ{ Q } \right) \mcup \left( \occ{ P } \mplus \occ{ Q } \right)
%\\
=
 \left( \left\{\left\{0\right\}\right\} \mcup \occ{ P } \right) \mplus \occ{ Q }
%\\=&
%\end{array}
$.



Consider the case for the \textit{medial1} rule, where $P \not\equiv \cunit$ and $Q \not\equiv\cunit$.
By distributivity of $\discup$ over $\mcup$,
%the following multiset equality holds:
$
%\begin{array}{rl}
\occ{ \mathopen{\forall x.} \left( P \cseq Q \right) }
=
\left\{\left\{ 0 \right\}\right\} \mcup \left( \occ{ P } \discup \occ{ Q } \right)
=
\left( \left\{\left\{ 0 \right\}\right\} \mcup \occ{ P } \right) \discup \left( \left\{\left\{ 0 \right\}\right\} \mcup \occ{ Q } \right)
=
\occ{ \forall x. P \cseq \forall x. Q }
%\end{array}
$.
Also $\wensize{ \mathopen{\forall x.} \left( P \cseq Q \right) } = \wensize{ \forall x. P \cseq \forall x. Q }$
and $\newsize{ \mathopen{\forall x.} \left( P \cseq Q \right) } = \newsize{ \forall x. P \cseq \forall x. Q }$.



For the \textit{select} rule,
$
 \occ{ \exists x. P }
   = \left\{\left\{ 0 \right\}\right\} \mcup \occ{ P }
   \mstrict \occ{ P }
   = \occ{ P \sub{x}{t} }
$,
by Lemma~\ref{lemma:size-sub}.



Consider the case for the \textit{switch} rule, where $P \not\equiv \cunit$ and $R \not\equiv \cunit$.
If $Q \not\equiv \cunit$, then, since $R \not\equiv \cunit$ we have $\left\{\left\{ 0 \right\}\right\} \mstrict \occ{R}$ and hence $\occ{P} = \occ{P} \mplus \left\{\left\{ 0 \right\}\right\} \mstrict \occ{P} \mplus \occ{R}$; and therefore the following holds since $\discup$ distributes over $\mplus$.
\[
\begin{array}{rl}
\occ{P \tensor \left(Q \cpar R\right)}
=& \occ{P} \discup \left( \occ{Q} \mplus \occ{R} \right)
\\
\mstrict& \left(\occ{P} \mplus \occ{R}\right) \discup \left( \occ{Q} \mplus \occ{R} \right)
\\
=& \left(\occ{P} \discup \occ{Q}\right) \mplus \occ{R}
= \occ{\left(P \tensor Q\right) \cpar R}
\end{array}
\]
If $Q \equiv \cunit$ then, since $\left\{\left\{0\right\}\right\} \mstrict \occ{P}$ and $\left\{\left\{0\right\}\right\} \mstrict \occ{R}$, the following hold.
\[
\occ{P \tensor \left( \cunit \cpar R\right)}
= \occ{P} \discup \occ{R}
\mstrict \occ{P} \mplus \occ{R}
= \occ{\left(P \tensor \cunit\right) \cpar R}
\]



Consider the case of the \textit{sequence} rule, where $P \not\equiv \cunit$ and $S \not\equiv \cunit$.
If $Q \not\equiv \cunit$ and $R \not\equiv \cunit$, then the following holds since $\discup$ distributes over $\mplus$.
\[
\!\!
\begin{array}{rl}
\occ{\left(P \cpar R\right) \cseq \left(Q \cpar S\right)}
=&\!\!\left( \occ{P} \mplus \occ{R} \right) \discup \left( \occ{Q} \mplus \occ{S} \right)
\\
\mstrict&\!\!\left( \occ{P} \mplus \occ{R} \right) \discup \left( \occ{Q} \mplus \occ{S} \right) \discup \left( \occ{P} \mplus \occ{S} \right) \discup \left( \occ{Q} \mplus \occ{R} \right)
\\
=&\!\!\left( \occ{P} \discup \occ{Q} \right) \mplus \left( \occ{R} \discup \occ{S} \right)
= \occ{\left(P \cseq Q\right) \cpar \left(R \cpar S\right)}
\end{array}
\]
If $Q \equiv \cunit$ and $R \not\equiv \cunit$, then, since $\left\{\left\{0\right\}\right\} \mstrict \occ{R}$, and hence $\occ{S} = \occ{S} \mplus \left\{\left\{0\right\}\right\} \mstrict \occ{S} \mplus \occ{R}$, therefore since $\discup$ distributes over $\mplus$.
\[
\begin{array}{rl}
\!\!
\occ{\left(P \cpar R\right) \cseq \left(\cunit \cpar S\right)}
=\left( \occ{P} \mplus \occ{R} \right) \discup \occ{S}
\mstrict&\!\!\left( \occ{P} \mplus \occ{R} \right) \discup \left(\occ{P} \mplus \occ{S}\right)
\\
=&\!\! \occ{P} \mplus \left( \occ{R} \discup \occ{S} \right)
=
\occ{\left(P \cseq \cunit \right) \cpar \left(R \cpar S\right)}
\end{array}
\]
A symmetric argument holds when $Q \not\equiv \cunit$ and $R \equiv \cunit$.

If $Q \equiv \cunit$ and $R \equiv \cunit$, then $\left\{\left\{0\right\}\right\} \mstrict \occ{P}$ and $\left\{\left\{0\right\}\right\} \mstrict \occ{S}$; hence the following strict inequality holds:
$
%\begin{array}{rl}
\occ{\left(P \cpar \cunit\right) \cseq \left(\cunit \cpar S\right)}
= \occ{P} \discup \occ{S}
\mstrict \occ{P} \mplus \occ{S}
= \occ{\left(P \cseq \cunit \right) \cpar \left(\cunit \cseq S\right)}
%\end{array}
$.




Consider the case of the \textit{medial new} rule where $P \not\equiv \cunit$ and $Q \not\equiv \cunit$. For the occurrence count the equality $\occ{ \mathopen{\new x.} \left( P \cseq Q \right) } = \occ{ P } \mplus \occ{ Q } = \occ{ \new x. P \cseq \new x.Q }$ holds.
For the wen count,
$\wensize{ \mathopen{\new x.} \left( P \cseq Q \right) } = 
\wensize{ P } \wensize{ Q } = 
\wensize{ \new x. P \cseq \new x. Q }$.
For the new count the following equality holds:
$
\newsize{ \mathopen{\new x.} \left( P \cseq Q \right) }
=
1 + \mmax{\newsize{ P }, \newsize{ Q }}
=
\mmax{ 1 + \newsize{ P }, 1 + \newsize{ Q }}
=
\newsize{ \new x. P \cseq \new x. Q }
$.



Consider the case for the \textit{medial} rule, where either $P \not\equiv \cunit$ or $R \not\equiv \cunit$ and also either $Q \not\equiv \cunit$ or $S \not\equiv \cunit$.
When all of $P$, $Q$, $R$ and $S$ are not equivalent to the unit, we have the following.
\[
\begin{array}{rl}
\!\!\!
\occ{
 \left(P \wwith R\right) \cseq \left(Q \wwith S\right)
}
=&\!\!\!
\left( \occ{P} \mcup \occ{R} \right) \discup \left( \occ{Q} \mcup \occ{S} \right)
\\
\mstrict
&\!\!\!
\left( \occ{P} \mcup \occ{R} \right) \discup \left( \occ{Q} \mcup \occ{S} \right) \discup \left( \occ{P} \mcup \occ{S} \right) \discup \left( \occ{Q} \mcup \occ{R} \right)
\\
=
&\!\!\!
\left( \occ{P} \discup \occ{Q} \right) \mcup \left( \occ{R} \discup \occ{S} \right)
=
\occ{
 \left(P \cseq Q\right) \wwith \left(R \cseq S\right)
}
\end{array}
\]
For when exactly one of $P$, $Q$, $R$ and $S$ is equivalent to the unit, all cases are symmetric. Without loss of generality suppose that $S \equiv \cunit$ (and possibly also $Q \equiv \cunit$, but $R \not\equiv \cunit$). By distributivity of $\discup$ over $\mcup$ the following holds.
\[
\begin{array}{rl}
\occ{
 \left(P \wwith R\right) \cseq \left(Q \wwith \cunit\right)
}
=&
\left( \occ{P} \mcup \occ{R} \right) \discup \left( \occ{Q} \mcup \left\{\left\{ 0 \right\}\right\} \right)
\\
\mstrict&
\left( \occ{P} \mcup \occ{R} \right) \discup \left( \occ{Q} \mcup \occ{R} \right)
\\
=&
\left( \occ{P} \discup \occ{Q} \right) \mcup \occ{R}
=
\occ{
 \left(P \cseq Q\right) \wwith \left(R \cseq \cunit \right)
}
\end{array}
\]
There is one more form of case to consider for the medial: either $P \not\equiv \cunit$, $Q \equiv \cunit$, $R \equiv \cunit$ and $S \not\equiv \cunit$; or $P \equiv \cunit$, $Q \not\equiv \cunit$, $R \not\equiv \cunit$ and $S \equiv \cunit$. We consider only the former case. The later case, can be treated symmetrically.
Since $P \not\equiv \cunit$ and $S \not\equiv \cunit$, $\left\{\left\{ 0 \right\}\right\} \mstrict \occ{P}$ and $\left\{\left\{ 0 \right\}\right\} \mstrict \occ{S}$. Therefore, $\occ{P} \mcup \left\{\left\{ 0 \right\}\right\} \mstrict \occ{P} \mcup \occ{S}$ and $\occ{Q} \mcup \left\{\left\{ 0 \right\}\right\} \mstrict \occ{P} \mcup \occ{S}$. Hence, we have established that $\left(\occ{P} \mcup \left\{\left\{ 0 \right\}\right\}\right) \discup \left( \occ{Q} \mcup \left\{\left\{ 0 \right\}\right\} \right) \mstrict \occ{P} \mcup \occ{S}$.
%
Note that the restriction on the \textit{medial} rule, either $P \not\equiv \cunit$ or $R \not\equiv \cunit$ and also either $Q \not\equiv \cunit$ or $S \not\equiv \cunit$, excludes any further cases. Hence we have established that $
\occ{\left(P \wwith R\right) \cseq \left(Q \wwith S\right)}
\mstrict
\occ{\left(P \cseq Q\right) \wwith \left(R \cseq S\right)}$. %Wen new counts need not be considered due to lexicographical ordering.
%\ross{This is why we could simplify the measures!}

For the \textit{with name} rule $\occ{\quantifier x.P \wwith \quantifier x.Q} = \occ{P} \mcup \occ{Q} = \occ{\mathopen{\quantifier x.}\left(P \wwith Q\right)}$,
where $\quantifier \in \left\{ \new, \wen \right\}$. For the new count $\newsize{\new x.P \wwith \new x.Q} = 2 + \newsize{P} + \newsize{Q} > 1 + \newsize{P} + \newsize{Q} = \newsize{\mathopen{\new x.}\left(P \wwith Q\right)}$
and $\newsize{\wen x.P \wwith \wen x.Q} = \newsize{\mathopen{\wen x.}\left(P \wwith Q\right)}$.
Similarly, $\wensize{\wen x.P \wwith \wen x.Q} > \wensize{\mathopen{\wen x.}\left(P \wwith Q\right)}$.
For \textit{left name}, \textit{right name} and \textit{all name}, the size of formulae are invariant.

The cases for the rules \textit{tidy}, \textit{tidy name}, \textit{left}, \textit{right}, \textit{atomic interact} are established by the following inequalities:
$\occ{ \cunit } \mstrict \occ{ \cunit \wwith \cunit }$,
$\occ{ \cunit } \mstrict \occ{ \new x. \cunit }$,
$\occ{ \cunit } \mstrict \occ{ \co{a} \cpar a }$,
$\occ{ P } \mstrict \occ{ P \ooplus Q }$
and
$\occ{ Q } \mstrict \occ{ P \ooplus Q }$.

Hence the lemma holds by induction on the length of the derivation.
\end{proof}

%% Existential lemma

\begin{lemma}
[Lemma~\ref{lemma:split-exists}]
If $\vdash \exists x. P \cpar Q$, then there exist formulae $V_i$ and values $v_i$ such that $\vdash P\sub{x}{v_i} \cpar V_i$, where $1 \leq i \leq n$, and $n$-ary killing context $\tcontext{}$ such that $\vcenter{\infer{Q}{ \tcontext{V_1, V_2, \hdots, V_n}}}$
and if $\tcontext{}$ binds $y$ then $\nfv{y}{\left(\exists x . P\right)}$.
\end{lemma}

% --- begin split-exists.tex ---
\begin{proof}
The proof proceeds by induction on the size of the proof in Definition~\ref{definition:size},
until the principal \textit{exists} operator is removed from the proof, according to the base case.
In the base case, the bottommost rule in a proof is an instance of the \textit{select} rule of the form 
$
\vcenter{
\infer[,]{
\exists x. T \cpar U 
}{
 T\sub{x}{v} \cpar U
}
}$ where $\vdash T\sub{x}{v} \cpar V$ holds; hence splitting is immediately satisfied.
%
As in every splitting lemma, there are commutative cases for \textit{new}, \textit{wen}, \textit{all}, \textit{with}, \textit{times} and two for \textit{seq}.
%These case follow a pattern hence are omitted in the interest of space.

Consider the commutative case induced by the \textit{external} rule.
% where $\exists x. T$ is the principal formula.
The bottommost rule is the form
\[
\infer[]{
\exists x. T \cpar \left( U \wwith V \right) \cpar W \cpar P
}{
\left(\exists x. T \cpar U \cpar W \wwith
\exists x. T \cpar V \cpar W\right)
\cpar P
}
\]
where it holds that $\vdash \left(\left(\exists x. T \cpar U \cpar W\right) \wwith \left(\exists x. T \cpar V \cpar W\right)\right) \cpar P$.
By Lemma~\ref{lemma:split-times}, $\vdash \exists x. T \cpar U \cpar W \cpar P$ and $\vdash \exists x. T \cpar V \cpar W \cpar P$; and furthermore 
%the two strict multiset inequalities 
$\size{\exists x. T \cpar U \cpar W \cpar P} \mstrict \size{ \exists x. T \cpar \left( U \wwith V \right) \cpar W \cpar P }$
and
$\size{\exists x. T \cpar V \cpar W \cpar P} \mstrict \size{ \exists x. T \cpar \left( U \wwith V \right) \cpar W \cpar P }$ hold.
%
Hence, by the induction hypothesis, there exist $Q_i$ and $u_i$ such that $\vdash T\sub{x}{u_i} \cpar Q_i$, for $1 \leq i \leq m$, and $R_j$ and $v_j$ such that
$\vdash T\sub{x}{v_j} \cpar R_j$, for $1 \leq j \leq n$; and $m$-ary and $n$-ary killing contexts $\tcontextn{0}{}$ and $\tcontextn{1}{}$ such that 
the derivations $(1)$ and $(2)$ below hold. 
\[
\begin{array}{ccc}
\infer[]{
U \cpar W \cpar P
}{
 \tcontextn{0}{Q_1, \hdots, Q_m}
}
&
\infer[]{
V \cpar W \cpar P
}{
 \tcontextn{1}{R_1, \hdots, R_n}
}
&
\infer[]{
\left( U \wwith V \right) \cpar W \cpar P
}{
\infer[]{
\left(U \cpar W \cpar P\right) \wwith \left(V \cpar W \cpar P\right) 
}{
\tcontextn{0}{Q_1, \hdots, Q_m} \wwith \tcontextn{1}{R_1, \hdots, R_n}
}} \\
(1) & (2) & (3)
\end{array}
\]
Thus the derivation $(3)$ above can be constructed.
%\bogdan{the $R_j, \hdots, R_n$ above should be replaced by $R_1, \hdots, R_n$}
Notice that $\tcontextn{0}{} \wwith \tcontextn{1}{}$ is an $m+n$-ary killing context, as required.



Consider the commutative case induced by the \textit{extrude1} rule. In this case, the bottommost rule is
\[
\infer[]{
\exists x. T \cpar \forall y. U \cpar V \cpar W
}{
\mathopen{\forall y.}\left( \exists x. T \cpar U \cpar V \right) \cpar W
}
\]
 assuming $\nfv{y}{(\exists x. T \cpar V)}$
where $\vdash \mathopen{\forall y.}\left( \exists x. T \cpar U \cpar V \right) \cpar W$ holds.
By Lemma~\ref{lemma:universal}, for every variable $z$, $\vdash \left( \exists x. T \cpar U \cpar V \right)\sub{y}{z} \cpar W$ holds.
Furthermore, by definition of substitution $\left( \exists x. T \cpar U \cpar V \right)\sub{y}{z} \cpar W \equiv \exists x. T \cpar U\sub{y}{z} \cpar V \cpar W$,  since $\nfv{y}{(\exists x. T \cpar V)}$.
%
Now observe the strict multiset inequality $
\size{ \exists x. T \cpar U\sub{y}{z} \cpar V \cpar W } \mstrict \size{ \exists x. T \cpar \forall y. U \cpar V \cpar W }
$ holds;
hence, by the induction hypothesis, for every variable $z$, there exist formulae $P^z_i$ and values $v^z_i$ such that $\vdash T\sub{x}{v^z_i} \cpar P^z_i$ holds, for $1 \leq i \leq n$, and $n$-ary killing context $\tcontext{}$ such that derivation $(4)$ below can be constructed. Hence, for $\nfv{z}{(\forall y. U \cpar V \cpar W)}$, the  derivation $(5)$
below can be constructed:
\[
\begin{array}{cc}
\infer[]{
U\sub{y}{z} \cpar V \cpar W
}{
\tcontext{P^z_1, \hdots, P^z_n}
}
&
\infer[]{
\forall y. U \cpar V \cpar W
}{
\infer[]{
\mathopen{\forall z.} \left( U\sub{y}{z} \cpar V \cpar W \right) 
}{
\mathopen{\forall z.} \tcontext{P^z_1, \hdots, P^z_n} 
}} \\
(4) & (5) 
\end{array}
\]
Notice that $\mathopen{\forall z.} \tcontext{}$ is a $n$-ary killing context as required.

Consider the commutative cases involving the \textit{sequence} rule.
We present the scenario where the principal formula $\exists x. U$ moves entirely to the left hand side of \textit{seq} operator. The cases where the principal formula moves entirely to the right hand side of the \textit{seq} operator and the commutative case for \textit{times}, are similar to the cases presented below.
In the scenario we consider, the bottommost rule in a proof is of the following form:
\[
\infer[]{
\exists x. U \cpar \left( V \cseq P \right) \cpar W \cpar Q
}{
\left( \left( \exists x. U \cpar V \cpar W\right) \cseq P \right) \cpar Q
}
\]
such that $\vdash \left( \left( \exists x. U \cpar V \cpar W\right) \cseq P \right) \cpar Q$ holds.
%
By Lemma~\ref{lemma:split-times}, there exist $R_i$ and $S_i$ such that $\vdash \exists x. U \cpar V \cpar W \cpar R_i$ and $\vdash P \cpar S_i$ hold, for $1 \leq i \leq n$, and $n$-ary killing context $\tcontext{}$ such that the derivation
$
\vcenter{
\infer[]{
 Q 
}{
 \tcontext{ R_1 \cseq S_1, \hdots, R_n \cseq S_n }
}}
$ holds, 
and furthermore the size of the proof of ${ \exists x. U \cpar V \cpar W \cpar R_i }$ is bounded above by the size of the proof of $\left( \left( \exists x. U \cpar V \cpar W\right) \cseq P \right) \cpar Q$ hence strictly bounded above by the size of the proof of $\exists x. U \cpar \left( V \cseq P \right) \cpar W \cpar Q$.
By the induction hypothesis, for $1 \leq i \leq n$, there exist formulae $P^i_j$ and terms $t^i_j$ such that $\vdash U\sub{x}{t^i_j} \cpar P^i_j$, for $1 \leq j \leq m_i$, and killing contexts $\tcontextn{i}{}$ such that the derivation $
\vcenter{
\infer[]{
V \cpar W \cpar R_i
}{
 \tcontextn{i}{P^i_1, \hdots, P^i_{m_i}}
}}$ holds.
Hence the following derivation can be constructed, as required.
\[
\infer[]{
\left( V \cseq P \right) \cpar W \cpar Q
}{
\infer[]{
\left( V \cseq P \right) \cpar W \cpar \tcontext{ R_1 \cseq S_1, \hdots, R_n \cseq S_n } 
}{
\infer[]{
\tcontext{ \left( V \cseq P \right) \cpar W \cpar R_i \cseq S_i \colon 1 \leq i \leq n } 
}{
\infer[]{
\tcontext{ \left( V \cpar W \cpar R_i \right) \cseq \left( P \cpar S_i \right) \colon 1 \leq i \leq n }
}{
\infer[]{
\tcontext{ V \cpar W \cpar R_i \colon 1 \leq i \leq n } 
}{
\tcontext{ \tcontextn{1}{P^1_1, \hdots, P^1_{m_1}}, \hdots, \tcontextn{n}{P^n_1, \hdots, P^n_{m_n}} } 
}}}}}
\]
Notice that $\tcontext{ \tcontextn{1}{}, \hdots, \tcontextn{n}{} }$ is a $\sum_{i=1}^n{m_i}$-ary killing context as required.


Consider the commutative case induced by the \textit{extrude new} rule.
In this case, the bottommost rule of a proof is of the form 
\[
\infer[\mbox{, where $\nfv{y}{\exists x. P \cpar R}$ and $\vdash \mathopen{\new y.} \left( \exists x. P \cpar Q \cpar R\right) \cpar S$ holds.}]{
\exists x. P \cpar \new y. Q \cpar R \cpar S
}{
 \mathopen{\new y.} \left( \exists x. P \cpar Q \cpar R\right) \cpar S
}
\] 
By Lemma~\ref{lemma:split-times}, there exist $T$ and $U$ such that $\vdash \exists x. P \cpar Q \cpar R \cpar U$, $\nfv{y}{T}$ holds and either $T = U$ or $T = \wen y. U$, and also
$\vcenter{\infer{S}{T}}$.
Furthermore, the size of the proof of $\exists x. P \cpar Q \cpar R \cpar U$ is bounded above by the size of the proof of $\mathopen{\new y.} \left( \exists x. P \cpar Q \cpar R\right) \cpar S$ and hence strictly bounded above by the size of the proof of $\exists x. P \cpar \new y. Q \cpar R \cpar S$, enabling the induction hypothesis.
Hence, by the induction hypothesis, there exist formulae $V_i$ and terms $t_i$ such that $\vdash P\sub{x}{t_i} \cpar V_i$ holds, for $1 \leq i \leq n$, and $n$-ary killing context $\tcontext{}$ such that the derivation $
\vcenter{
\infer{
Q \cpar R \cpar U
}{
\tcontext{ V_1, \hdots, V_n }
}}$ holds.
%
Observe that, either $T = U$ and $\nfv{y}{U}$, and hence we have
derivation $(6)$ below; 
or $T = \wen y. U$ and hence we have derivation $(7)$ below. Thereby we can construct the derivation 
$(8)$ below.
\[
\begin{array}{ccc}
\infer[]{
\new y. Q \cpar R \cpar T
}{
 \mathopen{\new y.} \left(Q \cpar R \cpar U\right)
}
&
\infer[]{
\new y. Q \cpar R \cpar \wen y. U
}{
\infer[]{
 \mathopen{\new y.} \left(Q \cpar R\right) \cpar \wen y. U
}{
 \mathopen{\new y.} \left(Q \cpar R \cpar U\right)
}}
&
\infer[]{
\new y. Q \cpar R \cpar S
}{
\infer[]{
\new y. Q \cpar R \cpar T
}{
\infer[]{
\mathopen{\new y.} \left( Q \cpar R \cpar U \right)
}{
\new y. \tcontext{ V_1, \hdots, V_n }
}}} \\
(6) & (7) & (8)
\end{array}
\]
Observe that $\new y. \tcontext{ }$ is a $n$-ary killing context as required.


Consider the commutative case induced by the \textit{right wen} rule.
In this case, the bottommost rule of a proof is of the form
\[
\infer[\mbox{, where $\nfv{y}{\exists x. P \cpar R}$.}]{
\exists x. P \cpar \wen y. Q \cpar R \cpar S
}{
 \wen y \left( \exists x. P \cpar Q \cpar R\right) \cpar S
}
\]
By Lemma~\ref{lemma:split-times}, there exist $T$ and $U$ such that 
$\vdash \exists x. P \cpar Q \cpar R \cpar U$, $\nfv{y}{T}$ and either $T = U$ or $T = \new y. U$, and also
$\vcenter{\infer{S}{T}}$.
Furthermore, the size of the proof of $\exists x.P \cpar Q \cpar R \cpar U$ is bounded above by the size of the proof of $\mathopen{\wen y.} \left( \exists x. P \cpar Q \cpar R\right) \cpar S$ and hence strictly bounded above by the size of the proof of $\exists x. P \cpar \wen y. Q \cpar R \cpar S$, enabling the induction hyothesis.
Hence, by the induction hypothesis, there exist formulae $V_i$ and terms $t_i$ such that $\vdash P\sub{x}{t_i} \cpar V_i$, for $1 \leq i \leq n$, and $n$-ary killing context $\tcontext{}$ such that $
\vcenter{
\infer{
Q \cpar R \cpar U
}{
 \tcontext{ V_1, \hdots, V_n }
}}$.
%
Observe that either $T = U$ and $\nfv{y}{U}$ hence the derivation $(9)$ below holds; 
or $T = \new y. U$ hence the derivation $(10)$ below holds. Hence the derivation $(11)$ below can be constructed:
\[
\begin{array}{ccc}
\infer{
\wen y. Q \cpar R \cpar T
}{
\infer{
 \mathopen{\wen y.} \left(Q \cpar R \cpar U\right)
}{
 \mathopen{\new y.} \left(Q \cpar R \cpar U\right)
}}
&
\infer{
\wen y. Q \cpar R \cpar \new y. U
}{
\infer[]{
 \mathopen{\wen y.} \left(Q \cpar R\right) \cpar \new y. U 
}{
 \mathopen{\new y.} \left(Q \cpar R \cpar U\right)
}}
&
\infer[]{
\wen y. Q \cpar R \cpar S
}{
\infer[]{
\wen y. Q \cpar R \cpar T
}{
\infer[]{
\mathopen{\new y.} \left( Q \cpar R \cpar U \right)
}{
\new y. \tcontext{ V_1, \hdots, V_n }
}}} \\
(9) & (10) & (11)
\end{array}
\]
Observe that $\new y. \tcontext{ }$ is a $n$-ary killing context as required.

In many commutative cases, the bottommost rule does not interfere with the principal formula. Consider when a rule is applied outside the scope of the principal formula.
%
In this case, the bottommost rule in a proof is of the form
$
\vcenter{
\infer[]{
 \exists x. U \cpar \context{ V }
}{
 \exists x. U \cpar \context{ W }
}}$
such that $\vdash \exists x. U \cpar \context{ W }$.
By the induction hypothesis, there exist formulae $P_i$ and terms $t_i$, for $1 \leq i \leq n$ such that $\vdash U\sub{x}{t_i} \cpar P_i$, for $1 \leq i \leq n$, and $n$-ary killing context $\tcontext{}$ such that 
$
\vcenter{
\infer[]{
\context{ W }
}{
 \tcontext{ P_1, \ldots, P_n }
}}
$.
Hence $
\vcenter{
\infer[]{
\context{ V }
}{
\infer[]{
 \context{ W } 
}{
 \tcontext{ P_1, \ldots, P_n }
}}}
$ as required.

%In the remaining commutative cases, the bottommost rule may appear inside the context of principal formula without affecting the connective itself. 
Consider the following application of any rule 
$
\vcenter{
\infer[]{
\exists x. \context{T} \cpar W
}{
 \exists x. \context{U} \cpar W
}}
$
such that $\vdash \exists x. \context{U} \cpar W$. 
By the induction hypothesis, there exist formulae $P_i$ and terms $t_i$ where $\vdash \context{U}\sub{x}{t_i} \cpar P_i$, for $1 \leq i \leq n$, and $n$-ary killing context $\tcontext{}$ such that 
%the derivation 
$
\vcenter{
\infer[]{
W
}{
 \tcontext{ P_1, \ldots, P_n }
}}
$.
Hence, by Lemma~\ref{lemma:substitution}, the proof 
$
\vcenter{
\infer[]{
\context{T}\sub{x}{v_i} \cpar P_i
}{
\infer[]{
 \context{U}\sub{x}{v_i} \cpar P_i 
}{
 \cunit
}}}
$ holds.

All cases have been considered hence the lemma holds by induction on the size of a proof.
\end{proof}%end of proof

% --- end split-exists.tex ---

%% Context invariant lemma

% --- begin context-invariant.tex ---
\begin{lemma}
[Lemma~\ref{lemma:invariant}]
If $\vdash \context{T}$, then there exist formulae $U_i$ and substitutions $\sigma_i$, for $1 \leq i \leq n$, and $n$-ary killing context $\tcontext{}$ such that $\vdash T\sigma_i \cpar U_i$; and, for any formula $V$ there exist $W_i$ such that either $W_i = V\sigma_i \cpar U_i$ or $W_i = \cunit$ and the following holds:
$
\vcenter{
\infer[]{
 \context{V}
}{
 \tcontext{ W_1, W_2, \hdots, W_n }
}}
$.
\end{lemma}

\begin{proof}
The proof proceeds by induction on the size of the formula part of the context (n.b.\ not counting the size of atoms). The base case concerning one hole is immediate.

Consider the case for a context of the form $\exists x. \context{} \cpar P$, where $\vdash \exists x. \context{ T } \cpar P$.
By Lemma~\ref{lemma:split-exists},
there exist formulae $Q_i$ and values $v_i$ such that $\vdash \context{ T } \sub{x}{v_i} \cpar Q_i$, for $1 \leq i \leq n$;
and $n$-ary killing context $\tcontext{}$ such that the following derivation holds.
\[
\infer{P}{
 \tcontext{ Q_1, Q_2, \hdots, Q_n }
}
\]
%
For context $\context{}$ and any formula $U$, let $\contextn{i}{}$ and $\sigma_{i}$ be such that $\context{U}\sub{x}{v_i} \equiv \contextn{i}{U\sigma_{i}}$.
Notice that for first-order quantifiers, the substitutions does not increase the size of the formula part of the context. It can only increases the size of terms in atoms, which are not counted in this induction.
%
Since $\vdash \context{ T } \sub{x}{v_i} \cpar Q_i$ holds, then $\vdash \contextn{i}{ T\sigma_i }  \cpar Q_i$ holds.
Therefore, by the induction hypothesis, there exists formula $V^i_j$ such that either $V^i_j = \cunit$ or $V^i_j = \left(U\sigma_i\right) \sigma^i_j \cpar W^i_j$, where $\vdash \left(T \sigma_i\right) \sigma^i_j \cpar W^i_j$, for $1 \leq j \leq m_i$;
and $m_i$-ary killing context $\tcontextn{i}{}$
%\bogdan{shouldn't be $\tcontextn{i}{}$ here?}
such that 
$
\context{ U } \sub{x}{v_i} \cpar Q_i
\equiv
\contextn{i}{ U\sigma_i } \cpar Q_i
$
and the following derivation holds:
\[
\infer[.]{
\contextn{i}{ U\sigma_i } \cpar Q_i
}{
\tcontextn{i}{ V^i_1, V^i_2, \hdots, V^i_{m_i} }
}
\]
Hence the following derivation can be constructed for all formulae $U$.
\[
\infer[]{
\exists x. \context{ U } \cpar P
}{
\infer[]{
\exists x. \context{ U } \cpar \tcontext{ Q_1, \hdots, Q_n }
}{
\infer[]{
\exists x. \context{ U } \cpar \tcontext{ Q_i \colon 1 \leq i \leq n }
}{
\infer[]{
\tcontext{ \exists x. \context{ U } \cpar Q_i \colon 1 \leq i \leq n }
}{
\infer[]{
\tcontext{ \context{ U } \sub{x}{v_i} \cpar Q_i \colon 1 \leq i \leq n }
}{
\tcontext{ \tcontextn{i}{ V^i_j \colon 1 \leq j \leq m_i } \colon 1 \leq i \leq n }
}}}}}
\]
Observe $V^i_j = \cunit$ or $V^i_j = U \mathclose{\left(\sigma_i \cdot \sigma^{i}_j\right)} \cpar W^i_j$, such that $\vdash T \mathclose{\left(\sigma_i \cdot \sigma^{i}_j\right)} \cpar W^i_j$, for all $i$ and $j$, as required.



\begin{comment}
Consider the case for a context of the form $\left( \context{ } \ooplus U \right) \cpar P$ where $\vdash \left( \context{T} \ooplus U \right) \cpar P$.
By Lemma~\ref{lemma:split-plus}, there exist killing context $\tcontext{}$ and $V_i$ for $1 \leq i \leq n$ such that either $\vdash \context{T} \cpar V_i$ or $\vdash U \cpar V_i$ and the following derivation holds:
\[
 P \longrightarrow \tcontext{ V_1, \hdots, V_n }
\]
Now, by the induction hypothesis, if $\vdash \context{T} \cpar V_i$ holds, then there exist killing contexts $\tcontextn{i}{}$, predicates $W^i_j$ and substitutions $\sigma^i_j$ for $1 \leq j \leq m_i$ such that $\vdash T\sigma^i_j \cpar W^i_j$; and, for any predicate $Q$ for $1 \leq j \leq m_i$, there exists $S^i_j$ where wither $S^i_j = \cunit$ or $S^i_j = Q\sigma^i_j \cpar W^i_j$ and the following derivation holds:
\[
 \context{Q} \cpar V_i 
   \longrightarrow \tcontextn{i}{ S^i_1, \hdots, S^i_{m_i} }
\]
Hence following derivation can be constructed for any predicate $Q$.
\[
\begin{array}{rl}
\left( \context{Q} \ooplus U \right) \cpar P
  \longrightarrow& \left( \context{Q} \ooplus U \right) \cpar \tcontext{ V_1, \hdots, V_n } \\
  \longrightarrow& \tcontext{ \left( \context{Q} \ooplus U \right) \cpar V_i \colon 1 \leq i \leq n } \\
  \longrightarrow& \tcontext{ R_1, \hdots, R_n }
\end{array}
\]
where $R_i$ is defined as follows.
\[
R_i = 
\left\{
\begin{array}{lr}
\cunit                                                     & \mbox{if $\vdash U \cpar V_i$} \\
\tcontextn{i}{ S^i_1, \hdots, S^i_{m_i} }  & \mbox{otherwise}
\end{array}
\right.
\] 
The above is well-defined since if $\vdash U \cpar V_i$ holds, then 
\[
\left( \context{Q} \ooplus U \right) \cpar V_i
 \longrightarrow U \cpar V_i
 \longrightarrow \cunit = R_i
\]
and, if $\vdash U \cpar V_i$ does not hold, then $\vdash \context{T} \cpar V_i$ must hold;
hence the following derivation can be applied:
\[
\left( \context{Q} \ooplus U \right) \cpar V_i
  \longrightarrow \context{Q} \cpar V_i 
  \longrightarrow \tcontextn{i}{ S^i_1, \hdots, S^i_{m_i} } = R_i
\]
Hence the induction invariant is satisfied.

\end{comment}


Consider the case for a context of the form $\new x. \context{} \cpar P$, where $\vdash \new x. \context{ T } \cpar P$.
By Lemma~\ref{lemma:split-times},
there exist formulae $Q$ and $\hat{Q}$ such that $\vdash \context{ T } \cpar \hat{Q}$ and either $Q = \hat{Q}$ or $Q$ and $\wen x.\hat{Q}$, and also 
$
\vcenter{
\infer{P}{Q}}
%P \longrightarrow Q
$.
Therefore, by the induction hypothesis,
there exist formulae $V_i$ and $W_i$ and substitutions $\sigma_i$ such that either $V_i = \cunit$ or $V_i = U \sigma_i \cpar W_i$, where $\vdash T \sigma_i \cpar W_i$, for $1 \leq i \leq n$;
and $n$-ary killing context $\tcontext{}$
such that 
\[
\infer[.]{
\context{ U } \cpar \hat{Q}
}{
\tcontext{ V_1, V_2, \hdots, V_n }
}
\]
Hence the following derivation 
\[
\infer[]{
\new x. \context{ U } \cpar P
}{
\infer[]{
\new x. \context{ U } \cpar Q
}{
\infer[]{
\mathopen{\new x.} \left( \context{ U } \cpar \hat{Q} \right)
}{
\new x. \tcontextn{i}{ V_i \colon 1 \leq i \leq n }
}}}
\]
can be constructed for all formulae $U$, as required.

\begin{comment}
Consider the case for a context of the form $\new x \context{} \cpar P$, where $\vdash \new x \context{ T } \cpar P$ and $\nfv{x}{P}$.
By Lemma~\ref{lemma:split-times},
there exist predicates $Q_i$ and $\hat{Q}_i$ such that $\vdash \context{ T } \cpar \hat{Q}_i$ and either $Q_i = \hat{Q}_i$ or $Q_i$ and $\wen x \hat{Q}_i$, for $1 \leq i \leq n$;
and $n$-ary killing context $\tcontext{}$ such that 
$
P \longrightarrow \tcontext{ Q_i \colon 1 \leq i \leq n }
$.
Therefore, by induction on the size of the context,
there exist predicates $V^i_j$ and $W^i_j$ and substitutions $\sigma^i_j$ such that either $V^i_j = \cunit$ or $V^i_j = U \sigma^i_j \cpar W^i_j$, where $\vdash T \sigma^i_j \cpar W^i_j$, for $1 \leq j \leq k_i$;
and $k_i$-ary killing context $\tcontextn{i}{}$
such that 
$
\context{ U } \cpar \hat{Q}_i
\longrightarrow
\tcontextn{i}{ V^i_1, V^i_2, \hdots, V^i_{k_i} }
$.
Hence the following derivation can be constructed for all predicates $U$.
\[
\begin{array}{rl}
\new x \context{ U } \cpar P
\longrightarrow&
\new x \context{ U } \cpar \tcontext{ Q_i \colon 1 \leq i \leq n }
\\
\longrightarrow&
\tcontext{ \new x \context{ U } \cpar Q_i \colon 1 \leq i \leq n }
\\
\longrightarrow&
\tcontext{ \new x \left( \context{ U } \cpar \hat{Q}_i \right) \colon 1 \leq i \leq n }
\longrightarrow
\tcontext{ \new x \tcontextn{i}{ V^i_j \colon 1 \leq j \leq k_i } \colon 1 \leq i \leq n }
\end{array}
\]
Thereby the induction invariant is satisfied.
\end{comment}



Consider the case for a context of the form $\wen x. \context{} \cpar P$, where $\vdash \wen x. \context{ T } \cpar P$. %and $\nfv{x}{P}$.
By Lemma~\ref{lemma:split-times},
there exist formulae $Q$ and $R$ such that $\nfv{x}{Q}$ and $\vdash \context{T} \cpar R$ and either $Q = R$ or $Q = \new x. R$, and also
$
\vcenter{\infer{P}{Q}}
$.
Therefore, by the induction hypothesis,
there exist formulae $V_i$ and $W_i$ and substitutions $\sigma_i$ such that either $V_i = \cunit$ or $V_i = U \sigma_i \cpar W_i$, where $\vdash T \sigma_i \cpar W_i$, for $1 \leq i \leq n$;
and $n$-ary killing context $\tcontext{}$
such that 
\[
\infer[.]{
\context{ U } \cpar R
}{
\tcontext{ V_1, V_2, \hdots, V_n }
}
\]
%
In the former case that $Q = R$, since $\nfv{x}{Q}$, the derivation
\[
\infer[]{
\wen x. \context{ U } \cpar R
}{
\infer[]{
\new x. \context{ U } \cpar R
}{
\new x. \left( \context{ U } \cpar R \right)
}}
\] holds.
In the case, $Q = \new x. R$ the derivation 
\[
\infer[]{
\wen x .\context{ U } \cpar \new x. R
}{
\mathopen{\new x.} \left( \context{ U } \cpar R \right)
}
\] holds.
Hence, for all formulae $U$,
%by applying one of the above derivations we can construct the following derivation :
\[
\infer[]{
\wen x. \context{ U } \cpar P
}{
\infer[.]{
\wen x. \context{ U } \cpar Q
}{
\infer[]{
\mathopen{\new x.} \left( \context{ U } \cpar R \right)
}{
\new x. \tcontext{ V_1, V_2, \hdots, V_n}
}}}
\]
%Thereby the induction invariant is satisfied.
\begin{comment}
Consider the case for a context of the form $\wen x \context{} \cpar P$, where $\vdash \wen x \context{ T } \cpar P$. %and $\nfv{x}{P}$.
By Lemma~\ref{lemma:split-times},
there exist predicates $Q_i$ and $R_i$ such that $\vdash \context{T} \cpar R_i$ and either $Q_i = R_i$ or $Q_i = \new x R_i$ and , for $1 \leq i \leq n$,
and $n$-ary killing context $\tcontext{}$ such that 
$
P \longrightarrow \tcontext{ Q_1, Q_2, \hdots, Q_n }
$.
Therefore, by induction on the size of the context,
there exist predicates $V^i_j$ and $W^i_j$ and substitutions $\sigma^i_j$ such that either $V^i_j = \cunit$ or $V^i_j = U \sigma^i_j \cpar W^i_j$, where $\vdash T \sigma^i_j \cpar W^i_j$, for $1 \leq j \leq k_i$;
and $k_i$-ary killing context $\tcontextn{i}{}$
such that 
$
\context{ U } \cpar R_i
\longrightarrow
\tcontextn{i}{ V^i_1, V^i_2, \hdots, V^i_{k_i} }
$.

In the former case that $Q_i = R_i$, \ross{Check freshness conditions.} since $\nfv{x}{Q_i}$, the derivation
$\wen x \context{ U } \cpar R_i
\longrightarrow
\new x \context{ U } \cpar R_i
\longrightarrow
\new x \left( \context{ U } \cpar R_i \right)
$ holds.
In the case, $Q_i = \new x R_i$ the derivation $\wen x \context{ U } \cpar \new x R_i
\longrightarrow
\new x \left( \context{ U } \cpar R_i \right)$ holds.
Hence, by applying one of the above derivations for each $i$ we can construct the following derivation for all predicates $U$.
\[
\begin{array}{l}
\wen x \context{ U } \cpar P
\\\quad
\begin{array}{rl}
\longrightarrow&
\wen x \context{ U } \cpar \tcontext{ Q_1, Q_2, \hdots, Q_n }
\\
\longrightarrow&
\tcontext{ \wen x \context{ U } \cpar Q_i \colon 1 \leq i \leq n }
\\
\longrightarrow&
\tcontext{ \new x \left( \context{ U } \cpar R_i \right) \colon 1 \leq i \leq n }
\\
\longrightarrow&
\tcontext{ \new x \tcontextn{i}{ V^i_1, V^i_2, \hdots, V^i_{k_i} } \colon 1 \leq i \leq n }
\end{array}
\end{array}
\]
Thereby the induction invariant is satisfied.
\end{comment}



Consider the case of a context of the form $\forall x. \context{ } \cpar P$, where $\vdash \forall x. \context{ T } \cpar$ holds. By Lemma~\ref{lemma:universal}, for any variable $y$, $\vdash \context{ T } \sub{x}{y} \cpar P$ holds.
%
For name $y$, let $\contextn{y}{}$ be such that for any formula $U$, $\context{U} \sub{x}{y} \equiv \contextn{y}{U\sub{x}{y}}$.
%
For any $y$, by the induction hypothesis, for any formula $U$, there exist formulae $V^y_i$ such that either $V^y_i = \cunit$ or $V^y_i = U\sub{x}{y}\sigma^y_i \cpar W^y_i$, where $\vdash T\sub{x}{y}\sigma^y_i \cpar W^y_i$ holds, for $1 \leq i \leq n$; and $n$-ary killing context $\tcontextn{y}{}$ such that 
$
\context{ U } \sub{x}{y} \cpar P
\equiv
\contextn{y}{ U \sub{x}{y}}  \cpar P$
and the following derivation can be constructed:
\[
\infer[.]{
\contextn{y}{ U \sub{x}{y}}  \cpar P
}{
\tcontextn{y}{ V^{y}_i \colon 1 \leq i \leq n  }
}
\]
Therefore, for  $\nfv{y}{(\forall x. \context{ U } \cpar P)}$ and any $U$,
derivation
\[
\infer[]{
\forall x. \context{ U } \cpar P
}{
\infer[]{
\mathopen{\forall y.} \left(\context{ U } \sub{x}{y} \cpar P \right)
}{
\forall y. \tcontextn{y}{ V^{y}_i \colon 1 \leq i \leq n }
}}
\] 
holds.
In the above $V^{y}_i = \cunit$ or $V^y_i = U\sub{x}{y}\sigma^y_i \cpar W^y_i$, where $\vdash T\sub{x}{y}\sigma^y_i \cpar W^y_i$ holds, for all $i$, as required.

The cases for \textit{plus}, \textit{with}, \textit{tensor} and \textit{seq} do not differ significantly from \textsf{MAV}~\cite{Horne2015}.
\end{proof}
% --- end context-invariant.tex ---



%% Co-rules elimination

%\begin{lemma}[co-close: Lemma~\ref{lemma:co-close}]
%If $\vdash \context{ \wen x. P \tensor \new x. Q }$ holds
%then $\vdash \context{ \mathopen{\wen x.} \left( P \tensor Q \right) }$ holds.
%\end{lemma}
%\input{coclose}

\begin{comment}
\begin{lemma}[co-tidy name: Lemma~\ref{lemma:co-tidynew}]
If $\vdash \context{ \wen x. \cunit }$ holds
then $\vdash \context{ \cunit }$ holds.
\end{lemma}
\begin{proof}
Assume that $\vdash \wen x. \cunit \cpar P$ holds.
By Lemma~\ref{lemma:split-times}, there exists $Q$ such that $\vdash Q$ and $\infer{P}{Q}$.
Hence the following proof of $P$ can be constructed:
$
\vcenter{
\infer{
P
}{
\infer[]{
Q
}{
\cunit
}}}
$.
Therefore, by Lemma~\ref{lemma:context}, for any context $\context{}$, if $\vdash \context{\wen{x}. \cunit}$ then $\vdash \context{ \cunit }$, as required.
\end{proof}
\end{comment}

%
%\begin{lemma}[co-extrude1: Lemma~\ref{lemma:co-extrude}]
%If  $\nfv{x}{Q}$ and $\vdash \context{ \exists x. P \tensor Q }$ holds then $\vdash \context{ \mathopen{\exists x.} \left( P \tensor Q \right) }$ holds.
%\end{lemma}
%\tohide{
%\input{coextrude}
%}{}
%

\begin{comment}
\begin{lemma}[co-tidy1: Lemma~\ref{lemma:co-tidy1}]
If $\vdash \context{ \exists x. \cunit }$ holds
then $\vdash \context{ \cunit }$ holds.
\end{lemma}
\tohide{
\begin{proof}
Assume that $\vdash \exists x. \cunit \cpar T$ holds.
By Lemma~\ref{lemma:split-exists}, there exists $U_i$ such that $\vdash U_i$, for $1 \leq i \leq n$, and $n$-ary killing context $\tcontext{}$ such that 
$\vcenter{\infer{T}{\tcontext{U_1, \hdots, U_n}}}$.
Hence the following proof of $T$ can be constructed:
\[
\infer[.]{
\cunit \cpar T 
}{
\infer[]{
\tcontext{U_1, \hdots, U_n}
}{
\infer[]{
\tcontext{\cunit, \hdots, \cunit}
}{
\cunit
}}}
\]
Therefore, by Lemma~\ref{lemma:context}, if $\vdash \context{\exists{x} \cunit}$ then $\vdash \context{ \cunit }$, as required.
\end{proof}}{}
%
\end{comment}

\begin{lemma}[co-left and co-right: Lemma~\ref{lemma:co-branching}]
If $\vdash \context{ P \wwith Q }$ holds
then both $\vdash \context{ P }$ and $\vdash \context{ Q }$ hold.
\end{lemma}
% --- begin cobranching.tex ---
\begin{proof}
Assume that $\vdash \left(P \wwith Q\right)\mathclose\sigma \cpar R$ holds.
By Lemma~\ref{lemma:split-times}, $\vdash P\mathclose\sigma \cpar R$ and $\vdash Q\mathclose\sigma \cpar R$ hold.
Hence by Lemma~\ref{lemma:context}, for any context $\context{}$, if $\vdash \context{ P \wwith Q }$ then $\vdash \context{ P }$ and $\vdash \context{ Q }$.
\end{proof}
% --- end cobranching.tex ---
%

\begin{lemma}[co-external: Lemma~\ref{lemma:co-external}]
If $\vdash \context{ P \tensor  \left( Q \ooplus R \right) }$ holds
 then $\vdash \context{ \left( P \tensor Q \right) \ooplus \left( P \tensor R \right) }$ holds.
\end{lemma}
\tohide{
% --- begin coexternal.tex ---
\begin{proof}
Assume that $\vdash \left( \left( P \ooplus Q \right) \tensor R\right)\mathclose{\sigma} \cpar W$ holds, for some substitution $\sigma$.
By Lemma~\ref{lemma:split-times}, there exist formulae $T_i$ and $U_i$ such that $\vdash \left( P \ooplus Q \right)\mathclose{\sigma} \cpar T_i$ and $\vdash R\sigma \cpar U_i$, for $1 \leq i \leq n$, and killing context $\tcontext{}$ such that 
\[
\infer[.]{
 W
}{
 \tcontext{ T_1 \cpar U_1, \hdots, T_n \cpar U_n }
}
\]
%
Now, by Lemma~\ref{lemma:split-plus}, for every $i$, there exists killing context $\tcontextn{i}{}$ and types $V^i_j$ such that either $\vdash P\sigma \cpar V^i_j$ or $\vdash Q\sigma \cpar V^i_j$ holds, for $1 \leq j \leq m_i$, and the derivation
\[
\infer[]{
 T_i 
}{
 \tcontextn{i}{ V^i_1, V^i_2, \hdots, V^i_{m_i} }
}
\] holds.

Notice that if $\vdash P\sigma \cpar V^i_j$ holds then the following derivation can be constructed.
\[
\infer[]{
\left( \left(P \tensor R\right) \ooplus \left(Q \tensor R\right) \right)\mathclose{\sigma} \cpar V^i_j  \cpar U_i
}{
\infer[]{
 \left( P \tensor R \right)\mathclose{\sigma} \cpar V^i_j  \cpar U_i 
}{
\infer[]{
 \left( P\sigma  \cpar  V^i_j \right) \tensor \left( R\sigma  \cpar U_i \right) 
}{
 \cunit
}}}
\]
Otherwise $\vdash Q \cpar V^i_j$ holds, hence the following derivation can be constructed.
\[
\infer[]{
\left( \left(P \tensor R\right) \ooplus \left(Q \tensor R \right)\right)\mathclose{\sigma} \cpar V^i_j  \cpar U_i
}{
\infer[]{
 \left( Q \tensor R \right)\mathclose{\sigma} \cpar V^i_j  \cpar U_i 
}{
\infer[]{
 \left( Q\sigma  \cpar  V^i_j \right) \tensor \left( R\sigma  \cpar U_i \right)
}{ \cunit
}}}
\]
Hence by applying one of the above proofs for each $i$ and $j$ we can construct the following proof.
\[
\infer[]{
\left( \left(P \tensor R\right) \ooplus \left(Q \tensor R\right) \right)\mathclose{\sigma} \cpar W 
}{
\infer[]{
\left(\left( P \tensor R\right) \ooplus \left(Q \tensor R\right) \right)\mathclose{\sigma} \cpar \tcontext{ T_1 \cpar U_1, \hdots, T_n \cpar U_n } 
}{
\infer[]{
\left( \left(P \tensor R\right) \ooplus \left(Q \tensor R\right) \right)\mathclose{\sigma} \cpar \tcontext{ \tcontextn{i}{ V^i_1, V^i_2, \hdots, V^i_{m_i} } \cpar U_i \colon 1 \leq i \leq n } 
}{
\infer[]{
\left( \left(P \tensor R\right) \ooplus \left(Q \tensor R\right) \right)\mathclose{\sigma} \cpar \tcontext{ \tcontextn{i}{ V^i_j  \cpar U_i \colon 1 \leq j \leq m_i } \colon 1 \leq i \leq n }
}{
\infer[]{
 \tcontext{ \left( \left(P \tensor R\right) \ooplus \left(Q \tensor R \right)\right)\mathclose{\sigma} \cpar \tcontextn{i}{ V^i_j  \cpar U_i \colon 1 \leq j \leq m_i } \colon 1 \leq i \leq n } 
}{
\infer[]{
 \tcontext{ \tcontextn{i}{ \left(\left( P \tensor R\right) \ooplus \left(Q \tensor R \right)\right)\mathclose{\sigma} \cpar V^i_j  \cpar U_i \colon 1 \leq j \leq m_i } \colon 1 \leq i \leq n } 
%   \longrightarrow \tcontext{ \left( P \tensor R \ooplus Q \tensor R \right) \cpar  \tcontextn{i}{ V^i_1  \cpar U_1, V^i_2  \cpar U_i, \hdots, V^i_{m_i}  \cpar U_i} \colon 1 \leq i \leq n } \\
%   \longrightarrow \tcontext{ \tcontextn{i}{ \left( P \tensor R \ooplus Q \tensor R \right) \cpar V^i_j  \cpar U_i  \colon 1 \leq j \leq m_i }  \colon 1 \leq i \leq n } \\
}{
\infer[]{
 \tcontext{ \tcontextn{i}{ \cunit \colon 1 \leq j \leq m_i } \colon 1 \leq i \leq n } 
}{
 \cunit
}}}}}}}
\]
Hence $\vdash \left(\left( P \tensor R\right) \ooplus \left(Q \tensor R\right) \right) \cpar W$.
Therefore, by Lemma~\ref{lemma:context}, for any context $\vdash \context{ \left( P \ooplus Q \right) \tensor R }$ yields $\vdash \context{ \left(P \tensor R\right) \ooplus \left(Q \tensor R\right) }$, as required.
\end{proof}
% --- end coexternal.tex ---
}{}

\begin{lemma}[co-sequence: Lemma~\ref{lemma:co-sequence}]
If $\vdash \context{ \left( P \cseq Q \right) \tensor \left( R \andthen S \right) }$ holds
 then $\vdash \context{ \left( P \tensor R \right) \cseq \left( Q \tensor S \right) }$ holds.
\end{lemma}
\tohide{
% --- begin coseq.tex ---
\begin{proof}
Assume that $\vdash \left( \left( P \cseq Q \right) \tensor \left( R \cseq S \right) \right)\mathclose{\sigma} \cpar U$ holds, for some substitution $\sigma$.
By Lemma~\ref{lemma:split-times}, there exist $n$-ary killing context $\tcontext{}$ and $U^0_i$ and $U^1_i$, for $1 \leq i \leq n$, such that $\vdash \left(P \cseq Q\right)\mathclose{\sigma} \cpar U^0_i$ and $\vdash \left(R \cseq S\right)\sigma \cpar U^1_i$ and the derivation
\[
\infer{
 U
}{
 \tcontext{ U^0_1 \cpar U^1_1, U^0_2 \cpar U^1_2, \hdots }
}
\]  holds.

Hence by Lemma~\ref{lemma:split-times}, for $k \in \left\{0, 1\right\}$ there exists $m^k_i$-ary killing context $\tcontextmn{k}{i}{}$ and types $V^k_{i,j}$, $W^k_{i,j}$ for $1 \leq j \leq m_i^k$, such that 
$\vdash P\sigma \cpar V^0_{i,j}$ and $\vdash Q \sigma\cpar W^0_{i,j}$ and $\vdash R\sigma \cpar V^1_{i,j}$ and $\vdash S\sigma \cpar W^1_{i,j}$
and the following derivation
\[
\infer{
U^k_i
}{
 \tcontextmn{k}{i}{V^k_{i,1}~\andthen~W^k_{i,1}, V^k_{i,2}~\andthen~W^k_{i,2} \hdots}
}
\]  holds.

Hence we can construct the following proof.
\[
\infer[]{
\left(\left(P \tensor R\right) \cseq \left(Q \tensor S\right)\right)\mathclose{\sigma} \cpar U
}{
\infer[]{
 \left(\left(P \tensor R\right) \cseq \left(Q \tensor S\right)\right)\mathclose{\sigma} \cpar
 \tcontext{ U^0_1 \cpar U^1_1, U^0_2 \cpar U^1_2, \hdots } 
}{
\infer[]{
 \left(\left(P \tensor R\right) \cseq \left(Q \tensor S\right)\right)\mathclose{\sigma} \cpar
 \tcontext{
 \begin{array}{rl}
      &\tcontextmn{0}{i}{V^0_{i,j} \cseq W^0_{i,j} \colon 1 \leq j \leq m^0_i } \\
 \cpar&\tcontextmn{1}{i}{V^1_{i,k} \cseq W^1_{i,k} \colon 1 \leq k \leq m^1_i } 
 \end{array}
   \colon 1 \leq i \leq n                                               
 }  
}{
\infer[]{
 \left(\left(P \tensor R\right) \cseq \left(Q \tensor S\right)\right)\mathclose{\sigma}
   \cpar
   \tcontext{
\begin{array}{r}
     \tcontextmn{1}{i}{
       \begin{array}{r}
       \tcontextmn{0}{i}{
          V^0_{i,j} \cseq W^0_{i,j}
           \colon 1 \leq j \leq m^0_i
       }
       \cpar
       \left(V^1_{i,k} \cseq W^1_{i,k}\right) %\\
\\
       \colon 1 \leq k \leq m^1_i
       \end{array}
       }
\\
  \colon 1 \leq i \leq n 
\end{array}
                                       }  
}{
\infer[]{
 \left(\left(P \tensor R\right) \cseq \left(Q \tensor S\right)\right)\mathclose{\sigma}
   \cpar
   \tcontext{
     \begin{array}{r}
     \tcontextmn{1}{i}{
       \begin{array}{r}
       \tcontextmn{0}{i}{
          \begin{array}{l}
          \left(V^0_{i,j} \cseq W^0_{i,j}\right) \cpar
       \left(V^1_{i,k} \cseq W^1_{i,k}\right)
           \colon 1 \leq j \leq m^0_i
          \end{array}
       }
       \\
       \colon 1 \leq k \leq m^1_i
       \end{array} }
     \\ \colon 1 \leq i \leq n 
     \end{array}
                                       }  
}{
\infer[]{
 \left(\left(P \tensor R\right) \cseq \left(Q \tensor S\right)\right)\sigma
   \cpar
   \tcontext{
     \begin{array}{r}
     \tcontextmn{1}{i}{
       \begin{array}{r}
       \tcontextmn{0}{i}{
          \begin{array}{l}
          \left(V^0_{i,j} \cpar V^1_{i,k}\right) \cseq
       \left(W^0_{i,j} \cpar W^1_{i,k}\right)
           \colon 1 \leq j \leq m^0_i
          \end{array}
       }
       \\
       \colon 1 \leq k \leq m^1_i
       \end{array} }
     \\ \colon 1 \leq i \leq n 
     \end{array}
                                       }  
}{
\infer[]{
   \tcontext{
%     \begin{array}{l}
     \tcontextmn{1}{i}{
 %      \begin{array}{l}
       \tcontextmn{0}{i}{
          \begin{array}{l}
   \left(\left(P \tensor R\right) \cseq \left(Q \tensor S\right)\right)\mathclose{\sigma}
   \cpar \\
      \left(\left(V^0_{i,j} \cpar V^1_{i,k}\right) \cseq
       \left(W^0_{i,j} \cpar W^1_{i,k}\right)\right)
          \end{array} \colon 1 \leq j \leq m^0_i
       }
       %\\
       \colon 1 \leq k \leq m^1_i
  %     \end{array}
      }
     %\\
     \colon 1 \leq i \leq n 
%     \end{array}
  }  
}{
\infer[]{
   \tcontext{
     %\begin{array}{l}
     \tcontextmn{1}{i}{
       %\begin{array}{l}
       \tcontextmn{0}{i}{
         \begin{array}{l}
         \left(\left(P \tensor R\right)\sigma \cpar V^0_{i,j} \cpar V^1_{i,k}\right)
         \cseq \\ 
         \left(
          \left(Q \tensor S\right)\sigma \cpar
          W^0_{i,j} \cpar W^1_{i,k}\right)
        \end{array}
           \colon 1 \leq j \leq m^0_i 
       }
       \colon 1 \leq k \leq m^1_i
       %\end{array}
     }
     \colon 1 \leq i \leq n 
   } 
}{
\infer[]{
   \tcontext{
%     \begin{array}{l}
     \tcontextmn{1}{i}{
 %      \begin{array}{l}
       \tcontextmn{0}{i}{
          \begin{array}{l}
\left(\left(P\sigma \cpar V^0_{i,j}\right) \tensor \left(R\sigma \cpar V^1_{i,k}\right)\right)
   \cseq \\
      \left(
          \left(Q\sigma \cpar W^0_{i,j}\right) \tensor
       \left(S\sigma \cpar W^1_{i,k}\right)\right)
\end{array}
           \colon 1 \leq j \leq m^0_i          
       }
       %\\
       \colon 1 \leq k \leq m^1_i
%       \end{array}
       }
     %\\ 
     \colon 1 \leq i \leq n 
%     \end{array}
                                       }  
}{
\infer[]{
   \tcontext{
     \begin{array}{l}
     \tcontextmn{1}{i}{
       \begin{array}{l}
       \tcontextmn{0}{i}{
          \begin{array}{l}
          \cunit
           \colon 1 \leq j \leq m^0_i
          \end{array}
       }
       \colon 1 \leq k \leq m^1_i
       \end{array} }
      \colon 1 \leq i \leq n 
     \end{array}
                                       }  
}{
 \cunit
}}}}}}}}}}
\]
%Hence $\vdash \left(\left(P \tensor R\right) \cseq \left(Q \tensor S\right)\right) \cpar P$ holds.
Therefore, by Lemma~\ref{lemma:context}, for any context $\vdash \context{\left(P \cseq Q\right) \tensor \left(R \cseq S\right)}$ yields $\vdash \context{ \left(P \tensor R\right) \cseq \left(Q \tensor S\right) }$.
\end{proof}

% --- end coseq.tex ---
}{}
%

\begin{lemma}[co-tidy: Lemma~\ref{lemma:co-tidy}]
If $\vdash \context{ \cunit \ooplus \cunit }$ holds, then $\vdash \context{ \cunit }$ holds.
\end{lemma}
\begin{proof}
Assume that $\vdash \left( \cunit \ooplus \cunit \right) \cpar P$ holds. By Lemma~\ref{lemma:split-plus}, there exist killing context $\tcontext{}$
and formulae $U_i$ for $1 \leq i \leq n$ such that $\vdash \cunit \cpar U_i$ or $\vdash \cunit \cpar U_i$ hold, hence $\vdash U_i$ holds, and the following derivation can be constructed.
\[
\infer[.]{
P
}{
 \tcontext{ U_1, \hdots, U_n }
}
\]
Thereby the following proof can be constructed:
\[
%\begin{array}{rl}
\infer[.]{
 P
}{
\infer[]{
 \tcontext{ U_1, \hdots, U_n } 
}{
\infer[]{
 \tcontext{ \cunit, \cunit, \hdots } 
}{ \cunit
}}}
%\end{array}
\]
Therefore, by Lemma~\ref{lemma:context}, for any context $\vdash \context{ \cunit \ooplus \cunit }$ yields $\vdash \context{ \cunit }$, as required.
\end{proof}
%
\begin{lemma}[atomic co-interaction: Lemma~\ref{lemma:co-atoms}]
If $\vdash \context{ \alpha \tensor \co{\alpha} }$ holds
then $\vdash \context{ \cunit }$ holds.
\end{lemma}
% --- begin coatoms.tex ---
\begin{proof}
Assume for atom $\alpha$ that %and $\beta$, where $\alpha \leq \beta$, 
$\vdash \left( \alpha \tensor \co{\alpha} \right)\sigma \cpar P$, for some formula $P$ and some substitution $\sigma$.
By Lemma~\ref{lemma:split-times}, there exist $n$-ary killing context $\tcontext{}$ and formulae $U_i$ and $V_i$ such that $\vdash \alpha \sigma \cpar U_i$ and $\vdash \co{\alpha \sigma} \cpar ~V_i$, for $1 \leq i \leq n$, such that 
\[
\infer[.]{P}{
\tcontext{ U_1 \cpar V_1, U_2 \cpar  V_2, \hdots}
}
\]
By Lemma~\ref{lemma:split-atoms}, for every $i$, there exist $m_i^0$-ary killing contexts $\tcontextmn{0}{i}{}$
%and atoms $\gamma_j^i$ such that $\gamma^i_j \leq \alpha \sigma$ for $1 \leq j \leq m_i^0$
such that 
\[
\infer[.]{U_i}{
\tcontextmn{0}{i}{\co{\alpha \sigma}, \hdots, \co{\alpha \sigma}}
}
\]
By Lemma~\ref{lemma:split-atoms}, for every $i$, there exist $m_i^1$-ary killing contexts $\tcontextmn{1}{i}{}$
%and atoms $\delta_i^k$ such that $\beta \sigma \leq \delta_i^k$ for $1 \leq k \leq m_i^1$
such that 
\[
\infer[.]{
V_i}{
 \tcontextmn{1}{i}{\alpha \sigma, \hdots, \alpha \sigma}
}
\]
%Note that, where we assume that if $\alpha \leq \beta$ then $\alpha\sigma \leq \beta\sigma$, due to closure of the preorder with respect to substitution.
%Also, since $\gamma_i^j \leq \alpha \sigma$ and $\alpha \sigma \leq \beta \sigma$ and $\beta \sigma \leq \delta_i^k$, by the transitivity of $\leq$ for atoms, $\gamma_i^j \leq \delta_i^k$.
Thereby the following proof can be constructed.
\[
\begin{array}{rl}
\infer[]{
P
}{
\infer[]{
 \tcontext{U_1 \cpar V_1, U_2 \cpar V_2, \hdots} 
}{
\infer[]{
 \tcontext{
                               \tcontextmn{0}{i}{\co{\alpha \sigma} \colon 1 \leq j \leq m_i^0} 
                         \cpar \tcontextmn{1}{i}{\alpha \sigma \colon 1 \leq k \leq m_i^1 }
                             \colon 1 \leq i \leq n 
}
}{
\infer[]{
 \tcontext{
    \tcontextmn{1}{i}{\tcontextmn{0}{i}{\co{\alpha \sigma} \colon 1 \leq j \leq m_i^0} \cpar \alpha \sigma \colon 1 \leq k \leq m_i^1 }
    \colon 1 \leq i \leq n
}
}{
\infer[]{
 \tcontext{
    \tcontextmn{1}{i}{\tcontextmn{0}{i}{\co{\alpha \sigma} \cpar \alpha \sigma \colon 1 \leq j \leq m_i^0} \colon 1 \leq k \leq m_i^1 } 
    \colon 1 \leq i \leq n 
}
}{
\infer[]{
 \tcontext{
    \tcontextmn{1}{i}{\tcontextmn{0}{i}{\cunit \colon 1 \leq j \leq m_i^0} \colon 1 \leq k \leq m_i^1 }
\colon 1 \leq i \leq n }
}{
 \cunit
}}}}}}
\end{array}
\]
Therefore, by Lemma~\ref{lemma:context}, for any context $\context{}$, $\vdash \context{\alpha \tensor \co{\alpha}}$ yields that $\vdash \context{ \cunit }$, as required.
\end{proof}
% --- end coatoms.tex ---

% --- end appendix.tex ---

\end{document}


